\documentclass[13pt,openany]{book}

% === GEOMETRY ===
\usepackage[
    inner=0.75in,
    outer=0.65in,
    top=0.8in,
    bottom=0.8in,
    headheight=15pt,
    includefoot
]{geometry}

% === FONTS & ENCODING (pdfLaTeX optimized) ===
\usepackage[T1]{fontenc}
\usepackage[utf8]{inputenc}
\usepackage{lmodern}           % High-quality Latin Modern fonts (much better than default CM)

\usepackage{cjhebrew}          % Reliable Hebrew for pdfLaTeX

% === OTHER PACKAGES (Careful order) ===
\usepackage{xcolor}
\usepackage{titlesec}
\usepackage{array}
\usepackage{amsmath}
\DeclareMathOperator{\Sel}{Sel}
\usepackage{amscd}
\usepackage{amsthm}
\usepackage{amssymb}
\usepackage{enumitem}
\usepackage{bm}
\usepackage{mathtools}
\usepackage{physics}
\usepackage{booktabs}
\usepackage{verbatim}
\usepackage{float}
\usepackage{hyperref}
\usepackage{cleveref}
\usepackage{longtable}
\usepackage{listings}
\lstset{
  basicstyle=\ttfamily\small,
  breaklines=true,
  escapechar=|,
  literate=
    {χ}{{$\chi$}}1
    {≈}{{$\approx$}}1
    {Φ}{{$\Phi$}}1
    {Ĵ}{{$\hat{J}$}}1
    {Ê}{{$\hat{K}$}}1
    {Ψ}{{$\Psi$}}1
    {Ω}{{$\Omega$}}1
}

% === TATE-SHAFAREVICH GROUP SYMBOL ===
\DeclareFontFamily{U}{wncy}{}
\DeclareFontShape{U}{wncy}{m}{n}{<->wncyr10}{}
\DeclareSymbolFont{mcy}{U}{wncy}{m}{n}
\DeclareMathSymbol{\Sha}{\mathord}{mcy}{"58}

\newcommand{\hb}[1]{\cjRL{#1}}
\usepackage{parskip}

% Add this to your document preamble
\setcounter{secnumdepth}{4}
\setcounter{tocdepth}{4}

% === THEOREM STYLES ===
\theoremstyle{plain}
\newtheorem{theorem}{Theorem}[chapter]
\newtheorem{postulate}{Postulate}[chapter]
\newtheorem{corollary}{Corollary}[theorem]
\newtheorem{lemma}[theorem]{Lemma}
\newtheorem{remark}{Remark}[chapter]

\theoremstyle{definition}
\newtheorem{definition}{Definition}[chapter]
\newtheorem{protocol}{Protocol}[chapter]
\newtheorem{proposition}{Proposition}[chapter]
\newtheorem{principle}{Principle}[chapter]

\renewcommand{\qedsymbol}{$\blacksquare$}

% === HEADER/FOOTER SETUP ===
\usepackage{fancyhdr}
\pagestyle{fancy}
\fancyhf{} % Clear default headers and footers

% [LE,RO] means Left on Even pages, Right on Odd pages (outside edges)
\fancyhead[LE,RO]{\thepage} 

% [RE] means Right on Even pages (inside edge for Chapter title)
\fancyhead[RE]{\nouppercase{\leftmark}} 

% [LO] means Left on Odd pages (inside edge for Section title)
\fancyhead[LO]{\nouppercase{\rightmark}} 

\renewcommand{\headrulewidth}{0.4pt}

\pagenumbering{roman}

\usepackage[titles]{tocloft}

% Make Part entries and their page numbers bold in the ToC
\renewcommand{\cftpartfont}{\bfseries}
\renewcommand{\cftpartpagefont}{\bfseries}

% Make Chapter entries and their page numbers normal (not bold)
\renewcommand{\cftchapfont}{\normalfont}
\renewcommand{\cftchappagefont}{\normalfont}

\begin{document}

% === TITLE PAGE ===
\begin{titlepage}
    \centering
    % The top fill pushes the content down so it centers vertically
    \vspace*{\fill} 

    % Top horizontal line
    \noindent\rule{\textwidth}{1pt} \\[1.5em]

    {\Huge \textbf{Continuation Mathematics}} \\[1.5em]

    {\Large \textit{A New Classical Foundation for Mathematics}} \\[1.2em]

    % Bottom horizontal line
    \noindent\rule{\textwidth}{1pt} \\[3cm]

    {\Large \textbf{Samir Amier Saliem Boulos}} \\
    \vspace{1cm}
    {\large June 2026}

    % The bottom fill pushes up against the top fill to perfectly center the block
    \vspace*{\fill} 
\end{titlepage}


\frontmatter

% --- Dedication ---
\cleardoublepage
\thispagestyle{empty}
\vspace*{0.3\textheight}
\begin{center}
    {\Large\itshape To my Lord, Saviour, King, and God, Jesus Christ.}
\end{center}
\clearpage

% --- Epigraph ---
\cleardoublepage
\thispagestyle{empty}
\vspace*{0.3\textheight}
\begin{flushright}
    \begin{minipage}{0.7\textwidth}
        \raggedleft
        \Large\itshape
        ``In the beginning was the Word, and the Word was with God, and the Word was God... And the Word became flesh and dwelt among us.''
        
        \vspace{0.5em}
        \normalsize\normalfont\textsc{— John 1:1, 14}
    \end{minipage}
\end{flushright}
\clearpage

% === TABLE OF CONTENTS ===
\tableofcontents

% --- ABSTRACT (Technical) ---

\chapter*{Abstract}
\addcontentsline{toc}{chapter}{Abstract}

This work develops a new classical foundation for mathematics in which
\emph{admissibility} replaces mathematical objects as the primitive
mathematical notion. Rather than beginning with sets, spaces, algebraic
operations, categories, geometric entities, or any other established
mathematical structures, the theory begins by determining what is
intrinsically admissible. From this single primitive there emerges, by
necessary generation, a hierarchy of mathematical structures including
constraints, relations, continuation systems, completion structures,
distinguishability, information, and realization.

Beginning from admissibility alone, the theory develops an autonomous branch of
classical mathematics whose constructions are generated rather than postulated.
Continuation spaces, completion spaces, continuation algebra, continuation
geometry, universal continuation theory, and admissible mathematics are shown
to arise as successive stages of one canonical structural development. Every
construction is carried out within ordinary classical mathematics, without
appeal to non-classical logic, alternative foundational systems, or physically
motivated assumptions.

A central theme throughout the work is the distinction between
\emph{intrinsic mathematical structure} and its \emph{realizations}. The theory
shows that many familiar mathematical disciplines arise as faithful
realizations of common admissibility structures. In particular, classical set
theory, topology, algebra, geometry, analysis, probability theory, information
theory, category theory, and the mathematical framework underlying quantum
theory are recovered as distinct realizations of a single foundational
architecture. Their apparent diversity reflects differences of realization
rather than differences of mathematical foundation.

The development culminates in two general theories. The first is the
\emph{Theory of Admissible Mathematics}, which establishes that every
mathematical theory is generated from an underlying admissibility structure.
The second is the \emph{Mathematics of Mathematical Discovery}, in which
canonical mathematical investigation, theorem spaces, theorem closure,
structural fixed points, and discovery itself are developed as mathematical
objects. Within this framework, mathematical discovery becomes a canonical
structural process governed by precise mathematical laws rather than by
heuristic invention.

The resulting framework presents mathematics not as the study of isolated
objects but as the study of admissibility and of the canonical generation of
mathematical structure. It provides both a new classical foundation for
mathematics and a constitutional theory describing how mathematical knowledge
itself is systematically discovered.

This volume is devoted exclusively to the development of this foundational
framework. It is entirely independent of any particular mathematical
application. Applications to specific mathematical problems are developed
separately, where classical theories are reconstructed through the principles
of admissibility, continuation, canonical investigation, and structural
closure established in the present work.

\mainmatter

\part{The Constitution of Continuation}

%%%%%%%%%%%%%%%%%%%%%%%%%%%%%%%%%%%%%%%%%%%%%%%%%%%%%%%%%%%%%%%%%%%%%%%%%%%%%%
\chapter{The Insufficiency of Completed Mathematics}
%%%%%%%%%%%%%%%%%%%%%%%%%%%%%%%%%%%%%%%%%%%%%%%%%%%%%%%%%%%%%%%%%%%%%%%%%%%%%%

\section{Introduction}

Modern mathematics studies mathematical objects after they have been specified.
Groups, rings, topological spaces, manifolds, graphs, algorithms, formal proofs,
and dynamical systems are all treated as completed entities possessing
well-defined internal structure.

Yet the construction of these objects almost never proceeds in a single step.
Proofs are assembled statement by statement, graphs edge by edge,
algorithms instruction by instruction, and algebraic structures relation by
relation. Even when a mathematical object is ultimately finite, it is ordinarily
encountered through a sequence of partial constructions whose intermediate
stages possess genuine mathematical content.

These intermediate stages are ubiquitous throughout mathematics.
Nevertheless, they do not presently constitute an autonomous mathematical
subject. They are usually regarded merely as temporary artefacts of a
construction rather than as mathematical objects worthy of independent study.

The purpose of this chapter is to demonstrate that this omission is structural
rather than cosmetic. The absence of an autonomous theory of partial
mathematical realization obscures a common phenomenon that appears across
numerous branches of mathematics.

%%%%%%%%%%%%%%%%%%%%%%%%%%%%%%%%%%%%%%%%%%%%%%%%%%%%%%%%%%%%%%%%%%%%%%%%%%%%%%

\section{Completed objects and partial constructions}

A completed mathematical object is one that satisfies the defining conditions of
its ambient theory.

Examples include:

\begin{itemize}
\item a completed proof in formal logic;
\item a group satisfying the group axioms;
\item a graph with its entire edge set specified;
\item a complete colouring of a graph;
\item a fully specified algorithm;
\item a complete solution of a system of equations.
\end{itemize}

Each of these objects admits numerous intermediate stages during its
construction.

For example, a proof consists of finitely many individual deductions.
Removing the final deduction generally produces a mathematically meaningful
partial proof. Likewise, removing several edges from a graph or several
relations from a group presentation generally produces another mathematically
meaningful object that is incomplete relative to the original construction.

The mathematical significance of these intermediate stages depends neither on
their order of discovery nor on the process by which they were produced.
Rather, they possess intrinsic structural properties of their own.

%%%%%%%%%%%%%%%%%%%%%%%%%%%%%%%%%%%%%%%%%%%%%%%%%%%%%%%%%%%%%%%%%%%%%%%%%%%%%%

\section{The asymmetry of existing mathematical theories}

Existing mathematical disciplines overwhelmingly emphasize completed
structures.

Group theory studies groups.

Topology studies topological spaces.

Measure theory studies measurable spaces.

Graph theory studies graphs.

Model theory studies models.

Proof theory studies completed derivations.

In each case, the completed object forms the primary subject of investigation.

The intermediate stages leading to those objects are usually treated only as
auxiliary devices employed during proofs or constructions. Their mathematical
properties are seldom investigated independently of the completed object that
eventually contains them.

This asymmetry is striking.

Suppose two finite constructions share identical initial stages but diverge later.
Classical mathematics typically regards these common initial stages merely as
subsets or prefixes of larger objects.

However, those initial stages already possess structural information concerning
their possible future development. Existing mathematical language contains no
general framework for studying that information independently of the completed
objects into which it may eventually evolve.

%%%%%%%%%%%%%%%%%%%%%%%%%%%%%%%%%%%%%%%%%%%%%%%%%%%%%%%%%%%%%%%%%%%%%%%%%%%%%%

\section{The extension problem}

Every partial mathematical construction gives rise to a fundamental question.

\begin{quote}
Given the present stage of construction, which further constructions remain
possible?
\end{quote}

This question appears in numerous mathematical settings.

A partial proof may admit several valid next deductions.

A partially coloured graph may permit certain colourings while forbidding
others.

A partially specified algebraic presentation may admit additional defining
relations without becoming inconsistent.

A partially constructed combinatorial object may extend in many different ways,
or perhaps in none.

In every case, the central mathematical question concerns the collection of
admissible extensions of the current object.

This collection is itself a mathematical object.

Nevertheless, no general mathematical theory presently studies such collections
independently of the particular discipline in which they arise.

%%%%%%%%%%%%%%%%%%%%%%%%%%%%%%%%%%%%%%%%%%%%%%%%%%%%%%%%%%%%%%%%%%%%%%%%%%%%%%

\section{A structural omission}

The preceding observations reveal a common structural feature shared by many
apparently unrelated areas of mathematics.

Each possesses

\begin{enumerate}
\item partial constructions;
\item admissible extensions of those constructions;
\item forbidden extensions;
\item completed realizations;
\item maximal constructions beyond which no further extension is possible.
\end{enumerate}

These notions recur throughout mathematics with remarkable consistency.

Despite their ubiquity, they have no common mathematical language.
Instead, each discipline develops its own specialised terminology and methods.

The repeated appearance of the same structural pattern suggests that these
objects deserve investigation independently of any particular application.

%%%%%%%%%%%%%%%%%%%%%%%%%%%%%%%%%%%%%%%%%%%%%%%%%%%%%%%%%%%%%%%%%%%%%%%%%%%%%%

\section{The objective of this work}

The aim of the present work is not to replace existing mathematical theories.

Rather, its purpose is to isolate and develop the mathematical structures
associated with partial realization, admissible extension, and completed
construction as autonomous objects of study.

The theory developed in the following chapters begins with no assumptions about
numbers, geometry, algebra, topology, computation, or logic beyond those
already available in classical mathematics.

Instead, it seeks to identify the minimal mathematical structures required for
the study of partial constructions and their admissible continuations.

Only after these structures have been established abstractly will particular
areas of mathematics appear as concrete realizations of the general theory.

%%%%%%%%%%%%%%%%%%%%%%%%%%%%%%%%%%%%%%%%%%%%%%%%%%%%%%%%%%%%%%%%%%%%%%%%%%%%%%

\section{The constitutional principle}

The development of this book is guided by a single methodological principle.

\begin{quote}
Whenever a mathematical phenomenon appears independently across multiple
disciplines and admits an intrinsic structural description, it should be studied
as an autonomous mathematical object.
\end{quote}

The repeated appearance of partial constructions, admissible extensions,
realizations, and maximal completions satisfies precisely this criterion.

The remainder of this work develops the mathematical theory that follows from
taking these structures as primary.

%%%%%%%%%%%%%%%%%%%%%%%%%%%%%%%%%%%%%%%%%%%%%%%%%%%%%%%%%%%%%%%%%%%%%%%%%%%%%%
\chapter{Partial Mathematical Objects}
%%%%%%%%%%%%%%%%%%%%%%%%%%%%%%%%%%%%%%%%%%%%%%%%%%%%%%%%%%%%%%%%%%%%%%%%%%%%%%

\begin{quote}
\emph{Classical mathematics ordinarily studies completed objects.
Continuation Mathematics begins one step earlier and studies objects
that have not yet completed.}
\end{quote}

\vspace{1em}

Every branch of mathematics contains objects that are naturally
constructed by finite stages. Sequences are extended one term at a time.
Groups are generated by adjoining generators. Graphs grow by adding
vertices and edges. Differential equations evolve through time.
Algorithms reveal progressively more information as computation proceeds.

Classically, these evolving constructions are usually regarded merely as
methods for producing the completed object. The intermediate stages are
treated as temporary approximations possessing no independent
mathematical status.

Continuation Mathematics rejects this viewpoint.

The intermediate stages are themselves mathematical objects. They possess
their own structure, their own morphisms, their own invariants, and their
own laws of evolution. Completion is not primitive. Completion is a
property that certain partial objects may eventually possess.

The purpose of this chapter is to isolate the notion of a
\emph{partial mathematical object} independently of any particular
mathematical discipline.

%%%%%%%%%%%%%%%%%%%%%%%%%%%%%%%%%%%%%%%%%%%%%%%%%%%%%%%%%%%%%%%%%%%%%%%%%%%%%%
\section{The Classical Paradigm}
%%%%%%%%%%%%%%%%%%%%%%%%%%%%%%%%%%%%%%%%%%%%%%%%%%%%%%%%%%%%%%%%%%%%%%%%%%%%%%

Most mathematical theories begin by specifying completed objects.

Typical examples include

\begin{itemize}
\item a complete graph,
\item a completed group,
\item a finished manifold,
\item an infinite sequence,
\item an infinite series,
\item a complete proof,
\item an entire dynamical orbit.
\end{itemize}

Finite constructions appear only as devices for reaching these objects.

For example,

\[
G_0
\subseteq
G_1
\subseteq
G_2
\subseteq
\cdots
\subseteq
G
\]

is interpreted merely as a construction of the final graph \(G\).

Likewise,

\[
a_1,a_2,\ldots,a_n,\ldots
\]

is regarded as one infinite sequence rather than an infinite family of
finite sequences.

Continuation Mathematics reverses this perspective.

The primary objects are

\[
G_0,\;
G_1,\;
G_2,\;
\dots
\]

rather than the eventual limit.

%%%%%%%%%%%%%%%%%%%%%%%%%%%%%%%%%%%%%%%%%%%%%%%%%%%%%%%%%%%%%%%%%%%%%%%%%%%%%%
\section{The Principle of Partiality}
%%%%%%%%%%%%%%%%%%%%%%%%%%%%%%%%%%%%%%%%%%%%%%%%%%%%%%%%%%%%%%%%%%%%%%%%%%%%%%

The central observation is remarkably simple.

\begin{quote}
Every completed mathematical object is preceded by a family of partial
objects.
\end{quote}

These partial objects already possess mathematical structure.

They may

\begin{itemize}
\item satisfy identities,
\item possess symmetries,
\item admit morphisms,
\item carry invariants,
\item support measures,
\item admit extensions,
\item fail to admit extensions.
\end{itemize}

Consequently they deserve independent mathematical study.

This leads to the first primitive concept of Continuation Mathematics.

%%%%%%%%%%%%%%%%%%%%%%%%%%%%%%%%%%%%%%%%%%%%%%%%%%%%%%%%%%%%%%%%%%%%%%%%%%%%%%
\section{Partial Objects}
%%%%%%%%%%%%%%%%%%%%%%%%%%%%%%%%%%%%%%%%%%%%%%%%%%%%%%%%%%%%%%%%%%%%%%%%%%%%%%

\begin{definition}[Partial mathematical object]
A \emph{partial mathematical object} is an object whose mathematical
description is not assumed to be complete but which admits at least one
legal continuation.
\end{definition}

Nothing in this definition refers to

\begin{itemize}
\item topology,
\item algebra,
\item geometry,
\item arithmetic,
\item logic,
\item computation,
\item dynamics.
\end{itemize}

The definition is intentionally universal.

Only one primitive notion appears:

\[
\boxed{\text{continuation}.}
\]

Everything else will be derived from it.

%%%%%%%%%%%%%%%%%%%%%%%%%%%%%%%%%%%%%%%%%%%%%%%%%%%%%%%%%%%%%%%%%%%%%%%%%%%%%%
\section{Partiality versus Approximation}
%%%%%%%%%%%%%%%%%%%%%%%%%%%%%%%%%%%%%%%%%%%%%%%%%%%%%%%%%%%%%%%%%%%%%%%%%%%%%%

Partiality should not be confused with approximation.

An approximation attempts to estimate a completed object.

A partial object need not estimate anything.

It merely records the information presently available.

For example,

\[
(3,5,8)
\]

is not an approximation to

\[
(3,5,8,13,21,\ldots).
\]

It is simply a shorter object.

Likewise,

a finite proof is not an approximate proof.

It is an incomplete proof.

Likewise,

a finite path inside a graph is not an approximate infinite path.

It is literally an initial segment.

Continuation Mathematics therefore separates

\[
\text{approximation}
\]

from

\[
\text{partiality}.
\]

These concepts are logically independent.

%%%%%%%%%%%%%%%%%%%%%%%%%%%%%%%%%%%%%%%%%%%%%%%%%%%%%%%%%%%%%%%%%%%%%%%%%%%%%%
\section{Completion is a Property}
%%%%%%%%%%%%%%%%%%%%%%%%%%%%%%%%%%%%%%%%%%%%%%%%%%%%%%%%%%%%%%%%%%%%%%%%%%%%%%

The traditional viewpoint implicitly assumes that completion is primitive.

Continuation Mathematics reverses this order.

Completion is merely a property possessed by certain partial objects.

\begin{definition}[Completed object]
A partial object is called
\emph{completed}
when no further continuation is required by the mathematical context under
consideration.
\end{definition}

Notice that this definition depends upon the continuation system itself.

A graph may be complete inside one continuation system while remaining
extendable inside another.

Completion is therefore relative.

Continuation is fundamental.

%%%%%%%%%%%%%%%%%%%%%%%%%%%%%%%%%%%%%%%%%%%%%%%%%%%%%%%%%%%%%%%%%%%%%%%%%%%%%%
\section{Examples}
%%%%%%%%%%%%%%%%%%%%%%%%%%%%%%%%%%%%%%%%%%%%%%%%%%%%%%%%%%%%%%%%%%%%%%%%%%%%%%

The same primitive idea appears throughout mathematics.

\subsection*{Finite sequences}

Every finite sequence

\[
(a_1,\ldots,a_n)
\]

is a partial object.

Appending another term produces a continuation.

\vspace{0.5em}

\subsection*{Finite graphs}

A graph together with its current vertices and edges is partial.

Adding one edge is a continuation.

Adding one vertex is another continuation.

\vspace{0.5em}

\subsection*{Generated groups}

A subgroup generated by finitely many generators is partial.

Adjoining another generator is a continuation.

\vspace{0.5em}

\subsection*{Proofs}

A proof under construction is partial.

Appending another valid inference is a continuation.

\vspace{0.5em}

\subsection*{Collatz}

A finite Collatz trajectory is partial.

Applying one further Collatz step is a continuation.

Notice that nothing in these examples depends upon the internal algebra of
the objects.

Only continuation matters.

%%%%%%%%%%%%%%%%%%%%%%%%%%%%%%%%%%%%%%%%%%%%%%%%%%%%%%%%%%%%%%%%%%%%%%%%%%%%%%
\section{Intrinsic Questions}
%%%%%%%%%%%%%%%%%%%%%%%%%%%%%%%%%%%%%%%%%%%%%%%%%%%%%%%%%%%%%%%%%%%%%%%%%%%%%%

Once partial objects are regarded as genuine mathematical objects,
entirely new questions arise.

Given a partial object \(X\),

one may ask

\begin{enumerate}
\item Does \(X\) admit any continuation?
\item How many continuations exist?
\item Are continuations unique?
\item Are different continuations compatible?
\item Can continuation terminate?
\item Can continuation continue forever?
\item Is there a maximal continuation?
\item Are all maximal continuations equivalent?
\item Does every infinite continuation converge?
\item Which invariants survive continuation?
\end{enumerate}

Observe that these questions precede every classical theory built upon the
completed objects.

%%%%%%%%%%%%%%%%%%%%%%%%%%%%%%%%%%%%%%%%%%%%%%%%%%%%%%%%%%%%%%%%%%%%%%%%%%%%%%
\section{The Fundamental Shift}
%%%%%%%%%%%%%%%%%%%%%%%%%%%%%%%%%%%%%%%%%%%%%%%%%%%%%%%%%%%%%%%%%%%%%%%%%%%%%%

The central philosophical change may now be stated mathematically.

Classical mathematics studies

\[
\boxed{\text{objects}.}
\]

Continuation Mathematics studies

\[
\boxed{\text{objects together with their continuation structure}.}
\]

The object itself becomes only one component of a larger mathematical
entity.

The continuation relation carries independent mathematical information.

Entirely new structures therefore become visible that disappear once
attention is restricted to completed objects alone.

%%%%%%%%%%%%%%%%%%%%%%%%%%%%%%%%%%%%%%%%%%%%%%%%%%%%%%%%%%%%%%%%%%%%%%%%%%%%%%
\section{Consequences}
%%%%%%%%%%%%%%%%%%%%%%%%%%%%%%%%%%%%%%%%%%%%%%%%%%%%%%%%%%%%%%%%%%%%%%%%%%%%%%

The introduction of partial mathematical objects has several immediate
consequences.

\begin{enumerate}
\item Mathematical existence separates from mathematical completion.

\item Completion becomes derivable rather than primitive.

\item Infinite objects become limits of continuation systems rather than
primitive entities.

\item Every branch of mathematics acquires an associated continuation
theory.

\item Questions concerning extension, obstruction, rigidity, maximality,
and compatibility become universal mathematical phenomena rather than
problem-specific constructions.
\end{enumerate}

These observations justify the introduction of an entirely new
mathematical framework.

Partial mathematical objects are not special cases arising inside existing
theories.

They are the primitive inhabitants of Continuation Mathematics.

The next chapter introduces the second primitive notion upon which the
entire theory is built: the continuation relation itself.

%%%%%%%%%%%%%%%%%%%%%%%%%%%%%%%%%%%%%%%%%%%%%%%%%%%%%%%%%%%%%%%%%%%%%%%%%%%%%%
\chapter{Continuation Relations}
%%%%%%%%%%%%%%%%%%%%%%%%%%%%%%%%%%%%%%%%%%%%%%%%%%%%%%%%%%%%%%%%%%%%%%%%%%%%%%

\begin{quote}
\emph{Partial objects possess no mathematics until one specifies how they
may continue. The continuation relation is therefore the first genuine
structure of Continuation Mathematics.}
\end{quote}

\vspace{1em}

The preceding chapter isolated partial mathematical objects independently
of every particular branch of mathematics. Such objects, however, do not
exist in isolation. Their mathematical significance lies entirely in the
ways they may be extended.

The purpose of this chapter is to introduce the primitive notion of
\emph{continuation}. We shall derive the minimal properties that every
continuation relation must satisfy before any further mathematical
structure can be developed.

The resulting theory is completely independent of arithmetic, topology,
algebra, geometry, or logic. It applies uniformly to every mathematical
discipline.

%%%%%%%%%%%%%%%%%%%%%%%%%%%%%%%%%%%%%%%%%%%%%%%%%%%%%%%%%%%%%%%%%%%%%%%%%%%%%%
\section{The Primitive Relation}
%%%%%%%%%%%%%%%%%%%%%%%%%%%%%%%%%%%%%%%%%%%%%%%%%%%%%%%%%%%%%%%%%%%%%%%%%%%%%%

Suppose that \(\mathcal{P}\) denotes a collection of partial mathematical
objects.

Whenever one object may legally be extended into another, there exists a
fundamental relation between them.

\begin{definition}[Continuation relation]
A \emph{continuation relation} on a class of partial mathematical objects
is a binary relation

\[
\leadsto
\;\subseteq\;
\mathcal P\times\mathcal P
\]

whose interpretation is

\[
X\leadsto Y
\]

if and only if \(Y\) is obtained from \(X\) by one admissible
continuation.
\end{definition}

The relation \(\leadsto\) is the primitive structure upon which the entire
theory rests.

Unlike equality, order, or algebraic operations, continuation expresses
possibility rather than state.

%%%%%%%%%%%%%%%%%%%%%%%%%%%%%%%%%%%%%%%%%%%%%%%%%%%%%%%%%%%%%%%%%%%%%%%%%%%%%%
\section{The Principle of Information Growth}
%%%%%%%%%%%%%%%%%%%%%%%%%%%%%%%%%%%%%%%%%%%%%%%%%%%%%%%%%%%%%%%%%%%%%%%%%%%%%%

Every continuation enlarges mathematical information.

Nothing is discarded.

Nothing is forgotten.

Only new structure may be added.

This observation becomes the first structural axiom.

\begin{postulate}[Information Growth]
If

\[
X\leadsto Y,
\]

then every mathematical property already determined by \(X\) remains
determined within \(Y\).
\end{postulate}

The postulate deliberately avoids specifying what ``property'' means.

That notion depends upon the mathematical category under consideration.

The principle merely asserts that continuation is additive rather than
destructive.

%%%%%%%%%%%%%%%%%%%%%%%%%%%%%%%%%%%%%%%%%%%%%%%%%%%%%%%%%%%%%%%%%%%%%%%%%%%%%%
\section{One-Step Continuation}
%%%%%%%%%%%%%%%%%%%%%%%%%%%%%%%%%%%%%%%%%%%%%%%%%%%%%%%%%%%%%%%%%%%%%%%%%%%%%%

Continuation is fundamentally local.

Large extensions are built from elementary ones.

\begin{definition}[Elementary continuation]
A continuation

\[
X\leadsto Y
\]

is called \emph{elementary} when no intermediate continuation exists,

\[
X
\leadsto
Z
\leadsto
Y.
\]

\end{definition}

Elementary continuations play the same role that generators play in group
theory.

Every larger continuation will eventually be decomposed into elementary
ones.

%%%%%%%%%%%%%%%%%%%%%%%%%%%%%%%%%%%%%%%%%%%%%%%%%%%%%%%%%%%%%%%%%%%%%%%%%%%%%%
\section{Finite Continuations}
%%%%%%%%%%%%%%%%%%%%%%%%%%%%%%%%%%%%%%%%%%%%%%%%%%%%%%%%%%%%%%%%%%%%%%%%%%%%%%

Repeated continuation produces finite chains.

\begin{definition}[Continuation chain]
A finite continuation chain is a sequence

\[
X_0
\leadsto
X_1
\leadsto
\cdots
\leadsto
X_n.
\]

Its length is \(n\).
\end{definition}

The object \(X_n\) is said to extend \(X_0\).

Notice that extension itself is no longer primitive.

It is generated by repeated elementary continuation.

%%%%%%%%%%%%%%%%%%%%%%%%%%%%%%%%%%%%%%%%%%%%%%%%%%%%%%%%%%%%%%%%%%%%%%%%%%%%%%
\section{The Derived Extension Relation}
%%%%%%%%%%%%%%%%%%%%%%%%%%%%%%%%%%%%%%%%%%%%%%%%%%%%%%%%%%%%%%%%%%%%%%%%%%%%%%

Finite continuation naturally induces a second relation.

\begin{definition}[Extension]
For partial objects \(X\) and \(Y\),

\[
X\preceq Y
\]

means that there exists a finite continuation chain beginning at \(X\) and
ending at \(Y\).
\end{definition}

The relation \(\preceq\) is the reflexive-transitive closure of
\(\leadsto\).

Unlike continuation,

extension ignores the intermediate stages.

%%%%%%%%%%%%%%%%%%%%%%%%%%%%%%%%%%%%%%%%%%%%%%%%%%%%%%%%%%%%%%%%%%%%%%%%%%%%%%
\section{The Fundamental Properties}
%%%%%%%%%%%%%%%%%%%%%%%%%%%%%%%%%%%%%%%%%%%%%%%%%%%%%%%%%%%%%%%%%%%%%%%%%%%%%%

The extension relation immediately satisfies familiar structural laws.

\begin{theorem}
The extension relation satisfies

\begin{enumerate}
\item Reflexivity:
\[
X\preceq X.
\]

\item Transitivity:
\[
X\preceq Y,\qquad
Y\preceq Z
\Longrightarrow
X\preceq Z.
\]
\end{enumerate}
\end{theorem}

\begin{proof}
Reflexivity follows from the empty continuation chain.

Transitivity follows by concatenating finite continuation chains.
\end{proof}

Thus every continuation system carries a natural preorder.

No further assumptions have yet been imposed.

%%%%%%%%%%%%%%%%%%%%%%%%%%%%%%%%%%%%%%%%%%%%%%%%%%%%%%%%%%%%%%%%%%%%%%%%%%%%%%
\section{Branching}
%%%%%%%%%%%%%%%%%%%%%%%%%%%%%%%%%%%%%%%%%%%%%%%%%%%%%%%%%%%%%%%%%%%%%%%%%%%%%%

Continuation need not be unique.

A partial object may admit many distinct continuations.

\begin{definition}[Branching]
The branching degree of a partial object \(X\) is

\[
\deg(X)
=
\#\{Y:X\leadsto Y\},
\]

possibly infinite.
\end{definition}

Several important cases occur.

\[
\deg(X)=0
\]

means continuation terminates.

\[
\deg(X)=1
\]

means continuation is deterministic.

\[
\deg(X)>1
\]

means genuine mathematical choice exists.

Branching is therefore an intrinsic invariant of continuation systems.

%%%%%%%%%%%%%%%%%%%%%%%%%%%%%%%%%%%%%%%%%%%%%%%%%%%%%%%%%%%%%%%%%%%%%%%%%%%%%%
\section{Dead Ends}
%%%%%%%%%%%%%%%%%%%%%%%%%%%%%%%%%%%%%%%%%%%%%%%%%%%%%%%%%%%%%%%%%%%%%%%%%%%%%%

Termination appears naturally.

\begin{definition}[Terminal object]
A partial object \(X\) is terminal when

\[
\deg(X)=0.
\]
\end{definition}

Terminality is entirely relative to the continuation relation.

An object may be terminal in one continuation system while remaining
extendable in another.

%%%%%%%%%%%%%%%%%%%%%%%%%%%%%%%%%%%%%%%%%%%%%%%%%%%%%%%%%%%%%%%%%%%%%%%%%%%%%%
\section{Infinite Continuations}
%%%%%%%%%%%%%%%%%%%%%%%%%%%%%%%%%%%%%%%%%%%%%%%%%%%%%%%%%%%%%%%%%%%%%%%%%%%%%%

Nothing in the preceding definitions restricts continuation to finite
length.

\begin{definition}[Infinite continuation]
An infinite continuation is an infinite chain

\[
X_0
\leadsto
X_1
\leadsto
X_2
\leadsto
\cdots.
\]
\end{definition}

Such objects will become fundamental later.

Every major unsolved problem involving infinite processes may now be
rephrased as a statement concerning the existence or non-existence of
certain infinite continuation chains.

%%%%%%%%%%%%%%%%%%%%%%%%%%%%%%%%%%%%%%%%%%%%%%%%%%%%%%%%%%%%%%%%%%%%%%%%%%%%%%
\section{Continuation Systems}
%%%%%%%%%%%%%%%%%%%%%%%%%%%%%%%%%%%%%%%%%%%%%%%%%%%%%%%%%%%%%%%%%%%%%%%%%%%%%%

The primitive objects and the primitive relation together determine the
first genuine mathematical structure of the theory.

\begin{definition}[Continuation system]
A \emph{continuation system} is a pair

\[
(\mathcal P,\leadsto),
\]

where

\begin{itemize}
\item \(\mathcal P\) is a class of partial mathematical objects,

\item \(\leadsto\) is a continuation relation on \(\mathcal P\).
\end{itemize}
\end{definition}

Every subsequent construction in this book will arise from continuation
systems.

Nothing further is assumed.

%%%%%%%%%%%%%%%%%%%%%%%%%%%%%%%%%%%%%%%%%%%%%%%%%%%%%%%%%%%%%%%%%%%%%%%%%%%%%%
\section{Examples}
%%%%%%%%%%%%%%%%%%%%%%%%%%%%%%%%%%%%%%%%%%%%%%%%%%%%%%%%%%%%%%%%%%%%%%%%%%%%%%

Many familiar mathematical structures already possess hidden continuation
systems.

\subsection*{Finite words}

Appending one symbol defines continuation.

\vspace{0.5em}

\subsection*{Finite graphs}

Adding one edge defines continuation.

\vspace{0.5em}

\subsection*{Generated groups}

Adjoining one generator defines continuation.

\vspace{0.5em}

\subsection*{Proofs}

Adding one valid inference defines continuation.

\vspace{0.5em}

\subsection*{Collatz}

Applying one accelerated step defines continuation.

Each of these examples belongs to a completely different branch of
mathematics.

Yet the continuation relation has identical formal behaviour in every
case.

%%%%%%%%%%%%%%%%%%%%%%%%%%%%%%%%%%%%%%%%%%%%%%%%%%%%%%%%%%%%%%%%%%%%%%%%%%%%%%
\section{The Fundamental Discovery}
%%%%%%%%%%%%%%%%%%%%%%%%%%%%%%%%%%%%%%%%%%%%%%%%%%%%%%%%%%%%%%%%%%%%%%%%%%%%%%

The preceding definitions reveal the first genuinely new phenomenon.

Classical mathematics ordinarily studies completed mathematical objects.

Continuation Mathematics studies

\[
(\mathcal P,\leadsto),
\]

where the object and its admissible continuations are inseparable.

This shifts mathematical attention away from states and toward
possibilities.

Completion becomes merely one possible outcome of continuation rather than
its starting assumption.

%%%%%%%%%%%%%%%%%%%%%%%%%%%%%%%%%%%%%%%%%%%%%%%%%%%%%%%%%%%%%%%%%%%%%%%%%%%%%%
\section{Toward Continuation Spaces}
%%%%%%%%%%%%%%%%%%%%%%%%%%%%%%%%%%%%%%%%%%%%%%%%%%%%%%%%%%%%%%%%%%%%%%%%%%%%%%

A continuation system contains considerably more information than its
individual objects.

Every object possesses

\begin{itemize}
\item descendants,
\item ancestors,
\item branches,
\item dead ends,
\item infinite extensions,
\item maximal extensions.
\end{itemize}

These are not properties of isolated objects.

They are properties of the global continuation structure.

Accordingly, the next chapter introduces the first global object of the
theory.

Rather than studying isolated continuation chains, we shall study the
entire space generated by all possible continuations simultaneously.

This leads naturally to the notion of a \emph{Continuation Space}, the
first genuinely geometric object of Continuation Mathematics.

%%%%%%%%%%%%%%%%%%%%%%%%%%%%%%%%%%%%%%%%%%%%%%%%%%%%%%%%%%%%%%%%%%%%%%%%%%%%%%
\chapter{Continuation Systems}
%%%%%%%%%%%%%%%%%%%%%%%%%%%%%%%%%%%%%%%%%%%%%%%%%%%%%%%%%%%%%%%%%%%%%%%%%%%%%%

\begin{quote}
\emph{A continuation relation describes individual extensions. A
continuation system describes the entire universe generated by those
extensions.}
\end{quote}

\vspace{1em}

The previous chapter introduced the primitive continuation relation

\[
\leadsto,
\]

which determines when one partial object may be extended into another.

Although the relation is local, it immediately generates a much richer
global structure. Every object possesses ancestors, descendants,
alternative futures, maximal extensions, and possible obstructions.
Collectively these form an entire mathematical universe.

The purpose of this chapter is to define that universe.

%%%%%%%%%%%%%%%%%%%%%%%%%%%%%%%%%%%%%%%%%%%%%%%%%%%%%%%%%%%%%%%%%%%%%%%%%%%%%%
\section{From Relations to Systems}
%%%%%%%%%%%%%%%%%%%%%%%%%%%%%%%%%%%%%%%%%%%%%%%%%%%%%%%%%%%%%%%%%%%%%%%%%%%%%%

A continuation relation never exists in isolation.

Whenever

\[
X\leadsto Y
\]

is defined, there simultaneously exist

\begin{itemize}
\item all objects reachable from \(X\),
\item all objects extending \(Y\),
\item all alternative continuations of \(X\),
\item all common descendants of different objects,
\item all maximal continuation chains.
\end{itemize}

The continuation relation therefore generates an entire mathematical
structure.

This motivates the following definition.

\begin{definition}[Continuation system]
A \emph{continuation system} is a pair

\[
\mathfrak C=(\mathcal P,\leadsto),
\]

where

\begin{itemize}

\item \(\mathcal P\) is a class of partial mathematical objects;

\item \(\leadsto\) is a continuation relation on \(\mathcal P\).

\end{itemize}

The class \(\mathcal P\) is called the \emph{continuation universe} of
\(\mathfrak C\).
\end{definition}

%%%%%%%%%%%%%%%%%%%%%%%%%%%%%%%%%%%%%%%%%%%%%%%%%%%%%%%%%%%%%%%%%%%%%%%%%%%%%%
\section{Reachability}
%%%%%%%%%%%%%%%%%%%%%%%%%%%%%%%%%%%%%%%%%%%%%%%%%%%%%%%%%%%%%%%%%%%%%%%%%%%%%%

Every continuation system possesses a natural notion of reachability.

\begin{definition}[Reachability]
Let \(X,Y\in\mathcal P\).

The object \(Y\) is said to be
\emph{reachable}
from \(X\) if

\[
X\preceq Y.
\]

The set

\[
\operatorname{Reach}(X)
=
\{Y\in\mathcal P:X\preceq Y\}
\]

is called the
\emph{reachability set}
of \(X\).
\end{definition}

Reachability describes every possible mathematical future of a partial
object.

Unlike a single continuation, it is a genuinely global construction.

%%%%%%%%%%%%%%%%%%%%%%%%%%%%%%%%%%%%%%%%%%%%%%%%%%%%%%%%%%%%%%%%%%%%%%%%%%%%%%
\section{Ancestry}
%%%%%%%%%%%%%%%%%%%%%%%%%%%%%%%%%%%%%%%%%%%%%%%%%%%%%%%%%%%%%%%%%%%%%%%%%%%%%%

Reachability admits a dual notion.

\begin{definition}[Ancestry]
The ancestry of a partial object \(X\) is

\[
\operatorname{Anc}(X)
=
\{Y\in\mathcal P:Y\preceq X\}.
\]
\end{definition}

The continuation system therefore carries two canonical directions.

Forward continuation.

Backward ancestry.

Many classical mathematical structures possess only one of these
explicitly.

Continuation Mathematics studies both simultaneously.

%%%%%%%%%%%%%%%%%%%%%%%%%%%%%%%%%%%%%%%%%%%%%%%%%%%%%%%%%%%%%%%%%%%%%%%%%%%%%%
\section{Continuation Components}
%%%%%%%%%%%%%%%%%%%%%%%%%%%%%%%%%%%%%%%%%%%%%%%%%%%%%%%%%%%%%%%%%%%%%%%%%%%%%%

Reachability naturally partitions the continuation universe.

\begin{definition}[Continuation component]
A continuation component is a maximal subset

\[
\mathcal C
\subseteq
\mathcal P
\]

such that every pair of objects is connected by a finite zig-zag of
continuation and reverse continuation.
\end{definition}

The continuation universe decomposes uniquely into continuation
components.

Each component may therefore be studied independently.

%%%%%%%%%%%%%%%%%%%%%%%%%%%%%%%%%%%%%%%%%%%%%%%%%%%%%%%%%%%%%%%%%%%%%%%%%%%%%%
\section{Roots}
%%%%%%%%%%%%%%%%%%%%%%%%%%%%%%%%%%%%%%%%%%%%%%%%%%%%%%%%%%%%%%%%%%%%%%%%%%%%%%

Every continuation process begins somewhere.

\begin{definition}[Root]
A root of a continuation system is an object having no proper ancestor.

Equivalently,

\[
Y\preceq X
\Longrightarrow
Y=X.
\]
\end{definition}

Roots generalize

\begin{itemize}

\item empty words,

\item trivial graphs,

\item identity elements,

\item initial states,

\item axioms.

\end{itemize}

Every continuation system need not possess a root.

When one exists, however, it provides a canonical origin.

%%%%%%%%%%%%%%%%%%%%%%%%%%%%%%%%%%%%%%%%%%%%%%%%%%%%%%%%%%%%%%%%%%%%%%%%%%%%%%
\section{Leaves}
%%%%%%%%%%%%%%%%%%%%%%%%%%%%%%%%%%%%%%%%%%%%%%%%%%%%%%%%%%%%%%%%%%%%%%%%%%%%%%

Dual to roots are terminal objects.

\begin{definition}[Leaf]
A leaf is a partial object admitting no continuation,

\[
\deg(X)=0.
\]
\end{definition}

Leaves need not represent completed mathematics.

They merely represent continuation systems in which no admissible
extension exists.

%%%%%%%%%%%%%%%%%%%%%%%%%%%%%%%%%%%%%%%%%%%%%%%%%%%%%%%%%%%%%%%%%%%%%%%%%%%%%%
\section{Maximal Continuations}
%%%%%%%%%%%%%%%%%%%%%%%%%%%%%%%%%%%%%%%%%%%%%%%%%%%%%%%%%%%%%%%%%%%%%%%%%%%%%%

Finite termination is only one possible endpoint.

A continuation chain may instead become maximal.

\begin{definition}[Maximal continuation]
A continuation chain is maximal if it cannot be extended further.

A maximal continuation may be

\begin{enumerate}

\item finite,

\item infinite.

\end{enumerate}

\end{definition}

Thus maximality and infinitude are distinct notions.

Every finite leaf determines a maximal continuation.

Infinite continuations are automatically maximal whenever no further
continuation exists beyond them.

%%%%%%%%%%%%%%%%%%%%%%%%%%%%%%%%%%%%%%%%%%%%%%%%%%%%%%%%%%%%%%%%%%%%%%%%%%%%%%
\section{Continuation Trees}
%%%%%%%%%%%%%%%%%%%%%%%%%%%%%%%%%%%%%%%%%%%%%%%%%%%%%%%%%%%%%%%%%%%%%%%%%%%%%%

Beginning from any object, all continuations assemble into a tree-like
structure.

\begin{definition}[Continuation tree]
For \(X\in\mathcal P\), the continuation tree of \(X\) is the directed
graph whose

\begin{itemize}

\item vertices are the elements of \(\operatorname{Reach}(X)\),

\item edges are the elementary continuations.

\end{itemize}

\end{definition}

Continuation trees constitute the first geometric objects generated by the
theory.

Their shape reflects the intrinsic branching behaviour of continuation.

%%%%%%%%%%%%%%%%%%%%%%%%%%%%%%%%%%%%%%%%%%%%%%%%%%%%%%%%%%%%%%%%%%%%%%%%%%%%%%
\section{Obstructions}
%%%%%%%%%%%%%%%%%%%%%%%%%%%%%%%%%%%%%%%%%%%%%%%%%%%%%%%%%%%%%%%%%%%%%%%%%%%%%%

Not every continuation succeeds.

Failure itself is mathematical information.

\begin{definition}[Continuation obstruction]
An obstruction is any mathematical condition preventing an admissible
continuation.
\end{definition}

Examples include

\begin{itemize}

\item violated congruences,

\item incompatibility,

\item positivity failure,

\item loss of integrality,

\item violation of defining axioms.

\end{itemize}

Continuation Mathematics studies obstructions directly rather than merely
recording their consequences.

%%%%%%%%%%%%%%%%%%%%%%%%%%%%%%%%%%%%%%%%%%%%%%%%%%%%%%%%%%%%%%%%%%%%%%%%%%%%%%
\section{Deterministic and Branching Systems}
%%%%%%%%%%%%%%%%%%%%%%%%%%%%%%%%%%%%%%%%%%%%%%%%%%%%%%%%%%%%%%%%%%%%%%%%%%%%%%

Continuation systems naturally divide into two broad classes.

\begin{definition}

A continuation system is

\begin{itemize}

\item \emph{deterministic}
if every object has branching degree at most one;

\item \emph{branching}
otherwise.

\end{itemize}

\end{definition}

Many dynamical systems belong to the deterministic class.

Most construction problems belong to the branching class.

Collatz exhibits an interesting intermediate behaviour: the forward map is
deterministic, but the continuation space generated by admissible
valuation sequences branches.

%%%%%%%%%%%%%%%%%%%%%%%%%%%%%%%%%%%%%%%%%%%%%%%%%%%%%%%%%%%%%%%%%%%%%%%%%%%%%%
\section{Universal Questions}
%%%%%%%%%%%%%%%%%%%%%%%%%%%%%%%%%%%%%%%%%%%%%%%%%%%%%%%%%%%%%%%%%%%%%%%%%%%%%%

Every continuation system immediately gives rise to universal mathematical
questions.

\begin{enumerate}

\item Does every continuation terminate?

\item Can infinite continuations exist?

\item Are maximal continuations unique?

\item Are roots unique?

\item What is the branching spectrum?

\item What obstructions govern continuation?

\item Which invariants survive every continuation?

\item What distinguishes finite from infinite continuation?

\end{enumerate}

These questions make no reference to arithmetic, topology, geometry,
analysis, or algebra.

They belong to continuation theory itself.

%%%%%%%%%%%%%%%%%%%%%%%%%%%%%%%%%%%%%%%%%%%%%%%%%%%%%%%%%%%%%%%%%%%%%%%%%%%%%%
\section{The First Structural Principle}
%%%%%%%%%%%%%%%%%%%%%%%%%%%%%%%%%%%%%%%%%%%%%%%%%%%%%%%%%%%%%%%%%%%%%%%%%%%%%%

We may now formulate the first structural principle governing continuation
systems.

\begin{theorem}[Structural Principle]
Every continuation relation canonically generates a continuation system,
and every continuation system canonically determines

\begin{enumerate}

\item its reachability sets,

\item its ancestry sets,

\item its continuation trees,

\item its continuation components,

\item its maximal continuation chains,

\item its obstruction structure.

\end{enumerate}

These constructions depend only upon the continuation relation itself.
\end{theorem}

\begin{proof}
Each construction is obtained directly from the reflexive-transitive
closure of the continuation relation together with elementary graph
theoretic notions. No additional mathematical structure is required.
\end{proof}

%%%%%%%%%%%%%%%%%%%%%%%%%%%%%%%%%%%%%%%%%%%%%%%%%%%%%%%%%%%%%%%%%%%%%%%%%%%%%%
\section{The Hidden Mathematical Object}
%%%%%%%%%%%%%%%%%%%%%%%%%%%%%%%%%%%%%%%%%%%%%%%%%%%%%%%%%%%%%%%%%%%%%%%%%%%%%%

The preceding chapter introduced partial mathematical objects.

This chapter has shown that no partial object should be studied alone.

Its mathematical identity is inseparable from the continuation system in
which it lives.

Consequently the true primitive object of Continuation Mathematics is not

\[
X,
\]

but rather

\[
(\mathfrak C,X),
\]

that is,

an object together with its continuation universe.

This shift parallels earlier developments in mathematics.

Geometry eventually replaced isolated points by spaces.

Algebra replaced isolated elements by algebraic structures.

Topology replaced isolated sets by topological spaces.

Continuation Mathematics replaces isolated partial objects by continuation
systems.

%%%%%%%%%%%%%%%%%%%%%%%%%%%%%%%%%%%%%%%%%%%%%%%%%%%%%%%%%%%%%%%%%%%%%%%%%%%%%%
\section{Toward the Continuation Order}
%%%%%%%%%%%%%%%%%%%%%%%%%%%%%%%%%%%%%%%%%%%%%%%%%%%%%%%%%%%%%%%%%%%%%%%%%%%%%%

Although continuation systems already possess remarkable structure, their
most fundamental feature has not yet appeared.

The extension relation

\[
\preceq
\]

behaves much like an order.

The question naturally arises:

\begin{quote}
Can continuation itself be viewed as a new kind of ordering principle?
\end{quote}

The answer is affirmative.

The next chapter develops the order-theoretic foundations of Continuation
Mathematics by deriving the \emph{Continuation Order}. This order will
become the backbone upon which continuation spaces, continuation algebra,
and eventually universal continuation are constructed.

%%%%%%%%%%%%%%%%%%%%%%%%%%%%%%%%%%%%%%%%%%%%%%%%%%%%%%%%%%%%%%%%%%%%%%%%%%%%%%
\part{Continuation Systems}
%%%%%%%%%%%%%%%%%%%%%%%%%%%%%%%%%%%%%%%%%%%%%%%%%%%%%%%%%%%%%%%%%%%%%%%%%%%%%%

%%%%%%%%%%%%%%%%%%%%%%%%%%%%%%%%%%%%%%%%%%%%%%%%%%%%%%%%%%%%%%%%%%%%%%%%%%%%%%
\chapter{Continuation Systems}
%%%%%%%%%%%%%%%%%%%%%%%%%%%%%%%%%%%%%%%%%%%%%%%%%%%%%%%%%%%%%%%%%%%%%%%%%%%%%%

\begin{quote}
\emph{The Constitution established what continuation is. We now begin the
mathematics that follows from it.}
\end{quote}

\vspace{1em}

Part I introduced the primitive notions upon which Continuation
Mathematics is founded:

\begin{itemize}
\item partial mathematical objects;
\item continuation relations;
\item continuation systems.
\end{itemize}

No substantial mathematics was developed there. The objective was only to
identify the primitive language of the theory.

Beginning with the present chapter, the development changes character.

Continuation systems will now be treated as mathematical objects in their
own right. We shall prove structural results that hold for every
continuation system, independently of the nature of the objects being
continued.

The philosophy is identical to that of modern algebra.

Groups are not studied because integers form a group.

Rather, integers are one realization of an abstract group.

Likewise, Collatz trajectories, proof systems, graph constructions,
algebraic extensions, and dynamical processes are not the theory itself.
They are realizations of a more fundamental continuation system.

%%%%%%%%%%%%%%%%%%%%%%%%%%%%%%%%%%%%%%%%%%%%%%%%%%%%%%%%%%%%%%%%%%%%%%%%%%%%%%
\section{The Mathematical Object}
%%%%%%%%%%%%%%%%%%%%%%%%%%%%%%%%%%%%%%%%%%%%%%%%%%%%%%%%%%%%%%%%%%%%%%%%%%%%%%

A continuation system consists of two inseparable pieces of information:

\[
\mathfrak C=(\mathcal P,\leadsto),
\]

where

\begin{itemize}
\item \(\mathcal P\) is the universe of partial objects;
\item \(\leadsto\) specifies the admissible elementary continuations.
\end{itemize}

Everything else must be derived from these data.

This is the constitutional principle of the theory.

%%%%%%%%%%%%%%%%%%%%%%%%%%%%%%%%%%%%%%%%%%%%%%%%%%%%%%%%%%%%%%%%%%%%%%%%%%%%%%
\section{Internal Versus External Structure}
%%%%%%%%%%%%%%%%%%%%%%%%%%%%%%%%%%%%%%%%%%%%%%%%%%%%%%%%%%%%%%%%%%%%%%%%%%%%%%

Every continuation system possesses two distinct kinds of structure.

\begin{definition}[Internal structure]
The \emph{internal structure} of a continuation system consists of the
partial objects themselves together with whatever mathematical data they
already carry.
\end{definition}

\begin{definition}[Continuation structure]
The \emph{continuation structure} consists entirely of the admissible
continuations among those objects.
\end{definition}

The distinction is fundamental.

Two continuation systems may possess identical internal objects while
having completely different continuation structures.

Conversely, continuation systems arising from entirely different branches
of mathematics may possess isomorphic continuation structures.

Continuation Mathematics studies the latter.

%%%%%%%%%%%%%%%%%%%%%%%%%%%%%%%%%%%%%%%%%%%%%%%%%%%%%%%%%%%%%%%%%%%%%%%%%%%%%%
\section{Structural Equivalence}
%%%%%%%%%%%%%%%%%%%%%%%%%%%%%%%%%%%%%%%%%%%%%%%%%%%%%%%%%%%%%%%%%%%%%%%%%%%%%%

The identity of a continuation system is determined not by the names of
its objects but by the way continuation operates.

\begin{definition}[Continuation isomorphism]
Let

\[
\mathfrak C_1=(\mathcal P_1,\leadsto_1),
\qquad
\mathfrak C_2=(\mathcal P_2,\leadsto_2).
\]

A \emph{continuation isomorphism} is a bijection

\[
\Phi:\mathcal P_1\longrightarrow\mathcal P_2
\]

such that

\[
X\leadsto_1Y
\quad\Longleftrightarrow\quad
\Phi(X)\leadsto_2\Phi(Y)
\]

for every pair of objects.
\end{definition}

Thus two continuation systems are mathematically identical precisely when
their continuation behaviour is identical.

%%%%%%%%%%%%%%%%%%%%%%%%%%%%%%%%%%%%%%%%%%%%%%%%%%%%%%%%%%%%%%%%%%%%%%%%%%%%%%
\section{Intrinsic Invariants}
%%%%%%%%%%%%%%%%%%%%%%%%%%%%%%%%%%%%%%%%%%%%%%%%%%%%%%%%%%%%%%%%%%%%%%%%%%%%%%

The previous definition immediately raises an important question.

Which properties survive continuation isomorphism?

Such quantities are the intrinsic invariants of continuation systems.

Examples already encountered include

\begin{itemize}
\item branching degree;
\item existence of roots;
\item existence of leaves;
\item finite versus infinite continuation;
\item reachability;
\item ancestry;
\item maximal continuation chains.
\end{itemize}

Every future invariant introduced in this book must be preserved under
continuation isomorphism.

%%%%%%%%%%%%%%%%%%%%%%%%%%%%%%%%%%%%%%%%%%%%%%%%%%%%%%%%%%%%%%%%%%%%%%%%%%%%%%
\section{Generated Subsystems}
%%%%%%%%%%%%%%%%%%%%%%%%%%%%%%%%%%%%%%%%%%%%%%%%%%%%%%%%%%%%%%%%%%%%%%%%%%%%%%

Every partial object determines its own local continuation universe.

\begin{definition}[Generated continuation subsystem]
Let \(X\in\mathcal P\).

The continuation subsystem generated by \(X\) is

\[
\mathfrak C(X)
=
\left(
\operatorname{Reach}(X),
\leadsto
\right),
\]

where the continuation relation is restricted to the reachable objects.
\end{definition}

Thus every object generates its own mathematical universe.

Global continuation systems are assembled from these local systems.

%%%%%%%%%%%%%%%%%%%%%%%%%%%%%%%%%%%%%%%%%%%%%%%%%%%%%%%%%%%%%%%%%%%%%%%%%%%%%%
\section{Restriction}
%%%%%%%%%%%%%%%%%%%%%%%%%%%%%%%%%%%%%%%%%%%%%%%%%%%%%%%%%%%%%%%%%%%%%%%%%%%%%%

Continuation systems admit natural restriction.

\begin{definition}[Continuation subsystem]
A continuation subsystem of

\[
(\mathcal P,\leadsto)
\]

is any subset

\[
\mathcal Q\subseteq\mathcal P
\]

that is closed under continuation, together with the induced continuation
relation.
\end{definition}

Closure under continuation means that whenever

\[
X\in\mathcal Q,
\qquad
X\leadsto Y,
\]

one also has

\[
Y\in\mathcal Q.
\]

Subsystems therefore inherit all continuation structure.

%%%%%%%%%%%%%%%%%%%%%%%%%%%%%%%%%%%%%%%%%%%%%%%%%%%%%%%%%%%%%%%%%%%%%%%%%%%%%%
\section{Universality}
%%%%%%%%%%%%%%%%%%%%%%%%%%%%%%%%%%%%%%%%%%%%%%%%%%%%%%%%%%%%%%%%%%%%%%%%%%%%%%

The preceding constructions reveal an important phenomenon.

Nothing introduced so far depends upon

\begin{itemize}
\item arithmetic;
\item topology;
\item geometry;
\item algebra;
\item analysis;
\item logic.
\end{itemize}

Only continuation is used.

Consequently every mathematical discipline possesses an associated
continuation theory obtained simply by identifying

\begin{enumerate}
\item its partial objects;
\item its admissible continuations.
\end{enumerate}

This universality is one of the central discoveries of the present work.

%%%%%%%%%%%%%%%%%%%%%%%%%%%%%%%%%%%%%%%%%%%%%%%%%%%%%%%%%%%%%%%%%%%%%%%%%%%%%%
\section{The Classification Problem}
%%%%%%%%%%%%%%%%%%%%%%%%%%%%%%%%%%%%%%%%%%%%%%%%%%%%%%%%%%%%%%%%%%%%%%%%%%%%%%

Having identified continuation systems as genuine mathematical objects,
the first major problem naturally appears.

\begin{quote}
\emph{When should two continuation systems be regarded as the same?}
\end{quote}

Equivalently,

\begin{quote}
\emph{Which properties belong to the continuation system itself, and which
depend merely upon the particular realization?}
\end{quote}

This is the foundational classification problem of Continuation
Mathematics.

Every later theorem will contribute to its solution.

%%%%%%%%%%%%%%%%%%%%%%%%%%%%%%%%%%%%%%%%%%%%%%%%%%%%%%%%%%%%%%%%%%%%%%%%%%%%%%
\section{Toward the Continuation Order}
%%%%%%%%%%%%%%%%%%%%%%%%%%%%%%%%%%%%%%%%%%%%%%%%%%%%%%%%%%%%%%%%%%%%%%%%%%%%%%

The continuation relation possesses an unexpected feature.

Repeated continuation induces the extension relation

\[
\preceq.
\]

This relation behaves much like an order, but it carries additional
information absent from ordinary order theory.

The next chapter develops this observation into the first major structural
theory of the subject.

We shall show that every continuation system canonically generates a
continuation order, and that many of its deepest properties arise from the
interaction between continuation and order rather than from either
structure individually.

%%%%%%%%%%%%%%%%%%%%%%%%%%%%%%%%%%%%%%%%%%%%%%%%%%%%%%%%%%%%%%%%%%%%%%%%%%%%%%
\chapter{Continuation Morphisms}
%%%%%%%%%%%%%%%%%%%%%%%%%%%%%%%%%%%%%%%%%%%%%%%%%%%%%%%%%%%%%%%%%%%%%%%%%%%%%%

A continuation system is not merely a collection of partial objects together
with continuation relations. It is itself a mathematical object, and therefore
admits structure-preserving maps.

The purpose of this chapter is to identify the correct notion of such maps.

The guiding principle is identical to that of every successful branch of
mathematics.

Groups are studied through homomorphisms.

Topological spaces are studied through continuous maps.

Vector spaces are studied through linear transformations.

Continuation systems must therefore be studied through maps that preserve
continuation.

The resulting notion is the continuation morphism.

%%%%%%%%%%%%%%%%%%%%%%%%%%%%%%%%%%%%%%%%%%%%%%%%%%%%%%%%%%%%%%%%%%%%%%%%%%%%%%

\section{Preservation of continuation}

Let

\[
\mathcal C=(P,\leadsto)
\]

and

\[
\mathcal D=(Q,\rightsquigarrow)
\]

be continuation systems.

A map

\[
f:P\rightarrow Q
\]

should not merely send objects to objects.

It must preserve the continuation structure itself.

%%%%%%%%%%%%%%%%%%%%%%%%%%%%%%%%%%%%%%%%%%%%%%%%%%%%%%%%%%%%%%%%%%%%%%%%%%%%%%

\begin{definition}[Continuation morphism]

A map

\[
f:\mathcal C\rightarrow\mathcal D
\]

is called a continuation morphism if

\[
x\leadsto y
\quad\Longrightarrow\quad
f(x)\rightsquigarrow f(y)
\]

for every continuation relation in
\(\mathcal C\).

\end{definition}

Thus every admissible continuation remains admissible after applying the map.

%%%%%%%%%%%%%%%%%%%%%%%%%%%%%%%%%%%%%%%%%%%%%%%%%%%%%%%%%%%%%%%%%%%%%%%%%%%%%%

The definition deliberately imposes no converse.

A continuation morphism preserves existing continuation.

It need not preserve every forbidden continuation.

Consequently continuation morphisms are naturally order-preserving rather than
bijective.

%%%%%%%%%%%%%%%%%%%%%%%%%%%%%%%%%%%%%%%%%%%%%%%%%%%%%%%%%%%%%%%%%%%%%%%%%%%%%%

\section{Examples}

Every identity map

\[
\operatorname{id}_{\mathcal C}
:
\mathcal C\rightarrow\mathcal C
\]

is a continuation morphism.

Indeed,

\[
x\leadsto y
\]

immediately implies

\[
x\leadsto y.
\]

%%%%%%%%%%%%%%%%%%%%%%%%%%%%%%%%%%%%%%%%%%%%%%%%%%%%%%%%%%%%%%%%%%%%%%%%%%%%%%

Whenever

\[
\mathcal C
\subseteq
\mathcal D
\]

is a continuation subsystem, the inclusion

\[
i:\mathcal C\hookrightarrow\mathcal D
\]

is a continuation morphism.

%%%%%%%%%%%%%%%%%%%%%%%%%%%%%%%%%%%%%%%%%%%%%%%%%%%%%%%%%%%%%%%%%%%%%%%%%%%%%%

Constant maps generally fail to be continuation morphisms unless the image
possesses a self-continuation relation.

Thus continuation behaves differently from ordinary set theory.

%%%%%%%%%%%%%%%%%%%%%%%%%%%%%%%%%%%%%%%%%%%%%%%%%%%%%%%%%%%%%%%%%%%%%%%%%%%%%%

\section{Composition}

Continuation morphisms compose.

\begin{theorem}

Suppose

\[
f:\mathcal C\rightarrow\mathcal D
\]

and

\[
g:\mathcal D\rightarrow\mathcal E
\]

are continuation morphisms.

Then

\[
g\circ f
:
\mathcal C
\rightarrow
\mathcal E
\]

is again a continuation morphism.

\end{theorem}

\begin{proof}

Assume

\[
x\leadsto y.
\]

Since \(f\) preserves continuation,

\[
f(x)\rightsquigarrow f(y).
\]

Applying preservation for \(g\),

\[
g(f(x))
\Longrightarrow
g(f(y)).
\]

Hence

\[
g\circ f
\]

preserves continuation.

\end{proof}

%%%%%%%%%%%%%%%%%%%%%%%%%%%%%%%%%%%%%%%%%%%%%%%%%%%%%%%%%%%%%%%%%%%%%%%%%%%%%%

Identity morphisms together with composition satisfy the usual associative
laws.

Therefore continuation systems naturally form a category.

%%%%%%%%%%%%%%%%%%%%%%%%%%%%%%%%%%%%%%%%%%%%%%%%%%%%%%%%%%%%%%%%%%%%%%%%%%%%%%

\begin{theorem}

Continuation systems together with continuation morphisms form a category.

\end{theorem}

\begin{proof}

Identity morphisms exist.

Composition is associative because composition of functions is associative.

Identity maps act as left and right identities.

Therefore the category axioms hold.

\end{proof}

%%%%%%%%%%%%%%%%%%%%%%%%%%%%%%%%%%%%%%%%%%%%%%%%%%%%%%%%%%%%%%%%%%%%%%%%%%%%%%

This category will be denoted

\[
\mathbf{Cont}.
\]

Unlike the classical categories of sets or topological spaces, the primitive
structure being preserved is not membership, topology, or algebra.

It is continuation.

%%%%%%%%%%%%%%%%%%%%%%%%%%%%%%%%%%%%%%%%%%%%%%%%%%%%%%%%%%%%%%%%%%%%%%%%%%%%%%

\section{Continuation equivalence}

Not every continuation morphism preserves all continuation information.

Those that do should be regarded as isomorphisms.

\begin{definition}

A continuation morphism

\[
f:\mathcal C\rightarrow\mathcal D
\]

is a continuation isomorphism if there exists a continuation morphism

\[
g:\mathcal D\rightarrow\mathcal C
\]

such that

\[
g\circ f
=
\operatorname{id}_{\mathcal C}
\]

and

\[
f\circ g
=
\operatorname{id}_{\mathcal D}.
\]

\end{definition}

When such a map exists we write

\[
\mathcal C
\cong
\mathcal D.
\]

Thus two continuation systems are mathematically identical whenever they
possess exactly the same continuation structure, regardless of how they are
presented.

%%%%%%%%%%%%%%%%%%%%%%%%%%%%%%%%%%%%%%%%%%%%%%%%%%%%%%%%%%%%%%%%%%%%%%%%%%%%%%

\section{The Principle of Continuation Invariance}

The purpose of mathematics is not to study particular presentations.

It is to study intrinsic structure.

Continuation mathematics therefore adopts the following constitutional
principle.

\begin{postulate}[Continuation Invariance]

Every intrinsic property of a continuation system is preserved under
continuation isomorphism.

\end{postulate}

Consequently every theorem developed throughout this work must ultimately
depend only upon continuation structure itself and not upon any accidental
presentation of that structure.

This principle will govern every subsequent construction in the theory.

In particular, all continuation spaces, continuation algebras, continuation
geometries, and universal continuation structures introduced in later parts of
this book will be required to respect continuation morphisms and continuation
equivalence.


%%%%%%%%%%%%%%%%%%%%%%%%%%%%%%%%%%%%%%%%%%%%%%%%%%%%%%%%%%%%%%%%%%%%%%%%%%%%%%
\chapter{Continuation Completion}
%%%%%%%%%%%%%%%%%%%%%%%%%%%%%%%%%%%%%%%%%%%%%%%%%%%%%%%%%%%%%%%%%%%%%%%%%%%%%%

The central question of every continuation system is whether its partial objects
can be completed.

Classical mathematics possesses many different notions of completion.
Metric spaces admit Cauchy completions.
Ordered sets admit Dedekind completions.
Fields admit algebraic closures.
Topological spaces admit compactifications.

Each of these constructions is highly specialized.
They depend upon additional structures introduced after the objects themselves
have already been defined.

Continuation Mathematics begins from a different principle.

Completion is not attached to a particular category.

Completion is an intrinsic property of continuation itself.

The purpose of this chapter is to construct this notion.

%%%%%%%%%%%%%%%%%%%%%%%%%%%%%%%%%%%%%%%%%%%%%%%%%%%%%%%%%%%%%%%%%%%%%%%%%%%%%%

\section{Partiality}

Every continuation system contains objects that are not yet complete.

The notion of incompleteness here is entirely structural.

\begin{definition}[Partial object]
Let $(X,\leadsto)$ be a continuation system.

An element $x\in X$ is called \emph{partial} if there exists at least one
non-trivial continuation beginning at $x$.

Equivalently,

\[
x\leadsto y,
\qquad
y\neq x.
\]

The collection of all partial objects is denoted

\[
P(X).
\]
\end{definition}

Partiality is therefore defined solely by the existence of further admissible
continuations.

Nothing external is required.

%%%%%%%%%%%%%%%%%%%%%%%%%%%%%%%%%%%%%%%%%%%%%%%%%%%%%%%%%%%%%%%%%%%%%%%%%%%%%%

\section{Terminal Objects}

Not every object admits further continuation.

\begin{definition}[Terminal object]
An object $x\in X$ is terminal if

\[
x\leadsto y
\quad\Longrightarrow\quad
y=x.
\]

The collection of terminal objects is denoted

\[
T(X).
\]
\end{definition}

Terminality is therefore the complete absence of further admissible extension.

Notice that

\[
P(X)\cap T(X)=\varnothing,
\]

although either collection may be empty.

%%%%%%%%%%%%%%%%%%%%%%%%%%%%%%%%%%%%%%%%%%%%%%%%%%%%%%%%%%%%%%%%%%%%%%%%%%%%%%

\section{Completion Chains}

A continuation need not terminate immediately.

Instead it may proceed through many intermediate stages.

\begin{definition}[Completion chain]
A completion chain is a continuation

\[
x_0
\leadsto
x_1
\leadsto
x_2
\leadsto
\cdots
\]

whose successive terms are related by admissible continuation.

Its length may be finite or infinite.
\end{definition}

Completion chains are the basic dynamical objects of Continuation Mathematics.

They replace trajectories, derivations,
proof sequences,
algorithmic executions,
orbits,
and iterative constructions.

%%%%%%%%%%%%%%%%%%%%%%%%%%%%%%%%%%%%%%%%%%%%%%%%%%%%%%%%%%%%%%%%%%%%%%%%%%%%%%

\section{Completion}

Completion is defined entirely by the existence of terminal endpoints.

\begin{definition}[Completion]
A completion of $x$ is a terminal object $t$ satisfying

\[
x\leadsto t.
\]

The set of all completions of $x$ is denoted

\[
C(x).
\]
\end{definition}

Unlike classical completion,
we do not assume existence.

We merely define the object.

%%%%%%%%%%%%%%%%%%%%%%%%%%%%%%%%%%%%%%%%%%%%%%%%%%%%%%%%%%%%%%%%%%%%%%%%%%%%%%

\section{Completable Objects}

Some partial objects possess completions.

Others do not.

\begin{definition}[Completable object]
An object $x$ is completable if

\[
C(x)\neq\varnothing.
\]
\end{definition}

Otherwise $x$ is called
\emph{non-completable}.

%%%%%%%%%%%%%%%%%%%%%%%%%%%%%%%%%%%%%%%%%%%%%%%%%%%%%%%%%%%%%%%%%%%%%%%%%%%%%%

\section{The Completion Operator}

Completion itself determines an operator.

\begin{definition}[Completion operator]
The completion operator is

\[
C:X
\longrightarrow
\mathcal P(T(X)),
\]

defined by

\[
C(x)
=
\{t\in T(X):x\leadsto t\}.
\]
\end{definition}

Unlike classical closure operators,
this operator is set-valued.

Multiple completions are permitted.

%%%%%%%%%%%%%%%%%%%%%%%%%%%%%%%%%%%%%%%%%%%%%%%%%%%%%%%%%%%%%%%%%%%%%%%%%%%%%%

\section{Existence versus Uniqueness}

Continuation Mathematics separates two questions.

\begin{enumerate}
\item Does a completion exist?
\item If it exists, is it unique?
\end{enumerate}

Classically these are often merged together.

Here they are fundamentally independent.

%%%%%%%%%%%%%%%%%%%%%%%%%%%%%%%%%%%%%%%%%%%%%%%%%%%%%%%%%%%%%%%%%%%%%%%%%%%%%%

\section{The Completion Spectrum}

Different continuation systems exhibit different completion behaviour.

\begin{definition}[Completion spectrum]
The completion spectrum of a continuation system is the partition

\[
X
=
T(X)
\sqcup
C_1(X)
\sqcup
C_m(X)
\sqcup
N(X),
\]

where

\begin{itemize}
\item $T(X)$ consists of terminal objects,

\item $C_1(X)$ consists of objects possessing a unique completion,

\item $C_m(X)$ consists of objects possessing multiple completions,

\item $N(X)$ consists of non-completable objects.
\end{itemize}
\end{definition}

This decomposition is intrinsic to every continuation system.

%%%%%%%%%%%%%%%%%%%%%%%%%%%%%%%%%%%%%%%%%%%%%%%%%%%%%%%%%%%%%%%%%%%%%%%%%%%%%%

\section{The Completion Problem}

Every continuation system therefore carries a fundamental mathematical problem.

\begin{definition}[Completion Problem]
Determine the completion spectrum of a continuation system.
\end{definition}

This is the universal form of many classical questions.

For example,

\begin{itemize}
\item convergence becomes completion,

\item solvability becomes completion,

\item termination becomes completion,

\item existence becomes completion,

\item classification becomes completion.
\end{itemize}

The diverse completion problems of mathematics are therefore manifestations of a
single continuation-theoretic problem.

%%%%%%%%%%%%%%%%%%%%%%%%%%%%%%%%%%%%%%%%%%%%%%%%%%%%%%%%%%%%%%%%%%%%%%%%%%%%%%

\section{The Completion Principle}

The chapter culminates in the first universal structural principle.

\begin{theorem}[Completion Principle]
Every continuation system determines canonically

\[
(T(X),\,P(X),\,C(X),\,N(X)),
\]

consisting respectively of

\begin{enumerate}
\item terminal objects,

\item partial objects,

\item completion sets,

\item non-completable objects.
\end{enumerate}

These four collections are determined solely by the continuation relation and
require no additional mathematical structure.
\end{theorem}

\begin{proof}
Each collection is defined entirely from the continuation relation
$\leadsto$.

No topology,
metric,
order,
algebra,
or geometry enters the construction.

Hence each is an intrinsic invariant of the continuation system.
\end{proof}

The notion of completion has therefore been recovered directly from continuation
itself.

No appeal has been made to limits, convergence, completeness axioms, or
category-specific constructions.

Completion has become a universal structural object.

The next chapter will show that among all possible completions there frequently
exists a distinguished one, forced by the continuation relation itself. This
leads to the theory of \emph{canonical continuations}.

%%%%%%%%%%%%%%%%%%%%%%%%%%%%%%%%%%%%%%%%%%%%%%%%%%%%%%%%%%%%%%%%%%%%%%%%%%%%%%
\chapter{Canonical Continuations}
%%%%%%%%%%%%%%%%%%%%%%%%%%%%%%%%%%%%%%%%%%%%%%%%%%%%%%%%%%%%%%%%%%%%%%%%%%%%%%

The preceding chapter established that completion is an intrinsic notion of every
continuation system. Existence of completions, however, is not sufficient for a
mathematical theory.

Whenever several completions are possible, one naturally asks whether one of
them is distinguished by the continuation structure itself.

The purpose of this chapter is to isolate such distinguished completions.

Unlike arbitrary completions, canonical continuations are forced entirely by the
continuation relation.

%%%%%%%%%%%%%%%%%%%%%%%%%%%%%%%%%%%%%%%%%%%%%%%%%%%%%%%%%%%%%%%%%%%%%%%%%%%%%%

\section{The Need for Canonicality}

Many branches of mathematics encounter non-uniqueness.

A polynomial may have several roots.

A topological space may possess several compactifications.

A graph may admit many spanning trees.

An equation may possess multiple solutions.

The mere existence of several continuations does not imply that they are equally
natural.

Continuation Mathematics therefore separates

\[
\text{existence}
\qquad\text{from}\qquad
\text{canonicity}.
\]

%%%%%%%%%%%%%%%%%%%%%%%%%%%%%%%%%%%%%%%%%%%%%%%%%%%%%%%%%%%%%%%%%%%%%%%%%%%%%%

\section{Canonical Completion}

\begin{definition}[Canonical completion]
Let $(X,\leadsto)$ be a continuation system.

A completion

\[
t\in C(x)
\]

is called \emph{canonical} if it is uniquely determined by the continuation
structure itself.

The canonical completion of $x$, when it exists, is denoted

\[
\widehat{x}.
\]
\end{definition}

Canonicality is therefore an intrinsic structural property.

It is not imposed externally by convention.

%%%%%%%%%%%%%%%%%%%%%%%%%%%%%%%%%%%%%%%%%%%%%%%%%%%%%%%%%%%%%%%%%%%%%%%%%%%%%%

\section{Canonical Continuation Maps}

Whenever every completable object possesses a canonical completion, one obtains
a distinguished mapping.

\begin{definition}[Canonical continuation]
The canonical continuation operator is the map

\[
\widehat{(\cdot)}:
X_c
\longrightarrow
T(X),
\]

where

\[
X_c
=
\{
x\in X:
\widehat{x}
\text{ exists}
\}.
\]
\end{definition}

The domain consists precisely of those objects admitting canonical completion.

%%%%%%%%%%%%%%%%%%%%%%%%%%%%%%%%%%%%%%%%%%%%%%%%%%%%%%%%%%%%%%%%%%%%%%%%%%%%%%

\section{Structural Characterization}

Canonical completions are characterized entirely by continuation.

\begin{theorem}[Structural Characterization]
If $\widehat{x}$ exists, then

\[
x\leadsto\widehat{x},
\]

and every continuation of $x$ compatible with the continuation structure factors
through $\widehat{x}$.
\end{theorem}

\begin{proof}
The defining property of canonicality is that no competing completion possesses
equal structural status.

Consequently every admissible continuation compatible with the continuation
relation necessarily passes through the distinguished terminal object.

Thus $\widehat{x}$ is universal among completions of $x$.
\end{proof}

%%%%%%%%%%%%%%%%%%%%%%%%%%%%%%%%%%%%%%%%%%%%%%%%%%%%%%%%%%%%%%%%%%%%%%%%%%%%%%

\section{Canonicality as Rigidity}

Canonical continuation represents structural rigidity.

\begin{definition}[Rigid continuation]
A continuation is rigid if every admissible continuation compatible with the
system determines the same canonical completion.
\end{definition}

Rigidity excludes structural ambiguity.

%%%%%%%%%%%%%%%%%%%%%%%%%%%%%%%%%%%%%%%%%%%%%%%%%%%%%%%%%%%%%%%%%%%%%%%%%%%%%%

\section{Stable Objects}

Some objects coincide with their own canonical completion.

\begin{definition}[Stable object]
An object $x$ is stable if

\[
\widehat{x}=x.
\]
\end{definition}

Every terminal object is stable.

The converse need not hold until additional axioms are introduced.

%%%%%%%%%%%%%%%%%%%%%%%%%%%%%%%%%%%%%%%%%%%%%%%%%%%%%%%%%%%%%%%%%%%%%%%%%%%%%%

\section{Canonical Fibres}

Canonical continuation partitions the system.

\begin{definition}[Canonical fibre]
Let

\[
t\in T(X).
\]

Its canonical fibre is

\[
F(t)
=
\{
x:
\widehat{x}=t
\}.
\]
\end{definition}

The fibres partition every canonically completable object.

%%%%%%%%%%%%%%%%%%%%%%%%%%%%%%%%%%%%%%%%%%%%%%%%%%%%%%%%%%%%%%%%%%%%%%%%%%%%%%

\section{Canonical Equivalence}

Canonical completion induces an equivalence relation.

\begin{definition}[Canonical equivalence]
Two objects satisfy

\[
x\sim_c y
\]

whenever

\[
\widehat{x}
=
\widehat{y}.
\]
\end{definition}

Objects are therefore equivalent precisely when they possess the same canonical
completion.

%%%%%%%%%%%%%%%%%%%%%%%%%%%%%%%%%%%%%%%%%%%%%%%%%%%%%%%%%%%%%%%%%%%%%%%%%%%%%%

\section{Canonical Quotients}

The canonical equivalence classes determine a quotient system.

\begin{definition}[Canonical quotient]
The canonical quotient of a continuation system is

\[
X/{\sim_c}.
\]
\end{definition}

The quotient identifies precisely those objects having identical canonical
continuations.

%%%%%%%%%%%%%%%%%%%%%%%%%%%%%%%%%%%%%%%%%%%%%%%%%%%%%%%%%%%%%%%%%%%%%%%%%%%%%%

\section{The Canonical Principle}

The existence of canonical continuation fundamentally simplifies mathematical
structure.

\begin{theorem}[Canonical Principle]
Whenever canonical continuations exist, every continuation system admits a
canonical decomposition into fibres indexed by terminal objects.
\end{theorem}

\begin{proof}
Every canonically completable object belongs to exactly one canonical fibre.

Distinct fibres are disjoint because canonical completion is unique.

Their union is the domain of the canonical continuation operator.

Hence the fibres form a partition indexed by terminal objects.
\end{proof}

%%%%%%%%%%%%%%%%%%%%%%%%%%%%%%%%%%%%%%%%%%%%%%%%%%%%%%%%%%%%%%%%%%%%%%%%%%%%%%

\section{Canonical Mathematics}

Many familiar mathematical constructions may now be viewed through the language
of canonical continuation.

Examples include:

\begin{itemize}
\item canonical representatives of equivalence classes,

\item normal forms,

\item reduced words,

\item Jordan canonical form,

\item reduced fractions,

\item prime decomposition,

\item canonical models,

\item normal proofs,

\item terminating rewrite systems.
\end{itemize}

Although these objects arise in diverse mathematical settings, they share a
common structural feature: they are canonical continuations.

Continuation Mathematics therefore identifies canonicity itself as a fundamental
mathematical phenomenon.

%%%%%%%%%%%%%%%%%%%%%%%%%%%%%%%%%%%%%%%%%%%%%%%%%%%%%%%%%%%%%%%%%%%%%%%%%%%%%%

\section{The Canonical Completion Theorem}

The chapter culminates in the central observation.

\begin{theorem}[Canonical Completion Theorem]
Canonical continuation is determined entirely by the continuation relation.

No metric,
order,
topology,
algebra,
category,
or geometric structure is required to define it.

Whenever canonical continuations exist, they are intrinsic invariants of the
continuation system.
\end{theorem}

\begin{proof}
Every definition introduced in this chapter depends solely upon admissible
continuation.

All additional mathematical structures are absent.

Hence canonical continuation is an invariant of continuation itself.
\end{proof}

Canonical continuation transforms completion from an existential notion into a
structural one.

Instead of merely asking whether an object can be completed, Continuation
Mathematics asks whether the continuation relation itself forces a unique
completion.

This distinction will become increasingly important in the remainder of the
book, culminating in the construction of universal continuation systems.

%%%%%%%%%%%%%%%%%%%%%%%%%%%%%%%%%%%%%%%%%%%%%%%%%%%%%%%%%%%%%%%%%%%%%%%%%%%%%%
\chapter{Universal Continuation Systems}
%%%%%%%%%%%%%%%%%%%%%%%%%%%%%%%%%%%%%%%%%%%%%%%%%%%%%%%%%%%%%%%%%%%%%%%%%%%%%%

The preceding chapters developed the theory of individual continuation systems.
They established continuation relations, continuation morphisms, completion,
and canonical continuation as intrinsic mathematical structures.

The next natural question is whether these individual systems themselves belong
to a larger mathematical universe.

The purpose of this chapter is to construct that universe.

%%%%%%%%%%%%%%%%%%%%%%%%%%%%%%%%%%%%%%%%%%%%%%%%%%%%%%%%%%%%%%%%%%%%%%%%%%%%%%

\section{The Class of Continuation Systems}

Every continuation system consists of a set equipped with an admissible
continuation relation.

It is therefore natural to regard continuation systems themselves as
mathematical objects.

\begin{definition}[Continuation universe]
The \emph{continuation universe}, denoted

\[
\mathbf{Cont},
\]

is the collection of all continuation systems together with their continuation
morphisms.
\end{definition}

Thus the objects of $\mathbf{Cont}$ are continuation systems, while its arrows
are continuation morphisms.

%%%%%%%%%%%%%%%%%%%%%%%%%%%%%%%%%%%%%%%%%%%%%%%%%%%%%%%%%%%%%%%%%%%%%%%%%%%%%%

\section{Identity Systems}

Every continuation system possesses an intrinsic identity morphism.

\begin{proposition}
For every continuation system

\[
(X,\leadsto),
\]

the identity map

\[
\mathrm{id}_X:X\to X
\]

is a continuation morphism.
\end{proposition}

\begin{proof}

If

\[
x\leadsto y,
\]

then

\[
\mathrm{id}_X(x)=x
\leadsto
y=\mathrm{id}_X(y),
\]

so continuation is preserved.
\end{proof}

%%%%%%%%%%%%%%%%%%%%%%%%%%%%%%%%%%%%%%%%%%%%%%%%%%%%%%%%%%%%%%%%%%%%%%%%%%%%%%

\section{Composition}

Continuation morphisms compose naturally.

\begin{theorem}
If

\[
f:X\to Y
\]

and

\[
g:Y\to Z
\]

are continuation morphisms, then

\[
g\circ f:X\to Z
\]

is also a continuation morphism.
\end{theorem}

\begin{proof}

Suppose

\[
x\leadsto x'.
\]

Since $f$ preserves continuation,

\[
f(x)\leadsto f(x').
\]

Applying preservation for $g$ gives

\[
g(f(x))
\leadsto
g(f(x')).
\]

Hence $g\circ f$ preserves continuation.
\end{proof}

%%%%%%%%%%%%%%%%%%%%%%%%%%%%%%%%%%%%%%%%%%%%%%%%%%%%%%%%%%%%%%%%%%%%%%%%%%%%%%

\section{The Category of Continuation}

The preceding results immediately yield the first universal structure.

\begin{theorem}[Category of Continuation Systems]

Continuation systems and continuation morphisms form a category.

\end{theorem}

\begin{proof}

Identity morphisms exist.

Composition exists.

Associativity follows from ordinary composition of maps.

Hence the category axioms hold.
\end{proof}

This category will also be denoted by

\[
\mathbf{Cont}.
\]

%%%%%%%%%%%%%%%%%%%%%%%%%%%%%%%%%%%%%%%%%%%%%%%%%%%%%%%%%%%%%%%%%%%%%%%%%%%%%%

\section{Universal Properties}

The significance of $\mathbf{Cont}$ is not categorical formalism.

Rather, it is structural universality.

Every mathematical theory that admits continuation embeds naturally into
$\mathbf{Cont}$.

Examples include

\begin{itemize}

\item iterative dynamical systems,

\item rewriting systems,

\item proof systems,

\item recursive constructions,

\item search trees,

\item decision procedures,

\item algebraic reduction systems,

\item computational executions.

\end{itemize}

Continuation therefore appears as the common structural language underlying
these seemingly unrelated theories.

%%%%%%%%%%%%%%%%%%%%%%%%%%%%%%%%%%%%%%%%%%%%%%%%%%%%%%%%%%%%%%%%%%%%%%%%%%%%%%

\section{Intrinsic Invariants}

Every continuation morphism preserves certain quantities.

\begin{definition}[Intrinsic continuation invariant]

A property

\[
I(X)
\]

is an intrinsic continuation invariant if

\[
I(X)=I(Y)
\]

whenever

\[
X
\cong
Y
\]

as continuation systems.

\end{definition}

Examples include

\begin{itemize}

\item existence of terminal objects,

\item existence of canonical completion,

\item completion spectrum,

\item continuation depth,

\item continuation height,

\item branching behaviour.

\end{itemize}

The classification of continuation systems therefore reduces to the discovery of
their intrinsic invariants.

%%%%%%%%%%%%%%%%%%%%%%%%%%%%%%%%%%%%%%%%%%%%%%%%%%%%%%%%%%%%%%%%%%%%%%%%%%%%%%

\section{Universal Constructions}

The category $\mathbf{Cont}$ naturally raises universal questions.

For example,

\begin{itemize}

\item products,

\item coproducts,

\item quotients,

\item pullbacks,

\item pushouts,

\item limits,

\item colimits,

\item free continuation systems,

\item universal completions.

\end{itemize}

The present work establishes the foundational language only.

The investigation of these constructions is deferred to later developments.

%%%%%%%%%%%%%%%%%%%%%%%%%%%%%%%%%%%%%%%%%%%%%%%%%%%%%%%%%%%%%%%%%%%%%%%%%%%%%%

\section{The Universal Continuation Principle}

The entire development of Parts I and II culminates in the following principle.

\begin{theorem}[Universal Continuation Principle]

Continuation systems form a universal mathematical framework whose fundamental
objects, morphisms, completion theory, and canonical structures are determined
solely by continuation.

Every mathematical discipline possessing admissible extension may therefore be
studied as a specialization of continuation theory.

\end{theorem}

\begin{proof}

Parts I and II established:

\begin{enumerate}

\item continuation relations;

\item continuation systems;

\item continuation morphisms;

\item completion theory;

\item canonical continuation;

\item composition of continuation morphisms.

Together these determine the category $\mathbf{Cont}$.

Every theory whose objects admit admissible continuation defines an object of
this category.

Hence continuation is universal.

\end{enumerate}

\end{proof}

%%%%%%%%%%%%%%%%%%%%%%%%%%%%%%%%%%%%%%%%%%%%%%%%%%%%%%%%%%%%%%%%%%%%%%%%%%%%%%

\section*{Closing Remarks}

Parts I and II have established the constitutional foundations of Continuation
Mathematics.

Beginning only with the primitive notion of admissible continuation, we have
constructed continuation relations, systems, morphisms, completion, canonical
continuation, and the universal category of continuation systems.

Nothing in this development has depended upon topology, algebra, geometry,
analysis, probability, logic, or computation.

Continuation has been shown to be an autonomous mathematical object.

The remainder of the book investigates the internal mathematics that this new
object necessarily generates.


\part{Continuation Spaces}

%%%%%%%%%%%%%%%%%%%%%%%%%%%%%%%%%%%%%%%%%%%%%%%%%%%%%%%%%%%%%%%%%%%%%%%%%%%%%%
\chapter{The Emergence of Continuation Spaces}
%%%%%%%%%%%%%%%%%%%%%%%%%%%%%%%%%%%%%%%%%%%%%%%%%%%%%%%%%%%%%%%%%%%%%%%%%%%%%%

The preceding parts established continuation as a primitive mathematical
relation and constructed the algebra generated by admissible continuation.

Everything developed thus far has been fundamentally local.

Given an object, one may ask which continuations are admissible.

One may study morphisms between continuation systems.

One may investigate completion and canonical continuation.

Yet an essential question remains unanswered.

What global mathematical object is determined by the totality of all
continuations?

The purpose of this chapter is to answer that question.

Continuation itself generates a mathematical space.

%%%%%%%%%%%%%%%%%%%%%%%%%%%%%%%%%%%%%%%%%%%%%%%%%%%%%%%%%%%%%%%%%%%%%%%%%%%%%%

\section{From Relations to Spaces}

A continuation relation

\[
\leadsto
\]

records only immediate admissibility.

It does not describe the overall organization produced by repeated
continuation.

Just as a graph is more than its individual edges, a continuation system is more
than its elementary continuation relation.

The collection of all continuation paths generates a global mathematical
structure.

This structure is called a continuation space.

%%%%%%%%%%%%%%%%%%%%%%%%%%%%%%%%%%%%%%%%%%%%%%%%%%%%%%%%%%%%%%%%%%%%%%%%%%%%%%

\section{The Fundamental Principle}

The transition from continuation systems to continuation spaces is governed by a
single principle.

\begin{theorem}[Space Generation Principle]

Every continuation system determines a unique continuation space.

\end{theorem}

The theorem is presently a structural assertion.

The construction of the space occupies the remainder of this chapter.

%%%%%%%%%%%%%%%%%%%%%%%%%%%%%%%%%%%%%%%%%%%%%%%%%%%%%%%%%%%%%%%%%%%%%%%%%%%%%%

\section{Reachability}

Continuation naturally generates reachability.

\begin{definition}[Reachability]

Let $(X,\leadsto)$ be a continuation system.

For objects

\[
x,y\in X,
\]

we write

\[
x\rightsquigarrow y
\]

if there exists a finite continuation chain

\[
x=x_0
\leadsto
x_1
\leadsto
\cdots
\leadsto
x_n=y.
\]

The relation

\[
\rightsquigarrow
\]

is called the reachability relation.

\end{definition}

Reachability is the transitive closure of continuation.

%%%%%%%%%%%%%%%%%%%%%%%%%%%%%%%%%%%%%%%%%%%%%%%%%%%%%%%%%%%%%%%%%%%%%%%%%%%%%%

\section{Continuation Components}

Reachability partitions every continuation system.

\begin{definition}[Continuation component]

The continuation component of an object $x$ is

\[
\mathcal C(x)
=
\{
y\in X:
x\rightsquigarrow y
\text{ or }
y\rightsquigarrow x
\}.
\]

\end{definition}

Objects belong to the same continuation component precisely when they are linked
by continuation.

%%%%%%%%%%%%%%%%%%%%%%%%%%%%%%%%%%%%%%%%%%%%%%%%%%%%%%%%%%%%%%%%%%%%%%%%%%%%%%

\section{Continuation Space}

The collection of continuation components forms the first global object.

\begin{definition}[Continuation space]

The continuation space associated to a continuation system is

\[
\mathfrak C(X)
=
\{
\mathcal C(x):
x\in X
\}.
\]

\end{definition}

Unlike the original continuation system, which records individual admissible
steps, the continuation space records global continuation structure.

%%%%%%%%%%%%%%%%%%%%%%%%%%%%%%%%%%%%%%%%%%%%%%%%%%%%%%%%%%%%%%%%%%%%%%%%%%%%%%

\section{Interior Objects}

Some objects cannot be reached from outside their own component.

\begin{definition}[Interior object]

An object is interior if every continuation path passing through it remains
inside its continuation component.

\end{definition}

Interiority is therefore defined entirely by continuation.

%%%%%%%%%%%%%%%%%%%%%%%%%%%%%%%%%%%%%%%%%%%%%%%%%%%%%%%%%%%%%%%%%%%%%%%%%%%%%%

\section{Boundary Objects}

Other objects connect different regions of continuation.

\begin{definition}[Boundary object]

A boundary object is an object through which distinct continuation components
interact.

\end{definition}

Boundary is therefore not geometric.

It is generated by continuation.

%%%%%%%%%%%%%%%%%%%%%%%%%%%%%%%%%%%%%%%%%%%%%%%%%%%%%%%%%%%%%%%%%%%%%%%%%%%%%%

\section{Continuation Separation}

Continuation components are maximal connected regions.

\begin{theorem}

Distinct continuation components are disjoint.

\end{theorem}

\begin{proof}

Suppose two components intersect.

Their common object provides continuation chains connecting every object of one
component with every object of the other.

Hence the two components coincide.

\end{proof}

%%%%%%%%%%%%%%%%%%%%%%%%%%%%%%%%%%%%%%%%%%%%%%%%%%%%%%%%%%%%%%%%%%%%%%%%%%%%%%

\section{Global Structure}

The continuation space therefore decomposes every continuation system into
maximal regions of mutual continuation.

This decomposition is intrinsic.

No metric has been introduced.

No topology has been assumed.

No geometry has been imposed.

The space is generated entirely by continuation.

%%%%%%%%%%%%%%%%%%%%%%%%%%%%%%%%%%%%%%%%%%%%%%%%%%%%%%%%%%%%%%%%%%%%%%%%%%%%%%

\section{The Emergence Theorem}

The chapter culminates in the fundamental observation.

\begin{theorem}[Emergence Theorem]

Every continuation relation generates a canonical continuation space.

This space is determined solely by admissible continuation and is independent of
all additional mathematical structure.

\end{theorem}

\begin{proof}

The continuation relation determines reachability.

Reachability determines continuation components.

The collection of continuation components determines the continuation space.

Each construction depends only upon continuation.

Hence the resulting space is intrinsic.

\end{proof}

%%%%%%%%%%%%%%%%%%%%%%%%%%%%%%%%%%%%%%%%%%%%%%%%%%%%%%%%%%%%%%%%%%%%%%%%%%%%%%

\section*{Closing Remarks}

Continuation Mathematics has now passed from local mathematics to global
mathematics.

The primitive object is no longer merely a continuation relation.

It is now the continuation space generated by that relation.

This transition mirrors one of the deepest recurring themes throughout
mathematics: local laws generate global structure.

The remaining chapters of this part investigate the internal geometry of these
spaces, deriving notions analogous to neighbourhoods, boundaries, dimension,
and continuity directly from continuation itself, without assuming any prior
topological or geometric framework.

%%%%%%%%%%%%%%%%%%%%%%%%%%%%%%%%%%%%%%%%%%%%%%%%%%%%%%%%%%%%%%%%%%%%%%%%%%%%%%
\chapter{Local Continuation Structure}
%%%%%%%%%%%%%%%%%%%%%%%%%%%%%%%%%%%%%%%%%%%%%%%%%%%%%%%%%%%%%%%%%%%%%%%%%%%%%%

The preceding chapter established that every continuation system generates a
continuation space. The existence of this space is a global consequence of the
continuation relation.

The next objective is to understand the local organization of continuation.

Every global mathematical space derives its structure from the interaction of
its local regions.

Continuation Mathematics is no exception.

The purpose of this chapter is to identify the local structures forced by
continuation itself.

%%%%%%%%%%%%%%%%%%%%%%%%%%%%%%%%%%%%%%%%%%%%%%%%%%%%%%%%%%%%%%%%%%%%%%%%%%%%%%

\section{Locality}

Continuation is inherently local.

Every continuation step originates from a single object and extends only to its
immediate admissible continuations.

Global structure is therefore accumulated from local information.

%%%%%%%%%%%%%%%%%%%%%%%%%%%%%%%%%%%%%%%%%%%%%%%%%%%%%%%%%%%%%%%%%%%%%%%%%%%%%%

\section{Continuation Neighborhoods}

The first local object is determined immediately by admissibility.

\begin{definition}[Continuation neighborhood]

Let $(X,\leadsto)$ be a continuation system.

The continuation neighborhood of an object $x$ is

\[
N(x)
=
\{
y\in X:
x\leadsto y
\}.
\]

\end{definition}

Thus every object determines its own local continuation environment.

%%%%%%%%%%%%%%%%%%%%%%%%%%%%%%%%%%%%%%%%%%%%%%%%%%%%%%%%%%%%%%%%%%%%%%%%%%%%%%

\section{Incoming Neighborhoods}

Continuation possesses two directions.

The previous definition records future continuation.

The reverse direction records admissible origins.

\begin{definition}[Incoming neighborhood]

The incoming neighborhood of an object $x$ is

\[
N^{-}(x)
=
\{
y\in X:
y\leadsto x
\}.
\]

\end{definition}

The pair

\[
(N^{-}(x),N(x))
\]

constitutes the complete local continuation data of the object.

%%%%%%%%%%%%%%%%%%%%%%%%%%%%%%%%%%%%%%%%%%%%%%%%%%%%%%%%%%%%%%%%%%%%%%%%%%%%%%

\section{Local Degree}

Neighborhoods immediately determine numerical invariants.

\begin{definition}[Continuation degree]

The continuation degree of an object is

\[
d^{+}(x)
=
|N(x)|.
\]

The incoming continuation degree is

\[
d^{-}(x)
=
|N^{-}(x)|.
\]

\end{definition}

These quantities measure local branching.

%%%%%%%%%%%%%%%%%%%%%%%%%%%%%%%%%%%%%%%%%%%%%%%%%%%%%%%%%%%%%%%%%%%%%%%%%%%%%%

\section{Branching Objects}

Continuation is not necessarily deterministic.

\begin{definition}[Branch point]

An object is a branch point whenever

\[
d^{+}(x)\ge2.
\]

\end{definition}

Such objects admit multiple admissible continuations.

%%%%%%%%%%%%%%%%%%%%%%%%%%%%%%%%%%%%%%%%%%%%%%%%%%%%%%%%%%%%%%%%%%%%%%%%%%%%%%

\section{Merge Points}

The dual phenomenon is equally important.

\begin{definition}[Merge point]

An object is a merge point whenever

\[
d^{-}(x)\ge2.
\]

\end{definition}

Merge points identify distinct continuation histories that converge.

%%%%%%%%%%%%%%%%%%%%%%%%%%%%%%%%%%%%%%%%%%%%%%%%%%%%%%%%%%%%%%%%%%%%%%%%%%%%%%

\section{Regular Objects}

Many continuation systems contain objects exhibiting uniform local behaviour.

\begin{definition}[Regular object]

An object is regular if

\[
d^{+}(x)=1
\qquad\text{and}\qquad
d^{-}(x)=1.
\]

\end{definition}

Regular objects neither branch nor merge.

They constitute the locally simplest regions of continuation.

%%%%%%%%%%%%%%%%%%%%%%%%%%%%%%%%%%%%%%%%%%%%%%%%%%%%%%%%%%%%%%%%%%%%%%%%%%%%%%

\section{Local Profiles}

The neighborhood of an object contains more information than its degree alone.

\begin{definition}[Local continuation profile]

The local continuation profile of an object is

\[
P(x)
=
\left(
N^{-}(x),
N(x)
\right).
\]

\end{definition}

Two objects possess identical local structure precisely when their continuation
profiles are isomorphic.

%%%%%%%%%%%%%%%%%%%%%%%%%%%%%%%%%%%%%%%%%%%%%%%%%%%%%%%%%%%%%%%%%%%%%%%%%%%%%%

\section{Local Equivalence}

Continuation therefore induces an intrinsic equivalence relation.

\begin{definition}[Local equivalence]

Objects

\[
x,y\in X
\]

are locally equivalent if their continuation profiles are isomorphic.

\end{definition}

This relation depends only upon continuation.

%%%%%%%%%%%%%%%%%%%%%%%%%%%%%%%%%%%%%%%%%%%%%%%%%%%%%%%%%%%%%%%%%%%%%%%%%%%%%%

\section{Homogeneous Regions}

Local equivalence naturally partitions continuation spaces.

\begin{definition}[Homogeneous continuation region]

A homogeneous continuation region is a maximal subset of a continuation space in
which every object is locally equivalent.

\end{definition}

These regions represent areas of uniform continuation behaviour.

%%%%%%%%%%%%%%%%%%%%%%%%%%%%%%%%%%%%%%%%%%%%%%%%%%%%%%%%%%%%%%%%%%%%%%%%%%%%%%

\section{The Local Structure Principle}

The local organization of continuation is entirely intrinsic.

\begin{theorem}[Local Structure Principle]

Every continuation system canonically determines

\begin{enumerate}

\item continuation neighborhoods,

\item incoming neighborhoods,

\item continuation degrees,

\item branch points,

\item merge points,

\item regular objects,

\item homogeneous continuation regions.

\end{enumerate}

Each of these structures depends solely upon the continuation relation.

\end{theorem}

\begin{proof}

Each object is obtained directly from the continuation relation
\(
\leadsto
\).

No auxiliary mathematical structure enters any construction.

Hence every local continuation structure is intrinsic.

\end{proof}

%%%%%%%%%%%%%%%%%%%%%%%%%%%%%%%%%%%%%%%%%%%%%%%%%%%%%%%%%%%%%%%%%%%%%%%%%%%%%%

\section*{Closing Remarks}

The continuation space now possesses both global and local organization.

Globally it decomposes into continuation components.

Locally every object possesses an intrinsic continuation neighborhood together
with numerical and structural invariants describing its immediate continuation
behaviour.

These local structures provide the first indications that continuation spaces
possess an internal geometry.

The next chapter investigates the interfaces between distinct regions of
continuation, leading to the theory of continuation frontiers.

%%%%%%%%%%%%%%%%%%%%%%%%%%%%%%%%%%%%%%%%%%%%%%%%%%%%%%%%%%%%%%%%%%%%%%%%%%%%%%
\chapter{Continuation Frontiers}
%%%%%%%%%%%%%%%%%%%%%%%%%%%%%%%%%%%%%%%%%%%%%%%%%%%%%%%%%%%%%%%%%%%%%%%%%%%%%%

\section{The emergence of frontiers}

Every continuation space possesses an intrinsic distinction between those
objects that admit further continuation and those whose continuation has
become impossible.

This distinction is not imposed externally.

It is generated by the continuation relation itself.

The resulting structure is called the \emph{continuation frontier}.

Unlike the boundary of a topological space, which depends upon an ambient
topology, a continuation frontier depends only upon the continuation
structure.

It therefore exists prior to topology, geometry, or metric notions.

Its role is to separate regions of mathematical evolution from regions where
the continuation process terminates.

Throughout this chapter let
\[
(X,\rightsquigarrow)
\]
be a continuation space.

%%%%%%%%%%%%%%%%%%%%%%%%%%%%%%%%%%%%%%%%%%%%%%%%%%%%%%%%%%%%%%%%%%%%%%%%%%%%%%

\section{Continuation-maximal objects}

Some objects cannot be continued any further.

These form the first fundamental class.

\begin{definition}[Continuation-maximal object]
An object
\[
x\in X
\]
is called \emph{continuation-maximal} if there exists no object
\[
y\neq x
\]
such that
\[
x\rightsquigarrow y.
\]
\end{definition}

Thus a continuation-maximal object possesses no genuine extension.

Its continuation process has terminated.

%%%%%%%%%%%%%%%%%%%%%%%%%%%%%%%%%%%%%%%%%%%%%%%%%%%%%%%%%%%%%%%%%%%%%%%%%%%%%%

\section{Continuation-minimal objects}

Dually one obtains the notion of initial objects.

\begin{definition}[Continuation-minimal object]
An object
\[
x\in X
\]
is called \emph{continuation-minimal}
if no distinct object satisfies
\[
y\rightsquigarrow x.
\]
\end{definition}

Minimal objects represent the beginnings of continuation chains.

Maximal objects represent their termination.

%%%%%%%%%%%%%%%%%%%%%%%%%%%%%%%%%%%%%%%%%%%%%%%%%%%%%%%%%%%%%%%%%%%%%%%%%%%%%%

\section{Frontier objects}

The most interesting objects are those lying between the interior of the
continuation process and complete termination.

\begin{definition}[Frontier object]
An object
\[
x\in X
\]
is called a \emph{frontier object}
if

\begin{enumerate}
\item
there exists some
\[
y
\]
with
\[
y\rightsquigarrow x,
\]

\item
every continuation beginning at
\[
x
\]
is finite.
\end{enumerate}

\end{definition}

Thus frontier objects remain inside the continuation process while lying at
its edge.

No infinite continuation can pass beyond them.

%%%%%%%%%%%%%%%%%%%%%%%%%%%%%%%%%%%%%%%%%%%%%%%%%%%%%%%%%%%%%%%%%%%%%%%%%%%%%%

\section{The continuation frontier}

These objects assemble into a canonical subset.

\begin{definition}[Continuation frontier]
The \emph{continuation frontier} of a continuation space is

\[
\partial_C X
=
\{
x\in X:
x
\text{ is a frontier object}
\}.
\]

\end{definition}

This subset is determined entirely by the continuation relation.

No additional structure is required.

%%%%%%%%%%%%%%%%%%%%%%%%%%%%%%%%%%%%%%%%%%%%%%%%%%%%%%%%%%%%%%%%%%%%%%%%%%%%%%

\section{Interior of a continuation space}

The complement of the frontier consists of those objects capable of unlimited
further continuation.

\begin{definition}[Continuation interior]
The continuation interior is

\[
\operatorname{Int}_C(X)
=
X\setminus\partial_C X.
\]

\end{definition}

Objects in the continuation interior admit arbitrarily long continuation
chains.

They constitute the dynamically active region of the continuation space.

%%%%%%%%%%%%%%%%%%%%%%%%%%%%%%%%%%%%%%%%%%%%%%%%%%%%%%%%%%%%%%%%%%%%%%%%%%%%%%

\section{Frontiers are structural}

The continuation frontier is invariant under continuation isomorphism.

\begin{theorem}[Frontier invariance]
Let

\[
F:X\rightarrow Y
\]

be a continuation isomorphism.

Then

\[
F(\partial_C X)
=
\partial_C Y.
\]

\end{theorem}

\begin{proof}

Continuation isomorphisms preserve continuation chains in both directions.

An object possesses arbitrarily long continuations precisely when its image
does.

Likewise finite termination is preserved.

Hence frontier objects correspond bijectively.

\end{proof}

%%%%%%%%%%%%%%%%%%%%%%%%%%%%%%%%%%%%%%%%%%%%%%%%%%%%%%%%%%%%%%%%%%%%%%%%%%%%%%

\section{Nested frontiers}

Frontiers may themselves possess continuation structure.

Removing one frontier may reveal another.

\begin{definition}[Higher continuation frontier]

Define recursively

\[
\partial_C^{\,0}(X)=X
\]

and

\[
\partial_C^{\,n+1}(X)
=
\partial_C
\left(
\partial_C^{\,n}(X)
\right).
\]

\end{definition}

Thus every continuation space generates a hierarchy

\[
X,
\partial_C X,
\partial_C^2X,
\partial_C^3X,
\dots
\]

called its
\emph{frontier filtration}.

%%%%%%%%%%%%%%%%%%%%%%%%%%%%%%%%%%%%%%%%%%%%%%%%%%%%%%%%%%%%%%%%%%%%%%%%%%%%%%

\section{Frontier depth}

This filtration measures structural complexity.

\begin{definition}[Frontier depth]

The frontier depth of an object
\[
x
\]
is the smallest integer
\[
d(x)\ge0
\]
such that

\[
x
\in
\partial_C^{\,d(x)}X
\]

but

\[
x
\notin
\partial_C^{\,d(x)+1}X.
\]

If no such integer exists,
the frontier depth is defined to be infinite.

\end{definition}

Objects of infinite frontier depth belong to infinitely nested continuation
processes.

Finite depth measures proximity to structural termination.

%%%%%%%%%%%%%%%%%%%%%%%%%%%%%%%%%%%%%%%%%%%%%%%%%%%%%%%%%%%%%%%%%%%%%%%%%%%%%%

\section{Termination versus incompleteness}

The continuation frontier separates two fundamentally different phenomena.

The first is genuine termination.

The second is merely incomplete knowledge.

A mathematical object may lie beyond present understanding without lying on
the continuation frontier.

Likewise an object may be completely understood while still possessing
unbounded continuation.

Continuation Mathematics therefore distinguishes between

\begin{enumerate}

\item epistemic incompleteness,

\item structural termination.

\end{enumerate}

These notions coincide only accidentally.

%%%%%%%%%%%%%%%%%%%%%%%%%%%%%%%%%%%%%%%%%%%%%%%%%%%%%%%%%%%%%%%%%%%%%%%%%%%%%%

\section{The universal role of frontiers}

Every continuation process eventually encounters one of three possibilities.

Either

\begin{enumerate}

\item continuation terminates,

\item continuation becomes periodic,

\item continuation remains unbounded.

\end{enumerate}

The continuation frontier is precisely the structure that separates these
regimes.

Later chapters will show that many classical mathematical problems may be
reformulated as determining whether particular continuation frontiers are
empty, finite, infinite, connected, or canonically complete.

Thus continuation frontiers constitute the first genuinely geometric objects
generated by continuation alone.

They represent the visible edge of mathematical evolution.

No topology has yet been introduced.

Nevertheless a notion resembling boundary has already emerged purely from
continuation itself.

This phenomenon will become the foundation upon which continuation geometry is
constructed.

%%%%%%%%%%%%%%%%%%%%%%%%%%%%%%%%%%%%%%%%%%%%%%%%%%%%%%%%%%%%%%%%%%%%%%%%%%%%%%
\chapter{Continuation Closure}
%%%%%%%%%%%%%%%%%%%%%%%%%%%%%%%%%%%%%%%%%%%%%%%%%%%%%%%%%%%%%%%%%%%%%%%%%%%%%%

\section{The necessity of closure}

Continuation systems are generated by repeated application of continuation.
The first natural question is whether every continuation process remains
inside the same mathematical universe.

In classical mathematics, closure is usually imposed as an axiom.
Groups are required to be closed under multiplication.
Vector spaces are required to be closed under addition and scalar
multiplication.

Continuation Mathematics reverses this viewpoint.

Closure is not assumed.

Closure is itself an object that must be constructed.

%%%%%%%%%%%%%%%%%%%%%%%%%%%%%%%%%%%%%%%%%%%%%%%%%%%%%%%%%%%%%%%%%%%%%%%%%%%%%%

\section{Continuation closure}

\begin{definition}[Continuation closure]
Let
\[
(X,\rightsquigarrow)
\]
be a continuation space and let
\[
A\subseteq X.
\]

The \emph{continuation closure} of \(A\), denoted
\[
\operatorname{Cl}_C(A),
\]
is the smallest continuation subsystem of \(X\) containing \(A\).

Equivalently,

\[
\operatorname{Cl}_C(A)
=
\bigcap
\{
Y\subseteq X:
A\subseteq Y
\text{ and }
Y
\text{ is continuation closed}
\}.
\]

\end{definition}

Thus continuation closure is the minimal universe in which every continuation
of every object of \(A\) already exists.

%%%%%%%%%%%%%%%%%%%%%%%%%%%%%%%%%%%%%%%%%%%%%%%%%%%%%%%%%%%%%%%%%%%%%%%%%%%%%%

\section{Continuation-closed subsets}

\begin{definition}

A subset
\[
Y\subseteq X
\]
is called \emph{continuation closed}
whenever

\[
x\in Y,
\qquad
x\rightsquigarrow y
\quad\Longrightarrow\quad
y\in Y.
\]

\end{definition}

Continuation closure therefore means closure under future evolution.

%%%%%%%%%%%%%%%%%%%%%%%%%%%%%%%%%%%%%%%%%%%%%%%%%%%%%%%%%%%%%%%%%%%%%%%%%%%%%%

\section{Existence and uniqueness}

Continuation closure is always well-defined.

\begin{theorem}[Closure Theorem]

Every subset
\[
A\subseteq X
\]
possesses a unique continuation closure.

\end{theorem}

\begin{proof}

The entire space \(X\) is continuation closed.

Hence the family of continuation-closed subsets containing \(A\) is
non-empty.

The intersection of any family of continuation-closed subsets remains
continuation closed.

Therefore the displayed intersection exists and is itself continuation
closed.

Minimality follows immediately from construction.

\end{proof}

%%%%%%%%%%%%%%%%%%%%%%%%%%%%%%%%%%%%%%%%%%%%%%%%%%%%%%%%%%%%%%%%%%%%%%%%%%%%%%

\section{Generated continuation systems}

Closure generates new mathematical objects.

\begin{definition}

The continuation subsystem generated by \(A\) is

\[
\langle A\rangle_C
=
\operatorname{Cl}_C(A).
\]

\end{definition}

Thus every subset generates a canonical continuation system.

Generation is therefore no longer primitive.

It is recovered from continuation closure.

%%%%%%%%%%%%%%%%%%%%%%%%%%%%%%%%%%%%%%%%%%%%%%%%%%%%%%%%%%%%%%%%%%%%%%%%%%%%%%

\section{Finite and infinite closure}

Continuation closure may terminate after finitely many steps or continue
indefinitely.

\begin{definition}

The continuation closure of \(A\) is called

\begin{enumerate}

\item
\emph{finite}
if only finitely many continuation steps are required to generate it;

\item
\emph{infinite}
otherwise.

\end{enumerate}

\end{definition}

This distinction depends entirely upon the continuation relation.

%%%%%%%%%%%%%%%%%%%%%%%%%%%%%%%%%%%%%%%%%%%%%%%%%%%%%%%%%%%%%%%%%%%%%%%%%%%%%%

\section{Closure rank}

The complexity of closure may itself be measured.

\begin{definition}[Closure rank]

For
\[
A\subseteq X
\]
define recursively

\[
C_0(A)=A,
\]

\[
C_{n+1}(A)
=
C_n(A)
\cup
\{
y:
x\rightsquigarrow y
\text{ for some }
x\in C_n(A)
\}.
\]

If there exists an integer
\[
r
\]
such that

\[
C_r(A)=C_{r+1}(A),
\]

then the smallest such integer is called the
\emph{closure rank}
of \(A\).

If no such integer exists,
the closure rank is infinite.

\end{definition}

Closure rank measures the number of continuation generations required before
no genuinely new objects appear.

%%%%%%%%%%%%%%%%%%%%%%%%%%%%%%%%%%%%%%%%%%%%%%%%%%%%%%%%%%%%%%%%%%%%%%%%%%%%%%

\section{Closure operators}

Continuation closure satisfies the classical closure axioms, but derives them
rather than assuming them.

\begin{theorem}[Closure Operator]

For every
\[
A,B\subseteq X,
\]

the continuation closure satisfies

\begin{enumerate}

\item
\[
A
\subseteq
\operatorname{Cl}_C(A),
\]

\item
\[
A\subseteq B
\Longrightarrow
\operatorname{Cl}_C(A)
\subseteq
\operatorname{Cl}_C(B),
\]

\item
\[
\operatorname{Cl}_C
(\operatorname{Cl}_C(A))
=
\operatorname{Cl}_C(A).
\]

\end{enumerate}

\end{theorem}

\begin{proof}

The first follows immediately from construction.

The second follows because every continuation-closed subset containing
\(B\) also contains \(A\).

The third follows because the continuation closure is already continuation
closed.

\end{proof}

%%%%%%%%%%%%%%%%%%%%%%%%%%%%%%%%%%%%%%%%%%%%%%%%%%%%%%%%%%%%%%%%%%%%%%%%%%%%%%

\section{Closure versus completion}

Closure and completion are fundamentally different notions.

Closure asks whether every continuation remains inside the same universe.

Completion asks whether missing continuation objects have been adjoined.

Consequently,

\[
\text{completion}
\quad\Longrightarrow\quad
\text{closure},
\]

but generally

\[
\text{closure}
\not\Longrightarrow
\text{completion}.
\]

Closure concerns stability.

Completion concerns existence.

%%%%%%%%%%%%%%%%%%%%%%%%%%%%%%%%%%%%%%%%%%%%%%%%%%%%%%%%%%%%%%%%%%%%%%%%%%%%%%

\section{Universal property of closure}

Continuation closure satisfies a universal characterization.

\begin{theorem}[Universal Property]

Let
\[
A\subseteq X.
\]

If
\[
Y
\]
is any continuation-closed subsystem satisfying

\[
A\subseteq Y,
\]

then

\[
\operatorname{Cl}_C(A)
\subseteq
Y.
\]

\end{theorem}

\begin{proof}

This is immediate from the definition of continuation closure as the smallest
continuation-closed subsystem containing \(A\).

\end{proof}

%%%%%%%%%%%%%%%%%%%%%%%%%%%%%%%%%%%%%%%%%%%%%%%%%%%%%%%%%%%%%%%%%%%%%%%%%%%%%%

\section{Closure as mathematical inevitability}

Continuation closure is not merely a technical construction.

It expresses one of the central principles of Continuation Mathematics.

Whenever a mathematical object exists, every continuation forced by its own
structure must either

\begin{enumerate}

\item
already exist,

\item
be adjoined,

\item
or expose the incompleteness of the surrounding mathematical universe.

\end{enumerate}

Closure therefore measures the internal sufficiency of a mathematical system.

In the chapters that follow, closure will interact with continuation
dimension, continuation topology, and continuation compactness to produce the
first genuinely geometric invariants of continuation spaces.

Closure is thus the bridge between continuation dynamics and continuation
geometry.

%%%%%%%%%%%%%%%%%%%%%%%%%%%%%%%%%%%%%%%%%%%%%%%%%%%%%%%%%%%%%%%%%%%%%%%%%%%%%%
\chapter{Continuation Interior}
%%%%%%%%%%%%%%%%%%%%%%%%%%%%%%%%%%%%%%%%%%%%%%%%%%%%%%%%%%%%%%%%%%%%%%%%%%%%%%

\section{Beyond closure}

Closure identifies those objects that remain inside a continuation system.
The next question is fundamentally different.

Among the objects of a continuation space, which belong to the genuine
``heart'' of the continuation process, and which merely lie on its edge?

This distinction gives rise to the notion of continuation interior.

Unlike topological interior, which depends upon open sets, continuation
interior is generated directly by the continuation relation itself.

It therefore precedes topology and is independent of any metric or geometric
structure.

Throughout this chapter let
\[
(X,\rightsquigarrow)
\]
be a continuation space.

%%%%%%%%%%%%%%%%%%%%%%%%%%%%%%%%%%%%%%%%%%%%%%%%%%%%%%%%%%%%%%%%%%%%%%%%%%%%%%

\section{Interior points}

An object belongs to the continuation interior when its future is
structurally rich and locally stable under continuation.

\begin{definition}[Interior object]

An object
\[
x\in X
\]
is called an \emph{interior object} if

\begin{enumerate}

\item
there exists
\[
y\neq x
\]
with
\[
x\rightsquigarrow y,
\]

\item
every continuation of
\[
x
\]
can itself be further continued.

\end{enumerate}

\end{definition}

Thus an interior object never terminates immediately, nor does it continue
into a dead end after finitely many forced steps.

%%%%%%%%%%%%%%%%%%%%%%%%%%%%%%%%%%%%%%%%%%%%%%%%%%%%%%%%%%%%%%%%%%%%%%%%%%%%%%

\section{The continuation interior}

\begin{definition}[Continuation interior]

The continuation interior of \(X\) is

\[
\operatorname{Int}_C(X)
=
\{
x\in X:
x
\text{ is an interior object}
\}.
\]

\end{definition}

The continuation interior consists precisely of those objects that remain
inside the active continuation process.

%%%%%%%%%%%%%%%%%%%%%%%%%%%%%%%%%%%%%%%%%%%%%%%%%%%%%%%%%%%%%%%%%%%%%%%%%%%%%%

\section{Interior and frontier}

Continuation interior and continuation frontier are complementary notions.

\begin{theorem}[Interior--Frontier Decomposition]

Every continuation space satisfies

\[
X
=
\operatorname{Int}_C(X)
\cup
\partial_C X.
\]

Moreover,

\[
\operatorname{Int}_C(X)
\cap
\partial_C X
=
\varnothing.
\]

\end{theorem}

\begin{proof}

By definition every object either possesses indefinitely extendable
continuation or eventually reaches termination.

These alternatives are mutually exclusive.

Hence every object belongs to exactly one of the two classes.

\end{proof}

%%%%%%%%%%%%%%%%%%%%%%%%%%%%%%%%%%%%%%%%%%%%%%%%%%%%%%%%%%%%%%%%%%%%%%%%%%%%%%

\section{Interior subsystems}

The continuation interior naturally carries its own continuation structure.

\begin{proposition}

The restriction of
\[
\rightsquigarrow
\]
to

\[
\operatorname{Int}_C(X)
\]

defines a continuation subsystem.

\end{proposition}

\begin{proof}

If
\[
x\in\operatorname{Int}_C(X)
\]
and
\[
x\rightsquigarrow y,
\]
then every continuation beginning at
\[
y
\]
remains indefinitely extendable by the defining property of interior objects.

Thus

\[
y\in\operatorname{Int}_C(X).
\]

\end{proof}

%%%%%%%%%%%%%%%%%%%%%%%%%%%%%%%%%%%%%%%%%%%%%%%%%%%%%%%%%%%%%%%%%%%%%%%%%%%%%%

\section{Interior chains}

The continuation interior is characterized by infinite continuation chains.

\begin{definition}

An \emph{interior chain} is an infinite continuation sequence

\[
x_0
\rightsquigarrow
x_1
\rightsquigarrow
x_2
\rightsquigarrow
\cdots
\]

contained entirely inside

\[
\operatorname{Int}_C(X).
\]

\end{definition}

Interior chains represent unbounded mathematical evolution.

%%%%%%%%%%%%%%%%%%%%%%%%%%%%%%%%%%%%%%%%%%%%%%%%%%%%%%%%%%%%%%%%%%%%%%%%%%%%%%

\section{Interior rank}

Not every interior object possesses the same continuation richness.

\begin{definition}[Interior rank]

Let

\[
x\in\operatorname{Int}_C(X).
\]

The \emph{interior rank} of
\[
x
\]
is the supremum of the lengths of pairwise distinct continuation chains
originating at
\[
x.
\]

If arbitrarily many independent continuation chains originate at
\[
x,
\]
its interior rank is infinite.

\end{definition}

Interior rank measures the local continuation complexity of an object.

%%%%%%%%%%%%%%%%%%%%%%%%%%%%%%%%%%%%%%%%%%%%%%%%%%%%%%%%%%%%%%%%%%%%%%%%%%%%%%

\section{Persistence}

Continuation interior expresses persistence under mathematical evolution.

\begin{definition}[Persistent object]

An object is called \emph{persistent} if every finite continuation chain
beginning at that object may be extended further.

\end{definition}

Persistent objects necessarily belong to the continuation interior.

%%%%%%%%%%%%%%%%%%%%%%%%%%%%%%%%%%%%%%%%%%%%%%%%%%%%%%%%%%%%%%%%%%%%%%%%%%%%%%

\section{Interior morphisms}

Interior is preserved by continuation isomorphisms.

\begin{theorem}[Interior Invariance]

Let

\[
F:X\rightarrow Y
\]

be a continuation isomorphism.

Then

\[
F
\left(
\operatorname{Int}_C(X)
\right)
=
\operatorname{Int}_C(Y).
\]

\end{theorem}

\begin{proof}

Continuation isomorphisms preserve continuation chains in both directions.

An object admits arbitrarily extensible continuation precisely when its image
does.

Therefore interior objects correspond bijectively.

\end{proof}

%%%%%%%%%%%%%%%%%%%%%%%%%%%%%%%%%%%%%%%%%%%%%%%%%%%%%%%%%%%%%%%%%%%%%%%%%%%%%%

\section{Interior generation}

Continuation interior may be generated recursively.

Define

\[
I_0=X,
\]

and recursively

\[
I_{n+1}
=
\{
x\in I_n:
\exists\,y\in I_n
\text{ with }
x\rightsquigarrow y
\}.
\]

Each stage removes objects whose continuation has already terminated.

The limiting subsystem

\[
I_\infty
=
\bigcap_{n=0}^{\infty}
I_n
\]

is called the \emph{stable continuation interior}.

It consists precisely of those objects surviving every finite elimination of
terminal behaviour.

%%%%%%%%%%%%%%%%%%%%%%%%%%%%%%%%%%%%%%%%%%%%%%%%%%%%%%%%%%%%%%%%%%%%%%%%%%%%%%

\section{Interior as structural permanence}

Continuation interior captures one of the fundamental ideas of the theory.

A mathematical object is not intrinsically significant merely because it
exists.

Its significance depends upon the continuation that remains available beyond
it.

Objects lying in the continuation frontier represent structural exhaustion.

Objects lying in the continuation interior represent structural permanence.

This distinction is independent of topology, metric geometry, or analysis.

It is produced solely by continuation itself.

Later chapters will show that continuation compactness, connectedness,
dimension, and eventually continuation geometry are all governed by the
interaction between closure, frontier, and interior.

The continuation interior therefore completes the first intrinsic
decomposition of every continuation space into its stable core and its
terminating boundary.


%%%%%%%%%%%%%%%%%%%%%%%%%%%%%%%%%%%%%%%%%%%%%%%%%%%%%%%%%%%%%%%%%%%%%%%%%%%%%%
%%%%%%%%%%%%%%%%%%%%%%%%%%%%%%%%%%%%%%%%%%%%%%%%%%%%%%%%%%%%%%%%%%%%%%%%%%%%%%
\chapter{Continuation Dimension}
%%%%%%%%%%%%%%%%%%%%%%%%%%%%%%%%%%%%%%%%%%%%%%%%%%%%%%%%%%%%%%%%%%%%%%%%%%%%%%

\section{Dimension is not primitive}

Dimension has traditionally been introduced by counting independent coordinates,
independent vectors, or independent geometric directions.

Continuation Mathematics rejects all of these as primitive.

Dimension is not assumed.

Dimension is recovered from the intrinsic algebra of continuation.

The primitive object is not a dimension.

The primitive object is the family of possible continuations.

%%%%%%%%%%%%%%%%%%%%%%%%%%%%%%%%%%%%%%%%%%%%%%%%%%%%%%%%%%%%%%%%%%%%%%%%%%%%%%

\section{Continuation cones}

Every object possesses a collection of immediate continuations.

\begin{definition}[Continuation cone]

Let
\[
(X,\rightsquigarrow)
\]
be a continuation space.

For
\[
x\in X,
\]

the \emph{continuation cone} of \(x\) is

\[
\mathcal C(x)
=
\{
y\in X:
x\rightsquigarrow y
\}.
\]

\end{definition}

The continuation cone records every possible first continuation of an object.

It is the local continuation geometry surrounding that object.

%%%%%%%%%%%%%%%%%%%%%%%%%%%%%%%%%%%%%%%%%%%%%%%%%%%%%%%%%%%%%%%%%%%%%%%%%%%%%%

\section{Continuation dependence}

Not every continuation represents genuinely new mathematical information.

Some continuations arise by factoring through others.

\begin{definition}

A continuation

\[
x\rightsquigarrow z
\]

is said to be generated by another continuation

\[
x\rightsquigarrow y
\]

if there exists a continuation chain

\[
y
\rightsquigarrow
\cdots
\rightsquigarrow
z.
\]

\end{definition}

Generated continuations introduce no genuinely new continuation direction.

%%%%%%%%%%%%%%%%%%%%%%%%%%%%%%%%%%%%%%%%%%%%%%%%%%%%%%%%%%%%%%%%%%%%%%%%%%%%%%

\section{Continuation generators}

The continuation cone possesses distinguished generating families.

\begin{definition}

A subset

\[
G(x)
\subseteq
\mathcal C(x)
\]

is called a generating family if every continuation in

\[
\mathcal C(x)
\]

is generated by some member of \(G(x)\).

\end{definition}

%%%%%%%%%%%%%%%%%%%%%%%%%%%%%%%%%%%%%%%%%%%%%%%%%%%%%%%%%%%%%%%%%%%%%%%%%%%%%%

\section{Minimal generating families}

Among all generating families some are minimal.

\begin{definition}

A generating family

\[
G(x)
\]

is minimal if no proper subset still generates the continuation cone.

\end{definition}

%%%%%%%%%%%%%%%%%%%%%%%%%%%%%%%%%%%%%%%%%%%%%%%%%%%%%%%%%%%%%%%%%%%%%%%%%%%%%%

\section{Continuation rank}

The intrinsic size of a continuation cone is measured by its minimal generators.

\begin{definition}

The \emph{continuation rank} of

\[
x
\]

is the cardinality of a minimal generating family of

\[
\mathcal C(x).
\]

It is denoted

\[
\operatorname{rank}_C(x).
\]

\end{definition}

Unlike branching degree, continuation rank ignores redundant continuations.

%%%%%%%%%%%%%%%%%%%%%%%%%%%%%%%%%%%%%%%%%%%%%%%%%%%%%%%%%%%%%%%%%%%%%%%%%%%%%%

\section{Continuation dimension}

Dimension is now recovered from continuation rank.

\begin{definition}

The continuation dimension of

\[
x
\]

is

\[
\dim_C(x)
=
\operatorname{rank}_C(x).
\]

\end{definition}

Dimension is therefore not primitive.

It is a derived invariant of continuation geometry.

%%%%%%%%%%%%%%%%%%%%%%%%%%%%%%%%%%%%%%%%%%%%%%%%%%%%%%%%%%%%%%%%%%%%%%%%%%%%%%

\section{Dimension of a continuation space}

The continuation dimension of a space is

\[
\dim_C(X)
=
\sup_{x\in X}
\dim_C(x).
\]

%%%%%%%%%%%%%%%%%%%%%%%%%%%%%%%%%%%%%%%%%%%%%%%%%%%%%%%%%%%%%%%%%%%%%%%%%%%%%%

\section{Universality}

Every mathematical object possessing continuation possesses a continuation
cone.

Every continuation cone possesses generating families.

Every generating family determines continuation rank.

Continuation rank determines continuation dimension.

Thus dimension is recovered from continuation itself.

No reference has been made to coordinates,
linear algebra,
topology,
metrics,
or geometry.

Dimension has become a universal continuation invariant.


%%%%%%%%%%%%%%%%%%%%%%%%%%%%%%%%%%%%%%%%%%%%%%%%%%%%%%%%%%%%%%%%%%%%%%%%%%%%%%
\chapter{Continuation Compactness}
%%%%%%%%%%%%%%%%%%%%%%%%%%%%%%%%%%%%%%%%%%%%%%%%%%%%%%%%%%%%%%%%%%%%%%%%%%%%%%

The preceding chapters have developed the intrinsic geometry of continuation
spaces. We have defined closure, interior, boundary, frontiers, and dimension
purely from the behaviour of admissible continuations.

The next structural question is unavoidable.

Given an infinite continuation process, must some coherent continuation always
survive?

Classical topology answers an analogous question through open covers. That
approach is intentionally avoided here. Open covers describe a space from
outside. Continuation Mathematics derives compactness from the internal
behaviour of continuation itself.

The result is a new primitive notion.

%%%%%%%%%%%%%%%%%%%%%%%%%%%%%%%%%%%%%%%%%%%%%%%%%%%%%%%%%%%%%%%%%%%%%%%%%%%%%%%
\section{Infinite continuation}
%%%%%%%%%%%%%%%%%%%%%%%%%%%%%%%%%%%%%%%%%%%%%%%%%%%%%%%%%%%%%%%%%%%%%%%%%%%%%%%

Let $(X,\mathcal C)$ be a continuation space.

An infinite continuation is a sequence

\[
x_0,x_1,x_2,\ldots
\]

such that

\[
x_{i+1}\in\mathcal C(x_i)
\]

for every $i\ge0$.

Thus every step is locally admissible.

The sequence itself represents a continuation process rather than merely a
collection of points.

%%%%%%%%%%%%%%%%%%%%%%%%%%%%%%%%%%%%%%%%%%%%%%%%%%%%%%%%%%%%%%%%%%%%%%%%%%%%%%%
\section{Coherent subsequences}
%%%%%%%%%%%%%%%%%%%%%%%%%%%%%%%%%%%%%%%%%%%%%%%%%%%%%%%%%%%%%%%%%%%%%%%%%%%%%%%

Not every infinite continuation possesses global structure.

Some eventually become inconsistent.

Others repeatedly revisit coherent regions.

The distinction is fundamental.

\begin{definition}[Coherent continuation]
An infinite continuation is called coherent if every finite prefix extends to an
admissible continuation.
\end{definition}

Equivalently, no finite stage permanently blocks all future continuation.

Coherence is therefore an intrinsic extension property.

%%%%%%%%%%%%%%%%%%%%%%%%%%%%%%%%%%%%%%%%%%%%%%%%%%%%%%%%%%%%%%%%%%%%%%%%%%%%%%%
\section{Continuation accumulation}
%%%%%%%%%%%%%%%%%%%%%%%%%%%%%%%%%%%%%%%%%%%%%%%%%%%%%%%%%%%%%%%%%%%%%%%%%%%%%%%

Classical topology defines accumulation through neighbourhoods.

Continuation Mathematics defines accumulation through persistent
extendability.

\begin{definition}[Continuation accumulation point]
Let

\[
(x_i)_{i\ge0}
\]

be an infinite continuation.

A point

\[
p\in X
\]

is called a continuation accumulation point if every continuation neighbourhood
of $p$ contains infinitely many terms of the sequence.
\end{definition}

This definition depends only upon continuation neighbourhoods and therefore
requires no metric.

Accumulation measures persistent revisitation under continuation rather than
distance.

%%%%%%%%%%%%%%%%%%%%%%%%%%%%%%%%%%%%%%%%%%%%%%%%%%%%%%%%%%%%%%%%%%%%%%%%%%%%%%%
\section{Continuation compactness}
%%%%%%%%%%%%%%%%%%%%%%%%%%%%%%%%%%%%%%%%%%%%%%%%%%%%%%%%%%%%%%%%%%%%%%%%%%%%%%%

We now introduce the central notion.

\begin{definition}[Continuation compactness]
A continuation space is continuation compact if every infinite continuation
contains a coherent infinite subcontinuation possessing a continuation
accumulation point.
\end{definition}

This definition contains no reference to open covers.

Instead it states that continuation itself can never escape indefinitely without
generating persistent coherent behaviour.

Compactness is therefore an intrinsic stabilization principle.

%%%%%%%%%%%%%%%%%%%%%%%%%%%%%%%%%%%%%%%%%%%%%%%%%%%%%%%%%%%%%%%%%%%%%%%%%%%%%%%
\section{Local versus global continuation}
%%%%%%%%%%%%%%%%%%%%%%%%%%%%%%%%%%%%%%%%%%%%%%%%%%%%%%%%%%%%%%%%%%%%%%%%%%%%%%%

Continuation compactness separates two fundamentally different phenomena.

Local continuation means that every individual step is admissible.

Global continuation means that infinitely many local steps remain jointly
coherent.

These need not coincide.

A system may admit arbitrarily long finite continuations while possessing no
infinite coherent continuation.

Conversely, compact continuation spaces force coherent infinite behaviour to
reappear through accumulation.

Compactness therefore transforms local admissibility into global persistence.

%%%%%%%%%%%%%%%%%%%%%%%%%%%%%%%%%%%%%%%%%%%%%%%%%%%%%%%%%%%%%%%%%%%%%%%%%%%%%%%
\section{Finite spaces}
%%%%%%%%%%%%%%%%%%%%%%%%%%%%%%%%%%%%%%%%%%%%%%%%%%%%%%%%%%%%%%%%%%%%%%%%%%%%%%%

Finite continuation spaces satisfy compactness automatically.

\begin{theorem}[Finite Continuation Compactness]
Every finite continuation space is continuation compact.
\end{theorem}

\begin{proof}
Every infinite continuation visits only finitely many points.

Hence some point occurs infinitely often.

That point is a continuation accumulation point.

Passing to the corresponding infinite subsequence produces a coherent infinite
subcontinuation.

Therefore the space is continuation compact.
\end{proof}

Thus compactness generalizes the elementary pigeonhole principle into a
continuation principle.

%%%%%%%%%%%%%%%%%%%%%%%%%%%%%%%%%%%%%%%%%%%%%%%%%%%%%%%%%%%%%%%%%%%%%%%%%%%%%%%
\section{Compactness as continuation conservation}
%%%%%%%%%%%%%%%%%%%%%%%%%%%%%%%%%%%%%%%%%%%%%%%%%%%%%%%%%%%%%%%%%%%%%%%%%%%%%%%

Continuation compactness expresses a conservation law.

Infinite continuation cannot disperse forever.

Eventually some continuation behaviour must recur.

Persistence therefore replaces escape.

This principle is entirely independent of topology, metric, measure, or
algebra.

It belongs solely to continuation itself.

%%%%%%%%%%%%%%%%%%%%%%%%%%%%%%%%%%%%%%%%%%%%%%%%%%%%%%%%%%%%%%%%%%%%%%%%%%%%%%%
\section{Compactness and closure}
%%%%%%%%%%%%%%%%%%%%%%%%%%%%%%%%%%%%%%%%%%%%%%%%%%%%%%%%%%%%%%%%%%%%%%%%%%%%%%%

Continuation closure and continuation compactness are complementary.

Closure determines what continuations cannot leave.

Compactness determines what infinite continuations cannot avoid.

Closure governs existence.

Compactness governs persistence.

Together they form the two principal stabilization mechanisms of continuation
geometry.

%%%%%%%%%%%%%%%%%%%%%%%%%%%%%%%%%%%%%%%%%%%%%%%%%%%%%%%%%%%%%%%%%%%%%%%%%%%%%%%
\section{Universal characterization}
%%%%%%%%%%%%%%%%%%%%%%%%%%%%%%%%%%%%%%%%%%%%%%%%%%%%%%%%%%%%%%%%%%%%%%%%%%%%%%%

Continuation compactness is the first genuinely global property of a
continuation space.

Dimension measures available freedom.

Boundary measures obstruction.

Closure measures completion.

Compactness measures unavoidable persistence.

These four notions together constitute the structural backbone of continuation
geometry.

They will subsequently interact to produce continuation algebra and, ultimately,
the universal theory of continuation itself.

%%%%%%%%%%%%%%%%%%%%%%%%%%%%%%%%%%%%%%%%%%%%%%%%%%%%%%%%%%%%%%%%%%%%%%%%%%%%%%
\chapter{Continuation Connectivity}
%%%%%%%%%%%%%%%%%%%%%%%%%%%%%%%%%%%%%%%%%%%%%%%%%%%%%%%%%%%%%%%%%%%%%%%%%%%%%%

The preceding chapter established the principle of continuation compactness.
Compactness governs the persistence of infinite continuation.

The next structural question concerns communication.

When are two mathematical objects connected through continuation itself?

Classical topology introduces connectivity through open sets, paths, or
separation. Continuation Mathematics begins instead from admissible
continuation. Connectivity is not imposed externally; it is generated by the
existence of continuation chains.

%%%%%%%%%%%%%%%%%%%%%%%%%%%%%%%%%%%%%%%%%%%%%%%%%%%%%%%%%%%%%%%%%%%%%%%%%%%%%%%
\section{Continuation chains}
%%%%%%%%%%%%%%%%%%%%%%%%%%%%%%%%%%%%%%%%%%%%%%%%%%%%%%%%%%%%%%%%%%%%%%%%%%%%%%%

The primitive notion is the existence of successive admissible continuations.

\begin{definition}[Continuation chain]
Let $(X,\mathcal C)$ be a continuation space.

A continuation chain from $x$ to $y$ is a finite sequence

\[
x=x_0,x_1,\ldots,x_n=y
\]

such that

\[
x_{i+1}\in\mathcal C(x_i)
\]

for every $i=0,\ldots,n-1$.
\end{definition}

A continuation chain is therefore a finite admissible propagation through the
space.

Connectivity begins with the existence of such chains.

%%%%%%%%%%%%%%%%%%%%%%%%%%%%%%%%%%%%%%%%%%%%%%%%%%%%%%%%%%%%%%%%%%%%%%%%%%%%%%%
\section{Continuation-connected points}
%%%%%%%%%%%%%%%%%%%%%%%%%%%%%%%%%%%%%%%%%%%%%%%%%%%%%%%%%%%%%%%%%%%%%%%%%%%%%%%

\begin{definition}[Continuation-connected]
Two points

\[
x,y\in X
\]

are called continuation-connected if there exists a continuation chain from
$x$ to $y$.
\end{definition}

This relation depends only upon admissible continuation.

No metric, topology, graph structure, or algebra is assumed.

%%%%%%%%%%%%%%%%%%%%%%%%%%%%%%%%%%%%%%%%%%%%%%%%%%%%%%%%%%%%%%%%%%%%%%%%%%%%%%%
\section{Symmetric connectivity}
%%%%%%%%%%%%%%%%%%%%%%%%%%%%%%%%%%%%%%%%%%%%%%%%%%%%%%%%%%%%%%%%%%%%%%%%%%%%%%%

Continuation itself is generally directional.

A continuation from $x$ to $y$ need not admit one from $y$ to $x$.

This asymmetry is intrinsic.

Accordingly we distinguish directed and mutual continuation.

\begin{definition}[Strong continuation connectivity]
Two points are strongly continuation-connected if each is continuation-connected
to the other.
\end{definition}

Strong continuation connectivity is symmetric.

Ordinary continuation connectivity need not be.

%%%%%%%%%%%%%%%%%%%%%%%%%%%%%%%%%%%%%%%%%%%%%%%%%%%%%%%%%%%%%%%%%%%%%%%%%%%%%%%
\section{Continuation components}
%%%%%%%%%%%%%%%%%%%%%%%%%%%%%%%%%%%%%%%%%%%%%%%%%%%%%%%%%%%%%%%%%%%%%%%%%%%%%%%

Connectivity partitions a continuation space into maximal coherent regions.

\begin{definition}[Continuation component]
A continuation component is a maximal subset

\[
Y\subseteq X
\]

such that every pair of points of $Y$ is strongly continuation-connected.
\end{definition}

Each continuation component is therefore internally self-sustaining.

No admissible continuation leaves the component while preserving mutual
continuation.

%%%%%%%%%%%%%%%%%%%%%%%%%%%%%%%%%%%%%%%%%%%%%%%%%%%%%%%%%%%%%%%%%%%%%%%%%%%%%%%
\section{Disconnected spaces}
%%%%%%%%%%%%%%%%%%%%%%%%%%%%%%%%%%%%%%%%%%%%%%%%%%%%%%%%%%%%%%%%%%%%%%%%%%%%%%%

A continuation space need not possess only one component.

\begin{definition}[Disconnected continuation space]
A continuation space is disconnected if it possesses at least two distinct
continuation components.

Otherwise it is continuation-connected.
\end{definition}

Disconnectedness therefore measures the failure of continuation to propagate
through the entire mathematical object.

%%%%%%%%%%%%%%%%%%%%%%%%%%%%%%%%%%%%%%%%%%%%%%%%%%%%%%%%%%%%%%%%%%%%%%%%%%%%%%%
\section{Continuation barriers}
%%%%%%%%%%%%%%%%%%%%%%%%%%%%%%%%%%%%%%%%%%%%%%%%%%%%%%%%%%%%%%%%%%%%%%%%%%%%%%%

The boundary between two continuation components is not necessarily geometric.

Instead it is generated by obstruction.

\begin{definition}[Continuation barrier]
A continuation barrier is a collection of objects through which no admissible
continuation chain can pass.
\end{definition}

Continuation barriers produce disconnectedness.

Their existence is independent of topology.

They arise solely from the failure of continuation.

%%%%%%%%%%%%%%%%%%%%%%%%%%%%%%%%%%%%%%%%%%%%%%%%%%%%%%%%%%%%%%%%%%%%%%%%%%%%%%%
\section{Reachability}
%%%%%%%%%%%%%%%%%%%%%%%%%%%%%%%%%%%%%%%%%%%%%%%%%%%%%%%%%%%%%%%%%%%%%%%%%%%%%%%

Every point generates a natural region of influence.

\begin{definition}[Reachability set]
The reachability set of a point $x$ is

\[
\operatorname{Reach}(x)
=
\{
y\in X:
y
\text{ is continuation-connected to }
x
\}.
\]
\end{definition}

Reachability is generated recursively.

Beginning from a single object, one repeatedly applies admissible continuation.

The resulting collection is the smallest continuation-closed set containing the
starting point.

%%%%%%%%%%%%%%%%%%%%%%%%%%%%%%%%%%%%%%%%%%%%%%%%%%%%%%%%%%%%%%%%%%%%%%%%%%%%%%%
\section{The continuation preorder}
%%%%%%%%%%%%%%%%%%%%%%%%%%%%%%%%%%%%%%%%%%%%%%%%%%%%%%%%%%%%%%%%%%%%%%%%%%%%%%%

Continuation connectivity naturally induces an order relation.

\begin{definition}[Continuation preorder]
Define

\[
x\preceq y
\]

whenever $y$ is continuation-connected to $x$.
\end{definition}

\begin{theorem}
The relation $\preceq$ is a preorder.
\end{theorem}

\begin{proof}
Every object is connected to itself by the empty continuation chain, so the
relation is reflexive.

If

\[
x\preceq y
\]

and

\[
y\preceq z,
\]

then concatenating the two continuation chains produces a continuation chain
from $x$ to $z$.

Hence

\[
x\preceq z.
\]

Therefore $\preceq$ is transitive.
\end{proof}

Strong continuation connectivity is precisely the equivalence relation generated
by this preorder.

Thus continuation components are exactly the equivalence classes of mutual
continuation.

%%%%%%%%%%%%%%%%%%%%%%%%%%%%%%%%%%%%%%%%%%%%%%%%%%%%%%%%%%%%%%%%%%%%%%%%%%%%%%%
\section{Connectivity as propagation}
%%%%%%%%%%%%%%%%%%%%%%%%%%%%%%%%%%%%%%%%%%%%%%%%%%%%%%%%%%%%%%%%%%%%%%%%%%%%%%%

Continuation connectivity should not be interpreted merely as the existence of
paths.

Its deeper meaning is propagation.

A continuation component is a region inside which mathematical information may
propagate without obstruction.

Disconnectedness identifies absolute barriers to propagation.

Thus connectivity measures not geometric nearness but structural
communicability.

%%%%%%%%%%%%%%%%%%%%%%%%%%%%%%%%%%%%%%%%%%%%%%%%%%%%%%%%%%%%%%%%%%%%%%%%%%%%%%%
\section{Universal characterization}
%%%%%%%%%%%%%%%%%%%%%%%%%%%%%%%%%%%%%%%%%%%%%%%%%%%%%%%%%%%%%%%%%%%%%%%%%%%%%%%

Continuation connectivity completes the fundamental local geometry of
continuation spaces.

Dimension measures freedom.

Boundary measures obstruction.

Closure measures completion.

Compactness measures persistence.

Connectivity measures propagation.

Together these notions determine the intrinsic geometry generated by
continuation itself.

The remaining chapter of Part III will introduce Completion Spaces, where
continuation geometry reaches its maximal completion. Completion spaces will
serve as the bridge from continuation geometry into continuation algebra, where
continuation will become an algebraic operation rather than merely a geometric
possibility.

%%%%%%%%%%%%%%%%%%%%%%%%%%%%%%%%%%%%%%%%%%%%%%%%%%%%%%%%%%%%%%%%%%%%%%%%%%%%%%
\chapter{Completion Spaces}
%%%%%%%%%%%%%%%%%%%%%%%%%%%%%%%%%%%%%%%%%%%%%%%%%%%%%%%%%%%%%%%%%%%%%%%%%%%%%%

The preceding chapters have developed the intrinsic geometry of continuation.
Every continuation space possesses boundaries, interiors, dimensions,
compactness properties, and connectivity generated entirely by admissible
continuation.

A fundamental question nevertheless remains.

When continuation is pursued without arbitrary interruption, what mathematical
object is ultimately obtained?

The purpose of this chapter is to answer that question.

Completion is not defined by adjoining missing points, as in classical
analysis. Instead, completion is the intrinsic terminal object generated by
continuation itself.

%%%%%%%%%%%%%%%%%%%%%%%%%%%%%%%%%%%%%%%%%%%%%%%%%%%%%%%%%%%%%%%%%%%%%%%%%%%%%%%
\section{Completion}
%%%%%%%%%%%%%%%%%%%%%%%%%%%%%%%%%%%%%%%%%%%%%%%%%%%%%%%%%%%%%%%%%%%%%%%%%%%%%%%

Every continuation process enlarges mathematical information.

Some processes terminate naturally.

Others continue indefinitely.

Completion describes the mathematical object obtained after every admissible
continuation has been exhausted.

\begin{definition}[Completion]
Let $(X,\mathcal C)$ be a continuation space.

A completion of $X$ is a continuation space

\[
\widehat X
\]

equipped with an embedding

\[
\iota:X\hookrightarrow\widehat X
\]

such that

\begin{enumerate}
\item every admissible continuation in $X$ remains admissible in
$\widehat X$;

\item every admissible continuation beginning in $\iota(X)$ possesses a
terminal realization inside $\widehat X$;

\item no proper continuation subspace of $\widehat X$ satisfies the preceding
two conditions.
\end{enumerate}

\end{definition}

Thus completion is characterized simultaneously by extension,
realization, and minimality.

%%%%%%%%%%%%%%%%%%%%%%%%%%%%%%%%%%%%%%%%%%%%%%%%%%%%%%%%%%%%%%%%%%%%%%%%%%%%%%%
\section{Terminal continuation}
%%%%%%%%%%%%%%%%%%%%%%%%%%%%%%%%%%%%%%%%%%%%%%%%%%%%%%%%%%%%%%%%%%%%%%%%%%%%%%%

Completion is governed by a universal terminal property.

\begin{definition}[Terminal continuation]
A continuation is terminal if no admissible continuation extends it further.
\end{definition}

Completion therefore consists precisely of adjoining all terminal continuations
forced by the continuation system.

Nothing else is introduced.

%%%%%%%%%%%%%%%%%%%%%%%%%%%%%%%%%%%%%%%%%%%%%%%%%%%%%%%%%%%%%%%%%%%%%%%%%%%%%%%
\section{Completion points}
%%%%%%%%%%%%%%%%%%%%%%%%%%%%%%%%%%%%%%%%%%%%%%%%%%%%%%%%%%%%%%%%%%%%%%%%%%%%%%%

Completion introduces new mathematical objects.

These are not arbitrary.

Each represents an entire coherent continuation process.

\begin{definition}[Completion point]
A completion point is a maximal coherent continuation that admits no proper
extension.
\end{definition}

Ordinary objects describe finite stages.

Completion points describe completed continuation histories.

%%%%%%%%%%%%%%%%%%%%%%%%%%%%%%%%%%%%%%%%%%%%%%%%%%%%%%%%%%%%%%%%%%%%%%%%%%%%%%%
\section{The completion boundary}
%%%%%%%%%%%%%%%%%%%%%%%%%%%%%%%%%%%%%%%%%%%%%%%%%%%%%%%%%%%%%%%%%%%%%%%%%%%%%%%

Every completion separates original objects from newly completed ones.

\begin{definition}[Completion boundary]
The completion boundary is

\[
\partial_c X
=
\widehat X\setminus\iota(X).
\]
\end{definition}

The completion boundary contains precisely those objects created through
completion.

It measures incompleteness of the original continuation space.

%%%%%%%%%%%%%%%%%%%%%%%%%%%%%%%%%%%%%%%%%%%%%%%%%%%%%%%%%%%%%%%%%%%%%%%%%%%%%%%
\section{Completion closure}
%%%%%%%%%%%%%%%%%%%%%%%%%%%%%%%%%%%%%%%%%%%%%%%%%%%%%%%%%%%%%%%%%%%%%%%%%%%%%%%

The completion possesses a natural closure property.

\begin{theorem}[Completion Closure]
Every admissible continuation in a completion space terminates inside the
completion.
\end{theorem}

\begin{proof}
This is immediate from the defining property of completion.

Every admissible continuation possesses a terminal realization in
$\widehat X$.

Hence continuation cannot escape the completion.

Therefore the completion is continuation closed.
\end{proof}

%%%%%%%%%%%%%%%%%%%%%%%%%%%%%%%%%%%%%%%%%%%%%%%%%%%%%%%%%%%%%%%%%%%%%%%%%%%%%%%
\section{Minimality}
%%%%%%%%%%%%%%%%%%%%%%%%%%%%%%%%%%%%%%%%%%%%%%%%%%%%%%%%%%%%%%%%%%%%%%%%%%%%%%%

Completion is not merely closed.

It is the smallest continuation space having this property.

\begin{theorem}[Minimal Completion]
If

\[
Y
\]

is any continuation space satisfying the defining properties of completion,
then

\[
\widehat X
\]

embeds into $Y$.

Consequently the completion is unique up to continuation isomorphism.
\end{theorem}

\begin{proof}
Both spaces contain realizations of every admissible continuation.

Minimality forces each completed continuation to correspond uniquely to the
same terminal continuation in the other space.

The resulting correspondence preserves continuation and is therefore a
continuation isomorphism.
\end{proof}

Thus completion is canonical.

%%%%%%%%%%%%%%%%%%%%%%%%%%%%%%%%%%%%%%%%%%%%%%%%%%%%%%%%%%%%%%%%%%%%%%%%%%%%%%%
\section{Completion and dimension}
%%%%%%%%%%%%%%%%%%%%%%%%%%%%%%%%%%%%%%%%%%%%%%%%%%%%%%%%%%%%%%%%%%%%%%%%%%%%%%%

Completion preserves intrinsic structure.

\begin{theorem}[Dimension Preservation]
Completion does not increase continuation dimension.

Every independent continuation direction already appears before completion.

Completion merely realizes terminal continuations that were previously
incomplete.
\end{theorem}

\begin{proof}
Dimension measures independent continuation directions.

Completion adds no new directions.

It merely extends already existing ones to their maximal realizations.

Hence continuation dimension is preserved.
\end{proof}

Thus completion enlarges realization rather than freedom.

%%%%%%%%%%%%%%%%%%%%%%%%%%%%%%%%%%%%%%%%%%%%%%%%%%%%%%%%%%%%%%%%%%%%%%%%%%%%%%%
\section{Completion and compactness}
%%%%%%%%%%%%%%%%%%%%%%%%%%%%%%%%%%%%%%%%%%%%%%%%%%%%%%%%%%%%%%%%%%%%%%%%%%%%%%%

Completion and compactness are complementary principles.

Compactness guarantees persistence.

Completion guarantees realization.

Persistence without completion allows infinite unfinished continuation.

Completion without compactness allows isolated terminal objects without
persistent behaviour.

Together they describe stabilized continuation.

%%%%%%%%%%%%%%%%%%%%%%%%%%%%%%%%%%%%%%%%%%%%%%%%%%%%%%%%%%%%%%%%%%%%%%%%%%%%%%%
\section{Universal completion}
%%%%%%%%%%%%%%%%%%%%%%%%%%%%%%%%%%%%%%%%%%%%%%%%%%%%%%%%%%%%%%%%%%%%%%%%%%%%%%%

Among all continuation spaces there exists a distinguished class.

\begin{definition}[Complete continuation space]
A continuation space is complete if it is equal to its own completion.
\end{definition}

Equivalently,

\[
X=\widehat X.
\]

Such spaces already contain realizations of every admissible continuation.

Nothing further can be forced.

%%%%%%%%%%%%%%%%%%%%%%%%%%%%%%%%%%%%%%%%%%%%%%%%%%%%%%%%%%%%%%%%%%%%%%%%%%%%%%%
\section{The Completion Principle}
%%%%%%%%%%%%%%%%%%%%%%%%%%%%%%%%%%%%%%%%%%%%%%%%%%%%%%%%%%%%%%%%%%%%%%%%%%%%%%%

The preceding constructions establish the central geometric principle of
Continuation Mathematics.

\begin{theorem}[Completion Principle]
Every continuation space admits a unique completion, up to continuation
isomorphism.
\end{theorem}

\begin{proof}
Construct the collection of all maximal coherent continuations generated by
the continuation structure.

Adjoin these as completion points.

The resulting space is continuation closed.

Minimality follows because every completion must contain precisely these
terminal continuations.

Uniqueness therefore follows from the universal property.
\end{proof}

The Completion Principle is the geometric analogue of completion theorems
throughout classical mathematics, yet it does not depend upon metric,
topology, order, algebra, or measure.

Instead it arises solely from continuation itself.

%%%%%%%%%%%%%%%%%%%%%%%%%%%%%%%%%%%%%%%%%%%%%%%%%%%%%%%%%%%%%%%%%%%%%%%%%%%%%%%
\section{Completion as the geometry of possibility}
%%%%%%%%%%%%%%%%%%%%%%%%%%%%%%%%%%%%%%%%%%%%%%%%%%%%%%%%%%%%%%%%%%%%%%%%%%%%%%%

Completion concludes the geometric development of Continuation Mathematics.

Continuation spaces begin with admissible local extension.

They acquire frontiers, interiors, boundaries, dimensions, compactness,
and connectivity.

Completion gathers every coherent continuation into a single universal
geometric object.

The result is the maximal realization of mathematical possibility generated by
continuation.

Geometry therefore culminates not in distance or shape but in completion.

The next part of this work begins the algebraic study of continuation.

Where geometry asks which continuations exist, algebra asks how
continuations combine.

The transition from continuation spaces to continuation algebra is therefore
forced by the Completion Principle itself.

Every completed continuation naturally becomes an algebraic element whose
compositions, decompositions, identities, and symmetries constitute the
subject of the next part.

\part{Continuation Algebra}

%%%%%%%%%%%%%%%%%%%%%%%%%%%%%%%%%%%%%%%%%%%%%%%%%%%%%%%%%%%%%%%%%%%%%%%%%%%%%%
\chapter{Continuation Algebra}
%%%%%%%%%%%%%%%%%%%%%%%%%%%%%%%%%%%%%%%%%%%%%%%%%%%%%%%%%%%%%%%%%%%%%%%%%%%%%%

The preceding part established the geometry of continuation. Every continuation
space possesses intrinsic notions of boundary, interior, closure,
compactness, connectivity, and completion.

Geometry alone, however, cannot describe how continuations interact.

The purpose of this part is to develop the algebra generated by continuation
itself.

Classical algebra begins with operations defined on previously existing
objects. Continuation Mathematics reverses this order.

The primitive mathematical objects are continuations.

Algebra is recovered from their interaction.

%%%%%%%%%%%%%%%%%%%%%%%%%%%%%%%%%%%%%%%%%%%%%%%%%%%%%%%%%%%%%%%%%%%%%%%%%%%%%%%
\section{The algebraic viewpoint}
%%%%%%%%%%%%%%%%%%%%%%%%%%%%%%%%%%%%%%%%%%%%%%%%%%%%%%%%%%%%%%%%%%%%%%%%%%%%%%%

A continuation is not merely a process.

Whenever two continuations may be performed consecutively, they determine a
new continuation.

Thus continuation possesses an intrinsic composition.

The existence of composition is not postulated.

It is forced by the definition of continuation.

The first objective of continuation algebra is therefore to understand the
algebra generated by admissible composition.

%%%%%%%%%%%%%%%%%%%%%%%%%%%%%%%%%%%%%%%%%%%%%%%%%%%%%%%%%%%%%%%%%%%%%%%%%%%%%%%
\section{Composable continuations}
%%%%%%%%%%%%%%%%%%%%%%%%%%%%%%%%%%%%%%%%%%%%%%%%%%%%%%%%%%%%%%%%%%%%%%%%%%%%%%%

Let $(X,\mathcal C)$ be a continuation space.

Suppose

\[
\alpha:x\rightsquigarrow y
\]

and

\[
\beta:y\rightsquigarrow z
\]

are admissible continuations.

Since the terminal object of the first continuation is the initial object of
the second, they naturally determine a longer continuation.

%%%%%%%%%%%%%%%%%%%%%%%%%%%%%%%%%%%%%%%%%%%%%%%%%%%%%%%%%%%%%%%%%%%%%%%%%%%%%%%
\section{Continuation composition}
%%%%%%%%%%%%%%%%%%%%%%%%%%%%%%%%%%%%%%%%%%%%%%%%%%%%%%%%%%%%%%%%%%%%%%%%%%%%%%%

\begin{definition}[Continuation composition]
Let

\[
\alpha:x\rightsquigarrow y,
\qquad
\beta:y\rightsquigarrow z
\]

be admissible continuations.

The continuation composition

\[
\beta\circ\alpha
\]

is the continuation obtained by performing $\alpha$ followed by $\beta$.
\end{definition}

Composition is defined precisely when the terminal object of the first
continuation agrees with the initial object of the second.

Thus composition is intrinsically partial.

%%%%%%%%%%%%%%%%%%%%%%%%%%%%%%%%%%%%%%%%%%%%%%%%%%%%%%%%%%%%%%%%%%%%%%%%%%%%%%%
\section{Associativity}
%%%%%%%%%%%%%%%%%%%%%%%%%%%%%%%%%%%%%%%%%%%%%%%%%%%%%%%%%%%%%%%%%%%%%%%%%%%%%%%

Composition satisfies the fundamental associative law.

\begin{theorem}[Associativity]
Whenever all compositions are defined,

\[
(\gamma\circ\beta)\circ\alpha
=
\gamma\circ(\beta\circ\alpha).
\]
\end{theorem}

\begin{proof}
Both sides represent the same ordered execution of the three admissible
continuations.

Only the placement of parentheses differs.

Since continuation records order rather than grouping, both constructions
produce the identical continuation.

Therefore composition is associative.
\end{proof}

Associativity is therefore structural.

It follows from continuation itself rather than from any numerical operation.

%%%%%%%%%%%%%%%%%%%%%%%%%%%%%%%%%%%%%%%%%%%%%%%%%%%%%%%%%%%%%%%%%%%%%%%%%%%%%%%
\section{Identity continuations}
%%%%%%%%%%%%%%%%%%%%%%%%%%%%%%%%%%%%%%%%%%%%%%%%%%%%%%%%%%%%%%%%%%%%%%%%%%%%%%%

Every mathematical object possesses a trivial continuation.

\begin{definition}[Identity continuation]
For every object

\[
x\in X
\]

the identity continuation

\[
\operatorname{id}_x
\]

is the continuation that leaves $x$ unchanged.
\end{definition}

Identity continuations represent the absence of change rather than the absence
of mathematical content.

%%%%%%%%%%%%%%%%%%%%%%%%%%%%%%%%%%%%%%%%%%%%%%%%%%%%%%%%%%%%%%%%%%%%%%%%%%%%%%%
\section{Identity laws}
%%%%%%%%%%%%%%%%%%%%%%%%%%%%%%%%%%%%%%%%%%%%%%%%%%%%%%%%%%%%%%%%%%%%%%%%%%%%%%%

\begin{theorem}[Identity]
Whenever the compositions are defined,

\[
\alpha\circ\operatorname{id}_x
=
\alpha,
\]

and

\[
\operatorname{id}_y\circ\alpha
=
\alpha,
\]

for every continuation

\[
\alpha:x\rightsquigarrow y.
\]
\end{theorem}

\begin{proof}
Composition with an identity continuation contributes no additional
continuation.

Hence the original continuation is preserved.

Therefore both identities hold.
\end{proof}

%%%%%%%%%%%%%%%%%%%%%%%%%%%%%%%%%%%%%%%%%%%%%%%%%%%%%%%%%%%%%%%%%%%%%%%%%%%%%%%
\section{Continuation algebra}
%%%%%%%%%%%%%%%%%%%%%%%%%%%%%%%%%%%%%%%%%%%%%%%%%%%%%%%%%%%%%%%%%%%%%%%%%%%%%%%

The collection of all admissible continuations equipped with composition and
identity forms the fundamental algebraic object of Continuation Mathematics.

\begin{definition}[Continuation algebra]
The continuation algebra of a continuation space consists of

\begin{enumerate}
\item all admissible continuations,

\item the partial composition operation,

\item the family of identity continuations.
\end{enumerate}
\end{definition}

Unlike ordinary algebraic systems, the binary operation is not globally
defined.

Its domain is determined by continuation itself.

%%%%%%%%%%%%%%%%%%%%%%%%%%%%%%%%%%%%%%%%%%%%%%%%%%%%%%%%%%%%%%%%%%%%%%%%%%%%%%%
\section{Generated algebra}
%%%%%%%%%%%%%%%%%%%%%%%%%%%%%%%%%%%%%%%%%%%%%%%%%%%%%%%%%%%%%%%%%%%%%%%%%%%%%%%

Every continuation algebra is generated from elementary continuations.

Repeated composition constructs all finite continuations.

Thus the entire algebra is generated recursively from its primitive admissible
extensions.

The algebra therefore records every possible finite mathematical evolution
inside the continuation space.

%%%%%%%%%%%%%%%%%%%%%%%%%%%%%%%%%%%%%%%%%%%%%%%%%%%%%%%%%%%%%%%%%%%%%%%%%%%%%%%
\section{Completion and algebra}
%%%%%%%%%%%%%%%%%%%%%%%%%%%%%%%%%%%%%%%%%%%%%%%%%%%%%%%%%%%%%%%%%%%%%%%%%%%%%%%

Completion spaces established the maximal realization of continuation.

Continuation algebra studies the internal laws obeyed by those realizations.

Geometry determines existence.

Algebra determines interaction.

The two theories are therefore complementary manifestations of the same
underlying continuation structure.

%%%%%%%%%%%%%%%%%%%%%%%%%%%%%%%%%%%%%%%%%%%%%%%%%%%%%%%%%%%%%%%%%%%%%%%%%%%%%%%
\section{The Algebra Principle}
%%%%%%%%%%%%%%%%%%%%%%%%%%%%%%%%%%%%%%%%%%%%%%%%%%%%%%%%%%%%%%%%%%%%%%%%%%%%%%%

The constructions of this chapter establish the first fundamental principle of
continuation algebra.

\begin{theorem}[Algebra Principle]
Every continuation space canonically generates a continuation algebra through
admissible composition.
\end{theorem}

\begin{proof}
Every continuation space determines its admissible continuations.

Whenever two continuations are composable, their composition is uniquely
defined.

Identity continuations exist for every object.

Associativity has already been established.

Therefore every continuation space canonically generates a continuation
algebra.
\end{proof}

Continuation algebra is therefore not an additional mathematical structure.

It is the inevitable algebra generated by continuation itself.

The remaining chapters of this part will investigate increasingly refined
algebraic structures generated from continuation, culminating in universal
continuation algebras whose operations are entirely determined by continuation
rather than imposed externally.


%%%%%%%%%%%%%%%%%%%%%%%%%%%%%%%%%%%%%%%%%%%%%%%%%%%%%%%%%%%%%%%%%%%%%%%%%%%%%%
\chapter{Continuation Operations}
%%%%%%%%%%%%%%%%%%%%%%%%%%%%%%%%%%%%%%%%%%%%%%%%%%%%%%%%%%%%%%%%%%%%%%%%%%%%%%

The preceding chapter established that every continuation space canonically
generates an algebra through admissible composition.

The next question concerns the operations acting upon continuation.

Classically, an operation is simply a function between sets.

Continuation Mathematics requires considerably more.

An operation must preserve continuation itself.

The preservation of continuation is therefore more fundamental than the
evaluation of functions.

%%%%%%%%%%%%%%%%%%%%%%%%%%%%%%%%%%%%%%%%%%%%%%%%%%%%%%%%%%%%%%%%%%%%%%%%%%%%%%%
\section{Operations as continuation-preserving transformations}
%%%%%%%%%%%%%%%%%%%%%%%%%%%%%%%%%%%%%%%%%%%%%%%%%%%%%%%%%%%%%%%%%%%%%%%%%%%%%%%

Every mathematical operation transforms one collection of objects into another.

Not every transformation is mathematically admissible.

A transformation that destroys continuation destroys the intrinsic structure of
the objects upon which it acts.

Accordingly, admissible operations are defined by preservation rather than
evaluation.

%%%%%%%%%%%%%%%%%%%%%%%%%%%%%%%%%%%%%%%%%%%%%%%%%%%%%%%%%%%%%%%%%%%%%%%%%%%%%%%
\section{Continuation operations}
%%%%%%%%%%%%%%%%%%%%%%%%%%%%%%%%%%%%%%%%%%%%%%%%%%%%%%%%%%%%%%%%%%%%%%%%%%%%%%%

\begin{definition}[Continuation operation]
Let

\[
(X,\mathcal C_X)
\quad\text{and}\quad
(Y,\mathcal C_Y)
\]

be continuation spaces.

A continuation operation is a mapping

\[
T:X\longrightarrow Y
\]

such that whenever

\[
y\in\mathcal C_X(x),
\]

one has

\[
T(y)\in\mathcal C_Y(T(x)).
\]

\end{definition}

Thus admissible continuation is preserved by the operation.

%%%%%%%%%%%%%%%%%%%%%%%%%%%%%%%%%%%%%%%%%%%%%%%%%%%%%%%%%%%%%%%%%%%%%%%%%%%%%%%
\section{Preservation principle}
%%%%%%%%%%%%%%%%%%%%%%%%%%%%%%%%%%%%%%%%%%%%%%%%%%%%%%%%%%%%%%%%%%%%%%%%%%%%%%%

Continuation operations preserve the continuation structure rather than merely
individual objects.

\begin{theorem}[Continuation Preservation]
If

\[
x_0,x_1,\ldots,x_n
\]

is a continuation chain in \(X\), then

\[
T(x_0),T(x_1),\ldots,T(x_n)
\]

is a continuation chain in \(Y\).
\end{theorem}

\begin{proof}

Each consecutive pair satisfies

\[
x_{i+1}\in\mathcal C_X(x_i).
\]

Since \(T\) preserves continuation,

\[
T(x_{i+1})
\in
\mathcal C_Y(T(x_i))
\]

for every \(i\).

Hence the transformed sequence remains a continuation chain.

\end{proof}

Thus continuation operations preserve every finite propagation.

%%%%%%%%%%%%%%%%%%%%%%%%%%%%%%%%%%%%%%%%%%%%%%%%%%%%%%%%%%%%%%%%%%%%%%%%%%%%%%%
\section{Composition of operations}
%%%%%%%%%%%%%%%%%%%%%%%%%%%%%%%%%%%%%%%%%%%%%%%%%%%%%%%%%%%%%%%%%%%%%%%%%%%%%%%

Continuation operations compose naturally.

\begin{theorem}[Closure under composition]

Let

\[
T:X\longrightarrow Y,
\qquad
S:Y\longrightarrow Z
\]

be continuation operations.

Then

\[
S\circ T
\]

is again a continuation operation.

\end{theorem}

\begin{proof}

Suppose

\[
y\in\mathcal C_X(x).
\]

Then

\[
T(y)\in\mathcal C_Y(T(x)).
\]

Applying the continuation-preserving property of \(S\),

\[
S(T(y))
\in
\mathcal C_Z(S(T(x))).
\]

Therefore

\[
S\circ T
\]

preserves continuation.

\end{proof}

Hence continuation operations are closed under composition.

%%%%%%%%%%%%%%%%%%%%%%%%%%%%%%%%%%%%%%%%%%%%%%%%%%%%%%%%%%%%%%%%%%%%%%%%%%%%%%%
\section{Identity operation}
%%%%%%%%%%%%%%%%%%%%%%%%%%%%%%%%%%%%%%%%%%%%%%%%%%%%%%%%%%%%%%%%%%%%%%%%%%%%%%%

Every continuation space possesses a distinguished operation.

\begin{definition}[Identity operation]

The identity operation on a continuation space is

\[
\operatorname{id}_X:X\longrightarrow X,
\]

defined by

\[
\operatorname{id}_X(x)=x.
\]

\end{definition}

The identity preserves every continuation trivially.

%%%%%%%%%%%%%%%%%%%%%%%%%%%%%%%%%%%%%%%%%%%%%%%%%%%%%%%%%%%%%%%%%%%%%%%%%%%%%%%
\section{Invertible operations}
%%%%%%%%%%%%%%%%%%%%%%%%%%%%%%%%%%%%%%%%%%%%%%%%%%%%%%%%%%%%%%%%%%%%%%%%%%%%%%%

Some continuation operations admit inverses.

\begin{definition}[Continuation automorphism]

A continuation automorphism is a bijective continuation operation

\[
T:X\longrightarrow X
\]

whose inverse is also a continuation operation.

\end{definition}

Such operations preserve the continuation structure completely.

They represent intrinsic symmetries of continuation.

%%%%%%%%%%%%%%%%%%%%%%%%%%%%%%%%%%%%%%%%%%%%%%%%%%%%%%%%%%%%%%%%%%%%%%%%%%%%%%%
\section{Generated operations}
%%%%%%%%%%%%%%%%%%%%%%%%%%%%%%%%%%%%%%%%%%%%%%%%%%%%%%%%%%%%%%%%%%%%%%%%%%%%%%%

Elementary continuation operations generate more complicated operations through
composition.

Beginning with primitive continuation-preserving transformations,

\[
T_1,T_2,\ldots,T_n,
\]

one constructs

\[
T_n\circ\cdots\circ T_2\circ T_1.
\]

Thus every finite operation is recursively generated from elementary ones.

%%%%%%%%%%%%%%%%%%%%%%%%%%%%%%%%%%%%%%%%%%%%%%%%%%%%%%%%%%%%%%%%%%%%%%%%%%%%%%%
\section{Local and global operations}
%%%%%%%%%%%%%%%%%%%%%%%%%%%%%%%%%%%%%%%%%%%%%%%%%%%%%%%%%%%%%%%%%%%%%%%%%%%%%%%

Continuation operations naturally divide into two classes.

A local continuation operation preserves admissible continuation only within a
specified continuation region.

A global continuation operation preserves continuation throughout the entire
space.

This distinction parallels the difference between local and global continuation
developed geometrically in the preceding part.

%%%%%%%%%%%%%%%%%%%%%%%%%%%%%%%%%%%%%%%%%%%%%%%%%%%%%%%%%%%%%%%%%%%%%%%%%%%%%%%
\section{The operation algebra}
%%%%%%%%%%%%%%%%%%%%%%%%%%%%%%%%%%%%%%%%%%%%%%%%%%%%%%%%%%%%%%%%%%%%%%%%%%%%%%%

The continuation operations acting on a continuation space themselves possess
algebraic structure.

Composition is associative.

The identity operation exists.

Invertible operations form the symmetry structure of the continuation space.

Thus operations become algebraic objects in their own right.

The algebra of operations is therefore distinct from the algebra generated by
individual continuations.

%%%%%%%%%%%%%%%%%%%%%%%%%%%%%%%%%%%%%%%%%%%%%%%%%%%%%%%%%%%%%%%%%%%%%%%%%%%%%%%
\section{The Operation Principle}
%%%%%%%%%%%%%%%%%%%%%%%%%%%%%%%%%%%%%%%%%%%%%%%%%%%%%%%%%%%%%%%%%%%%%%%%%%%%%%%

The preceding constructions establish the second foundational principle of
continuation algebra.

\begin{theorem}[Operation Principle]

Every continuation-preserving transformation induces a transformation of the
entire continuation algebra.

\end{theorem}

\begin{proof}

A continuation operation preserves every admissible continuation chain.

Hence it preserves compositions, identities, reachability, continuation
components, and every algebraic construction generated from continuation.

Therefore the operation acts naturally on the continuation algebra itself.

\end{proof}

Continuation operations therefore preserve not merely mathematical objects but
entire systems of continuation.

This principle marks the beginning of the intrinsic symmetry theory of
Continuation Mathematics.

The subsequent chapters will show that the algebra generated by continuation
operations naturally organizes into increasingly rich algebraic structures,
beginning with continuation monoids.

%%%%%%%%%%%%%%%%%%%%%%%%%%%%%%%%%%%%%%%%%%%%%%%%%%%%%%%%%%%%%%%%%%%%%%%%%%%%%%
\chapter{Continuation Monoids}
%%%%%%%%%%%%%%%%%%%%%%%%%%%%%%%%%%%%%%%%%%%%%%%%%%%%%%%%%%%%%%%%%%%%%%%%%%%%%%

The first genuinely algebraic object forced by Continuation Mathematics is not
a group but a monoid.

This is not an arbitrary choice. A continuation may always be extended by
performing another continuation, yet an extension need not admit an inverse.
Continuation is therefore naturally directional. The algebra of continuation is
an algebra of construction rather than reversal.

The monoid therefore appears before every richer algebraic structure.

%%%%%%%%%%%%%%%%%%%%%%%%%%%%%%%%%%%%%%%%%%%%%%%%%%%%%%%%%%%%%%%%%%%%%%%%%%%%%%
\section{The continuation product}
%%%%%%%%%%%%%%%%%%%%%%%%%%%%%%%%%%%%%%%%%%%%%%%%%%%%%%%%%%%%%%%%%%%%%%%%%%%%%%

Let $(C,\preceq)$ be a continuation system.

Assume that continuation composition
\[
\circ:C\times C\longrightarrow C
\]
is defined whenever the terminal state of the first continuation coincides with
the initial state of the second.

Whenever this composition is everywhere defined, we obtain a binary operation

\[
*:C\times C\rightarrow C,
\]

called the \emph{continuation product}.

Intuitively,

\[
x*y
\]

means

\[
\text{perform }x\text{ followed by }y.
\]

No cancellation or reversibility is assumed.

%%%%%%%%%%%%%%%%%%%%%%%%%%%%%%%%%%%%%%%%%%%%%%%%%%%%%%%%%%%%%%%%%%%%%%%%%%%%%%
\section{Associativity}
%%%%%%%%%%%%%%%%%%%%%%%%%%%%%%%%%%%%%%%%%%%%%%%%%%%%%%%%%%%%%%%%%%%%%%%%%%%%%%

Continuation composition is forced to be associative.

\begin{theorem}[Associativity]
Whenever the compositions are defined,

\[
(x*y)*z
=
x*(y*z).
\]
\end{theorem}

\begin{proof}

Both expressions describe exactly the continuation obtained by performing
$x$, then $y$, then $z$.

The order of execution is identical.

Only the placement of parentheses differs.

Therefore the resulting continuation is the same.

\end{proof}

Associativity is therefore not an imposed algebraic law.

It is forced by the meaning of successive continuation.

%%%%%%%%%%%%%%%%%%%%%%%%%%%%%%%%%%%%%%%%%%%%%%%%%%%%%%%%%%%%%%%%%%%%%%%%%%%%%%
\section{Identity continuation}
%%%%%%%%%%%%%%%%%%%%%%%%%%%%%%%%%%%%%%%%%%%%%%%%%%%%%%%%%%%%%%%%%%%%%%%%%%%%%%

Every continuation system possesses a distinguished continuation that performs
no change.

\begin{definition}

A continuation

\[
e\in C
\]

is an \emph{identity continuation} if

\[
e*x=x,
\]

and

\[
x*e=x
\]

for every continuation $x$.

\end{definition}

The identity continuation represents persistence without extension.

%%%%%%%%%%%%%%%%%%%%%%%%%%%%%%%%%%%%%%%%%%%%%%%%%%%%%%%%%%%%%%%%%%%%%%%%%%%%%%
\section{Continuation monoids}
%%%%%%%%%%%%%%%%%%%%%%%%%%%%%%%%%%%%%%%%%%%%%%%%%%%%%%%%%%%%%%%%%%%%%%%%%%%%%%

\begin{definition}

A \emph{continuation monoid} is a triple

\[
(C,*,e)
\]

satisfying

\begin{enumerate}
\item associativity;

\item existence of an identity continuation.
\end{enumerate}

\end{definition}

No inverse operation is required.

%%%%%%%%%%%%%%%%%%%%%%%%%%%%%%%%%%%%%%%%%%%%%%%%%%%%%%%%%%%%%%%%%%%%%%%%%%%%%%
\section{Why inverses are exceptional}
%%%%%%%%%%%%%%%%%%%%%%%%%%%%%%%%%%%%%%%%%%%%%%%%%%%%%%%%%%%%%%%%%%%%%%%%%%%%%%

Most mathematical constructions are irreversible.

Examples include

\[
\begin{aligned}
&\text{proof extension},\\
&\text{adding information},\\
&\text{time evolution},\\
&\text{Collatz iteration},\\
&\text{algorithm execution},\\
&\text{language parsing},\\
&\text{cell division},\\
&\text{entropy increase}.
\end{aligned}
\]

Each admits continuation.

None naturally admits inversion.

Consequently groups arise only after imposing additional structure.

Monoids therefore constitute the natural algebra of continuation.

%%%%%%%%%%%%%%%%%%%%%%%%%%%%%%%%%%%%%%%%%%%%%%%%%%%%%%%%%%%%%%%%%%%%%%%%%%%%%%
\section{Generated continuation monoids}
%%%%%%%%%%%%%%%%%%%%%%%%%%%%%%%%%%%%%%%%%%%%%%%%%%%%%%%%%%%%%%%%%%%%%%%%%%%%%%

A continuation system is frequently generated by elementary continuation steps.

\begin{definition}

Let

\[
G\subseteq C.
\]

The continuation monoid generated by $G$ is the smallest continuation monoid
containing every element of $G$.

It is denoted

\[
\langle G\rangle.
\]

\end{definition}

Every continuation is therefore a finite product

\[
g_1g_2\cdots g_n,
\]

where each

\[
g_i\in G.
\]

Thus finite continuation is generated from elementary continuation primitives.

%%%%%%%%%%%%%%%%%%%%%%%%%%%%%%%%%%%%%%%%%%%%%%%%%%%%%%%%%%%%%%%%%%%%%%%%%%%%%%
\section{Length}
%%%%%%%%%%%%%%%%%%%%%%%%%%%%%%%%%%%%%%%%%%%%%%%%%%%%%%%%%%%%%%%%%%%%%%%%%%%%%%

Generation immediately forces a natural measure.

\begin{definition}

For

\[
x\in\langle G\rangle,
\]

the \emph{continuation length}

\[
\ell(x)
\]

is the minimum number of generators required to construct $x$.

\end{definition}

Unlike classical word length, continuation length measures constructive depth
rather than merely symbolic complexity.

%%%%%%%%%%%%%%%%%%%%%%%%%%%%%%%%%%%%%%%%%%%%%%%%%%%%%%%%%%%%%%%%%%%%%%%%%%%%%%
\section{Prefix order}
%%%%%%%%%%%%%%%%%%%%%%%%%%%%%%%%%%%%%%%%%%%%%%%%%%%%%%%%%%%%%%%%%%%%%%%%%%%%%%

The continuation product induces a canonical order.

\begin{definition}

For continuations $x,y$ define

\[
x\leq_p y
\]

whenever there exists

\[
z
\]

such that

\[
xz=y.
\]

The relation $\leq_p$ is called the \emph{prefix order}.

\end{definition}

The prefix order records constructive ancestry.

One continuation precedes another precisely when the latter extends the former.

%%%%%%%%%%%%%%%%%%%%%%%%%%%%%%%%%%%%%%%%%%%%%%%%%%%%%%%%%%%%%%%%%%%%%%%%%%%%%%
\section{Cancellation}
%%%%%%%%%%%%%%%%%%%%%%%%%%%%%%%%%%%%%%%%%%%%%%%%%%%%%%%%%%%%%%%%%%%%%%%%%%%%%%

Cancellation becomes a structural question rather than an axiom.

\begin{definition}

A continuation monoid is

\begin{enumerate}
\item left cancellative if

\[
ax=ay
\Longrightarrow
x=y;
\]

\item right cancellative if

\[
xa=ya
\Longrightarrow
x=y.
\]

\end{enumerate}

\end{definition}

Many continuation monoids fail cancellation.

Failure of cancellation reflects the existence of distinct histories that
produce identical continuations.

Such phenomena will become central in later applications.

%%%%%%%%%%%%%%%%%%%%%%%%%%%%%%%%%%%%%%%%%%%%%%%%%%%%%%%%%%%%%%%%%%%%%%%%%%%%%%
\section{Continuation congruences}
%%%%%%%%%%%%%%%%%%%%%%%%%%%%%%%%%%%%%%%%%%%%%%%%%%%%%%%%%%%%%%%%%%%%%%%%%%%%%%

Not every distinction between continuations is structurally significant.

\begin{definition}

A relation

\[
\sim
\]

on a continuation monoid is a
\emph{continuation congruence} if

\[
x\sim y
\]

implies

\[
ax\sim ay,
\qquad
xa\sim ya
\]

for every continuation $a$.

\end{definition}

Continuation congruences identify continuations while preserving every future
extension.

%%%%%%%%%%%%%%%%%%%%%%%%%%%%%%%%%%%%%%%%%%%%%%%%%%%%%%%%%%%%%%%%%%%%%%%%%%%%%%
\section{The universal viewpoint}
%%%%%%%%%%%%%%%%%%%%%%%%%%%%%%%%%%%%%%%%%%%%%%%%%%%%%%%%%%%%%%%%%%%%%%%%%%%%%%

Every continuation system possesses an intrinsic monoid.

Every continuation space carries one.

Every continuation morphism preserves one.

The continuation monoid is therefore not an additional algebraic object.

It is the universal algebra extracted from the possibility of finite extension
itself.

Later chapters will enrich this monoid with additional operations, leading to
continuation semirings, continuation algebras, and ultimately the algebraic
structures underlying completion theory.

%%%%%%%%%%%%%%%%%%%%%%%%%%%%%%%%%%%%%%%%%%%%%%%%%%%%%%%%%%%%%%%%%%%%%%%%%%%%%%
\chapter{Continuation Semigroups}
%%%%%%%%%%%%%%%%%%%%%%%%%%%%%%%%%%%%%%%%%%%%%%%%%%%%%%%%%%%%%%%%%%%%%%%%%%%%%%

The preceding chapter introduced continuation monoids as continuation systems
whose admissible continuations are closed under composition and possess a
distinguished neutral continuation. The existence of an identity continuation,
however, is an additional structural property rather than a necessary one.

Many continuation systems admit composition without admitting a universal
initial continuation from which every continuation may begin. Such systems give
rise naturally to semigroups.

The purpose of this chapter is to recover semigroups as intrinsic continuation
objects rather than as primitive algebraic structures.

%%%%%%%%%%%%%%%%%%%%%%%%%%%%%%%%%%%%%%%%%%%%%%%%%%%%%%%%%%%%%%%%%%%%%%%%%%%%%%%
\section{Incomplete continuation}
%%%%%%%%%%%%%%%%%%%%%%%%%%%%%%%%%%%%%%%%%%%%%%%%%%%%%%%%%%%%%%%%%%%%%%%%%%%%%%%

Continuation is fundamentally directional.

Whenever one continuation may legitimately follow another, composition becomes
possible. Nothing in this observation requires the existence of a distinguished
empty continuation.

Consequently one may possess an algebra of continuation while lacking an
identity element.

This is precisely the situation described by semigroups.

%%%%%%%%%%%%%%%%%%%%%%%%%%%%%%%%%%%%%%%%%%%%%%%%%%%%%%%%%%%%%%%%%%%%%%%%%%%%%%%
\section{Continuation semigroups}
%%%%%%%%%%%%%%%%%%%%%%%%%%%%%%%%%%%%%%%%%%%%%%%%%%%%%%%%%%%%%%%%%%%%%%%%%%%%%%%

\begin{definition}[Continuation Semigroup]
A \emph{continuation semigroup} is a pair
\[
(\mathcal{S},\circ)
\]
consisting of a continuation system together with an associative continuation
composition
\[
\circ:\mathcal{S}\times\mathcal{S}\longrightarrow\mathcal{S}
\]
satisfying

\begin{enumerate}
\item closure,
\item associativity,
\item continuation admissibility.
\end{enumerate}

No identity continuation is assumed.
\end{definition}

Thus semigroups arise whenever continuation exists but has no canonical
beginning.

%%%%%%%%%%%%%%%%%%%%%%%%%%%%%%%%%%%%%%%%%%%%%%%%%%%%%%%%%%%%%%%%%%%%%%%%%%%%%%%
\section{Semigroups as broken monoids}
%%%%%%%%%%%%%%%%%%%%%%%%%%%%%%%%%%%%%%%%%%%%%%%%%%%%%%%%%%%%%%%%%%%%%%%%%%%%%%%

Every continuation monoid determines a continuation semigroup by forgetting its
identity.

The converse need not hold.

\begin{theorem}[Identity Completion]
Every continuation monoid determines a continuation semigroup.

Not every continuation semigroup admits a continuation identity.
\end{theorem}

\begin{proof}
Removing the identity from a monoid preserves closure and associativity.

Conversely, adjoining an identity requires the existence of an element acting
neutrally upon every continuation. Such an element need not exist.
\end{proof}

Identity is therefore an additional continuation property rather than a basic
algebraic necessity.

%%%%%%%%%%%%%%%%%%%%%%%%%%%%%%%%%%%%%%%%%%%%%%%%%%%%%%%%%%%%%%%%%%%%%%%%%%%%%%%
\section{Continuation generation}
%%%%%%%%%%%%%%%%%%%%%%%%%%%%%%%%%%%%%%%%%%%%%%%%%%%%%%%%%%%%%%%%%%%%%%%%%%%%%%%

Let
\[
G\subseteq\mathcal S
\]
be a collection of primitive continuations.

Repeated continuation composition produces

\[
G,
\qquad
G^2,
\qquad
G^3,
\qquad
\ldots
\]

whose union forms the continuation semigroup generated by $G$.

\begin{definition}
The smallest continuation semigroup containing $G$ is called the
\emph{continuation semigroup generated by $G$}.
\end{definition}

Thus generation is itself a continuation process.

%%%%%%%%%%%%%%%%%%%%%%%%%%%%%%%%%%%%%%%%%%%%%%%%%%%%%%%%%%%%%%%%%%%%%%%%%%%%%%%
\section{Continuation depth}
%%%%%%%%%%%%%%%%%%%%%%%%%%%%%%%%%%%%%%%%%%%%%%%%%%%%%%%%%%%%%%%%%%%%%%%%%%%%%%%

Without an identity, every element possesses an intrinsic depth.

\begin{definition}
The \emph{continuation depth} of
\[
x\in\mathcal S
\]
is the minimum number of primitive continuations required to construct $x$.
\end{definition}

Depth measures how far a continuation lies from primitive generation.

Unlike classical word length, continuation depth depends only upon admissible
continuation rather than upon arbitrary presentations.

%%%%%%%%%%%%%%%%%%%%%%%%%%%%%%%%%%%%%%%%%%%%%%%%%%%%%%%%%%%%%%%%%%%%%%%%%%%%%%%
\section{Principal continuation ideals}
%%%%%%%%%%%%%%%%%%%%%%%%%%%%%%%%%%%%%%%%%%%%%%%%%%%%%%%%%%%%%%%%%%%%%%%%%%%%%%%

Every continuation generates descendants.

\begin{definition}
For
\[
x\in\mathcal S,
\]
define the principal continuation ideal

\[
\mathcal I(x)
=
\{x\circ y:y\in\mathcal S\}.
\]
\end{definition}

This consists precisely of every continuation obtainable after $x$.

Thus ideals become forward continuation cones.

Likewise,

\[
\mathcal J(x)
=
\{y\circ x:y\in\mathcal S\}
\]

collects every continuation capable of producing $x$.

%%%%%%%%%%%%%%%%%%%%%%%%%%%%%%%%%%%%%%%%%%%%%%%%%%%%%%%%%%%%%%%%%%%%%%%%%%%%%%%
\section{Continuation order}
%%%%%%%%%%%%%%%%%%%%%%%%%%%%%%%%%%%%%%%%%%%%%%%%%%%%%%%%%%%%%%%%%%%%%%%%%%%%%%%

Continuation induces a natural preorder.

\begin{definition}
For
\[
x,y\in\mathcal S,
\]
define

\[
x\preceq y
\]

whenever

\[
y=x\circ z
\]

for some continuation
\[
z.
\]
\end{definition}

Thus $y$ is a continuation of $x$.

Whenever antisymmetry holds this preorder becomes a genuine partial order.

%%%%%%%%%%%%%%%%%%%%%%%%%%%%%%%%%%%%%%%%%%%%%%%%%%%%%%%%%%%%%%%%%%%%%%%%%%%%%%%
\section{Continuation growth}
%%%%%%%%%%%%%%%%%%%%%%%%%%%%%%%%%%%%%%%%%%%%%%%%%%%%%%%%%%%%%%%%%%%%%%%%%%%%%%%

Repeated continuation generates growth.

Let

\[
g(n)
=
\#G^n
\]

denote the number of distinct continuations obtainable after exactly $n$
primitive continuation steps.

\begin{definition}
The function

\[
g:\mathbb N\to\mathbb N
\]

is called the \emph{continuation growth function}.
\end{definition}

Growth therefore measures the expansion of admissible continuation rather than
merely counting algebraic words.

%%%%%%%%%%%%%%%%%%%%%%%%%%%%%%%%%%%%%%%%%%%%%%%%%%%%%%%%%%%%%%%%%%%%%%%%%%%%%%%
\section{Maximal continuation}
%%%%%%%%%%%%%%%%%%%%%%%%%%%%%%%%%%%%%%%%%%%%%%%%%%%%%%%%%%%%%%%%%%%%%%%%%%%%%%%

Some semigroups admit infinite continuation.

Others necessarily terminate.

\begin{definition}
A continuation semigroup possesses the
\emph{finite continuation property}
if every continuation chain

\[
x_1\preceq x_2\preceq x_3\preceq\cdots
\]

stabilizes after finitely many steps.
\end{definition}

Failure of stabilization signals the existence of genuine infinite continuation.

Thus infinite semigroup behaviour becomes a question about continuation rather
than merely cardinality.

%%%%%%%%%%%%%%%%%%%%%%%%%%%%%%%%%%%%%%%%%%%%%%%%%%%%%%%%%%%%%%%%%%%%%%%%%%%%%%%
\section{Semigroups as incomplete continuation algebra}
%%%%%%%%%%%%%%%%%%%%%%%%%%%%%%%%%%%%%%%%%%%%%%%%%%%%%%%%%%%%%%%%%%%%%%%%%%%%%%%

The continuation viewpoint reveals that semigroups are not primitive algebraic
objects.

They are continuation algebras in which composition exists but universal
initiation has not yet been recovered.

Monoids therefore arise as completed semigroups.

Semigroups occupy the intermediate position between arbitrary continuation
systems and fully completed continuation algebras.

This reinterpretation changes the conceptual status of semigroups completely.
Rather than being assumed at the outset, they emerge as one stage in the
structural evolution of continuation.

The next chapter studies groups, where continuation not only possesses an
identity but also admits reversible continuation. Groups therefore represent
continuation systems whose histories may always be traversed in both directions.

%%%%%%%%%%%%%%%%%%%%%%%%%%%%%%%%%%%%%%%%%%%%%%%%%%%%%%%%%%%%%%%%%%%%%%%%%%%%%%
\chapter{Continuation Lattices}
%%%%%%%%%%%%%%%%%%%%%%%%%%%%%%%%%%%%%%%%%%%%%%%%%%%%%%%%%%%%%%%%%%%%%%%%%%%%%%

The preceding chapters recovered the classical algebraic structures of monoids
and semigroups from the primitive notion of continuation. The next structure
arises from an entirely different source.

Continuation systems are naturally ordered.

Some continuations contain more information than others. Some represent
extensions of earlier continuations. Others admit common refinements or common
simplifications.

These relationships force an intrinsic order.

The purpose of this chapter is to show that lattices arise naturally from the
order structure induced by continuation.

%%%%%%%%%%%%%%%%%%%%%%%%%%%%%%%%%%%%%%%%%%%%%%%%%%%%%%%%%%%%%%%%%%%%%%%%%%%%%%
\section{Continuation refinement}
%%%%%%%%%%%%%%%%%%%%%%%%%%%%%%%%%%%%%%%%%%%%%%%%%%%%%%%%%%%%%%%%%%%%%%%%%%%%%%

Continuation is directional.

Whenever one continuation contains every admissible continuation of another,
the first may be regarded as a refinement of the second.

\begin{definition}[Continuation refinement]
Let \(A\) and \(B\) be continuation objects.

We say that \(A\) is a \emph{refinement} of \(B\), written

\[
B\preccurlyeq A,
\]

whenever every continuation admissible from \(B\) remains admissible from
\(A\).
\end{definition}

Thus refinement corresponds to increasing continuation information.

%%%%%%%%%%%%%%%%%%%%%%%%%%%%%%%%%%%%%%%%%%%%%%%%%%%%%%%%%%%%%%%%%%%%%%%%%%%%%%
\section{The continuation order}
%%%%%%%%%%%%%%%%%%%%%%%%%%%%%%%%%%%%%%%%%%%%%%%%%%%%%%%%%%%%%%%%%%%%%%%%%%%%%%

The refinement relation satisfies the basic properties of order.

\begin{theorem}
Continuation refinement is reflexive and transitive.

If distinct continuation objects are identified whenever they possess identical
continuation behaviour, then refinement becomes antisymmetric.
\end{theorem}

\begin{proof}
Every continuation refines itself.

If \(A\) refines \(B\), and \(B\) refines \(C\), then every continuation
admissible from \(C\) is admissible from \(A\), establishing transitivity.

After quotienting by continuation equivalence, antisymmetry follows
immediately.
\end{proof}

Thus continuation systems naturally determine partially ordered sets.

%%%%%%%%%%%%%%%%%%%%%%%%%%%%%%%%%%%%%%%%%%%%%%%%%%%%%%%%%%%%%%%%%%%%%%%%%%%%%%
\section{Greatest common continuation}
%%%%%%%%%%%%%%%%%%%%%%%%%%%%%%%%%%%%%%%%%%%%%%%%%%%%%%%%%%%%%%%%%%%%%%%%%%%%%%

Two continuation objects often possess a common structural core.

\begin{definition}
Let \(A\) and \(B\) be continuation objects.

A \emph{greatest common continuation} is a continuation object
\(A\wedge B\) satisfying

\[
A\wedge B\preccurlyeq A,
\qquad
A\wedge B\preccurlyeq B,
\]

and maximal with this property.
\end{definition}

This object records everything that both continuations necessarily possess.

It is the continuation-theoretic analogue of intersection.

%%%%%%%%%%%%%%%%%%%%%%%%%%%%%%%%%%%%%%%%%%%%%%%%%%%%%%%%%%%%%%%%%%%%%%%%%%%%%%
\section{Least common extension}
%%%%%%%%%%%%%%%%%%%%%%%%%%%%%%%%%%%%%%%%%%%%%%%%%%%%%%%%%%%%%%%%%%%%%%%%%%%%%%

Continuation objects may also admit a common enlargement.

\begin{definition}
The \emph{least common continuation} of two continuation objects is a
continuation object

\[
A\vee B
\]

satisfying

\[
A\preccurlyeq A\vee B,
\qquad
B\preccurlyeq A\vee B,
\]

and minimal among all such continuation objects.
\end{definition}

This object represents the smallest continuation capable of containing both
initial continuations.

%%%%%%%%%%%%%%%%%%%%%%%%%%%%%%%%%%%%%%%%%%%%%%%%%%%%%%%%%%%%%%%%%%%%%%%%%%%%%%
\section{Continuation lattices}
%%%%%%%%%%%%%%%%%%%%%%%%%%%%%%%%%%%%%%%%%%%%%%%%%%%%%%%%%%%%%%%%%%%%%%%%%%%%%%

\begin{definition}[Continuation lattice]
A continuation system is called a
\emph{continuation lattice}
whenever every pair of continuation objects admits both

\[
A\wedge B
\]

and

\[
A\vee B.
\]
\end{definition}

Thus a continuation lattice is a continuation system in which every pair of
continuations possesses both a maximal common refinement and a minimal common
extension.

%%%%%%%%%%%%%%%%%%%%%%%%%%%%%%%%%%%%%%%%%%%%%%%%%%%%%%%%%%%%%%%%%%%%%%%%%%%%%%
\section{Completion viewpoint}
%%%%%%%%%%%%%%%%%%%%%%%%%%%%%%%%%%%%%%%%%%%%%%%%%%%%%%%%%%%%%%%%%%%%%%%%%%%%%%

The lattice operations admit a continuation interpretation.

The meet

\[
A\wedge B
\]

removes incompatible continuation.

The join

\[
A\vee B
\]

adds precisely the continuation required to unify two systems.

Neither operation is primitive.

Both arise from continuation compatibility.

%%%%%%%%%%%%%%%%%%%%%%%%%%%%%%%%%%%%%%%%%%%%%%%%%%%%%%%%%%%%%%%%%%%%%%%%%%%%%%
\section{Continuation intervals}
%%%%%%%%%%%%%%%%%%%%%%%%%%%%%%%%%%%%%%%%%%%%%%%%%%%%%%%%%%%%%%%%%%%%%%%%%%%%%%

Given continuation objects

\[
A\preccurlyeq B,
\]

define the continuation interval

\[
[A,B]
=
\{
X:
A\preccurlyeq X\preccurlyeq B
\}.
\]

The interval consists of every continuation lying between two prescribed
continuation states.

Intervals measure the internal complexity of continuation.

%%%%%%%%%%%%%%%%%%%%%%%%%%%%%%%%%%%%%%%%%%%%%%%%%%%%%%%%%%%%%%%%%%%%%%%%%%%%%%
\section{Atomic continuation}
%%%%%%%%%%%%%%%%%%%%%%%%%%%%%%%%%%%%%%%%%%%%%%%%%%%%%%%%%%%%%%%%%%%%%%%%%%%%%%

Some continuation objects cannot be decomposed further.

\begin{definition}
A non-minimal continuation object is called
\emph{continuation atomic}
if it covers a unique immediate predecessor in the continuation order.
\end{definition}

Atoms represent the elementary increments of continuation.

Every larger continuation is obtained by successive refinement of such atomic
continuations whenever the lattice is atomic.

%%%%%%%%%%%%%%%%%%%%%%%%%%%%%%%%%%%%%%%%%%%%%%%%%%%%%%%%%%%%%%%%%%%%%%%%%%%%%%
\section{Distributive continuation}
%%%%%%%%%%%%%%%%%%%%%%%%%%%%%%%%%%%%%%%%%%%%%%%%%%%%%%%%%%%%%%%%%%%%%%%%%%%%%%

Some continuation systems admit independent refinement.

\begin{definition}
A continuation lattice is called
\emph{distributive}
whenever

\[
A\wedge(B\vee C)
=
(A\wedge B)\vee(A\wedge C)
\]

for every triple of continuation objects.
\end{definition}

Distributivity expresses the independence of separate continuation directions.

Failure of distributivity indicates interacting continuation phenomena that
cannot be separated.

%%%%%%%%%%%%%%%%%%%%%%%%%%%%%%%%%%%%%%%%%%%%%%%%%%%%%%%%%%%%%%%%%%%%%%%%%%%%%%
\section{Completion lattices}
%%%%%%%%%%%%%%%%%%%%%%%%%%%%%%%%%%%%%%%%%%%%%%%%%%%%%%%%%%%%%%%%%%%%%%%%%%%%%%

Every continuation lattice possesses a distinguished interpretation.

The meet removes continuation.

The join restores continuation.

Thus the lattice operations become complementary aspects of completion.

Rather than viewing lattices as abstract ordered sets equipped with two binary
operations, continuation mathematics interprets them as spaces in which
information may be both discarded and recovered through admissible
continuation.

The classical theory of lattices therefore becomes one manifestation of a much
more primitive phenomenon: the organization of continuation itself.

In subsequent chapters this viewpoint will expand from algebra into geometry,
where refinement, closure, dimension, compactness, and completion become
geometric properties of continuation spaces rather than independent axiomatic
constructions.

%%%%%%%%%%%%%%%%%%%%%%%%%%%%%%%%%%%%%%%%%%%%%%%%%%%%%%%%%%%%%%%%%%%%%%%%%%%%%%
\chapter{Continuation Categories}
%%%%%%%%%%%%%%%%%%%%%%%%%%%%%%%%%%%%%%%%%%%%%%%%%%%%%%%%%%%%%%%%%%%%%%%%%%%%%%

The preceding chapters recovered several classical algebraic structures from the
primitive notion of continuation. Composition produced monoids and semigroups,
while refinement produced lattices.

The next fundamental structure concerns not individual continuation objects but
the transformations that preserve continuation itself.

Classical category theory begins with objects and morphisms. Continuation
Mathematics proceeds in the opposite direction.

Continuation is primary.

Objects appear only as stable locations at which continuation may begin or end.

%%%%%%%%%%%%%%%%%%%%%%%%%%%%%%%%%%%%%%%%%%%%%%%%%%%%%%%%%%%%%%%%%%%%%%%%%%%%%%
\section{Continuation before objects}
%%%%%%%%%%%%%%%%%%%%%%%%%%%%%%%%%%%%%%%%%%%%%%%%%%%%%%%%%%%%%%%%%%%%%%%%%%%%%%

Suppose one has a collection of admissible continuations.

Whenever two continuations may be concatenated, a larger continuation is
produced.

Thus composition is primitive.

Only afterwards does one recognize that certain continuations possess common
sources and common targets.

These stable endpoints become the objects of the resulting category.

Objects are therefore not assumed.

They emerge.

%%%%%%%%%%%%%%%%%%%%%%%%%%%%%%%%%%%%%%%%%%%%%%%%%%%%%%%%%%%%%%%%%%%%%%%%%%%%%%
\section{Continuation objects}
%%%%%%%%%%%%%%%%%%%%%%%%%%%%%%%%%%%%%%%%%%%%%%%%%%%%%%%%%%%%%%%%%%%%%%%%%%%%%%

\begin{definition}[Continuation object]

A \emph{continuation object} is a maximal state that admits both incoming and
outgoing admissible continuations.

\end{definition}

Objects are equilibrium points within the continuation network.

They are not primitive entities but distinguished locations where continuation
may begin, terminate, or pass through.

%%%%%%%%%%%%%%%%%%%%%%%%%%%%%%%%%%%%%%%%%%%%%%%%%%%%%%%%%%%%%%%%%%%%%%%%%%%%%%
\section{Continuation morphisms}
%%%%%%%%%%%%%%%%%%%%%%%%%%%%%%%%%%%%%%%%%%%%%%%%%%%%%%%%%%%%%%%%%%%%%%%%%%%%%%

\begin{definition}

A \emph{continuation morphism}

\[
f:A\rightarrow B
\]

is an admissible continuation carrying the continuation object \(A\) into the
continuation object \(B\).

\end{definition}

Morphisms therefore represent continuation itself.

Objects merely organize the admissible continuations.

%%%%%%%%%%%%%%%%%%%%%%%%%%%%%%%%%%%%%%%%%%%%%%%%%%%%%%%%%%%%%%%%%%%%%%%%%%%%%%
\section{Composition}
%%%%%%%%%%%%%%%%%%%%%%%%%%%%%%%%%%%%%%%%%%%%%%%%%%%%%%%%%%%%%%%%%%%%%%%%%%%%%%

Whenever

\[
A
\xrightarrow{f}
B
\xrightarrow{g}
C
\]

are continuation morphisms, admissibility permits their composition

\[
g\circ f
:
A
\rightarrow
C.
\]

\begin{theorem}

Continuation composition is associative.

\end{theorem}

\begin{proof}

Associativity follows directly from the associativity of continuation
composition already established for continuation systems.

No additional structure is required.

\end{proof}

Thus categorical composition is inherited from continuation itself.

%%%%%%%%%%%%%%%%%%%%%%%%%%%%%%%%%%%%%%%%%%%%%%%%%%%%%%%%%%%%%%%%%%%%%%%%%%%%%%
\section{Identity continuation}
%%%%%%%%%%%%%%%%%%%%%%%%%%%%%%%%%%%%%%%%%%%%%%%%%%%%%%%%%%%%%%%%%%%%%%%%%%%%%%

Every continuation object possesses a trivial continuation.

\begin{definition}

The continuation

\[
\operatorname{id}_A:A\rightarrow A
\]

that performs no continuation is called the
\emph{identity continuation}.

\end{definition}

Identity continuation is the completed form of doing nothing.

It represents perfect continuation stability.

%%%%%%%%%%%%%%%%%%%%%%%%%%%%%%%%%%%%%%%%%%%%%%%%%%%%%%%%%%%%%%%%%%%%%%%%%%%%%%
\section{Continuation categories}
%%%%%%%%%%%%%%%%%%%%%%%%%%%%%%%%%%%%%%%%%%%%%%%%%%%%%%%%%%%%%%%%%%%%%%%%%%%%%%

\begin{definition}

A \emph{continuation category} consists of

\begin{enumerate}

\item continuation objects,

\item continuation morphisms,

\item associative continuation composition,

\item identity continuations.

\end{enumerate}

\end{definition}

Thus every continuation category is generated by admissible continuation.

%%%%%%%%%%%%%%%%%%%%%%%%%%%%%%%%%%%%%%%%%%%%%%%%%%%%%%%%%%%%%%%%%%%%%%%%%%%%%%
\section{Continuation diagrams}
%%%%%%%%%%%%%%%%%%%%%%%%%%%%%%%%%%%%%%%%%%%%%%%%%%%%%%%%%%%%%%%%%%%%%%%%%%%%%%

Continuation categories naturally organize themselves into diagrams.

A continuation diagram records every admissible continuation between a finite
collection of continuation objects.

Commutativity expresses the independence of continuation order.

Whenever two different continuation paths produce the same completed
continuation, the corresponding diagram commutes.

Thus commutativity expresses structural inevitability rather than algebraic
coincidence.

%%%%%%%%%%%%%%%%%%%%%%%%%%%%%%%%%%%%%%%%%%%%%%%%%%%%%%%%%%%%%%%%%%%%%%%%%%%%%%
\section{Universal continuation}
%%%%%%%%%%%%%%%%%%%%%%%%%%%%%%%%%%%%%%%%%%%%%%%%%%%%%%%%%%%%%%%%%%%%%%%%%%%%%%

Certain continuation objects are characterized entirely by their continuation
behavior.

\begin{definition}

A continuation object is called
\emph{universal}
whenever every admissible continuation toward (or away from) the object factors
uniquely through it.

\end{definition}

Universality therefore expresses maximal continuation efficiency.

Many classical universal constructions arise from this single continuation
principle.

%%%%%%%%%%%%%%%%%%%%%%%%%%%%%%%%%%%%%%%%%%%%%%%%%%%%%%%%%%%%%%%%%%%%%%%%%%%%%%
\section{Continuation equivalence}
%%%%%%%%%%%%%%%%%%%%%%%%%%%%%%%%%%%%%%%%%%%%%%%%%%%%%%%%%%%%%%%%%%%%%%%%%%%%%%

Two continuation categories need not be identical in order to possess identical
continuation structure.

\begin{definition}

Two continuation categories are called
\emph{continuation equivalent}
whenever there exist continuation-preserving transformations between them whose
compositions preserve every continuation up to continuation isomorphism.

\end{definition}

Equivalence therefore becomes preservation of continuation rather than
preservation of presentation.

%%%%%%%%%%%%%%%%%%%%%%%%%%%%%%%%%%%%%%%%%%%%%%%%%%%%%%%%%%%%%%%%%%%%%%%%%%%%%%
\section{Categories as continuation networks}
%%%%%%%%%%%%%%%%%%%%%%%%%%%%%%%%%%%%%%%%%%%%%%%%%%%%%%%%%%%%%%%%%%%%%%%%%%%%%%

Continuation Mathematics reverses the traditional logical order of category
theory.

Classically,

\[
\text{Objects}
\longrightarrow
\text{Morphisms}.
\]

Continuation Mathematics replaces this with

\[
\text{Continuation}
\longrightarrow
\text{Morphisms}
\longrightarrow
\text{Objects}
\longrightarrow
\text{Categories}.
\]

Objects are recovered from stable continuation.

Morphisms are recovered from admissible continuation.

Categories are recovered from the global organization of continuation.

Thus category theory itself becomes an emergent manifestation of continuation
rather than a primitive mathematical foundation.

This shift is representative of the philosophy of the present work: every
classical mathematical structure is reconstructed from a more primitive theory
of continuation, revealing continuation as the hidden organizational principle
underlying mathematical structure itself.

%%%%%%%%%%%%%%%%%%%%%%%%%%%%%%%%%%%%%%%%%%%%%%%%%%%%%%%%%%%%%%%%%%%%%%%%%%%%%%
\chapter{Continuation Functors}
%%%%%%%%%%%%%%%%%%%%%%%%%%%%%%%%%%%%%%%%%%%%%%%%%%%%%%%%%%%%%%%%%%%%%%%%%%%%%%

The previous chapter recovered categories from the primitive notion of
continuation. Objects emerged as stable locations within continuation systems,
while morphisms represented admissible continuation itself.

The next question is unavoidable.

When do two continuation categories describe the same continuation mathematics?

Classically this question is answered by functors.

Continuation Mathematics recovers functors from a more primitive principle.

A functor is not fundamentally a map that preserves objects and morphisms.

It is a transformation that preserves continuation.

%%%%%%%%%%%%%%%%%%%%%%%%%%%%%%%%%%%%%%%%%%%%%%%%%%%%%%%%%%%%%%%%%%%%%%%%%%%%%%
\section{Transporting continuation}
%%%%%%%%%%%%%%%%%%%%%%%%%%%%%%%%%%%%%%%%%%%%%%%%%%%%%%%%%%%%%%%%%%%%%%%%%%%%%%

Suppose two continuation categories describe different mathematical systems.

A meaningful comparison between them must preserve not merely individual
objects but the admissible continuations connecting them.

The fundamental requirement is therefore the preservation of continuation
structure.

%%%%%%%%%%%%%%%%%%%%%%%%%%%%%%%%%%%%%%%%%%%%%%%%%%%%%%%%%%%%%%%%%%%%%%%%%%%%%%
\section{Continuation-preserving transformations}
%%%%%%%%%%%%%%%%%%%%%%%%%%%%%%%%%%%%%%%%%%%%%%%%%%%%%%%%%%%%%%%%%%%%%%%%%%%%%%

\begin{definition}[Continuation Functor]

Let

\[
\mathcal C,
\qquad
\mathcal D
\]

be continuation categories.

A \emph{continuation functor}

\[
F:\mathcal C\longrightarrow\mathcal D
\]

is a transformation satisfying the following conditions.

\begin{enumerate}

\item Every continuation object of $\mathcal C$ is assigned a continuation
object of $\mathcal D$.

\item Every continuation morphism of $\mathcal C$ is assigned an admissible
continuation morphism of $\mathcal D$.

\item Every admissible continuation remains admissible after transport.

\item Composition is preserved:
\[
F(g\circ f)
=
F(g)\circ F(f).
\]

\item Identity continuations are preserved:
\[
F(\operatorname{id}_A)
=
\operatorname{id}_{F(A)}.
\]

\end{enumerate}

\end{definition}

Thus a continuation functor preserves the continuation architecture of a
category.

%%%%%%%%%%%%%%%%%%%%%%%%%%%%%%%%%%%%%%%%%%%%%%%%%%%%%%%%%%%%%%%%%%%%%%%%%%%%%%
\section{Preservation of admissibility}
%%%%%%%%%%%%%%%%%%%%%%%%%%%%%%%%%%%%%%%%%%%%%%%%%%%%%%%%%%%%%%%%%%%%%%%%%%%%%%

The defining feature of continuation functors is not merely preservation of
composition.

They preserve the admissibility relation itself.

Whenever

\[
A
\longrightarrow
B
\]

is an admissible continuation,

its image

\[
F(A)
\longrightarrow
F(B)
\]

must remain admissible.

Continuation therefore survives transportation between mathematical worlds.

%%%%%%%%%%%%%%%%%%%%%%%%%%%%%%%%%%%%%%%%%%%%%%%%%%%%%%%%%%%%%%%%%%%%%%%%%%%%%%
\section{Faithful continuation}
%%%%%%%%%%%%%%%%%%%%%%%%%%%%%%%%%%%%%%%%%%%%%%%%%%%%%%%%%%%%%%%%%%%%%%%%%%%%%%

Not every continuation functor preserves continuation equally well.

\begin{definition}

A continuation functor is called
\emph{faithful}
if distinct continuation morphisms remain distinct after transportation.

\end{definition}

Faithfulness prevents independent continuation histories from collapsing into a
single continuation.

%%%%%%%%%%%%%%%%%%%%%%%%%%%%%%%%%%%%%%%%%%%%%%%%%%%%%%%%%%%%%%%%%%%%%%%%%%%%%%
\section{Full continuation}
%%%%%%%%%%%%%%%%%%%%%%%%%%%%%%%%%%%%%%%%%%%%%%%%%%%%%%%%%%%%%%%%%%%%%%%%%%%%%%

\begin{definition}

A continuation functor is called
\emph{full}
if every admissible continuation between transported objects arises from an
admissible continuation before transportation.

\end{definition}

Thus fullness guarantees that no new continuation is artificially introduced.

%%%%%%%%%%%%%%%%%%%%%%%%%%%%%%%%%%%%%%%%%%%%%%%%%%%%%%%%%%%%%%%%%%%%%%%%%%%%%%
\section{Continuation embeddings}
%%%%%%%%%%%%%%%%%%%%%%%%%%%%%%%%%%%%%%%%%%%%%%%%%%%%%%%%%%%%%%%%%%%%%%%%%%%%%%

\begin{definition}

A continuation functor is called a
\emph{continuation embedding}
whenever it is both faithful and injective on continuation objects.

\end{definition}

Continuation embeddings preserve an entire continuation system inside a larger
one without altering its internal continuation structure.

%%%%%%%%%%%%%%%%%%%%%%%%%%%%%%%%%%%%%%%%%%%%%%%%%%%%%%%%%%%%%%%%%%%%%%%%%%%%%%
\section{Continuation equivalence}
%%%%%%%%%%%%%%%%%%%%%%%%%%%%%%%%%%%%%%%%%%%%%%%%%%%%%%%%%%%%%%%%%%%%%%%%%%%%%%

The correct notion of sameness in continuation mathematics is not equality.

It is preservation of continuation.

\begin{definition}

Two continuation categories are
\emph{continuation equivalent}
if there exist continuation functors

\[
F:\mathcal C\rightarrow\mathcal D,
\qquad
G:\mathcal D\rightarrow\mathcal C,
\]

whose compositions preserve every continuation object and every continuation
morphism up to continuation isomorphism.

\end{definition}

Equivalent continuation categories describe the same continuation mathematics,
even when their presentations differ.

%%%%%%%%%%%%%%%%%%%%%%%%%%%%%%%%%%%%%%%%%%%%%%%%%%%%%%%%%%%%%%%%%%%%%%%%%%%%%%
\section{Transport of completion}
%%%%%%%%%%%%%%%%%%%%%%%%%%%%%%%%%%%%%%%%%%%%%%%%%%%%%%%%%%%%%%%%%%%%%%%%%%%%%%

Continuation functors transport more than objects.

They transport every continuation invariant established throughout this work.

These include

\begin{itemize}

\item continuation depth,

\item continuation dimension,

\item continuation closure,

\item continuation interior,

\item continuation compactness,

\item continuation connectivity,

\item completion structure.

\end{itemize}

Thus continuation invariants are functorial.

Their numerical values may change only when continuation itself changes.

%%%%%%%%%%%%%%%%%%%%%%%%%%%%%%%%%%%%%%%%%%%%%%%%%%%%%%%%%%%%%%%%%%%%%%%%%%%%%%
\section{Functorial completion}
%%%%%%%%%%%%%%%%%%%%%%%%%%%%%%%%%%%%%%%%%%%%%%%%%%%%%%%%%%%%%%%%%%%%%%%%%%%%%%

Every completion process induces a continuation functor.

Beginning with a partial continuation system,

\[
\mathcal C,
\]

completion produces

\[
\widehat{\mathcal C}.
\]

The assignment

\[
\mathcal C
\longmapsto
\widehat{\mathcal C}
\]

extends naturally to continuation morphisms.

Consequently completion itself is functorial.

Completion is therefore not merely an operation.

It is a structural transformation preserving continuation.

%%%%%%%%%%%%%%%%%%%%%%%%%%%%%%%%%%%%%%%%%%%%%%%%%%%%%%%%%%%%%%%%%%%%%%%%%%%%%%
\section{Universality of continuation transport}
%%%%%%%%%%%%%%%%%%%%%%%%%%%%%%%%%%%%%%%%%%%%%%%%%%%%%%%%%%%%%%%%%%%%%%%%%%%%%%

Continuation functors reveal that continuation is independent of mathematical
presentation.

Different mathematical languages may encode precisely the same continuation
structure.

The role of the continuation functor is to expose this hidden identity.

Consequently functoriality is not fundamentally about preserving algebraic
operations.

It is about preserving the possibility of continuation itself.

From the perspective of Continuation Mathematics, functors become the universal
mechanism by which continuation is transported between mathematical worlds.

This concludes the algebraic development of continuation-preserving
transformations. The next chapter studies natural transformations, where the
relationships between continuation functors themselves become mathematical
objects and reveal a second level of continuation architecture.

%%%%%%%%%%%%%%%%%%%%%%%%%%%%%%%%%%%%%%%%%%%%%%%%%%%%%%%%%%%%%%%%%%%%%%%%%%%%%%
\chapter{Continuation Natural Transformations}
%%%%%%%%%%%%%%%%%%%%%%%%%%%%%%%%%%%%%%%%%%%%%%%%%%%%%%%%%%%%%%%%%%%%%%%%%%%%%%

The previous chapter established continuation functors as transformations that
preserve admissible continuation between continuation categories.

The next level of structure concerns the relationships between the functors
themselves.

Classically these relationships are called natural transformations.

Continuation Mathematics interprets them differently.

A natural transformation is itself a continuation.

Not of objects.

Not of morphisms.

But of continuation-preserving transformations.

Natural transformations therefore constitute second-order continuation.

%%%%%%%%%%%%%%%%%%%%%%%%%%%%%%%%%%%%%%%%%%%%%%%%%%%%%%%%%%%%%%%%%%%%%%%%%%%%%%
\section{Continuation of transport}
%%%%%%%%%%%%%%%%%%%%%%%%%%%%%%%%%%%%%%%%%%%%%%%%%%%%%%%%%%%%%%%%%%%%%%%%%%%%%%

Suppose

\[
F,G:\mathcal C\longrightarrow\mathcal D
\]

are continuation functors.

Each transports continuation from one mathematical world into another.

There may exist a systematic method for continuously deforming the transport
performed by \(F\) into that performed by \(G\).

Such a deformation is itself a continuation.

%%%%%%%%%%%%%%%%%%%%%%%%%%%%%%%%%%%%%%%%%%%%%%%%%%%%%%%%%%%%%%%%%%%%%%%%%%%%%%
\section{Second-order continuation}
%%%%%%%%%%%%%%%%%%%%%%%%%%%%%%%%%%%%%%%%%%%%%%%%%%%%%%%%%%%%%%%%%%%%%%%%%%%%%%

\begin{definition}[Continuation Natural Transformation]

Let

\[
F,G:\mathcal C\rightarrow\mathcal D
\]

be continuation functors.

A \emph{continuation natural transformation}

\[
\eta:F\Longrightarrow G
\]

assigns to every continuation object

\[
A\in\mathcal C
\]

a continuation morphism

\[
\eta_A:
F(A)
\longrightarrow
G(A)
\]

such that every admissible continuation is preserved coherently.

\end{definition}

Thus each object possesses its own continuation between the two transported
images.

%%%%%%%%%%%%%%%%%%%%%%%%%%%%%%%%%%%%%%%%%%%%%%%%%%%%%%%%%%%%%%%%%%%%%%%%%%%%%%
\section{Coherence}
%%%%%%%%%%%%%%%%%%%%%%%%%%%%%%%%%%%%%%%%%%%%%%%%%%%%%%%%%%%%%%%%%%%%%%%%%%%%%%

Continuation cannot depend upon arbitrary choices.

Whenever

\[
f:A\rightarrow B
\]

is a continuation morphism, the transported continuations must satisfy

\[
G(f)\circ\eta_A
=
\eta_B\circ F(f).
\]

\begin{theorem}[Continuation Coherence]

Every continuation natural transformation satisfies the coherence condition

\[
\begin{CD}
F(A) @>{F(f)}>> F(B)\\
@V{\eta_A}VV @VV{\eta_B}V\\
G(A) @>>{G(f)}> G(B)
\end{CD}
\]

for every continuation morphism \(f\).

\end{theorem}

\begin{proof}

Continuation may not depend upon whether one first transports continuation or
first transforms the transport itself.

The two admissible continuation histories therefore coincide.

\end{proof}

Thus coherence expresses the independence of continuation order.

%%%%%%%%%%%%%%%%%%%%%%%%%%%%%%%%%%%%%%%%%%%%%%%%%%%%%%%%%%%%%%%%%%%%%%%%%%%%%%
\section{Pointwise continuation}
%%%%%%%%%%%%%%%%%%%%%%%%%%%%%%%%%%%%%%%%%%%%%%%%%%%%%%%%%%%%%%%%%%%%%%%%%%%%%%

Every continuation object contributes one local continuation

\[
\eta_A.
\]

The natural transformation is therefore assembled from infinitely many local
continuations.

The coherence theorem guarantees that these local continuations combine into a
single global continuation.

Naturality is therefore a completion principle.

%%%%%%%%%%%%%%%%%%%%%%%%%%%%%%%%%%%%%%%%%%%%%%%%%%%%%%%%%%%%%%%%%%%%%%%%%%%%%%
\section{Identity continuation}
%%%%%%%%%%%%%%%%%%%%%%%%%%%%%%%%%%%%%%%%%%%%%%%%%%%%%%%%%%%%%%%%%%%%%%%%%%%%%%

Every continuation functor possesses a trivial self-continuation.

\begin{definition}

The family

\[
\operatorname{id}_F:
F
\Longrightarrow
F
\]

consisting of identity continuations

\[
\operatorname{id}_{F(A)}
\]

is called the identity continuation transformation.

\end{definition}

Identity continuation expresses complete stability of continuation transport.

%%%%%%%%%%%%%%%%%%%%%%%%%%%%%%%%%%%%%%%%%%%%%%%%%%%%%%%%%%%%%%%%%%%%%%%%%%%%%%
\section{Composition}
%%%%%%%%%%%%%%%%%%%%%%%%%%%%%%%%%%%%%%%%%%%%%%%%%%%%%%%%%%%%%%%%%%%%%%%%%%%%%%

Continuation transformations compose.

If

\[
F
\xLongrightarrow{\eta}
G
\xLongrightarrow{\mu}
H,
\]

their composition is defined pointwise,

\[
(\mu\circ\eta)_A
=
\mu_A\circ\eta_A.
\]

\begin{theorem}

Composition of continuation natural transformations is associative.

\end{theorem}

\begin{proof}

Associativity follows immediately from the associativity of continuation
composition inside the target continuation category.

\end{proof}

%%%%%%%%%%%%%%%%%%%%%%%%%%%%%%%%%%%%%%%%%%%%%%%%%%%%%%%%%%%%%%%%%%%%%%%%%%%%%%
\section{Invertible continuation}
%%%%%%%%%%%%%%%%%%%%%%%%%%%%%%%%%%%%%%%%%%%%%%%%%%%%%%%%%%%%%%%%%%%%%%%%%%%%%%

Some continuation transformations may be reversed.

\begin{definition}

A continuation natural transformation

\[
\eta:F\Longrightarrow G
\]

is called a
\emph{continuation isomorphism}
if every component

\[
\eta_A
\]

is an isomorphism.

\end{definition}

Such transformations preserve continuation completely.

No continuation information is lost.

%%%%%%%%%%%%%%%%%%%%%%%%%%%%%%%%%%%%%%%%%%%%%%%%%%%%%%%%%%%%%%%%%%%%%%%%%%%%%%
\section{Transformation spaces}
%%%%%%%%%%%%%%%%%%%%%%%%%%%%%%%%%%%%%%%%%%%%%%%%%%%%%%%%%%%%%%%%%%%%%%%%%%%%%%

Continuation transformations themselves form mathematical objects.

For fixed continuation functors

\[
F,G,
\]

write

\[
\operatorname{Nat}(F,G)
\]

for the collection of all continuation natural transformations from \(F\) to
\(G\).

These collections possess their own continuation structure through
transformation composition.

Thus continuation naturally organizes itself into successive structural levels.

%%%%%%%%%%%%%%%%%%%%%%%%%%%%%%%%%%%%%%%%%%%%%%%%%%%%%%%%%%%%%%%%%%%%%%%%%%%%%%
\section{Higher continuation}
%%%%%%%%%%%%%%%%%%%%%%%%%%%%%%%%%%%%%%%%%%%%%%%%%%%%%%%%%%%%%%%%%%%%%%%%%%%%%%

Natural transformations reveal that continuation is recursive.

Objects admit continuation.

Functors transport continuation.

Natural transformations continue continuation transport.

Nothing prevents repeating this process.

One may study continuations between continuation transformations, and then
continuations between those continuations.

Thus continuation generates an infinite hierarchy of mathematical structure.

This hierarchy is not postulated.

It is forced by the recursive nature of continuation itself.

%%%%%%%%%%%%%%%%%%%%%%%%%%%%%%%%%%%%%%%%%%%%%%%%%%%%%%%%%%%%%%%%%%%%%%%%%%%%%%
\section{The emergence of higher mathematics}
%%%%%%%%%%%%%%%%%%%%%%%%%%%%%%%%%%%%%%%%%%%%%%%%%%%%%%%%%%%%%%%%%%%%%%%%%%%%%%

Continuation Mathematics interprets natural transformations as the first
manifestation of higher-order continuation.

Rather than viewing them merely as morphisms between functors, they become
continuations acting upon continuation-preserving transformations.

This observation suggests a general principle.

Whenever mathematics develops transformations between existing structures,
Continuation Mathematics asks whether those transformations themselves admit
continuation.

Whenever the answer is affirmative, a new mathematical level emerges.

Natural transformations therefore represent not the end of category theory, but
the beginning of an infinite continuation hierarchy whose successive levels are
generated by repeatedly allowing continuation to act upon itself.

%%%%%%%%%%%%%%%%%%%%%%%%%%%%%%%%%%%%%%%%%%%%%%%%%%%%%%%%%%%%%%%%%%%%%%%%%%%%%%
\chapter{Continuation Colimits}
%%%%%%%%%%%%%%%%%%%%%%%%%%%%%%%%%%%%%%%%%%%%%%%%%%%%%%%%%%%%%%%%%%%%%%%%%%%%%%

The preceding chapter showed that limits describe the universal object that is
compatible with an entire continuation diagram. Equally fundamental is the
opposite problem.

Suppose several continuation systems possess overlapping information. Can they
be merged into a single continuation system without introducing arbitrary
choices? More generally, can every compatible family of partial continuation
systems be completed into one universal continuation object?

The answer gives rise to the notion of a continuation colimit.

%%%%%%%%%%%%%%%%%%%%%%%%%%%%%%%%%%%%%%%%%%%%%%%%%%%%%%%%%%%%%%%%%%%%%%%%%%%%%%%
\section{Dual universality}
%%%%%%%%%%%%%%%%%%%%%%%%%%%%%%%%%%%%%%%%%%%%%%%%%%%%%%%%%%%%%%%%%%%%%%%%%%%%%%%

Limits gather information flowing inward toward a universal receiver.

Colimits describe information flowing outward from a universal source.

Where a limit identifies the largest object compatible with every projection,
a colimit identifies the smallest object containing every compatible extension.

This duality is intrinsic to continuation mathematics.

A continuation process may therefore be viewed from two complementary
directions:

\[
\text{Completion by restriction}
\qquad\Longleftrightarrow\qquad
\text{Completion by extension.}
\]

%%%%%%%%%%%%%%%%%%%%%%%%%%%%%%%%%%%%%%%%%%%%%%%%%%%%%%%%%%%%%%%%%%%%%%%%%%%%%%%
\section{Continuation cocones}
%%%%%%%%%%%%%%%%%%%%%%%%%%%%%%%%%%%%%%%%%%%%%%%%%%%%%%%%%%%%%%%%%%%%%%%%%%%%%%%

\begin{definition}[Continuation cocone]
Let
\[
D:J\longrightarrow\mathbf{Cont}
\]
be a continuation diagram.

A \emph{continuation cocone} over \(D\) consists of

\[
(C,\{\iota_j\}),
\]

where

\[
\iota_j:D(j)\longrightarrow C
\]

is a continuation morphism for every object of the diagram, such that every
triangle commutes.
\end{definition}

Thus every object of the diagram extends naturally into one common continuation
system.

%%%%%%%%%%%%%%%%%%%%%%%%%%%%%%%%%%%%%%%%%%%%%%%%%%%%%%%%%%%%%%%%%%%%%%%%%%%%%%%
\section{Universal continuation extension}
%%%%%%%%%%%%%%%%%%%%%%%%%%%%%%%%%%%%%%%%%%%%%%%%%%%%%%%%%%%%%%%%%%%%%%%%%%%%%%%

\begin{definition}[Continuation colimit]
A continuation colimit of a diagram \(D\) is a continuation cocone

\[
(L,\{\lambda_j\})
\]

such that for every continuation cocone

\[
(C,\{\iota_j\})
\]

there exists a unique continuation morphism

\[
u:L\longrightarrow C
\]

making every triangle commute.
\end{definition}

Equivalently,

\[
L
\]

is the smallest continuation system into which every object of the diagram
extends.

%%%%%%%%%%%%%%%%%%%%%%%%%%%%%%%%%%%%%%%%%%%%%%%%%%%%%%%%%%%%%%%%%%%%%%%%%%%%%%%
\section{Universal minimality}
%%%%%%%%%%%%%%%%%%%%%%%%%%%%%%%%%%%%%%%%%%%%%%%%%%%%%%%%%%%%%%%%%%%%%%%%%%%%%%%

\begin{theorem}[Universal Extension Property]
Whenever a continuation colimit exists, it is unique up to unique continuation
isomorphism.
\end{theorem}

\begin{proof}
Suppose

\[
(L,\lambda_j)
\]

and

\[
(L',\lambda'_j)
\]

are both universal.

Universality of \(L\) produces

\[
f:L\rightarrow L'.
\]

Universality of \(L'\) produces

\[
g:L'\rightarrow L.
\]

The compositions satisfy

\[
g\circ f=\operatorname{id}_L,
\qquad
f\circ g=\operatorname{id}_{L'}
\]

by uniqueness.

Hence \(f\) is an isomorphism.
\end{proof}

%%%%%%%%%%%%%%%%%%%%%%%%%%%%%%%%%%%%%%%%%%%%%%%%%%%%%%%%%%%%%%%%%%%%%%%%%%%%%%%
\section{Completion by directed extension}
%%%%%%%%%%%%%%%%%%%%%%%%%%%%%%%%%%%%%%%%%%%%%%%%%%%%%%%%%%%%%%%%%%%%%%%%%%%%%%%

Continuation mathematics is fundamentally concerned with systems that grow by
successive admissible extensions.

Accordingly, the most important colimits arise from directed systems

\[
C_1
\longrightarrow
C_2
\longrightarrow
C_3
\longrightarrow
\cdots
\]

whose objects become progressively more complete.

\begin{definition}[Directed continuation system]
A continuation diagram is directed if every pair of objects possesses a common
continuation.
\end{definition}

The colimit of such a system represents the object obtained after every finite
continuation has been performed.

%%%%%%%%%%%%%%%%%%%%%%%%%%%%%%%%%%%%%%%%%%%%%%%%%%%%%%%%%%%%%%%%%%%%%%%%%%%%%%%
\section{Free continuation objects}
%%%%%%%%%%%%%%%%%%%%%%%%%%%%%%%%%%%%%%%%%%%%%%%%%%%%%%%%%%%%%%%%%%%%%%%%%%%%%%%

One of the principal uses of colimits is the construction of free objects.

\begin{definition}[Free continuation object]
Let

\[
X
\]

be a set of generators.

A continuation object

\[
F(X)
\]

is free on \(X\) if every function

\[
X\rightarrow C
\]

extends uniquely to a continuation morphism

\[
F(X)\rightarrow C.
\]
\end{definition}

Thus the free continuation object introduces no relations except those forced
by the continuation axioms themselves.

%%%%%%%%%%%%%%%%%%%%%%%%%%%%%%%%%%%%%%%%%%%%%%%%%%%%%%%%%%%%%%%%%%%%%%%%%%%%%%%
\section{Generation versus completion}
%%%%%%%%%%%%%%%%%%%%%%%%%%%%%%%%%%%%%%%%%%%%%%%%%%%%%%%%%%%%%%%%%%%%%%%%%%%%%%%

Continuation mathematics distinguishes two fundamentally different operations.

The first generates new objects.

The second completes existing ones.

Generation introduces new continuation directions.

Completion realizes directions already present.

Colimits formalize generation.

Limits formalize completion.

Both are indispensable.

%%%%%%%%%%%%%%%%%%%%%%%%%%%%%%%%%%%%%%%%%%%%%%%%%%%%%%%%%%%%%%%%%%%%%%%%%%%%%%%
\section{Completion as a colimit}
%%%%%%%%%%%%%%%%%%%%%%%%%%%%%%%%%%%%%%%%%%%%%%%%%%%%%%%%%%%%%%%%%%%%%%%%%%%%%%%

Many completion processes may themselves be viewed as directed colimits.

Suppose

\[
C_0
\subseteq
C_1
\subseteq
C_2
\subseteq
\cdots
\]

is an ascending chain of continuation systems.

Each stage contributes additional admissible continuations.

The completed object is naturally identified with

\[
\operatorname{colim} C_n.
\]

Thus completion itself is obtained as a universal extension.

%%%%%%%%%%%%%%%%%%%%%%%%%%%%%%%%%%%%%%%%%%%%%%%%%%%%%%%%%%%%%%%%%%%%%%%%%%%%%%%
\section{The continuation principle}
%%%%%%%%%%%%%%%%%%%%%%%%%%%%%%%%%%%%%%%%%%%%%%%%%%%%%%%%%%%%%%%%%%%%%%%%%%%%%%%

The interaction between limits and colimits expresses a central principle.

\begin{theorem}[Continuation Duality Principle]
Continuation limits describe universal compatibility.

Continuation colimits describe universal extension.

Every sufficiently developed continuation theory possesses both structures,
which are related by categorical duality.
\end{theorem}

%%%%%%%%%%%%%%%%%%%%%%%%%%%%%%%%%%%%%%%%%%%%%%%%%%%%%%%%%%%%%%%%%%%%%%%%%%%%%%%
\section{Toward algebraic completion}
%%%%%%%%%%%%%%%%%%%%%%%%%%%%%%%%%%%%%%%%%%%%%%%%%%%%%%%%%%%%%%%%%%%%%%%%%%%%%%%

The algebra developed thus far contains

\[
\begin{array}{c}
\text{operations},\\
\text{monoids},\\
\text{semigroups},\\
\text{lattices},\\
\text{categories},\\
\text{functors},\\
\text{natural transformations},\\
\text{limits},\\
\text{colimits}.
\end{array}
\]

Yet these structures remain external.

The next stage is to understand when algebraic continuation itself stabilizes.

Instead of asking whether individual objects admit further continuations, we
shall investigate whether entire algebraic systems eventually become complete
under repeated continuation.

This leads naturally to the notion of **Continuation Completion Algebras**, the
final chapter of Part IV.


%%%%%%%%%%%%%%%%%%%%%%%%%%%%%%%%%%%%%%%%%%%%%%%%%%%%%%%%%%%%%%%%%%%%%%%%%%%%%%
\chapter{Continuation Adjunctions}
%%%%%%%%%%%%%%%%%%%%%%%%%%%%%%%%%%%%%%%%%%%%%%%%%%%%%%%%%%%%%%%%%%%%%%%%%%%%%%

The preceding chapters established two universal constructions.

Limits determine universal compatibility.

Colimits determine universal extension.

Both describe optimal objects, but from opposite directions.

The deeper question is whether these opposite constructions are themselves
connected by a universal principle.

Continuation mathematics answers this affirmatively.

%%%%%%%%%%%%%%%%%%%%%%%%%%%%%%%%%%%%%%%%%%%%%%%%%%%%%%%%%%%%%%%%%%%%%%%%%%%%%%%
\section{Universal duality}
%%%%%%%%%%%%%%%%%%%%%%%%%%%%%%%%%%%%%%%%%%%%%%%%%%%%%%%%%%%%%%%%%%%%%%%%%%%%%%%

Many constructions occur naturally in complementary pairs.

Generation is paired with forgetting.

Restriction is paired with extension.

Locality is paired with completion.

Partiality is paired with totality.

Continuation mathematics seeks not merely these pairs individually, but the
universal mechanism that relates them.

%%%%%%%%%%%%%%%%%%%%%%%%%%%%%%%%%%%%%%%%%%%%%%%%%%%%%%%%%%%%%%%%%%%%%%%%%%%%%%%
\section{Continuation adjunctions}
%%%%%%%%%%%%%%%%%%%%%%%%%%%%%%%%%%%%%%%%%%%%%%%%%%%%%%%%%%%%%%%%%%%%%%%%%%%%%%%

\begin{definition}[Continuation adjunction]
Let

\[
F:\mathcal C\rightarrow\mathcal D,
\qquad
G:\mathcal D\rightarrow\mathcal C
\]

be continuation functors.

The pair

\[
F\dashv G
\]

is called a \emph{continuation adjunction} if there exists a natural bijection

\[
\operatorname{Hom}_{\mathcal D}(F(X),Y)
\cong
\operatorname{Hom}_{\mathcal C}(X,G(Y))
\]

for every pair of objects

\[
X\in\mathcal C,
\qquad
Y\in\mathcal D.
\]

\end{definition}

Thus every continuation from the generated object \(F(X)\) into \(Y\)
corresponds uniquely to a continuation from \(X\) into the recovered object
\(G(Y)\).

%%%%%%%%%%%%%%%%%%%%%%%%%%%%%%%%%%%%%%%%%%%%%%%%%%%%%%%%%%%%%%%%%%%%%%%%%%%%%%%
\section{Continuation units}
%%%%%%%%%%%%%%%%%%%%%%%%%%%%%%%%%%%%%%%%%%%%%%%%%%%%%%%%%%%%%%%%%%%%%%%%%%%%%%%

Every continuation adjunction determines a canonical comparison between an
object and its generated completion.

\begin{definition}[Continuation unit]
The natural transformation

\[
\eta:
\operatorname{Id}_{\mathcal C}
\Longrightarrow
GF
\]

is called the continuation unit.
\end{definition}

For every object

\[
X,
\]

the morphism

\[
\eta_X:
X\rightarrow GF(X)
\]

records the canonical continuation of \(X\) into its reconstructed image.

%%%%%%%%%%%%%%%%%%%%%%%%%%%%%%%%%%%%%%%%%%%%%%%%%%%%%%%%%%%%%%%%%%%%%%%%%%%%%%%
\section{Continuation counits}
%%%%%%%%%%%%%%%%%%%%%%%%%%%%%%%%%%%%%%%%%%%%%%%%%%%%%%%%%%%%%%%%%%%%%%%%%%%%%%%

Dually there exists a universal collapse.

\begin{definition}[Continuation counit]
The natural transformation

\[
\varepsilon:
FG
\Longrightarrow
\operatorname{Id}_{\mathcal D}
\]

is called the continuation counit.
\end{definition}

The counit records the universal reduction of a generated object back to its
essential continuation.

%%%%%%%%%%%%%%%%%%%%%%%%%%%%%%%%%%%%%%%%%%%%%%%%%%%%%%%%%%%%%%%%%%%%%%%%%%%%%%%
\section{The triangular identities}
%%%%%%%%%%%%%%%%%%%%%%%%%%%%%%%%%%%%%%%%%%%%%%%%%%%%%%%%%%%%%%%%%%%%%%%%%%%%%%%

The unit and counit cannot be chosen independently.

They satisfy two compatibility identities.

\begin{theorem}[Continuation Triangle Identities]

For every continuation adjunction,

\[
F
\xrightarrow{F\eta}
FGF
\xrightarrow{\varepsilon F}
F
\]

equals the identity on \(F\), while

\[
G
\xrightarrow{\eta G}
GFG
\xrightarrow{G\varepsilon}
G
\]

equals the identity on \(G\).

\end{theorem}

\begin{proof}

The proof is identical to the classical proof for adjoint functors and follows
from the defining natural bijection.

\end{proof}

These identities express that generation followed immediately by recovery
introduces no unnecessary continuation.

%%%%%%%%%%%%%%%%%%%%%%%%%%%%%%%%%%%%%%%%%%%%%%%%%%%%%%%%%%%%%%%%%%%%%%%%%%%%%%%
\section{Generation and recovery}
%%%%%%%%%%%%%%%%%%%%%%%%%%%%%%%%%%%%%%%%%%%%%%%%%%%%%%%%%%%%%%%%%%%%%%%%%%%%%%%

Continuation adjunctions formalize one of the central philosophical ideas of
continuation mathematics.

Generation creates new mathematical structure.

Recovery extracts only what is forced.

Neither operation is inverse to the other.

Instead they determine one another universally.

Adjunction therefore replaces equality by optimal correspondence.

%%%%%%%%%%%%%%%%%%%%%%%%%%%%%%%%%%%%%%%%%%%%%%%%%%%%%%%%%%%%%%%%%%%%%%%%%%%%%%%
\section{Free continuation constructions}
%%%%%%%%%%%%%%%%%%%%%%%%%%%%%%%%%%%%%%%%%%%%%%%%%%%%%%%%%%%%%%%%%%%%%%%%%%%%%%%

Every free continuation object naturally produces an adjunction.

Suppose

\[
U:\mathbf{ContAlg}
\rightarrow
\mathbf{ContSet}
\]

forgets algebraic structure.

Suppose

\[
F:\mathbf{ContSet}
\rightarrow
\mathbf{ContAlg}
\]

constructs free continuation algebras.

Then

\[
F\dashv U.
\]

Thus free continuation algebras arise universally rather than by arbitrary
construction.

%%%%%%%%%%%%%%%%%%%%%%%%%%%%%%%%%%%%%%%%%%%%%%%%%%%%%%%%%%%%%%%%%%%%%%%%%%%%%%%
\section{Completion adjunctions}
%%%%%%%%%%%%%%%%%%%%%%%%%%%%%%%%%%%%%%%%%%%%%%%%%%%%%%%%%%%%%%%%%%%%%%%%%%%%%%%

Many completion procedures are likewise adjoint.

Suppose

\[
C
\longmapsto
\widehat C
\]

assigns the canonical completion of every continuation space.

The inclusion

\[
i:
C
\hookrightarrow
\widehat C
\]

is frequently the unit of an adjunction.

Completion is therefore characterized not merely by adding missing points but
by satisfying a universal continuation property.

%%%%%%%%%%%%%%%%%%%%%%%%%%%%%%%%%%%%%%%%%%%%%%%%%%%%%%%%%%%%%%%%%%%%%%%%%%%%%%%
\section{Adjunction as optimal continuation}
%%%%%%%%%%%%%%%%%%%%%%%%%%%%%%%%%%%%%%%%%%%%%%%%%%%%%%%%%%%%%%%%%%%%%%%%%%%%%%%

Adjunction replaces exact inversion by optimal approximation.

Instead of demanding

\[
GF=\operatorname{Id},
\]

continuation mathematics asks for the strongest correspondence compatible with
the continuation structure.

This distinction is fundamental.

Most mathematical constructions are not reversible.

Many are nevertheless adjoint.

%%%%%%%%%%%%%%%%%%%%%%%%%%%%%%%%%%%%%%%%%%%%%%%%%%%%%%%%%%%%%%%%%%%%%%%%%%%%%%%
\section{Continuation reflection}
%%%%%%%%%%%%%%%%%%%%%%%%%%%%%%%%%%%%%%%%%%%%%%%%%%%%%%%%%%%%%%%%%%%%%%%%%%%%%%%

Certain continuation systems admit canonical reflections.

A reflection identifies the closest object lying inside a distinguished
subcategory while preserving all forced continuations.

Dually, coreflections identify the largest object compatible with a prescribed
restriction.

Both arise naturally from continuation adjunctions.

%%%%%%%%%%%%%%%%%%%%%%%%%%%%%%%%%%%%%%%%%%%%%%%%%%%%%%%%%%%%%%%%%%%%%%%%%%%%%%%
\section{Toward completion algebras}
%%%%%%%%%%%%%%%%%%%%%%%%%%%%%%%%%%%%%%%%%%%%%%%%%%%%%%%%%%%%%%%%%%%%%%%%%%%%%%%

Continuation algebra now contains the principal universal structures of modern
algebra.

\[
\begin{array}{c}
\text{operations},\\
\text{monoids},\\
\text{semigroups},\\
\text{lattices},\\
\text{categories},\\
\text{functors},\\
\text{natural transformations},\\
\text{limits},\\
\text{colimits},\\
\text{adjunctions}.
\end{array}
\]

These constructions describe how continuation systems are generated, compared,
extended, recovered, and universally related.

One final question remains.

Can an algebra itself possess an intrinsic notion of completion, independent
of any surrounding category?

The answer produces the final object of Part IV: the **Continuation Completion
Algebra**, where algebraic continuation reaches its own internal fixed point.

%%%%%%%%%%%%%%%%%%%%%%%%%%%%%%%%%%%%%%%%%%%%%%%%%%%%%%%%%%%%%%%%%%%%%%%%%%%%%%
\chapter{Continuation Completion Algebras}
%%%%%%%%%%%%%%%%%%%%%%%%%%%%%%%%%%%%%%%%%%%%%%%%%%%%%%%%%%%%%%%%%%%%%%%%%%%%%%

The preceding chapters developed the algebraic structures that arise from
continuation.

Operations permit continuation.

Monoids organize continuation.

Categories compare continuation.

Adjunctions relate generation and recovery.

A deeper question now presents itself.

Can an entire continuation algebra itself become complete?

This chapter answers that question.

%%%%%%%%%%%%%%%%%%%%%%%%%%%%%%%%%%%%%%%%%%%%%%%%%%%%%%%%%%%%%%%%%%%%%%%%%%%%%%%
\section{Algebraic completion}
%%%%%%%%%%%%%%%%%%%%%%%%%%%%%%%%%%%%%%%%%%%%%%%%%%%%%%%%%%%%%%%%%%%%%%%%%%%%%%%

Every algebra evolves by admitting progressively richer operations.

Initially only primitive operations exist.

Subsequently compositions become admissible.

Eventually identities, universal constructions, and higher-order operations
appear.

The process naturally suggests the existence of a terminal algebraic state in
which no genuinely new continuation operation can be added.

%%%%%%%%%%%%%%%%%%%%%%%%%%%%%%%%%%%%%%%%%%%%%%%%%%%%%%%%%%%%%%%%%%%%%%%%%%%%%%%
\section{Continuation completion algebras}
%%%%%%%%%%%%%%%%%%%%%%%%%%%%%%%%%%%%%%%%%%%%%%%%%%%%%%%%%%%%%%%%%%%%%%%%%%%%%%%

\begin{definition}[Continuation Completion Algebra]
A \emph{Continuation Completion Algebra} is a continuation algebra
\(\mathcal A\) satisfying the following properties.

\begin{enumerate}

\item
\textbf{Closure.}
Every admissible continuation operation on \(\mathcal A\) already belongs to
\(\mathcal A\).

\item
\textbf{Universality.}
Every continuation morphism into another complete continuation algebra factors
uniquely through \(\mathcal A\).

\item
\textbf{Maximality.}
No proper continuation extension of \(\mathcal A\) preserves the same
continuation structure.

\end{enumerate}

\end{definition}

Thus a continuation completion algebra is algebraically complete relative to
the continuation relation.

%%%%%%%%%%%%%%%%%%%%%%%%%%%%%%%%%%%%%%%%%%%%%%%%%%%%%%%%%%%%%%%%%%%%%%%%%%%%%%%
\section{Intrinsic completion}
%%%%%%%%%%%%%%%%%%%%%%%%%%%%%%%%%%%%%%%%%%%%%%%%%%%%%%%%%%%%%%%%%%%%%%%%%%%%%%%

Unlike classical algebraic completion, continuation completion is intrinsic.

No ambient algebra is required.

No external topology is assumed.

No completion process is imposed.

Completion is determined entirely by the continuation structure generated
internally.

%%%%%%%%%%%%%%%%%%%%%%%%%%%%%%%%%%%%%%%%%%%%%%%%%%%%%%%%%%%%%%%%%%%%%%%%%%%%%%%
\section{Stabilization}
%%%%%%%%%%%%%%%%%%%%%%%%%%%%%%%%%%%%%%%%%%%%%%%%%%%%%%%%%%%%%%%%%%%%%%%%%%%%%%%

\begin{definition}[Algebraic stabilization]
A continuation algebra is said to stabilize if every ascending continuation
chain

\[
A_1
\subseteq
A_2
\subseteq
A_3
\subseteq
\cdots
\]

eventually reaches a stage after which every subsequent continuation is
equivalent to one already present.
\end{definition}

Stabilization expresses the exhaustion of algebraic novelty.

%%%%%%%%%%%%%%%%%%%%%%%%%%%%%%%%%%%%%%%%%%%%%%%%%%%%%%%%%%%%%%%%%%%%%%%%%%%%%%%
\section{The stabilization theorem}
%%%%%%%%%%%%%%%%%%%%%%%%%%%%%%%%%%%%%%%%%%%%%%%%%%%%%%%%%%%%%%%%%%%%%%%%%%%%%%%

\begin{theorem}[Stabilization Criterion]

A continuation algebra is complete if and only if every admissible continuation
chain stabilizes.

\end{theorem}

\begin{proof}

Suppose first that the algebra is complete.

Any strictly increasing continuation chain would produce a new admissible
operation outside the algebra, contradicting closure.

Conversely, suppose every continuation chain stabilizes.

Then no new admissible continuation operation can be generated.

Hence the algebra already contains every continuation operation forced by its
own structure.

Therefore it is complete.

\end{proof}

The theorem identifies completion with exhaustion of admissible continuation.

%%%%%%%%%%%%%%%%%%%%%%%%%%%%%%%%%%%%%%%%%%%%%%%%%%%%%%%%%%%%%%%%%%%%%%%%%%%%%%%
\section{Completion rank}
%%%%%%%%%%%%%%%%%%%%%%%%%%%%%%%%%%%%%%%%%%%%%%%%%%%%%%%%%%%%%%%%%%%%%%%%%%%%%%%

The approach toward completion can itself be measured.

\begin{definition}[Completion rank]

The \emph{completion rank} of a continuation algebra is the least ordinal
\(\rho\) such that after \(\rho\) successive continuation extensions the
algebra stabilizes.

\end{definition}

Finite completion ranks describe rapidly stabilizing algebras.

Infinite ranks measure deeper continuation complexity.

Thus completion possesses an intrinsic quantitative invariant.

%%%%%%%%%%%%%%%%%%%%%%%%%%%%%%%%%%%%%%%%%%%%%%%%%%%%%%%%%%%%%%%%%%%%%%%%%%%%%%%
\section{Canonical completion}
%%%%%%%%%%%%%%%%%%%%%%%%%%%%%%%%%%%%%%%%%%%%%%%%%%%%%%%%%%%%%%%%%%%%%%%%%%%%%%%

\begin{theorem}[Existence of Canonical Completion]

Every continuation algebra admits, whenever stabilization occurs, a unique
completion algebra up to continuation isomorphism.

\end{theorem}

\begin{proof}

Take the directed system of all admissible continuation extensions.

Whenever stabilization occurs, its direct continuation colimit satisfies the
axioms of a continuation completion algebra.

Uniqueness follows from the universal property.

\end{proof}

%%%%%%%%%%%%%%%%%%%%%%%%%%%%%%%%%%%%%%%%%%%%%%%%%%%%%%%%%%%%%%%%%%%%%%%%%%%%%%%
\section{Absolute completion}
%%%%%%%%%%%%%%%%%%%%%%%%%%%%%%%%%%%%%%%%%%%%%%%%%%%%%%%%%%%%%%%%%%%%%%%%%%%%%%%

Completion algebras reveal an important distinction.

Some algebras become complete only relative to a chosen collection of
continuation operations.

Others remain complete under every admissible enlargement.

\begin{definition}[Absolute completion]

A continuation completion algebra is called \emph{absolute} if every admissible
continuation operation already belongs to the algebra independently of the
ambient continuation universe.

\end{definition}

Absolute completion is the strongest algebraic form of completion.

%%%%%%%%%%%%%%%%%%%%%%%%%%%%%%%%%%%%%%%%%%%%%%%%%%%%%%%%%%%%%%%%%%%%%%%%%%%%%%%
\section{The algebra of continuation}
%%%%%%%%%%%%%%%%%%%%%%%%%%%%%%%%%%%%%%%%%%%%%%%%%%%%%%%%%%%%%%%%%%%%%%%%%%%%%%%

The preceding chapters collectively establish the internal algebra of
continuation.

\[
\begin{array}{cccc}
\text{operations}
&
\longrightarrow
&
\text{monoids}
&
\\
&
&
\downarrow
&
\\
&
&
\text{semigroups}
&
\\
&
&
\downarrow
&
\\
&
&
\text{lattices}
&
\\
&
&
\downarrow
&
\\
&
&
\text{categories}
&
\\
&
&
\downarrow
&
\\
&
&
\text{functors}
&
\\
&
&
\downarrow
&
\\
&
&
\text{natural transformations}
&
\\
&
&
\downarrow
&
\\
&
&
\text{limits}
&
\\
&
&
\downarrow
&
\\
&
&
\text{colimits}
&
\\
&
&
\downarrow
&
\\
&
&
\text{adjunctions}
&
\\
&
&
\downarrow
&
\\
&
&
\boxed{\text{completion algebras}}
\end{array}
\]

Every layer is forced by the preceding one.

%%%%%%%%%%%%%%%%%%%%%%%%%%%%%%%%%%%%%%%%%%%%%%%%%%%%%%%%%%%%%%%%%%%%%%%%%%%%%%%
\section{Toward continuation geometry}
%%%%%%%%%%%%%%%%%%%%%%%%%%%%%%%%%%%%%%%%%%%%%%%%%%%%%%%%%%%%%%%%%%%%%%%%%%%%%%%

Part IV has developed the algebra generated by continuation.

The next question is no longer algebraic.

Instead we ask:

How does continuation organize space itself?

If algebra describes admissible operations,

geometry describes admissible organization.

The passage from algebra to geometry is therefore not a change of subject.

It is the emergence of global continuation structure from local continuation
relations.

Accordingly, the next part develops the foundations of
\emph{Continuation Geometry}, where neighborhoods, boundaries, dimension,
curvature, and ultimately mathematical space itself are reconstructed from the
continuation relation rather than assumed a priori.


\part{Continuation Geometry}

\chapter{Continuation Geometry}

The preceding Parts developed continuation as an algebraic phenomenon. Objects possess continuation systems; continuation systems possess algebraic operations; algebraic operations organize into categories and admit universal constructions.

A different viewpoint now becomes unavoidable.

Every continuation process determines not merely an algebraic object but an entire geometric landscape. At every stage there exists a collection of admissible continuations together with constraints that determine how those continuations may evolve. These collections possess neighborhoods, boundaries, dimensions, completion points, and large-scale organization.

Continuation therefore has an intrinsic geometry.

The purpose of this Part is to construct that geometry from first principles.

Unlike classical geometry, which begins from points together with externally imposed notions of distance or incidence, continuation geometry begins from admissibility itself. Geometry is recovered from the organization of possible futures.

Consequently no metric is assumed.

No topology is postulated.

No coordinate system is introduced.

Instead, each of these notions will emerge from continuation structure.

%%%%%%%%%%%%%%%%%%%%%%%%%%%%%%%%%%%%%%%%%%%%%%%%%%%%%%%%%%%%%%%%%%%%%

\section{Geometric Emergence}

Suppose $\mathcal C$ is a continuation system.

Every state $x\in\mathcal C$ determines its continuation space

\[
\operatorname{Cont}(x).
\]

These continuation spaces are not isolated.

Whenever

\[
x\rightsquigarrow y,
\]

the continuation space of $y$ is related to that of $x$ through the continuation morphisms developed in Part II.

The collection of all continuation spaces therefore possesses an intrinsic organization.

This organization is the geometry of the continuation system.

\begin{definition}[Continuation Geometry]
The \emph{continuation geometry} of a continuation system is the geometric structure induced by the interaction of all continuation spaces under admissible continuation.
\end{definition}

Thus geometry is not attached to points.

Geometry is attached to possible futures.

%%%%%%%%%%%%%%%%%%%%%%%%%%%%%%%%%%%%%%%%%%%%%%%%%%%%%%%%%%%%%%%%%%%%%

\section{Geometric Neighborhoods}

Classically, neighborhoods are defined through metrics or open sets.

Continuation mathematics reverses this order.

Neighborhoods arise because two states possess similar futures.

\begin{definition}[Geometric Neighborhood]
Let $x,y\in\mathcal C$.

We say that $y$ lies in the continuation neighborhood of $x$ whenever the continuation structures of $x$ and $y$ admit sufficiently large common continuation.

Symbolically,

\[
y\in N(x).
\]

\end{definition}

Thus proximity is determined by continuation similarity rather than external distance.

Future behavior replaces spatial location.

%%%%%%%%%%%%%%%%%%%%%%%%%%%%%%%%%%%%%%%%%%%%%%%%%%%%%%%%%%%%%%%%%%%%%

\section{Continuation Coordinates}

Coordinates are usually introduced before geometry.

Continuation mathematics proceeds oppositely.

Coordinates are recovered from continuation structure.

\begin{definition}[Continuation Coordinate]
A continuation coordinate is any invariant obtained from continuation behavior that uniquely identifies a state within a specified continuation geometry.
\end{definition}

Examples include

\[
\dim(x),
\]

completion depth,

frontier depth,

continuation entropy,

branching degree,

or any complete continuation invariant.

Coordinates therefore arise from dynamics.

%%%%%%%%%%%%%%%%%%%%%%%%%%%%%%%%%%%%%%%%%%%%%%%%%%%%%%%%%%%%%%%%%%%%%

\section{Continuation Paths}

Every continuation process determines a path.

\begin{definition}
A continuation path is a sequence

\[
x_0
\rightsquigarrow
x_1
\rightsquigarrow
x_2
\rightsquigarrow
\cdots
\]

such that each step is admissible.
\end{definition}

These paths play the role occupied by curves in classical geometry.

Instead of parameterizing motion through Euclidean space, they parameterize motion through continuation space.

%%%%%%%%%%%%%%%%%%%%%%%%%%%%%%%%%%%%%%%%%%%%%%%%%%%%%%%%%%%%%%%%%%%%%

\section{Geometric Objects}

Every continuation path traces a geometric object.

Finite paths determine finite geometric figures.

Infinite paths determine geometric rays.

Closed paths determine continuation cycles.

Completion paths determine terminal components.

Thus every geometric object is generated dynamically.

Geometry is no longer static.

Geometry is the visible trace of continuation.

%%%%%%%%%%%%%%%%%%%%%%%%%%%%%%%%%%%%%%%%%%%%%%%%%%%%%%%%%%%%%%%%%%%%%

\section{Intrinsic Geometry}

The geometry constructed here is intrinsic.

No embedding space is required.

No ambient manifold exists.

Everything is determined internally by continuation.

This parallels the transition from extrinsic to intrinsic differential geometry, but is considerably more general.

The surrounding universe plays no mathematical role.

%%%%%%%%%%%%%%%%%%%%%%%%%%%%%%%%%%%%%%%%%%%%%%%%%%%%%%%%%%%%%%%%%%%%%

\section{The Fundamental Principle of Continuation Geometry}

The previous developments force the following constitutional statement.

\begin{theorem}[Fundamental Principle of Continuation Geometry]
Every continuation system determines a unique intrinsic geometry.

Conversely, every geometric notion admitted by continuation mathematics must be recoverable from continuation structure alone.
\end{theorem}

\begin{proof}
Continuation spaces determine neighborhoods.

Neighborhoods determine local geometry.

Continuation paths determine geometric figures.

Completion determines global organization.

Every geometric construction therefore arises from continuation structure itself.

No additional geometric primitives are required.
\end{proof}

%%%%%%%%%%%%%%%%%%%%%%%%%%%%%%%%%%%%%%%%%%%%%%%%%%%%%%%%%%%%%%%%%%%%%

\section{Outlook}

The remainder of this Part develops the principal objects of continuation geometry.

Beginning only with admissibility and continuation, we shall derive

\[
\text{metrics},
\]

\[
\text{distance},
\]

\[
\text{geodesics},
\]

\[
\text{curvature},
\]

\[
\text{volume},
\]

\[
\text{entropy},
\]

and

\[
\text{rigidity},
\]

without assuming any of them in advance.

Geometry will appear not as an independent discipline but as a necessary consequence of continuation itself.

\chapter{Continuation Metrics}

The first problem of continuation geometry is the recovery of distance.

In classical mathematics, a metric is introduced as an external primitive satisfying a collection of axioms. Geometry is then developed relative to that metric.

Continuation mathematics proceeds in the opposite direction.

Continuation already determines which futures are possible, which futures are forbidden, and how different continuation systems evolve. Distance must therefore emerge from these continuation relations themselves.

The purpose of this chapter is to derive the concept of metric directly from continuation structure.

%%%%%%%%%%%%%%%%%%%%%%%%%%%%%%%%%%%%%%%%%%%%%%%%%%%%%%%%%%%%%%%%%%%%%

\section{The insufficiency of classical metrics}

A classical metric assigns to every pair of points a non-negative real number

\[
d(x,y).
\]

The assignment is assumed to satisfy positivity, symmetry, and the triangle inequality.

Nothing in these axioms explains why two objects should be close.

The metric merely records that they are.

Continuation mathematics rejects this starting point.

Distance must arise from structural continuation.

Objects are close precisely because their possible futures possess structural agreement.

%%%%%%%%%%%%%%%%%%%%%%%%%%%%%%%%%%%%%%%%%%%%%%%%%%%%%%%%%%%%%%%%%%%%%

\section{Continuation similarity}

The primitive geometric quantity is not distance but similarity of continuation.

\begin{definition}[Continuation Similarity]
Let $x$ and $y$ be states of a continuation system.

The continuation similarity of $x$ and $y$ is the degree to which their admissible continuation structures coincide.
\end{definition}

Similarity therefore measures future agreement rather than present position.

If

\[
\operatorname{Cont}(x)
=
\operatorname{Cont}(y),
\]

then the two states possess identical immediate continuation structure.

If the continuation spaces differ substantially, the similarity decreases.

Distance will eventually be recovered from similarity.

%%%%%%%%%%%%%%%%%%%%%%%%%%%%%%%%%%%%%%%%%%%%%%%%%%%%%%%%%%%%%%%%%%%%%

\section{Continuation equivalence levels}

Continuation agreement naturally possesses depth.

Two states may agree for one continuation step while diverging later.

Alternatively they may agree indefinitely.

\begin{definition}[Continuation depth]
The continuation depth of two states is the largest integer

\[
d\ge0
\]

such that every admissible continuation of length at most $d$ is identical for both states.
\end{definition}

Depth therefore measures the persistence of common future behavior.

Greater depth corresponds to greater geometric proximity.

%%%%%%%%%%%%%%%%%%%%%%%%%%%%%%%%%%%%%%%%%%%%%%%%%%%%%%%%%%%%%%%%%%%%%

\section{Derived distance}

Similarity determines distance.

\begin{definition}[Continuation metric]
A continuation metric is any function

\[
D(x,y)
\]

whose value depends only upon continuation similarity and decreases monotonically as continuation depth increases.
\end{definition}

Many such metrics may exist.

The geometry itself does not depend upon a particular numerical realization.

Instead, continuation determines an entire family of compatible metrics.

%%%%%%%%%%%%%%%%%%%%%%%%%%%%%%%%%%%%%%%%%%%%%%%%%%%%%%%%%%%%%%%%%%%%%

\section{Canonical metrics}

Among all continuation metrics, some are determined uniquely by the continuation system itself.

\begin{definition}[Canonical continuation metric]
A continuation metric is canonical whenever it is completely recoverable from continuation structure and requires no external normalization.
\end{definition}

Canonical metrics depend solely upon admissibility.

No coordinates appear.

No embedding space appears.

No arbitrary scaling constants appear.

%%%%%%%%%%%%%%%%%%%%%%%%%%%%%%%%%%%%%%%%%%%%%%%%%%%%%%%%%%%%%%%%%%%%%

\section{Metric generation}

Continuation systems generate metrics recursively.

Suppose

\[
x
\rightsquigarrow
x'
\]

and

\[
y
\rightsquigarrow
y'.
\]

The relationship between

\[
D(x,y)
\]

and

\[
D(x',y')
\]

is determined entirely by the continuation morphisms introduced in Part II.

Distance therefore evolves together with continuation.

Metrics become dynamic objects.

%%%%%%%%%%%%%%%%%%%%%%%%%%%%%%%%%%%%%%%%%%%%%%%%%%%%%%%%%%%%%%%%%%%%%

\section{Metric compatibility}

Not every numerical function qualifies as a continuation metric.

\begin{definition}[Metric compatibility]
A continuation metric is compatible whenever continuation-equivalent states have zero distance and every continuation morphism preserves the induced geometric structure.
\end{definition}

Compatibility guarantees that geometry reflects continuation rather than arbitrary numerical choices.

%%%%%%%%%%%%%%%%%%%%%%%%%%%%%%%%%%%%%%%%%%%%%%%%%%%%%%%%%%%%%%%%%%%%%

\section{Completion distance}

Completion spaces provide another source of metric information.

Suppose

\[
x
\rightsquigarrow
\omega .
\]

The complexity of reaching completion determines another intrinsic measure.

\begin{definition}[Completion distance]
The completion distance of a state is the minimal continuation complexity required to reach a completion object.
\end{definition}

Completion distance is entirely intrinsic.

No coordinates are involved.

Only continuation complexity matters.

%%%%%%%%%%%%%%%%%%%%%%%%%%%%%%%%%%%%%%%%%%%%%%%%%%%%%%%%%%%%%%%%%%%%%

\section{Metric rigidity}

Continuation metrics satisfy an important rigidity principle.

\begin{theorem}[Metric Rigidity]
If two continuation systems are continuation-isomorphic, then every canonical continuation metric is preserved.
\end{theorem}

\begin{proof}
Continuation isomorphisms preserve admissibility, continuation depth, completion structure, neighborhoods, and continuation morphisms.

Every canonical metric is recovered solely from these structures.

Therefore canonical metrics are invariant under continuation isomorphism.
\end{proof}

%%%%%%%%%%%%%%%%%%%%%%%%%%%%%%%%%%%%%%%%%%%%%%%%%%%%%%%%%%%%%%%%%%%%%

\section{The emergence of geometry}

The previous Parts recovered algebra from continuation.

The present chapter recovers distance.

Neither algebra nor geometry has been assumed.

Both have emerged from continuation itself.

The remainder of this Part develops increasingly refined geometric invariants.

Metrics give rise to shortest continuation paths.

Shortest paths give rise to geodesics.

The behavior of geodesics reveals curvature.

Curvature will ultimately measure the intrinsic obstruction to continuation itself.

Thus the next stage of continuation geometry is forced.

\chapter{Continuation Geodesics}

The existence of a continuation metric immediately forces the existence of preferred continuation paths.

In classical geometry, geodesics are shortest paths with respect to a prescribed metric. Their definition therefore depends upon geometric structure that has already been assumed.

Continuation mathematics reverses this dependency.

Continuation paths exist before metrics.

Metrics merely quantify the complexity of those paths.

Consequently geodesics are not primitive geometric objects.

They are distinguished continuation processes.

%%%%%%%%%%%%%%%%%%%%%%%%%%%%%%%%%%%%%%%%%%%%%%%%%%%%%%%%%%%%%%%%%%%%%

\section{Continuation paths}

Every continuation system possesses admissible paths.

\begin{definition}[Continuation Path]
A continuation path is a finite or infinite sequence

\[
x_0
\rightsquigarrow
x_1
\rightsquigarrow
x_2
\rightsquigarrow
\cdots
\]

such that every transition is an admissible continuation.
\end{definition}

The collection of all continuation paths forms the dynamical skeleton of the continuation geometry.

%%%%%%%%%%%%%%%%%%%%%%%%%%%%%%%%%%%%%%%%%%%%%%%%%%%%%%%%%%%%%%%%%%%%%

\section{Path complexity}

Not all continuation paths are equally efficient.

Some reach a destination after few continuation steps.

Others require many intermediate continuations.

Some continually branch before eventually returning.

The intrinsic complexity of a path therefore measures the continuation cost of reaching its terminal state.

\begin{definition}[Path Complexity]

The continuation complexity of a path is the total continuation cost accumulated along the path.

\end{definition}

The precise form of this cost depends upon the continuation metric introduced in the previous chapter.

%%%%%%%%%%%%%%%%%%%%%%%%%%%%%%%%%%%%%%%%%%%%%%%%%%%%%%%%%%%%%%%%%%%%%

\section{Minimal continuation}

Among all continuation paths joining two states, some possess least continuation complexity.

\begin{definition}[Minimal Continuation]

A continuation path is minimal if no admissible continuation joining the same endpoints possesses smaller continuation complexity.

\end{definition}

Minimality replaces shortest distance as the primitive concept.

Distance becomes a numerical expression of minimal continuation.

%%%%%%%%%%%%%%%%%%%%%%%%%%%%%%%%%%%%%%%%%%%%%%%%%%%%%%%%%%%%%%%%%%%%%

\section{Continuation geodesics}

\begin{definition}[Continuation Geodesic]

A continuation geodesic is a minimal continuation path between two continuation states.

\end{definition}

Thus geodesics are generated by admissibility.

They are not externally imposed curves.

%%%%%%%%%%%%%%%%%%%%%%%%%%%%%%%%%%%%%%%%%%%%%%%%%%%%%%%%%%%%%%%%%%%%%

\section{Local geodesics}

Global minimality may fail while local minimality remains.

\begin{definition}[Local Geodesic]

A continuation path is locally geodesic if every sufficiently small subpath is minimal.

\end{definition}

Local continuation optimality therefore precedes global optimality.

%%%%%%%%%%%%%%%%%%%%%%%%%%%%%%%%%%%%%%%%%%%%%%%%%%%%%%%%%%%%%%%%%%%%%

\section{Uniqueness}

Continuation systems naturally divide into two classes.

Some admit unique geodesics.

Others admit multiple competing geodesics.

\begin{definition}

A continuation system is geodesically rigid if every pair of continuation states is joined by at most one continuation geodesic.

\end{definition}

Failure of rigidity reflects branching continuation.

%%%%%%%%%%%%%%%%%%%%%%%%%%%%%%%%%%%%%%%%%%%%%%%%%%%%%%%%%%%%%%%%%%%%%

\section{Geodesic branching}

Branching has no analogue in classical Euclidean geometry.

Continuation geometry permits one continuation process to split naturally into several equally admissible minimal continuations.

\begin{definition}

A branching geodesic is a geodesic admitting more than one minimal continuation extension.

\end{definition}

Branching is therefore an intrinsic geometric phenomenon.

%%%%%%%%%%%%%%%%%%%%%%%%%%%%%%%%%%%%%%%%%%%%%%%%%%%%%%%%%%%%%%%%%%%%%

\section{Completion geodesics}

Completion objects occupy a distinguished geometric role.

\begin{definition}

A completion geodesic is a minimal continuation path terminating at a completion object.

\end{definition}

These paths describe the intrinsic flow of the continuation system toward completion.

%%%%%%%%%%%%%%%%%%%%%%%%%%%%%%%%%%%%%%%%%%%%%%%%%%%%%%%%%%%%%%%%%%%%%

\section{Geodesic invariance}

\begin{theorem}[Geodesic Invariance]

Continuation isomorphisms preserve continuation geodesics.

\end{theorem}

\begin{proof}

Continuation isomorphisms preserve admissibility.

They preserve continuation metrics.

Therefore they preserve minimal continuation paths.

Hence they preserve geodesics.

\end{proof}

%%%%%%%%%%%%%%%%%%%%%%%%%%%%%%%%%%%%%%%%%%%%%%%%%%%%%%%%%%%%%%%%%%%%%

\section{Toward curvature}

Geodesics reveal the large-scale organization of continuation geometry.

If nearby geodesics remain close, the continuation space behaves regularly.

If they diverge rapidly, continuation possesses intrinsic instability.

If they repeatedly converge, the continuation system exhibits geometric attraction.

Thus curvature is not an independent primitive.

Curvature measures the behavior of families of continuation geodesics.

The next chapter develops this principle.

\chapter{Continuation Curvature}

The preceding chapter established that every continuation system possesses distinguished paths, namely the continuation geodesics.

The next question is unavoidable.

How do nearby continuation geodesics behave?

In classical geometry this behavior is measured by curvature.

Continuation mathematics derives an analogous concept directly from admissible continuation.

Curvature therefore measures not the bending of space but the bending of possibility.

%%%%%%%%%%%%%%%%%%%%%%%%%%%%%%%%%%%%%%%%%%%%%%%%%%%%%%%%%%%%%%%%%%%%%

\section{The insufficiency of classical curvature}

Classical curvature depends upon differentiability.

One assumes coordinates.

One constructs tangent spaces.

One introduces derivatives.

One then measures the variation of tangent directions.

Continuation systems possess none of these primitives.

They possess only admissible continuation.

Consequently curvature must arise from continuation itself.

%%%%%%%%%%%%%%%%%%%%%%%%%%%%%%%%%%%%%%%%%%%%%%%%%%%%%%%%%%%%%%%%%%%%%

\section{Geodesic divergence}

Suppose two continuation geodesics begin at the same state,

\[
\gamma_1,\gamma_2:
x
\rightsquigarrow
\cdots
\]

Initially the two continuations may coincide.

Eventually they may separate.

Alternatively they may repeatedly return toward one another.

Their relative behavior is intrinsic.

\begin{definition}[Geodesic Divergence]

The divergence of two continuation geodesics is the rate at which their continuation distance changes under simultaneous continuation.

\end{definition}

Divergence is therefore a purely continuation-theoretic quantity.

%%%%%%%%%%%%%%%%%%%%%%%%%%%%%%%%%%%%%%%%%%%%%%%%%%%%%%%%%%%%%%%%%%%%%

\section{Flat continuation}

The simplest possibility is complete stability.

\begin{definition}[Flat Continuation]

A continuation region is flat whenever parallel continuation geodesics preserve their mutual continuation distance.

\end{definition}

Flatness therefore means that continuation introduces no additional structural distortion.

Future evolution behaves uniformly.

%%%%%%%%%%%%%%%%%%%%%%%%%%%%%%%%%%%%%%%%%%%%%%%%%%%%%%%%%%%%%%%%%%%%%

\section{Positive continuation curvature}

Continuation may instead possess an attracting tendency.

\begin{definition}[Positive Continuation Curvature]

A continuation region has positive continuation curvature whenever nearby continuation geodesics systematically approach one another under continuation.

\end{definition}

Positive curvature therefore measures intrinsic continuation attraction.

Distinct futures become progressively less distinguishable.

%%%%%%%%%%%%%%%%%%%%%%%%%%%%%%%%%%%%%%%%%%%%%%%%%%%%%%%%%%%%%%%%%%%%%

\section{Negative continuation curvature}

The opposite phenomenon is equally important.

\begin{definition}[Negative Continuation Curvature]

A continuation region possesses negative continuation curvature whenever nearby continuation geodesics systematically separate under continuation.

\end{definition}

Negative curvature therefore measures intrinsic continuation instability.

Small continuation differences produce increasingly different futures.

%%%%%%%%%%%%%%%%%%%%%%%%%%%%%%%%%%%%%%%%%%%%%%%%%%%%%%%%%%%%%%%%%%%%%

\section{Continuation singularities}

Not every continuation process survives indefinitely.

Some continuations terminate abruptly.

Others become impossible after finite continuation.

These produce geometric singularities.

\begin{definition}[Continuation Singularity]

A continuation singularity is a point at which every admissible continuation is obstructed.

\end{definition}

Singularities represent complete geometric obstruction.

They form the continuation analogue of boundaries beyond which continuation cannot proceed.

%%%%%%%%%%%%%%%%%%%%%%%%%%%%%%%%%%%%%%%%%%%%%%%%%%%%%%%%%%%%%%%%%%%%%

\section{Curvature and branching}

Continuation curvature is intimately related to branching.

Suppose

\[
|\operatorname{Cont}(x)|=1.
\]

Continuation is locally deterministic.

Suppose instead

\[
|\operatorname{Cont}(x)|\gg1.
\]

Many futures coexist.

Large branching tends to generate negative continuation curvature, while repeated merging tends to generate positive continuation curvature.

Thus branching is the combinatorial source of curvature.

%%%%%%%%%%%%%%%%%%%%%%%%%%%%%%%%%%%%%%%%%%%%%%%%%%%%%%%%%%%%%%%%%%%%%

\section{Curvature as obstruction}

Curvature measures more than geometric deformation.

It measures continuation obstruction.

Whenever admissible continuation becomes increasingly constrained, continuation curvature increases.

Whenever admissible continuation becomes increasingly unconstrained, curvature decreases.

Thus curvature quantifies the local resistance of continuation.

%%%%%%%%%%%%%%%%%%%%%%%%%%%%%%%%%%%%%%%%%%%%%%%%%%%%%%%%%%%%%%%%%%%%%

\section{Intrinsic character}

No coordinates appear in the preceding definitions.

No embedding manifold is assumed.

No differentiable structure exists.

Curvature is determined entirely by continuation behavior.

Consequently continuation curvature is intrinsic.

%%%%%%%%%%%%%%%%%%%%%%%%%%%%%%%%%%%%%%%%%%%%%%%%%%%%%%%%%%%%%%%%%%%%%

\section{Curvature invariance}

\begin{theorem}[Continuation Curvature Invariance]

Continuation isomorphisms preserve continuation curvature.

\end{theorem}

\begin{proof}

Continuation isomorphisms preserve admissibility.

They preserve continuation geodesics.

They preserve continuation metrics.

Hence they preserve geodesic divergence.

Since curvature is recovered solely from geodesic divergence, continuation curvature is invariant under continuation isomorphism.

\end{proof}

%%%%%%%%%%%%%%%%%%%%%%%%%%%%%%%%%%%%%%%%%%%%%%%%%%%%%%%%%%%%%%%%%%%%%

\section{The Geometry of Mathematical Difficulty}

Curvature supplies the first genuinely global geometric invariant of continuation mathematics.

Low curvature corresponds to stable continuation.

High curvature corresponds to structural obstruction.

Infinite curvature corresponds to impossible continuation.

Consequently mathematical problems themselves possess continuation geometry.

A theorem is easy whenever continuation remains geometrically flat.

A theorem is difficult whenever continuation encounters regions of large curvature.

A theorem becomes impossible precisely when continuation reaches intrinsic obstruction.

Thus continuation curvature provides a geometric measure of mathematical difficulty itself.

%%%%%%%%%%%%%%%%%%%%%%%%%%%%%%%%%%%%%%%%%%%%%%%%%%%%%%%%%%%%%%%%%%%%%

\section{Outlook}

Distance measures separation.

Geodesics measure optimal continuation.

Curvature measures deformation of continuation.

The next remaining global quantity is size.

Classical geometry measures size through volume.

Continuation mathematics must derive an intrinsic notion of volume directly from admissible continuation.

This will be the subject of the next chapter.

\chapter{Continuation Boundaries}

The concept of boundary is one of the oldest ideas in mathematics. In classical geometry, topology, and analysis, a boundary separates an object from its complement. It is defined relative to an ambient space and therefore depends upon an external notion of neighbourhood.

Continuation Mathematics proceeds differently.

Objects are generated by continuation. Consequently the natural notion of boundary is not separation from an exterior, but termination of admissible continuation.

Boundaries therefore arise internally.

%%%%%%%%%%%%%%%%%%%%%%%%%%%%%%%%%%%%%%%%%%%%%%%%%%%%%%%%%%%%%%%%%%%%%%%%%%%%%%
\section{Why classical boundaries are not primitive}
%%%%%%%%%%%%%%%%%%%%%%%%%%%%%%%%%%%%%%%%%%%%%%%%%%%%%%%%%%%%%%%%%%%%%%%%%%%%%%

Suppose an object has been generated by successive continuation.

At every stage there exists a collection of admissible continuations.

Eventually this collection may become smaller.

Some continuations disappear.

Others remain.

The point at which continuation can no longer proceed in a particular direction is what should properly be called a boundary.

Thus boundaries are generated by the continuation process itself.

No surrounding space is required.

%%%%%%%%%%%%%%%%%%%%%%%%%%%%%%%%%%%%%%%%%%%%%%%%%%%%%%%%%%%%%%%%%%%%%%%%%%%%%%
\section{Boundary points}
%%%%%%%%%%%%%%%%%%%%%%%%%%%%%%%%%%%%%%%%%%%%%%%%%%%%%%%%%%%%%%%%%%%%%%%%%%%%%%

Let $(X,\mathcal C)$ be a continuation space.

\begin{definition}[Boundary point]
A point
\[
x\in X
\]
is called a \emph{boundary point} if there exists at least one admissible continuation arriving at $x$, but every continuation extending beyond $x$ fails.
\end{definition}

Thus a boundary point is neither initial nor interior.

It is a maximal continuation point.

%%%%%%%%%%%%%%%%%%%%%%%%%%%%%%%%%%%%%%%%%%%%%%%%%%%%%%%%%%%%%%%%%%%%%%%%%%%%%%
\section{Boundary objects}
%%%%%%%%%%%%%%%%%%%%%%%%%%%%%%%%%%%%%%%%%%%%%%%%%%%%%%%%%%%%%%%%%%%%%%%%%%%%%%

Entire objects may themselves form boundaries.

\begin{definition}
A continuation subsystem
\[
B\subseteq X
\]
is called a \emph{boundary object} if every continuation entering $B$ terminates inside $B$.
\end{definition}

Boundary objects absorb continuation.

No admissible extension exists beyond them.

%%%%%%%%%%%%%%%%%%%%%%%%%%%%%%%%%%%%%%%%%%%%%%%%%%%%%%%%%%%%%%%%%%%%%%%%%%%%%%
\section{Frontiers and boundaries}
%%%%%%%%%%%%%%%%%%%%%%%%%%%%%%%%%%%%%%%%%%%%%%%%%%%%%%%%%%%%%%%%%%%%%%%%%%%%%%

Every boundary determines a frontier.

The converse need not hold.

A frontier merely separates distinct continuation behaviours.

A boundary completely terminates one.

Thus

\[
\text{Boundary}
\Longrightarrow
\text{Frontier}
\]

but not conversely.

%%%%%%%%%%%%%%%%%%%%%%%%%%%%%%%%%%%%%%%%%%%%%%%%%%%%%%%%%%%%%%%%%%%%%%%%%%%%%%
\section{Boundary rank}
%%%%%%%%%%%%%%%%%%%%%%%%%%%%%%%%%%%%%%%%%%%%%%%%%%%%%%%%%%%%%%%%%%%%%%%%%%%%%%

Not every obstruction is equally severe.

Some terminate only one continuation.

Others terminate infinitely many.

This motivates the following definition.

\begin{definition}
The \emph{boundary rank} of a boundary object is the number (finite or infinite) of independent continuation directions that terminate there.
\end{definition}

Boundary rank measures the strength of an obstruction.

%%%%%%%%%%%%%%%%%%%%%%%%%%%%%%%%%%%%%%%%%%%%%%%%%%%%%%%%%%%%%%%%%%%%%%%%%%%%%%
\section{Terminal boundaries}
%%%%%%%%%%%%%%%%%%%%%%%%%%%%%%%%%%%%%%%%%%%%%%%%%%%%%%%%%%%%%%%%%%%%%%%%%%%%%%

Some boundaries admit absolutely no continuation.

\begin{definition}
A boundary is called \emph{terminal} if every admissible continuation ending there is maximal.
\end{definition}

These represent genuine endpoints of continuation.

%%%%%%%%%%%%%%%%%%%%%%%%%%%%%%%%%%%%%%%%%%%%%%%%%%%%%%%%%%%%%%%%%%%%%%%%%%%%%%
\section{Boundary completion}
%%%%%%%%%%%%%%%%%%%%%%%%%%%%%%%%%%%%%%%%%%%%%%%%%%%%%%%%%%%%%%%%%%%%%%%%%%%%%%

Suppose continuation repeatedly approaches a boundary without reaching it.

If adjoining a single object completes every such continuation, that object is called the completion of the boundary.

Thus completion naturally fills missing boundaries.

Completion therefore appears as the dual notion to obstruction.

%%%%%%%%%%%%%%%%%%%%%%%%%%%%%%%%%%%%%%%%%%%%%%%%%%%%%%%%%%%%%%%%%%%%%%%%%%%%%%
\section{Boundary spectra}
%%%%%%%%%%%%%%%%%%%%%%%%%%%%%%%%%%%%%%%%%%%%%%%%%%%%%%%%%%%%%%%%%%%%%%%%%%%%%%

Boundaries need not occur individually.

They may organise into families.

\begin{definition}
The collection of all boundary objects of a continuation space is called its
\emph{boundary spectrum}.
\end{definition}

The spectrum records every possible obstruction.

%%%%%%%%%%%%%%%%%%%%%%%%%%%%%%%%%%%%%%%%%%%%%%%%%%%%%%%%%%%%%%%%%%%%%%%%%%%%%%
\section{Functorial behaviour}
%%%%%%%%%%%%%%%%%%%%%%%%%%%%%%%%%%%%%%%%%%%%%%%%%%%%%%%%%%%%%%%%%%%%%%%%%%%%%%

Continuation morphisms preserve terminality.

\begin{theorem}
Continuation morphisms map terminal boundaries to terminal boundaries.
\end{theorem}

\begin{proof}

A continuation morphism preserves admissible continuation.

If no continuation exists beyond a terminal boundary before applying the morphism, none can exist afterwards.

Therefore terminality is preserved.

\end{proof}

%%%%%%%%%%%%%%%%%%%%%%%%%%%%%%%%%%%%%%%%%%%%%%%%%%%%%%%%%%%%%%%%%%%%%%%%%%%%%%
\section{Boundary rigidity}
%%%%%%%%%%%%%%%%%%%%%%%%%%%%%%%%%%%%%%%%%%%%%%%%%%%%%%%%%%%%%%%%%%%%%%%%%%%%%%

The most important property of boundaries is rigidity.

\begin{theorem}[Boundary Rigidity]

Suppose every continuation reaching a boundary object terminates.

Then no admissible continuation structure can remove that boundary without enlarging the continuation universe.

\end{theorem}

\begin{proof}

Removing the boundary would require introducing new admissible continuations.

Those continuations were absent from the original continuation system.

Hence the resulting system is no longer the same continuation universe.

Therefore the boundary is intrinsic.

\end{proof}

%%%%%%%%%%%%%%%%%%%%%%%%%%%%%%%%%%%%%%%%%%%%%%%%%%%%%%%%%%%%%%%%%%%%%%%%%%%%%%
\section{The Boundary Principle}
%%%%%%%%%%%%%%%%%%%%%%%%%%%%%%%%%%%%%%%%%%%%%%%%%%%%%%%%%%%%%%%%%%%%%%%%%%%%%%

The previous results reveal the correct foundational interpretation.

Boundaries are not geometric surfaces.

They are intrinsic termination loci generated by continuation itself.

Geometry therefore appears only after the continuation structure has already determined where continuation necessarily stops.

In Continuation Mathematics, obstruction is primary.

Boundary is merely its geometric manifestation.

\chapter{Continuation Entropy}

Classical entropy measures uncertainty, disorder, complexity, or information.

Although these notions arise in different areas of mathematics, they share a common feature: they attempt to quantify the number of possible future developments of a system.

Continuation Mathematics identifies this common structure directly.

Entropy is not fundamentally uncertainty.

Entropy is continuation freedom.

%%%%%%%%%%%%%%%%%%%%%%%%%%%%%%%%%%%%%%%%%%%%%%%%%%%%%%%%%%%%%%%%%%%%%%%%%%%%%%
\section{The continuation viewpoint}
%%%%%%%%%%%%%%%%%%%%%%%%%%%%%%%%%%%%%%%%%%%%%%%%%%%%%%%%%%%%%%%%%%%%%%%%%%%%%%

Suppose an object admits several admissible continuations.

Each continuation represents one possible extension of the present object.

The larger the admissible continuation family, the greater the continuation freedom.

Conversely, if only one continuation exists, the future is completely determined.

If no continuation exists, the process terminates.

Thus entropy naturally measures the size of continuation space.

%%%%%%%%%%%%%%%%%%%%%%%%%%%%%%%%%%%%%%%%%%%%%%%%%%%%%%%%%%%%%%%%%%%%%%%%%%%%%%
\section{Continuation entropy}
%%%%%%%%%%%%%%%%%%%%%%%%%%%%%%%%%%%%%%%%%%%%%%%%%%%%%%%%%%%%%%%%%%%%%%%%%%%%%%

Let $X$ be a continuation space.

For each object $x\in X$, denote its admissible continuation set by

\[
\mathcal C(x).
\]

\begin{definition}[Local continuation entropy]

The \emph{local continuation entropy} of $x$ is

\[
H(x)
=
\Phi(|\mathcal C(x)|),
\]

where $\Phi$ is any strictly increasing normalization.

\end{definition}

Typical choices include

\[
\Phi(n)=n,
\]

or

\[
\Phi(n)=\log n.
\]

The theory itself is independent of normalization.

Entropy is fundamentally induced by continuation multiplicity.

%%%%%%%%%%%%%%%%%%%%%%%%%%%%%%%%%%%%%%%%%%%%%%%%%%%%%%%%%%%%%%%%%%%%%%%%%%%%%%
\section{Global continuation entropy}
%%%%%%%%%%%%%%%%%%%%%%%%%%%%%%%%%%%%%%%%%%%%%%%%%%%%%%%%%%%%%%%%%%%%%%%%%%%%%%

Local entropy measures freedom at one object.

Entire continuation spaces possess global entropy.

\begin{definition}

The \emph{global continuation entropy} of a continuation space is the aggregate continuation freedom over all admissible objects.

\end{definition}

Different aggregation procedures recover different classical entropy theories.

%%%%%%%%%%%%%%%%%%%%%%%%%%%%%%%%%%%%%%%%%%%%%%%%%%%%%%%%%%%%%%%%%%%%%%%%%%%%%%
\section{Entropy generation}
%%%%%%%%%%%%%%%%%%%%%%%%%%%%%%%%%%%%%%%%%%%%%%%%%%%%%%%%%%%%%%%%%%%%%%%%%%%%%%

Continuation may generate entropy.

Suppose

\[
x
\longrightarrow
x'
\]

is admissible.

If

\[
|\mathcal C(x')|
>
|\mathcal C(x)|,
\]

then entropy has increased.

If

\[
|\mathcal C(x')|
<
|\mathcal C(x)|,
\]

entropy has decreased.

Thus entropy evolution becomes an intrinsic geometric property.

%%%%%%%%%%%%%%%%%%%%%%%%%%%%%%%%%%%%%%%%%%%%%%%%%%%%%%%%%%%%%%%%%%%%%%%%%%%%%%
\section{Zero entropy}
%%%%%%%%%%%%%%%%%%%%%%%%%%%%%%%%%%%%%%%%%%%%%%%%%%%%%%%%%%%%%%%%%%%%%%%%%%%%%%

The smallest possible entropy occurs when continuation is unique.

\begin{definition}

A continuation object has zero entropy whenever

\[
|\mathcal C(x)|=1.
\]

\end{definition}

Such objects possess completely deterministic evolution.

No branching occurs.

%%%%%%%%%%%%%%%%%%%%%%%%%%%%%%%%%%%%%%%%%%%%%%%%%%%%%%%%%%%%%%%%%%%%%%%%%%%%%%
\section{Maximal entropy}
%%%%%%%%%%%%%%%%%%%%%%%%%%%%%%%%%%%%%%%%%%%%%%%%%%%%%%%%%%%%%%%%%%%%%%%%%%%%%%

A continuation object has maximal entropy whenever every admissible continuation remains possible.

Such objects possess maximal continuation freedom.

They represent regions of greatest structural flexibility.

%%%%%%%%%%%%%%%%%%%%%%%%%%%%%%%%%%%%%%%%%%%%%%%%%%%%%%%%%%%%%%%%%%%%%%%%%%%%%%
\section{Entropy collapse}
%%%%%%%%%%%%%%%%%%%%%%%%%%%%%%%%%%%%%%%%%%%%%%%%%%%%%%%%%%%%%%%%%%%%%%%%%%%%%%

Continuation entropy may decrease abruptly.

\begin{definition}

An \emph{entropy collapse} occurs whenever infinitely many admissible continuations reduce to finitely many.

\end{definition}

Collapse represents structural selection.

Rather than destroying information, continuation removes inadmissible futures.

%%%%%%%%%%%%%%%%%%%%%%%%%%%%%%%%%%%%%%%%%%%%%%%%%%%%%%%%%%%%%%%%%%%%%%%%%%%%%%
\section{Entropy and completion}
%%%%%%%%%%%%%%%%%%%%%%%%%%%%%%%%%%%%%%%%%%%%%%%%%%%%%%%%%%%%%%%%%%%%%%%%%%%%%%

Completion generally decreases entropy.

As continuation becomes more complete, fewer compatible extensions remain.

\begin{theorem}

Canonical completion never increases continuation entropy.

\end{theorem}

\begin{proof}

Completion adds structural constraints.

Additional constraints cannot create new admissible continuations.

Hence

\[
|\mathcal C(x_{\mathrm{completed}})|
\le
|\mathcal C(x)|.
\]

Therefore entropy cannot increase.

\end{proof}

%%%%%%%%%%%%%%%%%%%%%%%%%%%%%%%%%%%%%%%%%%%%%%%%%%%%%%%%%%%%%%%%%%%%%%%%%%%%%%
\section{Entropy and geometry}
%%%%%%%%%%%%%%%%%%%%%%%%%%%%%%%%%%%%%%%%%%%%%%%%%%%%%%%%%%%%%%%%%%%%%%%%%%%%%%

Continuation entropy is geometric.

Regions possessing large entropy exhibit rich continuation branching.

Regions possessing small entropy exhibit rigid continuation.

Boundaries typically coincide with entropy collapse.

Completion typically corresponds to entropy minimization.

Thus geometry becomes a visible manifestation of continuation entropy.

%%%%%%%%%%%%%%%%%%%%%%%%%%%%%%%%%%%%%%%%%%%%%%%%%%%%%%%%%%%%%%%%%%%%%%%%%%%%%%
\section{The Entropy Principle}
%%%%%%%%%%%%%%%%%%%%%%%%%%%%%%%%%%%%%%%%%%%%%%%%%%%%%%%%%%%%%%%%%%%%%%%%%%%%%%

Classical entropy theories measure uncertainty because they observe systems after continuation has already generated branching.

Continuation Mathematics reverses this viewpoint.

Branching is primary.

Entropy merely measures branching.

Consequently probability, information, statistical mechanics, dynamical entropy, and algorithmic complexity all arise as different numerical realizations of a deeper continuation-theoretic quantity.

Entropy is therefore not an independent mathematical primitive.

It is the observable shadow of continuation freedom.

\chapter{Continuation Rigidity}

The preceding chapters developed the geometric language of continuation.
Continuation paths describe evolution.
Metrics measure accessibility.
Curvature measures obstruction.
Volume measures structural abundance.
Entropy measures uncertainty.

A fundamental question now arises.

\begin{center}
\emph{When does local continuation determine the entire global object?}
\end{center}

This question is universal.

In Euclidean geometry, rigid bodies are determined by local distances.

In graph theory, certain graphs are uniquely determined by their local neighborhoods.

In differential geometry, curvature may determine global geometry.

Continuation Mathematics asks for the corresponding principle governing continuation spaces.

The answer leads naturally to the notion of continuation rigidity.

%%%%%%%%%%%%%%%%%%%%%%%%%%%%%%%%%%%%%%%%%%%%%%%%%%%%%%%%%%%%

\section{Failure of rigidity}

Most continuation spaces are not rigid.

Distinct global continuation systems may possess identical local continuation behaviour while differing globally.

For example, two continuation trees may agree completely for the first hundred continuation levels while diverging afterwards.

Likewise two algebraic continuation systems may admit identical local morphisms but possess different completion spaces.

Thus local continuation alone does not generally determine the whole object.

Rigidity therefore becomes a structural property rather than an automatic consequence.

%%%%%%%%%%%%%%%%%%%%%%%%%%%%%%%%%%%%%%%%%%%%%%%%%%%%%%%%%%%%

\section{Continuation determination}

\begin{definition}[Continuation determination]
Let $X$ be a continuation space.

A subset
\[
D\subseteq X
\]
is called a
\emph{determining continuation set}
if every continuation compatible with $D$
extends uniquely to the whole space.
\end{definition}

Thus determination means that no independent continuation choices remain.

Every future continuation is already forced.

%%%%%%%%%%%%%%%%%%%%%%%%%%%%%%%%%%%%%%%%%%%%%%%%%%%%%%%%%%%%

\section{Rigid continuation spaces}

\begin{definition}[Rigid continuation space]
A continuation space is called
\emph{rigid}
if every finite determining continuation set uniquely determines the entire continuation structure.
\end{definition}

Intuitively,

\[
\boxed{
\text{Local continuation}
\Longrightarrow
\text{Global continuation}.
}
\]

%%%%%%%%%%%%%%%%%%%%%%%%%%%%%%%%%%%%%%%%%%%%%%%%%%%%%%%%%%%%

\section{Propagation of determination}

Rigidity is fundamentally a propagation phenomenon.

Once sufficient continuation information is fixed,
every admissible continuation becomes forced.

There is therefore no remaining branching.

Continuation propagates deterministically throughout the space.

%%%%%%%%%%%%%%%%%%%%%%%%%%%%%%%%%%%%%%%%%%%%%%%%%%%%%%%%%%%%

\section{Rigidity operators}

Every continuation space possesses a natural closure operator.

Starting from any subset,

\[
A_0=D,
\]

define recursively

\[
A_{n+1}
=
\operatorname{Cont}(A_n).
\]

This produces

\[
D
\subseteq
A_1
\subseteq
A_2
\subseteq
\cdots
\]

whose union is the continuation closure generated by $D$.

%%%%%%%%%%%%%%%%%%%%%%%%%%%%%%%%%%%%%%%%%%%%%%%%%%%%%%%%%%%%

\begin{definition}[Rigidity closure]
The rigidity closure generated by $D$ is

\[
\operatorname{Rig}(D)
=
\bigcup_{n=0}^{\infty}A_n.
\]

\end{definition}

If

\[
\operatorname{Rig}(D)=X,
\]

then $D$ determines the entire continuation space.

%%%%%%%%%%%%%%%%%%%%%%%%%%%%%%%%%%%%%%%%%%%%%%%%%%%%%%%%%%%%

\section{Rigidity rank}

The speed of propagation measures the strength of rigidity.

\begin{definition}[Rigidity rank]

The rigidity rank of $D$ is the least ordinal (or natural number whenever finite)

\[
\rho(D)
\]

such that

\[
A_{\rho(D)}
=
A_{\rho(D)+1}.
\]

\end{definition}

Small rigidity rank corresponds to rapid propagation.

Large rigidity rank corresponds to weak determination.

Infinite rigidity rank indicates persistent structural freedom.

%%%%%%%%%%%%%%%%%%%%%%%%%%%%%%%%%%%%%%%%%%%%%%%%%%%%%%%%%%%%

\section{Minimal determining systems}

Determination should be economical.

\begin{definition}

A determining continuation set is
\emph{minimal}
if none of its proper subsets determine the entire space.

\end{definition}

Minimal determining systems generalize

\begin{itemize}
\item generating sets,
\item algebraic bases,
\item graph separators,
\item spanning trees,
\item coordinate systems.
\end{itemize}

Continuation Mathematics treats all of these as manifestations of one common phenomenon.

%%%%%%%%%%%%%%%%%%%%%%%%%%%%%%%%%%%%%%%%%%%%%%%%%%%%%%%%%%%%

\section{Rigidity spectrum}

Not every continuation space possesses the same degree of rigidity.

One therefore associates to every continuation space its rigidity spectrum,

\[
\mathcal R(X),
\]

consisting of all rigidity ranks arising from determining subsets.

This spectrum measures how easily global structure is forced.

Highly rigid spaces possess small spectra.

Flexible spaces possess large spectra.

%%%%%%%%%%%%%%%%%%%%%%%%%%%%%%%%%%%%%%%%%%%%%%%%%%%%%%%%%%%%

\section{Rigidity and completion}

Rigidity interacts naturally with completion.

Incomplete continuation systems often admit multiple inequivalent completions.

As continuation information accumulates,
the family of admissible completions decreases.

Eventually one of two situations occurs.

Either

\[
|\operatorname{Comp}(X)|>1,
\]

and genuine structural ambiguity remains,

or

\[
|\operatorname{Comp}(X)|=1,
\]

in which case continuation rigidity has forced a unique completion.

Thus rigidity may be interpreted as the collapse of completion ambiguity.

%%%%%%%%%%%%%%%%%%%%%%%%%%%%%%%%%%%%%%%%%%%%%%%%%%%%%%%%%%%%

\section{Structural rigidity principle}

The concepts introduced throughout this chapter suggest the following general principle.

\begin{theorem}[Structural Rigidity Principle]
Let $X$ be a continuation space.

Whenever every admissible continuation extending a finite determining set is uniquely forced at every stage, the continuation closure generated by that set equals the unique completion of $X$.
\end{theorem}

The theorem expresses the central geometric philosophy of rigidity.

Global structure is not postulated.

Global structure emerges because every local continuation eventually becomes forced.

%%%%%%%%%%%%%%%%%%%%%%%%%%%%%%%%%%%%%%%%%%%%%%%%%%%%%%%%%%%%

\section{The role of rigidity}

Rigidity occupies a distinguished position within Continuation Mathematics.

Earlier chapters measured continuation.

Rigidity explains why continuation eventually stops admitting alternatives.

Many classical uniqueness theorems are manifestations of rigidity.

The continuation framework therefore places rigidity beside curvature, entropy, volume, and dimension as one of the fundamental geometric invariants governing continuation spaces.

The following part will show that these geometric invariants are not merely abstract constructions. They naturally generate universal mathematical structures that appear throughout analysis, algebra, topology, probability, number theory, and ultimately the mathematical foundations underlying quantum theory.

\chapter{Universal Continuation Geometry}

The preceding chapters have progressively revealed that continuation possesses an intrinsic geometry.

Beginning with continuation paths, we introduced metrics that measure continuation distance, geodesics that describe optimal continuation, curvature that measures structural obstruction, volume that quantifies continuation abundance, entropy that measures continuation uncertainty, and rigidity that characterizes when local continuation uniquely determines global structure.

These notions were introduced independently.

The purpose of this chapter is to show that they are not independent constructions. Rather, they arise as different manifestations of a single universal geometric object.

This object will be called the \emph{Universal Continuation Geometry}.

%%%%%%%%%%%%%%%%%%%%%%%%%%%%%%%%%%%%%%%%%%%%%%%%%%%%%%%%%%%%

\section{The geometric unification problem}

Every mature branch of mathematics eventually discovers that seemingly unrelated structures are governed by a common geometry.

Euclidean geometry unifies distance and angle.

Riemannian geometry unifies curvature and geodesics.

Algebraic geometry unifies algebra with topology.

Continuation Mathematics seeks the analogous unification.

The question is therefore not

\begin{center}
"What is continuation distance?"

or

"What is continuation curvature?"
\end{center}

The correct question is

\begin{center}
\emph{What geometric object simultaneously generates every continuation invariant?}
\end{center}

%%%%%%%%%%%%%%%%%%%%%%%%%%%%%%%%%%%%%%%%%%%%%%%%%%%%%%%%%%%%

\section{The continuation geometry}

\begin{definition}[Continuation geometry]

A \emph{continuation geometry} consists of a continuation space

\[
X
\]

equipped with the following compatible structures:

\begin{enumerate}
\item a continuation metric;

\item continuation paths;

\item continuation geodesics;

\item continuation curvature;

\item continuation volume;

\item continuation entropy;

\item continuation rigidity;

\item continuation completion.
\end{enumerate}

These structures are required to arise from the same continuation relation.

\end{definition}

Thus geometry is not additional data.

Geometry is generated by continuation itself.

%%%%%%%%%%%%%%%%%%%%%%%%%%%%%%%%%%%%%%%%%%%%%%%%%%%%%%%%%%%%

\section{Derived geometry}

One of the central philosophical principles of Continuation Mathematics now becomes visible.

Classically, geometric structures are usually introduced independently.

Continuation Mathematics reverses this viewpoint.

Continuation is primitive.

Geometry is derived.

Consequently,

\[
\boxed{
\text{Continuation}
\Longrightarrow
\text{Geometry}.
}
\]

%%%%%%%%%%%%%%%%%%%%%%%%%%%%%%%%%%%%%%%%%%%%%%%%%%%%%%%%%%%%

\section{Universal geometric invariants}

Every continuation geometry possesses a canonical family of invariants.

These include

\[
\begin{aligned}
d(X)
&\qquad
\text{(metric)},\\
\Gamma(X)
&\qquad
\text{(geodesics)},\\
K(X)
&\qquad
\text{(curvature)},\\
V(X)
&\qquad
\text{(volume)},\\
H(X)
&\qquad
\text{(entropy)},\\
\rho(X)
&\qquad
\text{(rigidity)},\\
\operatorname{Comp}(X)
&\qquad
\text{(completion)}.
\end{aligned}
\]

These invariants are not arbitrary.

Each measures a different aspect of the same continuation structure.

%%%%%%%%%%%%%%%%%%%%%%%%%%%%%%%%%%%%%%%%%%%%%%%%%%%%%%%%%%%%

\section{Compatibility}

The universal invariants cannot vary independently.

For example,

high rigidity typically implies low entropy,

while large curvature often increases geodesic complexity,

and maximal completion frequently corresponds to maximal continuation volume.

The exact relationships depend upon the continuation system under consideration.

This motivates the study of universal compatibility laws.

%%%%%%%%%%%%%%%%%%%%%%%%%%%%%%%%%%%%%%%%%%%%%%%%%%%%%%%%%%%%

\section{Universal continuation tensor}

To encode the interaction of all continuation invariants simultaneously, we introduce a universal geometric object.

\begin{definition}[Universal continuation tensor]

The \emph{Universal Continuation Tensor}

\[
\mathbb{T}(X)
\]

is the total collection of all continuation invariants together with every algebraic relation connecting them.

\end{definition}

Unlike tensors in differential geometry, the continuation tensor is not initially numerical.

It is structural.

Its components consist of continuation invariants and their interaction laws.

%%%%%%%%%%%%%%%%%%%%%%%%%%%%%%%%%%%%%%%%%%%%%%%%%%%%%%%%%%%%

\section{Geometric universality}

\begin{definition}

A geometric invariant is called
\emph{universal}
if it is definable solely from continuation and is independent of any particular mathematical realization.

\end{definition}

This distinguishes intrinsic continuation geometry from geometry inherited from ambient mathematical structures.

%%%%%%%%%%%%%%%%%%%%%%%%%%%%%%%%%%%%%%%%%%%%%%%%%%%%%%%%%%%%

\section{Functoriality}

Continuation geometry is naturally functorial.

Every continuation morphism

\[
f:X\rightarrow Y
\]

induces transformations of all continuation invariants.

Distances,

geodesics,

curvature,

volume,

entropy,

rigidity,

and completion

are transported together.

Thus continuation geometry behaves as a single mathematical object.

%%%%%%%%%%%%%%%%%%%%%%%%%%%%%%%%%%%%%%%%%%%%%%%%%%%%%%%%%%%%

\section{Universality Principle}

The constructions developed throughout this part suggest the following guiding theorem.

\begin{theorem}[Universality Principle]

Every intrinsic geometric property of a continuation system is recoverable from its continuation relation.

Conversely, every continuation relation determines a unique intrinsic continuation geometry.

\end{theorem}

The theorem states that continuation and geometry determine one another.

Neither is fundamentally prior once the continuation relation has been specified.

%%%%%%%%%%%%%%%%%%%%%%%%%%%%%%%%%%%%%%%%%%%%%%%%%%%%%%%%%%%%

\section{The hierarchy of geometry}

Continuation geometry organizes itself into successive levels.

\[
\begin{aligned}
\text{Continuation Relation}
&\\
\Downarrow
\\
\text{Continuation Space}
&\\
\Downarrow
\\
\text{Metric}
&\\
\Downarrow
\\
\text{Geodesics}
&\\
\Downarrow
\\
\text{Curvature}
&\\
\Downarrow
\\
\text{Volume}
&\\
\Downarrow
\\
\text{Entropy}
&\\
\Downarrow
\\
\text{Rigidity}
&\\
\Downarrow
\\
\text{Completion}
\end{aligned}
\]

Each level is forced by those preceding it.

Nothing is introduced independently.

%%%%%%%%%%%%%%%%%%%%%%%%%%%%%%%%%%%%%%%%%%%%%%%%%%%%%%%%%%%%

\section{Universal geometric emergence}

One of the central discoveries of Continuation Mathematics is that geometry itself is an emergent phenomenon.

The primitive notion is neither point, distance, manifold, topology, nor algebra.

The primitive notion is continuation.

Geometry appears only after continuation has generated sufficient structural organization.

Thus geometry is not assumed.

Geometry is the large-scale manifestation of continuation.

%%%%%%%%%%%%%%%%%%%%%%%%%%%%%%%%%%%%%%%%%%%%%%%%%%%%%%%%%%%%

\section{The completion of geometric foundations}

The present part has established the geometric foundations of Continuation Mathematics.

Beginning from continuation alone we have recovered

\begin{itemize}

\item metric,

\item path,

\item geodesic,

\item curvature,

\item volume,

\item entropy,

\item rigidity,

\item completion,

\item universal geometry.

\end{itemize}

Each has been derived from the continuation relation without external geometric assumptions.

This completes the geometric layer of the theory.

The following part enlarges the scope considerably.

Rather than studying continuation inside individual mathematical systems, we investigate continuation as a universal mathematical principle capable of generating entire mathematical universes.

The transition from geometry to universality mirrors one of the deepest transitions in classical mathematics: local structure gives way to universal structure.

Continuation Mathematics now follows the same progression.




\part{Universal Continuation}

\chapter{Continuation Universes}

The preceding parts developed Continuation Mathematics as an intrinsic mathematical theory.

Beginning from continuation itself, we constructed partial mathematical objects, continuation systems, continuation spaces, continuation algebra, and continuation geometry.

A natural question now arises.

\begin{center}
\emph{What mathematical worlds can continuation generate?}
\end{center}

This question leads to the notion of a continuation universe.

Continuation Mathematics does not begin with a fixed mathematical universe.

Rather, mathematical universes emerge from continuation itself.

%%%%%%%%%%%%%%%%%%%%%%%%%%%%%%%%%%%%%%%%%%%%%%%%%%%%%%%%%%%%

\section{The problem of mathematical universes}

Classical mathematics is usually developed inside a single ambient universe.

Depending upon one's foundations, this universe may consist of

\begin{itemize}

\item sets,

\item classes,

\item types,

\item categories,

\item higher categories,

\item universes,

\item toposes,

or similar foundational objects.

\end{itemize}

These are ordinarily assumed before mathematics begins.

Continuation Mathematics reverses this order.

Instead of assuming a universe, we ask

\begin{center}
\emph{Which universes are generated by continuation?}
\end{center}

%%%%%%%%%%%%%%%%%%%%%%%%%%%%%%%%%%%%%%%%%%%%%%%%%%%%%%%%%%%%

\section{The continuation viewpoint}

Continuation is fundamentally prior to objects.

Objects arise only as stabilized continuation structures.

Consequently,

entire mathematical universes become secondary constructions.

The primitive entity is therefore not

\[
x\in X,
\]

but rather

\[
x
\rightsquigarrow
y,
\]

the possibility of continuation.

%%%%%%%%%%%%%%%%%%%%%%%%%%%%%%%%%%%%%%%%%%%%%%%%%%%%%%%%%%%%

\section{Continuation universes}

\begin{definition}[Continuation universe]

A
\emph{continuation universe}
is a maximal continuation system closed under every admissible continuation operation.

\end{definition}

Closure means that whenever a continuation is admissible inside the universe,

its realization also belongs to the universe.

Thus the universe continuously regenerates itself.

%%%%%%%%%%%%%%%%%%%%%%%%%%%%%%%%%%%%%%%%%%%%%%%%%%%%%%%%%%%%

\section{Internal closure}

Every continuation universe possesses an intrinsic closure operator.

Beginning with any collection

\[
A\subseteq U,
\]

one repeatedly applies every admissible continuation.

This produces

\[
A
\subseteq
\operatorname{Cont}(A)
\subseteq
\operatorname{Cont}^2(A)
\subseteq
\cdots
\]

whose union generates the smallest continuation universe containing the original collection.

%%%%%%%%%%%%%%%%%%%%%%%%%%%%%%%%%%%%%%%%%%%%%%%%%%%%%%%%%%%%

\section{Generation}

\begin{definition}[Generating family]

A subset

\[
G\subseteq U
\]

is called a
\emph{generating family}

if repeated continuation of elements of $G$

produces the entire universe.

\end{definition}

Generation therefore replaces construction.

Universes are grown rather than assembled.

%%%%%%%%%%%%%%%%%%%%%%%%%%%%%%%%%%%%%%%%%%%%%%%%%%%%%%%%%%%%

\section{Local and global universes}

Continuation naturally distinguishes two scales.

A local continuation universe is generated by finitely many continuation objects.

A global continuation universe may require infinitely many generators.

This distinction parallels the classical distinction between finite and infinite mathematical structures.

%%%%%%%%%%%%%%%%%%%%%%%%%%%%%%%%%%%%%%%%%%%%%%%%%%%%%%%%%%%%

\section{Minimal universes}

Not every continuation universe contains unnecessary information.

\begin{definition}

A continuation universe is
\emph{minimal}

if no proper continuation subuniverse generates the same continuation structure.

\end{definition}

Minimal universes represent irreducible mathematical worlds.

%%%%%%%%%%%%%%%%%%%%%%%%%%%%%%%%%%%%%%%%%%%%%%%%%%%%%%%%%%%%

\section{Maximal universes}

At the opposite extreme lie universal continuation universes.

\begin{definition}

A continuation universe is
\emph{universal}

if every admissible continuation system embeds into it through a continuation morphism.

\end{definition}

Universal continuation universes play the role occupied by universal objects throughout mathematics.

%%%%%%%%%%%%%%%%%%%%%%%%%%%%%%%%%%%%%%%%%%%%%%%%%%%%%%%%%%%%

\section{Hierarchy of universes}

Continuation universes organize naturally into a hierarchy.

\[
U_0
\subseteq
U_1
\subseteq
U_2
\subseteq
\cdots
\]

where each level admits strictly richer continuation phenomena than the previous one.

Some universes admit only finite continuation.

Others admit infinite continuation.

Still others admit universal completion.

The hierarchy measures increasing structural richness.

%%%%%%%%%%%%%%%%%%%%%%%%%%%%%%%%%%%%%%%%%%%%%%%%%%%%%%%%%%%%

\section{Emergence of mathematics}

One of the central principles of this book may now be stated.

\begin{theorem}[Emergence Principle]

Every mathematical structure appearing inside a continuation universe is generated by admissible continuation.

No mathematical object is fundamentally primitive.

Every object is the completion of a continuation process.

\end{theorem}

This theorem represents a foundational reversal.

Rather than constructing continuation from mathematics,

Continuation Mathematics constructs mathematics from continuation.

%%%%%%%%%%%%%%%%%%%%%%%%%%%%%%%%%%%%%%%%%%%%%%%%%%%%%%%%%%%%

\section{The plurality of mathematics}

Classically one often speaks of

"the universe of mathematics."

Continuation Mathematics suggests a broader perspective.

There need not exist a unique mathematical universe.

Distinct continuation laws generate distinct continuation universes.

Each possesses its own admissible objects,

its own continuation geometry,

its own algebra,

its own completion theory,

and ultimately its own mathematics.

The study of mathematics therefore becomes the study of the possible continuation universes.

%%%%%%%%%%%%%%%%%%%%%%%%%%%%%%%%%%%%%%%%%%%%%%%%%%%%%%%%%%%%

\section{Towards realization theory}

A continuation universe is not yet classical mathematics.

It is a universal structural object.

The next question is therefore inevitable.

How do familiar mathematical systems arise from continuation universes?

The answer requires a theory of realization.

A realization identifies those continuation structures that manifest themselves as sets, topological spaces, groups, manifolds, measure spaces, categories, or other classical mathematical objects.

The next chapter develops precisely this theory.

\chapter{Realization Theory}

The preceding chapter introduced continuation universes as maximal systems generated entirely by admissible continuation.

A continuation universe is an abstract mathematical object.

The purpose of the present chapter is to explain how familiar mathematical structures arise from continuation universes.

This passage from universal continuation to concrete mathematics will be called \emph{realization}.

Realization Theory forms the bridge between Continuation Mathematics and every existing branch of mathematics.

%%%%%%%%%%%%%%%%%%%%%%%%%%%%%%%%%%%%%%%%%%%%%%%%%%%%%%%%%%%%

\section{The realization problem}

Continuation Mathematics is intentionally primitive.

Its objects need not initially resemble

\begin{itemize}

\item sets,

\item groups,

\item manifolds,

\item topological spaces,

\item vector spaces,

\item measure spaces, or 

\item any other familiar mathematical structures.

\end{itemize}

Instead, they possess only continuation.

The central question therefore becomes

\begin{center}

\emph{When does a continuation universe appear as a classical mathematical structure?}

\end{center}

This is the realization problem.

%%%%%%%%%%%%%%%%%%%%%%%%%%%%%%%%%%%%%%%%%%%%%%%%%%%%%%%%%%%%

\section{Realizations}

\begin{definition}[Realization]

A realization of a continuation universe

\[
U
\]

is a structure-preserving assignment

\[
\mathfrak R :
U
\longrightarrow
\mathcal C
\]

into some mathematical category

\[
\mathcal C,
\]

such that admissible continuations correspond to admissible constructions inside $\mathcal C$.

\end{definition}

Thus realization translates continuation into another mathematical language while preserving its structural content.

%%%%%%%%%%%%%%%%%%%%%%%%%%%%%%%%%%%%%%%%%%%%%%%%%%%%%%%%%%%%

\section{The realization principle}

The essential philosophy of realization theory is remarkably simple.

Continuation does not imitate mathematics.

Continuation becomes mathematics.

The same continuation universe may admit many different realizations.

Each realization reveals one aspect of the underlying continuation structure.

%%%%%%%%%%%%%%%%%%%%%%%%%%%%%%%%%%%%%%%%%%%%%%%%%%%%%%%%%%%%

\section{Faithful realizations}

Not every realization preserves all continuation information.

\begin{definition}

A realization

\[
\mathfrak R
\]

is called faithful

if distinct continuation structures always remain distinct after realization.

\end{definition}

Faithful realizations preserve the entire continuation geometry.

%%%%%%%%%%%%%%%%%%%%%%%%%%%%%%%%%%%%%%%%%%%%%%%%%%%%%%%%%%%%

\section{Complete realizations}

Some realizations preserve not only objects but every continuation relation.

\begin{definition}

A realization is called complete if every continuation relation is represented exactly.

\end{definition}

Complete realizations preserve

\begin{itemize}

\item continuation paths,

\item continuation spaces,

\item continuation algebra,

\item continuation geometry,

\item continuation completion.

\end{itemize}

%%%%%%%%%%%%%%%%%%%%%%%%%%%%%%%%%%%%%%%%%%%%%%%%%%%%%%%%%%%%

\section{Equivalent realizations}

Different mathematical languages may represent exactly the same continuation universe.

\begin{definition}

Two realizations

\[
\mathfrak R_1,
\mathfrak R_2
\]

are equivalent whenever there exists an isomorphism between their images preserving every continuation relation.

\end{definition}

Thus realization concerns structural content rather than notation.

%%%%%%%%%%%%%%%%%%%%%%%%%%%%%%%%%%%%%%%%%%%%%%%%%%%%%%%%%%%%

\section{The realization spectrum}

Every continuation universe possesses many possible realizations.

\begin{definition}

The realization spectrum of a continuation universe

\[
U
\]

is the collection

\[
\operatorname{Spec}(U)
\]

of all inequivalent realizations of $U$.

\end{definition}

The realization spectrum measures the mathematical diversity contained within a single continuation universe.

%%%%%%%%%%%%%%%%%%%%%%%%%%%%%%%%%%%%%%%%%%%%%%%%%%%%%%%%%%%%

\section{Universal realization}

Certain continuation universes admit extraordinarily rich realization spectra.

\begin{definition}

A continuation universe is universally realizable if every sufficiently rich mathematical language appears as one of its realizations.

\end{definition}

Such universes simultaneously contain

\begin{itemize}

\item algebra,

\item topology,

\item geometry,

\item analysis,

\item probability,

\item information,

\item category theory,

\end{itemize}

and many other mathematical disciplines.

%%%%%%%%%%%%%%%%%%%%%%%%%%%%%%%%%%%%%%%%%%%%%%%%%%%%%%%%%%%%

\section{Emergent mathematical objects}

Objects appearing inside a realization are not primitive.

They emerge through the realization process.

Thus

points,

sets,

functions,

groups,

manifolds,

operators,

categories,

and measures

are realization-dependent manifestations of deeper continuation objects.

Different realizations may assign entirely different meanings to the same continuation structure.

%%%%%%%%%%%%%%%%%%%%%%%%%%%%%%%%%%%%%%%%%%%%%%%%%%%%%%%%%%%%

\section{Structural invariance}

Although realizations may differ dramatically,

certain properties remain unchanged.

These are continuation invariants.

Examples include

\begin{itemize}

\item continuation dimension,

\item continuation entropy,

\item continuation rigidity,

\item continuation completion,

\item continuation spectra.

\end{itemize}

Such invariants belong to the continuation universe itself rather than to any particular realization.

%%%%%%%%%%%%%%%%%%%%%%%%%%%%%%%%%%%%%%%%%%%%%%%%%%%%%%%%%%%%

\section{The Realization Theorem}

The preceding ideas suggest one of the central principles of Continuation Mathematics.

\begin{theorem}[Realization Principle]

Every mathematical theory obtained from a continuation universe is determined by its realization.

Distinct realizations of the same continuation universe describe different mathematical manifestations of the same underlying continuation structure.

\end{theorem}

This theorem expresses a profound shift in mathematical foundations.

Classically,

one begins with mathematical objects and studies their properties.

Continuation Mathematics reverses this order.

One begins with continuation,

and mathematical objects emerge only after a realization has been chosen.

%%%%%%%%%%%%%%%%%%%%%%%%%%%%%%%%%%%%%%%%%%%%%%%%%%%%%%%%%%%%

\section{The role of realization theory}

Realization Theory transforms Continuation Mathematics from a single mathematical discipline into a universal foundational framework.

Rather than competing with existing mathematics,

Continuation Mathematics explains why existing mathematics possesses the structures that it does.

Every classical mathematical theory becomes a realization of a deeper continuation universe.

The remaining chapters of this part develop this claim explicitly.

Beginning with classical set theory, we shall recover the major branches of mathematics as canonical realizations of continuation universes.

\chapter{Universal Mathematical Structures}

\section{The universal structural problem}

Continuation Mathematics does not begin with sets, elements, operations,
spaces, numbers, functions, or relations. These are all mathematical
objects whose existence must themselves be explained.

The fundamental question is therefore not

\begin{center}
``What mathematical structures exist?''
\end{center}

but rather

\begin{center}
``How can mathematical structure itself arise?''
\end{center}

The purpose of this chapter is to answer this question.

The central claim is that mathematical structure is not primitive.
Structure is generated by continuation.

%%%%%%%%%%%%%%%%%%%%%%%%%%%%%%%%%%%%%%%%%%%%%%%%%%%%%%%%%%%%%%%%%%%%%%%%%%%%%%

\section{Structural emergence}

Every continuation system possesses admissible continuations.

Those continuations generate continuation spaces.

Continuation spaces possess closure operators.

Closure generates completion.

Completion produces stable mathematical objects.

Thus mathematical structure is never assumed.
It emerges.

\begin{definition}[Structural emergence]
A mathematical structure is said to emerge whenever it appears as a stable
completion of an underlying continuation process.
\end{definition}

Thus permanence is always secondary.

Continuation is always primary.

%%%%%%%%%%%%%%%%%%%%%%%%%%%%%%%%%%%%%%%%%%%%%%%%%%%%%%%%%%%%%%%%%%%%%%%%%%%%%%

\section{Universal structural generators}

Every continuation system possesses certain canonical constructions.

These include

\[
\begin{aligned}
&\text{Objects},\\
&\text{Relations},\\
&\text{Morphisms},\\
&\text{Closure},\\
&\text{Interior},\\
&\text{Boundary},\\
&\text{Dimension},\\
&\text{Connectivity},\\
&\text{Completion}.
\end{aligned}
\]

These are not independent notions.

They are generated by continuation itself.

\begin{theorem}[Universal structural generators]
Every continuation system canonically generates each of the preceding
structural notions.
\end{theorem}

\begin{proof}
Each construction has already been established in the preceding Parts of
this work.

Objects generate admissible continuations.

Continuation generates spaces.

Spaces generate closure.

Closure generates completion.

Completion generates stable mathematical structure.

Therefore every continuation system canonically generates the stated
structures.
\end{proof}

%%%%%%%%%%%%%%%%%%%%%%%%%%%%%%%%%%%%%%%%%%%%%%%%%%%%%%%%%%%%%%%%%%%%%%%%%%%%%%

\section{Universal realization principle}

The preceding theorem suggests a remarkable possibility.

Suppose one wishes to construct an arbitrary branch of mathematics.

Instead of postulating its axioms directly, one seeks an underlying
continuation system whose completed structures reproduce precisely that
branch.

This motivates the central realization principle.

\begin{definition}[Universal realization]
A mathematical theory is said to be realizable if it is recovered as the
stable completion of a continuation system.
\end{definition}

Realization is therefore not interpretation.

It is generation.

%%%%%%%%%%%%%%%%%%%%%%%%%%%%%%%%%%%%%%%%%%%%%%%%%%%%%%%%%%%%%%%%%%%%%%%%%%%%%%

\section{The realization hierarchy}

Different continuation systems may generate different completed structures.

Some systems produce finite structures.

Others generate infinite structures.

Others generate metric structures.

Others generate algebraic structures.

Others generate topological structures.

Others generate probabilistic structures.

Others generate quantum structures.

Thus realization naturally forms a hierarchy.

\begin{definition}[Realization hierarchy]
The realization hierarchy is the partially ordered collection of all
mathematical structures ordered by the continuation systems from which they
are generated.
\end{definition}

The ordering is determined by structural derivability rather than by logical
strength.

%%%%%%%%%%%%%%%%%%%%%%%%%%%%%%%%%%%%%%%%%%%%%%%%%%%%%%%%%%%%%%%%%%%%%%%%%%%%%%

\section{Structural universality}

The existence of the realization hierarchy suggests a stronger principle.

Rather than constructing individual mathematical theories separately, one
may seek continuation systems capable of simultaneously generating entire
families of mathematical structures.

\begin{definition}[Universal continuation system]
A continuation system is universal if every member of some specified class
of mathematical structures is realizable from it.
\end{definition}

Universality therefore refers to generative capacity.

%%%%%%%%%%%%%%%%%%%%%%%%%%%%%%%%%%%%%%%%%%%%%%%%%%%%%%%%%%%%%%%%%%%%%%%%%%%%%%

\section{The architecture of mathematics}

The picture that now emerges is fundamentally different from the traditional
foundational hierarchy.

Classically one begins with

\[
\text{Sets}
\longrightarrow
\text{Relations}
\longrightarrow
\text{Functions}
\longrightarrow
\text{Structures}.
\]

Continuation Mathematics instead produces

\[
\text{Continuation}
\longrightarrow
\text{Admissibility}
\longrightarrow
\text{Completion}
\longrightarrow
\text{Structure}
\longrightarrow
\text{Realization}.
\]

Thus structure appears only after continuation has stabilized.

%%%%%%%%%%%%%%%%%%%%%%%%%%%%%%%%%%%%%%%%%%%%%%%%%%%%%%%%%%%%%%%%%%%%%%%%%%%%%%

\section{The universal recovery programme}

The remainder of this Part is devoted to demonstrating that the principal
branches of mathematics are realizable within this framework.

Each recovery follows the same pattern.

\begin{enumerate}
\item Construct an appropriate continuation universe.

\item Identify its admissible continuations.

\item Determine the corresponding continuation space.

\item Construct its completion.

\item Prove that the completed structure is isomorphic to the desired
mathematical theory.
\end{enumerate}

Thus every mathematical theory becomes a realization theorem.

%%%%%%%%%%%%%%%%%%%%%%%%%%%%%%%%%%%%%%%%%%%%%%%%%%%%%%%%%%%%%%%%%%%%%%%%%%%%%%

\section{The structural reduction theorem}

The preceding discussion may be summarized by the fundamental reduction
principle.

\begin{theorem}[Structural Reduction Theorem]
Every realizable mathematical theory is completely determined by the
continuation system from which it is generated together with its completion
operator.
\end{theorem}

\begin{proof}
The continuation system determines all admissible continuations.

These determine the continuation space.

The continuation space determines closure.

Closure determines completion.

Completion determines the resulting mathematical structure.

No additional primitive mathematical assumptions are required.
\end{proof}

This theorem establishes the universal programme of Continuation
Mathematics.

The remainder of this Part consists of progressively stronger realization
theorems showing how increasingly sophisticated branches of mathematics
emerge from continuation.

\chapter{Recovery of Classical Set Theory}

\section{The problem of collections}

Classical mathematics begins by postulating sets.

A set is regarded as a primitive collection whose members are simply assumed
to belong to it. Membership is therefore an undefined relation, and the
entire hierarchy of modern mathematics is erected upon this primitive notion.

Continuation Mathematics reverses this order.

Nothing is assumed to be a collection.

The primary object is always a continuation system.

The question is therefore not

\begin{center}
``What is a set?''
\end{center}

but rather

\begin{center}
``When does a collection become forced by continuation itself?''
\end{center}

The purpose of this chapter is to answer that question.

%%%%%%%%%%%%%%%%%%%%%%%%%%%%%%%%%%%%%%%%%%%%%%%%%%%%%%%%%%%%%%%%%%%%%%%%%%%%%%

\section{Continuation equivalence}

Different objects may possess different descriptions while exhibiting
identical continuation behaviour.

From the standpoint of continuation, such objects are indistinguishable.

\begin{definition}[Continuation equivalence]
Let $\mathcal C$ be a continuation system.

Two objects $x,y\in\mathcal C$ are said to be
\emph{continuation equivalent}, written

\[
x\sim y,
\]

whenever they possess identical continuation behaviour.

Equivalently,

\[
\operatorname{Cont}(x)
=
\operatorname{Cont}(y).
\]
\end{definition}

Continuation equivalence is intrinsic.

It depends only upon admissible continuation.

%%%%%%%%%%%%%%%%%%%%%%%%%%%%%%%%%%%%%%%%%%%%%%%%%%%%%%%%%%%%%%%%%%%%%%%%%%%%%%

\section{Completion classes}

Continuation equivalence partitions every continuation system into canonical
classes.

\begin{definition}[Completion class]
The completion class generated by an object $x$ is

\[
[x]
=
\{y:y\sim x\}.
\]

It consists of every object possessing the same continuation structure as
$x$.
\end{definition}

No collection has been postulated.

The collection is forced by continuation.

%%%%%%%%%%%%%%%%%%%%%%%%%%%%%%%%%%%%%%%%%%%%%%%%%%%%%%%%%%%%%%%%%%%%%%%%%%%%%%

\section{Observable mathematical objects}

The completion classes are precisely the stable objects visible from
continuation.

\begin{definition}[Observable object]
An observable mathematical object is a completion class of continuation
equivalent objects.
\end{definition}

Thus the primitive mathematical objects of Continuation Mathematics are not
individual elements.

They are completion classes.

Individual representatives become secondary.

%%%%%%%%%%%%%%%%%%%%%%%%%%%%%%%%%%%%%%%%%%%%%%%%%%%%%%%%%%%%%%%%%%%%%%%%%%%%%%

\section{Membership}

Classical membership is primitive.

Continuation membership is derived.

\begin{definition}[Continuation membership]
An object belongs to an observable object whenever it lies in its completion
class.

Thus

\[
x\in [y]
\]

means precisely

\[
x\sim y.
\]
\end{definition}

Membership is therefore nothing more than continuation equivalence.

%%%%%%%%%%%%%%%%%%%%%%%%%%%%%%%%%%%%%%%%%%%%%%%%%%%%%%%%%%%%%%%%%%%%%%%%%%%%%%

\section{Extensionality}

The classical Axiom of Extensionality now becomes a theorem.

\begin{theorem}[Continuation Extensionality]
Two observable objects are identical if and only if they contain precisely
the same continuation representatives.
\end{theorem}

\begin{proof}
Observable objects are completion classes.

Completion classes are uniquely determined by continuation equivalence.

Hence two observable objects coincide precisely when their representatives
coincide.
\end{proof}

Extensionality is therefore generated rather than assumed.

%%%%%%%%%%%%%%%%%%%%%%%%%%%%%%%%%%%%%%%%%%%%%%%%%%%%%%%%%%%%%%%%%%%%%%%%%%%%%%

\section{Subobjects}

Continuation itself naturally generates refinement.

Some continuation classes possess stronger admissibility conditions than
others.

\begin{definition}[Continuation subobject]
A completion class $A$ is a continuation subobject of $B$ whenever every
representative of $A$ is continuation equivalent to some representative
inside $B$ after restriction to a smaller continuation system.
\end{definition}

Subobjects therefore arise from refinement of continuation rather than from
primitive inclusion.

%%%%%%%%%%%%%%%%%%%%%%%%%%%%%%%%%%%%%%%%%%%%%%%%%%%%%%%%%%%%%%%%%%%%%%%%%%%%%%

\section{Structural operations}

Completion classes admit natural structural operations.

Two continuation classes may be merged by taking the smallest completion
class generated by their representatives.

Likewise they may intersect by taking the largest common continuation
completion generated by both.

These operations are intrinsic.

No external logic is required.

\begin{definition}[Structural union]
The structural union of two completion classes is the smallest completion
class generated by both.
\end{definition}

\begin{definition}[Structural intersection]
The structural intersection of two completion classes is the largest
completion class common to both.
\end{definition}

These operations are generated by continuation closure.

%%%%%%%%%%%%%%%%%%%%%%%%%%%%%%%%%%%%%%%%%%%%%%%%%%%%%%%%%%%%%%%%%%%%%%%%%%%%%%

\section{Higher completion classes}

Completion classes themselves possess continuation behaviour.

They therefore admit completion.

\begin{definition}[Second-order completion]
A second-order completion class is a completion class whose representatives
are themselves completion classes.
\end{definition}

Iterating this construction generates an entire hierarchy of observable
mathematical objects.

Unlike the cumulative hierarchy of classical set theory, this hierarchy is
generated dynamically by continuation.

%%%%%%%%%%%%%%%%%%%%%%%%%%%%%%%%%%%%%%%%%%%%%%%%%%%%%%%%%%%%%%%%%%%%%%%%%%%%%%

\section{Recovery theorem}

We may now recover the structural content of elementary set theory.

\begin{theorem}[Recovery of Classical Set Theory]
The category of completion classes generated by continuation possesses the
same structural behaviour as elementary classical set theory.

Membership, extensionality, inclusion, unions, intersections, and higher
collections are all recoverable as intrinsic consequences of continuation.
\end{theorem}

\begin{proof}
Every completion class is generated intrinsically from continuation
equivalence.

Membership is equivalence.

Extensionality follows immediately.

Structural unions and intersections arise from continuation closure.

Higher completion classes generate the analogue of higher-order collections.

Thus every elementary structural feature of classical set theory is
recovered without assuming sets as primitive objects.
\end{proof}

%%%%%%%%%%%%%%%%%%%%%%%%%%%%%%%%%%%%%%%%%%%%%%%%%%%%%%%%%%%%%%%%%%%%%%%%%%%%%%

\section{The foundational reversal}

The significance of this recovery is fundamental.

Classically one begins with

\[
\text{Elements}
\longrightarrow
\text{Sets}
\longrightarrow
\text{Mathematics}.
\]

Continuation Mathematics instead establishes

\[
\text{Continuation}
\longrightarrow
\text{Continuation Equivalence}
\longrightarrow
\text{Completion Classes}
\longrightarrow
\text{Observable Objects}
\longrightarrow
\text{Mathematics}.
\]

Collections therefore cease to be primitive.

They become stable manifestations of continuation.

Set theory is not rejected.

It is recovered.

The recovery of classical set theory is therefore the first instance of the
general realization programme developed throughout this Part.

\chapter{Recovery of Topology}

The purpose of this chapter is to recover classical topology as a realization of continuation theory rather than as an independently postulated mathematical discipline.

In classical mathematics a topology is introduced axiomatically as a distinguished family of subsets satisfying closure conditions under unions and finite intersections. Those axioms describe the final structure but do not explain why such a structure should exist.

Continuation Mathematics proceeds in the opposite direction.

The primitive notion is not openness.

The primitive notion is continuation.

Topology will emerge from the behavior of continuation classes.

%%%%%%%%%%%%%%%%%%%%%%%%%%%%%%%%%%%%%%%%%%%%%%%%%%%%%%%%%%%%%%%%%%%%%%%%%%%%%%%
\section{Topology as stable continuation}
%%%%%%%%%%%%%%%%%%%%%%%%%%%%%%%%%%%%%%%%%%%%%%%%%%%%%%%%%%%%%%%%%%%%%%%%%%%%%%%

A mathematical object is never observed directly.

One only observes the possible continuations that remain available after the object has been partially specified.

Consequently the local structure surrounding an object is determined by the family of objects sharing sufficiently similar continuation behavior.

This motivates the fundamental definition.

\begin{definition}[Continuation neighbourhood]
Let $X$ be a continuation universe and let $x\in X$.

A \emph{continuation neighbourhood} of $x$ is a continuation system
\[
N(x)
\]
satisfying the following property:

every object of $N(x)$ possesses all sufficiently short admissible continuations possessed by $x$.
\end{definition}

Intuitively, objects are close whenever their futures initially agree.

Thus neighbourhoods arise from agreement of continuation rather than agreement of coordinates.

%%%%%%%%%%%%%%%%%%%%%%%%%%%%%%%%%%%%%%%%%%%%%%%%%%%%%%%%%%%%%%%%%%%%%%%%%%%%%%%
\section{Intrinsic openness}
%%%%%%%%%%%%%%%%%%%%%%%%%%%%%%%%%%%%%%%%%%%%%%%%%%%%%%%%%%%%%%%%%%%%%%%%%%%%%%%

Continuation neighbourhoods immediately determine the notion of openness.

\begin{definition}[Open continuation class]

A continuation class
\[
U\subseteq X
\]
is called \emph{open} if every object of $U$ possesses a continuation neighbourhood entirely contained inside $U$.
\end{definition}

This definition contains no arbitrary axioms.

It is forced.

A set is open precisely when continuation cannot immediately escape it.

Openness therefore measures local continuation stability.

%%%%%%%%%%%%%%%%%%%%%%%%%%%%%%%%%%%%%%%%%%%%%%%%%%%%%%%%%%%%%%%%%%%%%%%%%%%%%%%
\section{Recovery of the topology axioms}
%%%%%%%%%%%%%%%%%%%%%%%%%%%%%%%%%%%%%%%%%%%%%%%%%%%%%%%%%%%%%%%%%%%%%%%%%%%%%%%

The classical topology axioms now become theorems.

\begin{theorem}[Empty and total continuation classes]

Both
\[
\varnothing
\]
and
\[
X
\]
are open.
\end{theorem}

\begin{proof}

Immediate.

The empty class contains no points.

Every continuation neighbourhood of every object lies inside $X$.

\end{proof}

\begin{theorem}[Arbitrary unions]

The union of any family of open continuation classes is open.
\end{theorem}

\begin{proof}

Suppose

\[
x\in\bigcup_{\alpha}U_\alpha .
\]

Then

\[
x\in U_\beta
\]

for some index $\beta$.

Since $U_\beta$ is open there exists a continuation neighbourhood of $x$ lying inside $U_\beta$, hence inside the union.

\end{proof}

\begin{theorem}[Finite intersections]

The intersection of finitely many open continuation classes is open.
\end{theorem}

\begin{proof}

Suppose

\[
x\in U_1\cap\cdots\cap U_n .
\]

Each $U_i$ contains a continuation neighbourhood of $x$.

Intersecting these finitely many neighbourhoods again produces a continuation neighbourhood.

Hence the intersection is open.

\end{proof}

Therefore the family of open continuation classes forms a topology.

Notice that no topology was assumed.

It has been recovered.

%%%%%%%%%%%%%%%%%%%%%%%%%%%%%%%%%%%%%%%%%%%%%%%%%%%%%%%%%%%%%%%%%%%%%%%%%%%%%%%
\section{Closure revisited}
%%%%%%%%%%%%%%%%%%%%%%%%%%%%%%%%%%%%%%%%%%%%%%%%%%%%%%%%%%%%%%%%%%%%%%%%%%%%%%%

Earlier we introduced continuation closure as a primitive continuation operation.

We now recover the classical topological interpretation.

\begin{theorem}

The continuation closure of a class equals the intersection of all closed continuation classes containing it.

\end{theorem}

\begin{proof}

Closed continuation classes are complements of open continuation classes.

The result follows directly from the universal property of continuation closure established previously.

\end{proof}

Thus closure acquires its classical interpretation without changing its intrinsic continuation definition.

%%%%%%%%%%%%%%%%%%%%%%%%%%%%%%%%%%%%%%%%%%%%%%%%%%%%%%%%%%%%%%%%%%%%%%%%%%%%%%%
\section{Interior revisited}
%%%%%%%%%%%%%%%%%%%%%%%%%%%%%%%%%%%%%%%%%%%%%%%%%%%%%%%%%%%%%%%%%%%%%%%%%%%%%%%

Likewise the continuation interior agrees with the union of all open continuation classes contained inside a given class.

Thus the continuation interior coincides with the classical interior operator after realization.

%%%%%%%%%%%%%%%%%%%%%%%%%%%%%%%%%%%%%%%%%%%%%%%%%%%%%%%%%%%%%%%%%%%%%%%%%%%%%%%
\section{Boundary}
%%%%%%%%%%%%%%%%%%%%%%%%%%%%%%%%%%%%%%%%%%%%%%%%%%%%%%%%%%%%%%%%%%%%%%%%%%%%%%%

The notion of boundary also emerges naturally.

\begin{definition}[Continuation boundary]

For any continuation class $A$

\[
\partial A
=
\overline{A}
\setminus
A^\circ .
\]

\end{definition}

Objects on the boundary possess admissible continuations entering both the interior and the exterior.

Boundaries therefore represent maximal continuation uncertainty.

%%%%%%%%%%%%%%%%%%%%%%%%%%%%%%%%%%%%%%%%%%%%%%%%%%%%%%%%%%%%%%%%%%%%%%%%%%%%%%%
\section{Dense continuation classes}
%%%%%%%%%%%%%%%%%%%%%%%%%%%%%%%%%%%%%%%%%%%%%%%%%%%%%%%%%%%%%%%%%%%%%%%%%%%%%%%

A continuation class is dense whenever every continuation neighbourhood intersects it.

Equivalently,

\[
\overline{A}=X.
\]

Density therefore expresses universal continuation accessibility.

%%%%%%%%%%%%%%%%%%%%%%%%%%%%%%%%%%%%%%%%%%%%%%%%%%%%%%%%%%%%%%%%%%%%%%%%%%%%%%%
\section{Basis}
%%%%%%%%%%%%%%%%%%%%%%%%%%%%%%%%%%%%%%%%%%%%%%%%%%%%%%%%%%%%%%%%%%%%%%%%%%%%%%%

The neighbourhood construction determines a natural basis.

\begin{definition}

A family

\[
\mathcal B
\]

of continuation neighbourhoods is called a continuation basis if every continuation neighbourhood contains one element of $\mathcal B$.

\end{definition}

Thus topology may be generated entirely from local continuation data.

%%%%%%%%%%%%%%%%%%%%%%%%%%%%%%%%%%%%%%%%%%%%%%%%%%%%%%%%%%%%%%%%%%%%%%%%%%%%%%%
\section{Recovery Theorem}
%%%%%%%%%%%%%%%%%%%%%%%%%%%%%%%%%%%%%%%%%%%%%%%%%%%%%%%%%%%%%%%%%%%%%%%%%%%%%%%

\begin{theorem}[Recovery of topology]

Every continuation universe possesses a canonical topology generated entirely by continuation neighbourhoods.

Conversely every classical topological space is realized by a continuation universe whose continuation neighbourhoods recover exactly the original topology.

\end{theorem}

\begin{proof}

The first statement follows from the previous sections.

The converse is obtained by taking admissible continuations to be local topological refinements.

The resulting continuation neighbourhoods coincide with the original neighbourhood system.

\end{proof}

%%%%%%%%%%%%%%%%%%%%%%%%%%%%%%%%%%%%%%%%%%%%%%%%%%%%%%%%%%%%%%%%%%%%%%%%%%%%%%%
\section{The meaning of topology}
%%%%%%%%%%%%%%%%%%%%%%%%%%%%%%%%%%%%%%%%%%%%%%%%%%%%%%%%%%%%%%%%%%%%%%%%%%%%%%%

Topology is therefore not fundamental.

It is a realization of continuation.

Open sets describe regions from which continuation cannot immediately escape.

Closed sets describe completion barriers.

Neighbourhoods describe local continuation agreement.

Closure records all possible completion limits.

Interior records unrestricted continuation.

Boundary records competing continuations.

Compactness records finite continuation completeness.

Connectedness records continuation indivisibility.

Every major concept of topology is recovered from continuation stability.

Thus topology becomes one realization of a deeper mathematical structure rather than an independent foundation.

\chapter{Recovery of Algebra}

The purpose of this chapter is to recover algebra as a realization of continuation theory.

Classically, algebra begins by postulating operations on sets.

Continuation Mathematics reverses this order.

Operations are not primitive.

Continuation is primitive.

Algebra arises when continuations themselves become composable.

%%%%%%%%%%%%%%%%%%%%%%%%%%%%%%%%%%%%%%%%%%%%%%%%%%%%%%%%%%%%%%%%%%%%%%%%%%%%%%%
\section{Composition before operation}
%%%%%%%%%%%%%%%%%%%%%%%%%%%%%%%%%%%%%%%%%%%%%%%%%%%%%%%%%%%%%%%%%%%%%%%%%%%%%%%

A continuation system possesses admissible extensions.

Suppose

\[
x
\longrightarrow
y
\]

and

\[
y
\longrightarrow
z.
\]

Whenever both continuations are admissible there exists a natural composite continuation

\[
x
\longrightarrow
z.
\]

Thus composition exists before any algebraic operation has been introduced.

Composition is therefore an intrinsic structural property of continuation.

%%%%%%%%%%%%%%%%%%%%%%%%%%%%%%%%%%%%%%%%%%%%%%%%%%%%%%%%%%%%%%%%%%%%%%%%%%%%%%%
\section{Continuation operations}
%%%%%%%%%%%%%%%%%%%%%%%%%%%%%%%%%%%%%%%%%%%%%%%%%%%%%%%%%%%%%%%%%%%%%%%%%%%%%%%

\begin{definition}

A \emph{continuation operation} on a continuation system is any rule

\[
\star :
\mathcal C\times\mathcal C
\rightarrow
\mathcal C
\]

whose output is the canonical continuation obtained by composing two compatible continuations.

\end{definition}

The operation is not arbitrarily chosen.

It is forced by admissibility.

Thus algebraic multiplication is recovered as continuation composition.

%%%%%%%%%%%%%%%%%%%%%%%%%%%%%%%%%%%%%%%%%%%%%%%%%%%%%%%%%%%%%%%%%%%%%%%%%%%%%%%
\section{Associativity}
%%%%%%%%%%%%%%%%%%%%%%%%%%%%%%%%%%%%%%%%%%%%%%%%%%%%%%%%%%%%%%%%%%%%%%%%%%%%%%%

Composition possesses an intrinsic associativity.

\begin{theorem}

Whenever three continuations are composable,

\[
(\alpha\star\beta)\star\gamma
=
\alpha\star(\beta\star\gamma).
\]

\end{theorem}

\begin{proof}

Both expressions represent the unique continuation obtained after performing the same three admissible extensions.

Since continuation depends only upon the completed propagation, both constructions coincide.

\end{proof}

Associativity therefore precedes abstract semigroup theory.

%%%%%%%%%%%%%%%%%%%%%%%%%%%%%%%%%%%%%%%%%%%%%%%%%%%%%%%%%%%%%%%%%%%%%%%%%%%%%%%
\section{Identity}
%%%%%%%%%%%%%%%%%%%%%%%%%%%%%%%%%%%%%%%%%%%%%%%%%%%%%%%%%%%%%%%%%%%%%%%%%%%%%%%

Every continuation system possesses a trivial continuation.

\begin{definition}

The identity continuation is the unique continuation leaving every partial object unchanged.

\end{definition}

It satisfies

\[
e\star x=x,
\qquad
x\star e=x.
\]

Thus identity arises automatically.

%%%%%%%%%%%%%%%%%%%%%%%%%%%%%%%%%%%%%%%%%%%%%%%%%%%%%%%%%%%%%%%%%%%%%%%%%%%%%%%
\section{Invertibility}
%%%%%%%%%%%%%%%%%%%%%%%%%%%%%%%%%%%%%%%%%%%%%%%%%%%%%%%%%%%%%%%%%%%%%%%%%%%%%%%

Not every continuation is reversible.

This leads naturally to invertibility.

\begin{definition}

A continuation is invertible whenever there exists another continuation returning every realized object to its previous continuation state.

\end{definition}

Invertibility therefore measures reversibility of continuation.

Groups become realizations of completely reversible continuation systems.

%%%%%%%%%%%%%%%%%%%%%%%%%%%%%%%%%%%%%%%%%%%%%%%%%%%%%%%%%%%%%%%%%%%%%%%%%%%%%%%
\section{Algebraic laws}
%%%%%%%%%%%%%%%%%%%%%%%%%%%%%%%%%%%%%%%%%%%%%%%%%%%%%%%%%%%%%%%%%%%%%%%%%%%%%%%

Every familiar algebraic law acquires a continuation interpretation.

\begin{center}
\begin{tabular}{ll}
Associativity & Independence of continuation bracketing\\
Identity & Trivial continuation\\
Inverse & Reversal of continuation\\
Closure & Completion stability\\
Cancellation & Continuation rigidity
\end{tabular}
\end{center}

Thus the traditional algebraic axioms become structural consequences rather than assumptions.

%%%%%%%%%%%%%%%%%%%%%%%%%%%%%%%%%%%%%%%%%%%%%%%%%%%%%%%%%%%%%%%%%%%%%%%%%%%%%%%
\section{Realization of semigroups}
%%%%%%%%%%%%%%%%%%%%%%%%%%%%%%%%%%%%%%%%%%%%%%%%%%%%%%%%%%%%%%%%%%%%%%%%%%%%%%%

Whenever every pair of compatible continuations admits a composite, one obtains a semigroup realization.

The semigroup operation is precisely continuation composition.

Nothing further is required.

%%%%%%%%%%%%%%%%%%%%%%%%%%%%%%%%%%%%%%%%%%%%%%%%%%%%%%%%%%%%%%%%%%%%%%%%%%%%%%%
\section{Realization of monoids}
%%%%%%%%%%%%%%%%%%%%%%%%%%%%%%%%%%%%%%%%%%%%%%%%%%%%%%%%%%%%%%%%%%%%%%%%%%%%%%%

If the continuation system possesses the identity continuation, the semigroup becomes a monoid.

The identity is therefore not an added algebraic object.

It is the empty continuation.

%%%%%%%%%%%%%%%%%%%%%%%%%%%%%%%%%%%%%%%%%%%%%%%%%%%%%%%%%%%%%%%%%%%%%%%%%%%%%%%
\section{Realization of groups}
%%%%%%%%%%%%%%%%%%%%%%%%%%%%%%%%%%%%%%%%%%%%%%%%%%%%%%%%%%%%%%%%%%%%%%%%%%%%%%%

If every continuation admits a reverse continuation preserving realization, the monoid becomes a group.

Groups therefore describe perfectly reversible continuation.

They are exceptional rather than generic.

%%%%%%%%%%%%%%%%%%%%%%%%%%%%%%%%%%%%%%%%%%%%%%%%%%%%%%%%%%%%%%%%%%%%%%%%%%%%%%%
\section{Realization of rings}
%%%%%%%%%%%%%%%%%%%%%%%%%%%%%%%%%%%%%%%%%%%%%%%%%%%%%%%%%%%%%%%%%%%%%%%%%%%%%%%

Two fundamentally different continuation mechanisms may coexist.

One composes continuations sequentially.

The other combines independent continuations.

The interaction between these two modes recovers ring structure.

Sequential continuation becomes multiplication.

Independent continuation becomes addition.

Distributivity expresses the compatibility of these two continuation mechanisms.

%%%%%%%%%%%%%%%%%%%%%%%%%%%%%%%%%%%%%%%%%%%%%%%%%%%%%%%%%%%%%%%%%%%%%%%%%%%%%%%
\section{Realization of fields}
%%%%%%%%%%%%%%%%%%%%%%%%%%%%%%%%%%%%%%%%%%%%%%%%%%%%%%%%%%%%%%%%%%%%%%%%%%%%%%%

When every non-trivial continuation admits normalization by inversion, the ring realizes a field.

Fields therefore represent continuation systems possessing complete algebraic solvability.

%%%%%%%%%%%%%%%%%%%%%%%%%%%%%%%%%%%%%%%%%%%%%%%%%%%%%%%%%%%%%%%%%%%%%%%%%%%%%%%
\section{Universal Recovery Theorem}
%%%%%%%%%%%%%%%%%%%%%%%%%%%%%%%%%%%%%%%%%%%%%%%%%%%%%%%%%%%%%%%%%%%%%%%%%%%%%%%

\begin{theorem}[Recovery of Algebra]

Every classical algebraic structure is realized by a continuation system whose admissible continuation composition reproduces its operations.

Conversely every continuation system possessing stable composition realizes an algebraic structure determined entirely by continuation.

\end{theorem}

\begin{proof}

The preceding sections recover semigroups, monoids, groups, rings, and fields from progressively stronger continuation principles.

Conversely every algebraic structure defines admissible compositions of partial objects.

Thus the correspondence is complete.

\end{proof}

%%%%%%%%%%%%%%%%%%%%%%%%%%%%%%%%%%%%%%%%%%%%%%%%%%%%%%%%%%%%%%%%%%%%%%%%%%%%%%%
\section{The meaning of algebra}
%%%%%%%%%%%%%%%%%%%%%%%%%%%%%%%%%%%%%%%%%%%%%%%%%%%%%%%%%%%%%%%%%%%%%%%%%%%%%%%

Algebra is therefore not fundamentally the study of operations on sets.

It is the study of stable composition among admissible continuations.

Operations are realized continuation laws.

Associativity expresses continuation coherence.

Identity expresses trivial continuation.

Inverse expresses reversibility.

Closure expresses completion stability.

Every classical algebraic object is recovered as one realization of a more primitive continuation structure.

Thus algebra is no longer foundational.

Continuation is.

%%%%%%%%%%%%%%%%%%%%%%%%%%%%%%%%%%%%%%%%%%%%%%%%%%%%%%%%%%%%%%%%%%%%%%%%%%%%%%
\chapter{Recovery of Geometry}
%%%%%%%%%%%%%%%%%%%%%%%%%%%%%%%%%%%%%%%%%%%%%%%%%%%%%%%%%%%%%%%%%%%%%%%%%%%%%%

Geometry is traditionally introduced by postulating points together with
primitive geometric relations such as incidence, distance, angle, or
neighbourhood. Different geometries arise by replacing one collection of
axioms with another.

Continuation Mathematics approaches geometry differently.

The primitive object is not a point but a continuation system. Geometry
appears only after the admissible continuations of partial objects have been
organized into sufficiently rich completion spaces.

The purpose of this chapter is to show that geometric objects are recovered as
realizations of continuation structure rather than assumed primitives.

%%%%%%%%%%%%%%%%%%%%%%%%%%%%%%%%%%%%%%%%%%%%%%%%%%%%%%%%%%%%%%%%%%%%%%%%%%%%%%
\section{Geometry as continuation organization}
%%%%%%%%%%%%%%%%%%%%%%%%%%%%%%%%%%%%%%%%%%%%%%%%%%%%%%%%%%%%%%%%%%%%%%%%%%%%%%

A continuation space contains information about which partial objects may
continue together and which cannot.

This information determines an intrinsic organization of the space.

\begin{definition}[Geometric realization]
A \emph{geometric realization} of a continuation space is a realization in
which continuation relations admit an intrinsic spatial interpretation.
\end{definition}

Thus geometry is not primitive.

It is one possible realization of continuation.

%%%%%%%%%%%%%%%%%%%%%%%%%%%%%%%%%%%%%%%%%%%%%%%%%%%%%%%%%%%%%%%%%%%%%%%%%%%%%%
\section{Neighbourhoods from continuation}
%%%%%%%%%%%%%%%%%%%%%%%%%%%%%%%%%%%%%%%%%%%%%%%%%%%%%%%%%%%%%%%%%%%%%%%%%%%%%%

Let $X$ be a continuation space.

Two elements should be regarded as geometrically close whenever they admit
many common continuations.

This motivates the following notion.

\begin{definition}[Continuation neighbourhood]
The continuation neighbourhood of an element $x\in X$ is the collection
\[
N(x)
=
\{y\in X:
x\text{ and }y
\text{ possess sufficiently compatible continuations}\}.
\]
\end{definition}

Neighbourhoods therefore emerge directly from admissible continuation rather
than from topology.

Topology itself was recovered in the previous chapter.

Geometry now enriches that topology by measuring continuation compatibility.

%%%%%%%%%%%%%%%%%%%%%%%%%%%%%%%%%%%%%%%%%%%%%%%%%%%%%%%%%%%%%%%%%%%%%%%%%%%%%%
\section{Distance from continuation effort}
%%%%%%%%%%%%%%%%%%%%%%%%%%%%%%%%%%%%%%%%%%%%%%%%%%%%%%%%%%%%%%%%%%%%%%%%%%%%%%

Suppose two objects possess different continuation histories.

One may ask how difficult it is to transform one history into another while
preserving admissibility.

This naturally produces a notion of distance.

\begin{definition}[Continuation distance]
A continuation distance is a function

\[
d:X\times X\rightarrow[0,\infty]
\]

whose value measures the minimal continuation effort required to pass from one
element to another.
\end{definition}

Different continuation systems generate different metrics.

The metric is therefore not postulated.

It is recovered.

%%%%%%%%%%%%%%%%%%%%%%%%%%%%%%%%%%%%%%%%%%%%%%%%%%%%%%%%%%%%%%%%%%%%%%%%%%%%%%
\section{Geodesics}
%%%%%%%%%%%%%%%%%%%%%%%%%%%%%%%%%%%%%%%%%%%%%%%%%%%%%%%%%%%%%%%%%%%%%%%%%%%%%%

Once continuation distance exists, shortest continuation paths become
meaningful.

\begin{definition}[Continuation geodesic]
A continuation geodesic is a continuation path realizing minimal continuation
effort between its endpoints.
\end{definition}

Classical geodesics are recovered whenever continuation effort coincides with
ordinary metric length.

Thus straight lines become optimal continuation processes.

%%%%%%%%%%%%%%%%%%%%%%%%%%%%%%%%%%%%%%%%%%%%%%%%%%%%%%%%%%%%%%%%%%%%%%%%%%%%%%
\section{Coordinate systems}
%%%%%%%%%%%%%%%%%%%%%%%%%%%%%%%%%%%%%%%%%%%%%%%%%%%%%%%%%%%%%%%%%%%%%%%%%%%%%%

Coordinates are often mistaken for geometry itself.

Continuation Mathematics distinguishes sharply between the two.

The continuation structure exists independently of every coordinate system.

Coordinates are merely realizations that encode continuation information.

Consequently every coordinate representation possesses intrinsic redundancy,
while the continuation structure remains invariant.

%%%%%%%%%%%%%%%%%%%%%%%%%%%%%%%%%%%%%%%%%%%%%%%%%%%%%%%%%%%%%%%%%%%%%%%%%%%%%%
\section{Intrinsic versus extrinsic geometry}
%%%%%%%%%%%%%%%%%%%%%%%%%%%%%%%%%%%%%%%%%%%%%%%%%%%%%%%%%%%%%%%%%%%%%%%%%%%%%%

Continuation Mathematics naturally separates two kinds of geometry.

\begin{definition}
An intrinsic geometric property depends only upon continuation relations.

An extrinsic geometric property depends upon a chosen realization.
\end{definition}

This distinction mirrors—but does not assume—the classical distinction between
intrinsic and embedded geometry.

%%%%%%%%%%%%%%%%%%%%%%%%%%%%%%%%%%%%%%%%%%%%%%%%%%%%%%%%%%%%%%%%%%%%%%%%%%%%%%
\section{Geometry as completion}
%%%%%%%%%%%%%%%%%%%%%%%%%%%%%%%%%%%%%%%%%%%%%%%%%%%%%%%%%%%%%%%%%%%%%%%%%%%%%%

Geometry becomes fully visible only after sufficiently many continuation
processes have been completed.

Partial continuation systems often exhibit only local geometric behaviour.

Completion reveals global geometry.

Thus geometric objects should be regarded as stable completion classes rather
than primitive configurations.

%%%%%%%%%%%%%%%%%%%%%%%%%%%%%%%%%%%%%%%%%%%%%%%%%%%%%%%%%%%%%%%%%%%%%%%%%%%%%%
\section{Recovering classical geometries}
%%%%%%%%%%%%%%%%%%%%%%%%%%%%%%%%%%%%%%%%%%%%%%%%%%%%%%%%%%%%%%%%%%%%%%%%%%%%%%

Different continuation systems generate different geometric realizations.

Examples include

\begin{itemize}
\item Euclidean geometry,
\item affine geometry,
\item projective geometry,
\item hyperbolic geometry,
\item spherical geometry,
\item Riemannian geometry,
\item symplectic geometry,
\item algebraic geometry.
\end{itemize}

Each is characterized by the continuation constraints imposed upon its
underlying completion space.

%%%%%%%%%%%%%%%%%%%%%%%%%%%%%%%%%%%%%%%%%%%%%%%%%%%%%%%%%%%%%%%%%%%%%%%%%%%%%%
\section{Universality}
%%%%%%%%%%%%%%%%%%%%%%%%%%%%%%%%%%%%%%%%%%%%%%%%%%%%%%%%%%%%%%%%%%%%%%%%%%%%%%

The fundamental object of geometry is therefore not the point.

Nor is it the manifold.

Nor is it the metric.

The primitive object is the continuation structure.

Points, neighbourhoods, metrics, curves, geodesics, dimensions, manifolds,
connections and geometric spaces are all realizations obtained after suitable
completion.

Thus Continuation Mathematics reverses the traditional logical order.

Classical mathematics begins with geometry and studies continuations within
it.

Continuation Mathematics begins with continuation itself and recovers geometry
as one of its realizations.

This completes the recovery of classical geometry from the universal theory of
continuation.

%%%%%%%%%%%%%%%%%%%%%%%%%%%%%%%%%%%%%%%%%%%%%%%%%%%%%%%%%%%%%%%%%%%%%%%%%%%%%%
\chapter{Recovery of Analysis}
%%%%%%%%%%%%%%%%%%%%%%%%%%%%%%%%%%%%%%%%%%%%%%%%%%%%%%%%%%%%%%%%%%%%%%%%%%%%%%

Analysis is traditionally founded upon limits.

From limits arise continuity, differentiation, integration, infinite series,
functional analysis, measure theory, and much of modern mathematics.

Continuation Mathematics proposes the opposite logical order.

Limits are not primitive.

They arise because continuation systems admit stable completion.

Analysis is therefore the mathematics of controlled completion.

%%%%%%%%%%%%%%%%%%%%%%%%%%%%%%%%%%%%%%%%%%%%%%%%%%%%%%%%%%%%%%%%%%%%%%%%%%%%%%
\section{The origin of limiting behaviour}
%%%%%%%%%%%%%%%%%%%%%%%%%%%%%%%%%%%%%%%%%%%%%%%%%%%%%%%%%%%%%%%%%%%%%%%%%%%%%%

A continuation process consists of a succession of admissible extensions

\[
x_0
\longrightarrow
x_1
\longrightarrow
x_2
\longrightarrow
\cdots.
\]

Such a process naturally raises a single question.

Does the process stabilize?

If stabilization occurs, then the continuation possesses a completion.

If stabilization fails, the continuation remains genuinely incomplete.

The entire subject of analysis begins with this distinction.

%%%%%%%%%%%%%%%%%%%%%%%%%%%%%%%%%%%%%%%%%%%%%%%%%%%%%%%%%%%%%%%%%%%%%%%%%%%%%%
\section{Limits as universal completions}
%%%%%%%%%%%%%%%%%%%%%%%%%%%%%%%%%%%%%%%%%%%%%%%%%%%%%%%%%%%%%%%%%%%%%%%%%%%%%%

The notion of limit is therefore secondary.

\begin{definition}[Continuation limit]
Let
\[
(x_n)_{n\ge0}
\]
be a continuation sequence inside a continuation space.

A \emph{continuation limit} is a universal completion of this sequence whenever
such a completion exists.
\end{definition}

Thus a limit is not defined by epsilon-delta inequalities.

It is defined structurally as the terminal object completing a continuation.

The classical epsilon-delta definition is one realization of this more
primitive notion.

%%%%%%%%%%%%%%%%%%%%%%%%%%%%%%%%%%%%%%%%%%%%%%%%%%%%%%%%%%%%%%%%%%%%%%%%%%%%%%
\section{Continuity}
%%%%%%%%%%%%%%%%%%%%%%%%%%%%%%%%%%%%%%%%%%%%%%%%%%%%%%%%%%%%%%%%%%%%%%%%%%%%%%

Functions preserve mathematical structure only when they preserve admissible
continuations.

\begin{definition}[Continuation continuity]
A realization map

\[
f:X\rightarrow Y
\]

is continuation-continuous whenever every convergent continuation process in
\(X\) is carried into a convergent continuation process in \(Y\), with
completion preserved.
\end{definition}

Continuity is therefore preservation of completion.

Topology supplied the notion of neighbourhood.

Analysis strengthens this by requiring neighbourhood behaviour to survive
arbitrary continuation.

%%%%%%%%%%%%%%%%%%%%%%%%%%%%%%%%%%%%%%%%%%%%%%%%%%%%%%%%%%%%%%%%%%%%%%%%%%%%%%
\section{Differentiation}
%%%%%%%%%%%%%%%%%%%%%%%%%%%%%%%%%%%%%%%%%%%%%%%%%%%%%%%%%%%%%%%%%%%%%%%%%%%%%%

Differentiation measures the first-order behaviour of continuation.

Suppose a continuation process approaches a completed object.

The derivative records the universal linear approximation governing the local
continuation.

\begin{definition}[Continuation derivative]
A continuation derivative is the universal first-order realization governing
the local completion of a continuation process.
\end{definition}

Thus the derivative is not fundamentally a quotient.

The quotient formula is one realization of the intrinsic continuation
derivative.

%%%%%%%%%%%%%%%%%%%%%%%%%%%%%%%%%%%%%%%%%%%%%%%%%%%%%%%%%%%%%%%%%%%%%%%%%%%%%%
\section{Integration}
%%%%%%%%%%%%%%%%%%%%%%%%%%%%%%%%%%%%%%%%%%%%%%%%%%%%%%%%%%%%%%%%%%%%%%%%%%%%%%

Differentiation analyzes local continuation.

Integration reconstructs global continuation.

\begin{definition}[Continuation integral]
A continuation integral is the universal completion obtained by assembling
compatible local continuations into a single completed object.
\end{definition}

The Fundamental Theorem of Calculus expresses the compatibility of these two
universal constructions.

Differentiation decomposes completion.

Integration reconstructs completion.

%%%%%%%%%%%%%%%%%%%%%%%%%%%%%%%%%%%%%%%%%%%%%%%%%%%%%%%%%%%%%%%%%%%%%%%%%%%%%%
\section{Infinite series}
%%%%%%%%%%%%%%%%%%%%%%%%%%%%%%%%%%%%%%%%%%%%%%%%%%%%%%%%%%%%%%%%%%%%%%%%%%%%%%

Infinite series are continuation processes generated by repeated composition.

Partial sums form finite continuation objects.

Their completion, when it exists, defines the value of the series.

Convergence therefore becomes a special case of completion.

Divergence records the failure of admissible completion.

%%%%%%%%%%%%%%%%%%%%%%%%%%%%%%%%%%%%%%%%%%%%%%%%%%%%%%%%%%%%%%%%%%%%%%%%%%%%%%
\section{Function spaces}
%%%%%%%%%%%%%%%%%%%%%%%%%%%%%%%%%%%%%%%%%%%%%%%%%%%%%%%%%%%%%%%%%%%%%%%%%%%%%%

Functions themselves admit continuation.

One may therefore consider spaces whose elements are continuation-preserving
maps.

Such spaces inherit continuation structure from their constituent maps.

Function spaces become higher-order continuation spaces.

Analysis of operators is therefore analysis of continuation between
continuations.

%%%%%%%%%%%%%%%%%%%%%%%%%%%%%%%%%%%%%%%%%%%%%%%%%%%%%%%%%%%%%%%%%%%%%%%%%%%%%%
\section{Measure}
%%%%%%%%%%%%%%%%%%%%%%%%%%%%%%%%%%%%%%%%%%%%%%%%%%%%%%%%%%%%%%%%%%%%%%%%%%%%%%

Classically, measure assigns size.

Continuation Mathematics identifies a more primitive concept.

Before assigning size, one must first determine which continuations are
admissible.

Measure quantifies the abundance of admissible continuation.

\begin{definition}[Continuation measure]
A continuation measure assigns quantitative weight to admissible continuation
families while remaining invariant under equivalent continuation
decompositions.
\end{definition}

Lebesgue measure, Hausdorff measure, and probability measures are recovered as
realizations of continuation measure.

%%%%%%%%%%%%%%%%%%%%%%%%%%%%%%%%%%%%%%%%%%%%%%%%%%%%%%%%%%%%%%%%%%%%%%%%%%%%%%
\section{Completeness}
%%%%%%%%%%%%%%%%%%%%%%%%%%%%%%%%%%%%%%%%%%%%%%%%%%%%%%%%%%%%%%%%%%%%%%%%%%%%%%

Completeness is one of the central ideas of analysis.

Continuation Mathematics explains its origin.

\begin{definition}[Analytic completeness]
A continuation space is analytically complete whenever every admissible Cauchy
continuation possesses a universal completion.
\end{definition}

Thus completeness is not an axiom.

It is a structural property of continuation spaces.

Metric completeness, Banach spaces, Hilbert spaces, and complete lattices are
all manifestations of this single continuation principle.

%%%%%%%%%%%%%%%%%%%%%%%%%%%%%%%%%%%%%%%%%%%%%%%%%%%%%%%%%%%%%%%%%%%%%%%%%%%%%%
\section{Recovery of classical analysis}
%%%%%%%%%%%%%%%%%%%%%%%%%%%%%%%%%%%%%%%%%%%%%%%%%%%%%%%%%%%%%%%%%%%%%%%%%%%%%%

The principal structures of classical analysis now appear naturally.

\begin{itemize}
\item Limits become universal completions.
\item Continuity becomes preservation of completion.
\item Differentiation becomes local continuation.
\item Integration becomes global completion.
\item Infinite series become repeated continuation.
\item Function spaces become higher-order continuation spaces.
\item Measure becomes quantitative continuation.
\item Completeness becomes universal completion.
\end{itemize}

Thus the entire logical architecture of analysis is recovered without assuming
limits as primitive objects.

%%%%%%%%%%%%%%%%%%%%%%%%%%%%%%%%%%%%%%%%%%%%%%%%%%%%%%%%%%%%%%%%%%%%%%%%%%%%%%
\section{The Continuation Principle of Analysis}
%%%%%%%%%%%%%%%%%%%%%%%%%%%%%%%%%%%%%%%%%%%%%%%%%%%%%%%%%%%%%%%%%%%%%%%%%%%%%%

The preceding development suggests a single foundational principle.

\begin{theorem}[Continuation Principle of Analysis]
Every analytic concept is generated by one of three universal continuation
operations:

\begin{enumerate}
    \item local continuation,
    \item global completion,
    \item quantitative continuation.
\end{enumerate}

Every branch of classical analysis is obtained by realizing one or more of
these operations within an appropriate continuation universe.
\end{theorem}

This principle explains why apparently unrelated branches of analysis possess
the same structural phenomena: convergence, approximation, stability,
compactness, completeness, continuity, and decomposition are all manifestations
of the mathematics of continuation.

Analysis therefore ceases to be the study of limits.

It becomes the universal mathematics of completion.

This completes the recovery of classical analysis from the theory of
continuation.

%%%%%%%%%%%%%%%%%%%%%%%%%%%%%%%%%%%%%%%%%%%%%%%%%%%%%%%%%%%%%%%%%%%%%%%%%%%%%%
\chapter{Recovery of Probability}
%%%%%%%%%%%%%%%%%%%%%%%%%%%%%%%%%%%%%%%%%%%%%%%%%%%%%%%%%%%%%%%%%%%%%%%%%%%%%%

Probability is traditionally introduced by assigning numerical weights to
events.

Different foundations begin with different primitive notions:

\begin{itemize}
    \item frequencies,
    \item equally likely outcomes,
    \item measure spaces,
    \item subjective belief,
    \item axiomatic probability.
\end{itemize}

Continuation Mathematics adopts none of these as primitive.

Instead, probability arises from incomplete continuation.

Whenever the future continuation of a partial mathematical object is not yet
uniquely determined, one obtains a continuation spectrum.

Probability quantifies that spectrum.

Thus probability is not primitive uncertainty.

It is the mathematics of quantitative continuation.

%%%%%%%%%%%%%%%%%%%%%%%%%%%%%%%%%%%%%%%%%%%%%%%%%%%%%%%%%%%%%%%%%%%%%%%%%%%%%%
\section{Continuation uncertainty}
%%%%%%%%%%%%%%%%%%%%%%%%%%%%%%%%%%%%%%%%%%%%%%%%%%%%%%%%%%%%%%%%%%%%%%%%%%%%%%

Suppose a continuation object possesses several admissible extensions

\[
x
\longrightarrow
x_1,
\qquad
x
\longrightarrow
x_2,
\qquad
\ldots,
\qquad
x
\longrightarrow
x_n.
\]

Nothing in the continuation structure requires one continuation to be selected
before another.

The object therefore possesses genuine continuation uncertainty.

\begin{definition}[Continuation uncertainty]
Continuation uncertainty is the existence of more than one admissible
continuation extending the same partial object.
\end{definition}

Probability begins precisely here.

%%%%%%%%%%%%%%%%%%%%%%%%%%%%%%%%%%%%%%%%%%%%%%%%%%%%%%%%%%%%%%%%%%%%%%%%%%%%%%
\section{Continuation spectra}
%%%%%%%%%%%%%%%%%%%%%%%%%%%%%%%%%%%%%%%%%%%%%%%%%%%%%%%%%%%%%%%%%%%%%%%%%%%%%%

Every partial object determines the collection of all admissible
continuations.

\begin{definition}[Continuation spectrum]
The continuation spectrum of a partial object is the collection of every
admissible continuation extending that object.
\end{definition}

The spectrum itself is entirely deterministic.

What is uncertain is which continuation will ultimately be realized.

Probability therefore measures the structure of the spectrum rather than the
existence of randomness.

%%%%%%%%%%%%%%%%%%%%%%%%%%%%%%%%%%%%%%%%%%%%%%%%%%%%%%%%%%%%%%%%%%%%%%%%%%%%%%
\section{Weights on continuation}
%%%%%%%%%%%%%%%%%%%%%%%%%%%%%%%%%%%%%%%%%%%%%%%%%%%%%%%%%%%%%%%%%%%%%%%%%%%%%%

Continuation theory itself merely determines admissibility.

It does not prescribe quantitative preference.

A realization may therefore assign weights.

\begin{definition}[Continuation weight]
A continuation weight assigns a non-negative numerical value to each
admissible continuation while respecting continuation equivalence.
\end{definition}

Different realizations produce different probability theories.

Uniform probability is only one realization among many.

%%%%%%%%%%%%%%%%%%%%%%%%%%%%%%%%%%%%%%%%%%%%%%%%%%%%%%%%%%%%%%%%%%%%%%%%%%%%%%
\section{Probability as normalized continuation}
%%%%%%%%%%%%%%%%%%%%%%%%%%%%%%%%%%%%%%%%%%%%%%%%%%%%%%%%%%%%%%%%%%%%%%%%%%%%%%

Whenever the total continuation weight is finite, normalization becomes
possible.

\begin{definition}[Continuation probability]
A continuation probability is a normalized continuation weight satisfying

\[
P(C)
=
\frac{w(C)}
{\displaystyle\sum_{D}w(D)},
\]

where the sum ranges over every admissible continuation of the same partial
object.
\end{definition}

Thus the probability axioms are consequences of normalization.

They are not primitive assumptions.

%%%%%%%%%%%%%%%%%%%%%%%%%%%%%%%%%%%%%%%%%%%%%%%%%%%%%%%%%%%%%%%%%%%%%%%%%%%%%%
\section{Conditional probability}
%%%%%%%%%%%%%%%%%%%%%%%%%%%%%%%%%%%%%%%%%%%%%%%%%%%%%%%%%%%%%%%%%%%%%%%%%%%%%%

As continuation proceeds, some continuations become impossible while others
remain admissible.

Probability therefore changes.

\begin{definition}[Conditional continuation]
Conditional probability is the renormalization of continuation weights after
restricting to the surviving admissible continuations.
\end{definition}

Bayesian updating becomes a special realization of this universal continuation
operation.

Learning is continuation refinement.

%%%%%%%%%%%%%%%%%%%%%%%%%%%%%%%%%%%%%%%%%%%%%%%%%%%%%%%%%%%%%%%%%%%%%%%%%%%%%%
\section{Expectation}
%%%%%%%%%%%%%%%%%%%%%%%%%%%%%%%%%%%%%%%%%%%%%%%%%%%%%%%%%%%%%%%%%%%%%%%%%%%%%%

Suppose every continuation carries a numerical observable.

The observable need not be deterministic.

Its average over the continuation spectrum defines expectation.

\begin{definition}[Continuation expectation]
The continuation expectation of an observable is its weighted average over the
entire continuation spectrum.
\end{definition}

Expectation therefore measures average continuation behaviour.

%%%%%%%%%%%%%%%%%%%%%%%%%%%%%%%%%%%%%%%%%%%%%%%%%%%%%%%%%%%%%%%%%%%%%%%%%%%%%%
\section{Random variables}
%%%%%%%%%%%%%%%%%%%%%%%%%%%%%%%%%%%%%%%%%%%%%%%%%%%%%%%%%%%%%%%%%%%%%%%%%%%%%%

A random variable is not fundamentally random.

It is an observable defined upon a continuation spectrum.

Different continuations produce different observable values.

The randomness belongs to continuation selection rather than to the observable
itself.

%%%%%%%%%%%%%%%%%%%%%%%%%%%%%%%%%%%%%%%%%%%%%%%%%%%%%%%%%%%%%%%%%%%%%%%%%%%%%%
\section{Independence}
%%%%%%%%%%%%%%%%%%%%%%%%%%%%%%%%%%%%%%%%%%%%%%%%%%%%%%%%%%%%%%%%%%%%%%%%%%%%%%

Classically, independence is defined probabilistically.

Continuation Mathematics explains its origin.

\begin{definition}[Continuation independence]
Two continuation processes are independent whenever the admissible
continuations of one process impose no continuation constraints upon the
other.
\end{definition}

Probabilistic independence is recovered after normalization.

The structural notion precedes the numerical one.

%%%%%%%%%%%%%%%%%%%%%%%%%%%%%%%%%%%%%%%%%%%%%%%%%%%%%%%%%%%%%%%%%%%%%%%%%%%%%%
\section{Stochastic processes}
%%%%%%%%%%%%%%%%%%%%%%%%%%%%%%%%%%%%%%%%%%%%%%%%%%%%%%%%%%%%%%%%%%%%%%%%%%%%%%

Repeated continuation generates stochastic processes.

Each successive continuation enlarges the continuation history while
restricting future admissible continuations.

Markov processes arise whenever future admissibility depends only upon the
present continuation state.

More general stochastic processes arise whenever continuation memory is
retained.

%%%%%%%%%%%%%%%%%%%%%%%%%%%%%%%%%%%%%%%%%%%%%%%%%%%%%%%%%%%%%%%%%%%%%%%%%%%%%%
\section{Recovery of classical probability}
%%%%%%%%%%%%%%%%%%%%%%%%%%%%%%%%%%%%%%%%%%%%%%%%%%%%%%%%%%%%%%%%%%%%%%%%%%%%%%

The principal concepts of probability are therefore recovered naturally.

\begin{itemize}

\item Sample spaces become continuation spectra.

\item Events become continuation subsets.

\item Probability measures become normalized continuation weights.

\item Conditional probability becomes continuation refinement.

\item Bayes' theorem becomes iterative continuation updating.

\item Random variables become continuation observables.

\item Expectation becomes average continuation behaviour.

\item Independence becomes structural continuation independence.

\item Stochastic processes become evolving continuation systems.

\end{itemize}

Thus probability is no longer the mathematics of randomness.

It is the mathematics of incomplete continuation.

%%%%%%%%%%%%%%%%%%%%%%%%%%%%%%%%%%%%%%%%%%%%%%%%%%%%%%%%%%%%%%%%%%%%%%%%%%%%%%
\section{Objective and epistemic probability}
%%%%%%%%%%%%%%%%%%%%%%%%%%%%%%%%%%%%%%%%%%%%%%%%%%%%%%%%%%%%%%%%%%%%%%%%%%%%%%

One of the oldest questions in probability concerns the meaning of
probability itself.

Is probability objective?

Or is it merely a measure of knowledge?

Continuation Mathematics separates these notions.

The continuation spectrum is objective.

It is determined entirely by the continuation structure.

The weights assigned to that spectrum belong to a realization.

Some realizations describe physical frequencies.

Others describe logical uncertainty.

Others describe information.

Still others describe quantum amplitudes.

The underlying continuation object is identical.

Only the realization differs.

Thus the classical philosophical debate concerning the interpretation of
probability is dissolved.

Different interpretations correspond to different realizations of the same
continuation spectrum.

%%%%%%%%%%%%%%%%%%%%%%%%%%%%%%%%%%%%%%%%%%%%%%%%%%%%%%%%%%%%%%%%%%%%%%%%%%%%%%
\section{The Continuation Principle of Probability}
%%%%%%%%%%%%%%%%%%%%%%%%%%%%%%%%%%%%%%%%%%%%%%%%%%%%%%%%%%%%%%%%%%%%%%%%%%%%%%

Probability is not primitive.

Nor is randomness.

The primitive object is admissible continuation.

Whenever continuation is not uniquely determined, a continuation spectrum
appears.

Whenever that spectrum admits quantitative realization, probability emerges.

Therefore probability is neither an axiom nor an interpretation.

It is a realization of continuation.

This completes the recovery of classical probability from the universal theory
of continuation.

\chapter{Recovery of Information Theory}

Information theory is ordinarily introduced by postulating alphabets, probability
distributions, messages, and entropy functions. Within Continuation Mathematics,
none of these notions are primitive.

Instead, information is recovered from the structure of continuation itself.

The fundamental observation is that continuation distinguishes possibilities.
Whenever an object possesses multiple admissible continuations, there exists
structural uncertainty. Whenever continuation collapses to a unique completion,
that uncertainty disappears.

Information is therefore not an external quantity attached to an object.
It is an intrinsic property of continuation spaces.

%%%%%%%%%%%%%%%%%%%%%%%%%%%%%%%%%%%%%%%%%%%%%%%%%%%%%%%%%%%%%%%%%%%%%%%%%%%%%%%
\section{Structural Information}
%%%%%%%%%%%%%%%%%%%%%%%%%%%%%%%%%%%%%%%%%%%%%%%%%%%%%%%%%%%%%%%%%%%%%%%%%%%%%%%

Continuation spaces naturally measure how many futures remain available.

\begin{definition}[Structural information]
Let \(X\) be a continuation space.

The \emph{structural information} of an object \(x\in X\) is the mathematical
content contained in the admissible continuation class
\[
\mathcal C(x).
\]

Information therefore measures the richness of admissible continuation rather
than any external encoding.
\end{definition}

Thus information is fundamentally geometric.

Objects possessing identical continuation classes contain identical information.

Objects whose continuation classes differ are informationally distinct.

%%%%%%%%%%%%%%%%%%%%%%%%%%%%%%%%%%%%%%%%%%%%%%%%%%%%%%%%%%%%%%%%%%%%%%%%%%%%%%%
\section{Distinguishability}
%%%%%%%%%%%%%%%%%%%%%%%%%%%%%%%%%%%%%%%%%%%%%%%%%%%%%%%%%%%%%%%%%%%%%%%%%%%%%%%

The first role of information is distinction.

\begin{definition}[Continuation distinguishability]
Two objects \(x,y\) are called
\emph{continuation distinguishable}
whenever
\[
\mathcal C(x)\neq\mathcal C(y).
\]
\end{definition}

No reference to probability has appeared.

Information begins simply as the ability of continuation to separate objects.

This is considerably more primitive than Shannon's communication viewpoint.

%%%%%%%%%%%%%%%%%%%%%%%%%%%%%%%%%%%%%%%%%%%%%%%%%%%%%%%%%%%%%%%%%%%%%%%%%%%%%%%
\section{Information Refinement}
%%%%%%%%%%%%%%%%%%%%%%%%%%%%%%%%%%%%%%%%%%%%%%%%%%%%%%%%%%%%%%%%%%%%%%%%%%%%%%%

Continuation proceeds by successive refinement.

Each admissible continuation either

\[
\mathcal C(x)
\longrightarrow
\mathcal C(y),
\]

or eliminates possible futures.

Thus continuation naturally induces an ordering of informational states.

\begin{definition}[Information refinement]
If every admissible continuation of \(y\) is also an admissible continuation of
\(x\), then \(y\) is called an
\emph{informational refinement}
of \(x\).
\end{definition}

Refinement decreases structural uncertainty.

Completion represents maximal refinement.

%%%%%%%%%%%%%%%%%%%%%%%%%%%%%%%%%%%%%%%%%%%%%%%%%%%%%%%%%%%%%%%%%%%%%%%%%%%%%%%
\section{Information Order}
%%%%%%%%%%%%%%%%%%%%%%%%%%%%%%%%%%%%%%%%%%%%%%%%%%%%%%%%%%%%%%%%%%%%%%%%%%%%%%%

The refinement relation defines an intrinsic order.

\[
x
\preceq_I
y
\]

whenever every continuation of \(y\) is already forced by \(x\).

This ordering satisfies

\begin{itemize}
\item reflexivity,
\item antisymmetry,
\item transitivity.
\end{itemize}

Hence information naturally forms a partially ordered system.

The order is recovered from continuation alone.

%%%%%%%%%%%%%%%%%%%%%%%%%%%%%%%%%%%%%%%%%%%%%%%%%%%%%%%%%%%%%%%%%%%%%%%%%%%%%%%
\section{Completion as Perfect Information}
%%%%%%%%%%%%%%%%%%%%%%%%%%%%%%%%%%%%%%%%%%%%%%%%%%%%%%%%%%%%%%%%%%%%%%%%%%%%%%%

Completion removes every remaining ambiguity.

\begin{definition}[Perfect information]
A completed object possesses
\emph{perfect information}
whenever its continuation class consists solely of itself.

Equivalently,

\[
\mathcal C(x)=\{x\}.
\]
\end{definition}

Perfect information therefore corresponds precisely to completed continuation.

Nothing remains undecided.

%%%%%%%%%%%%%%%%%%%%%%%%%%%%%%%%%%%%%%%%%%%%%%%%%%%%%%%%%%%%%%%%%%%%%%%%%%%%%%%
\section{Structural Entropy}
%%%%%%%%%%%%%%%%%%%%%%%%%%%%%%%%%%%%%%%%%%%%%%%%%%%%%%%%%%%%%%%%%%%%%%%%%%%%%%%

Entropy is not primitive.

Instead it measures the size or complexity of continuation.

\begin{definition}[Structural entropy]
A structural entropy is any invariant

\[
H:
\mathsf{Cont}
\rightarrow
[0,\infty]
\]

satisfying

\[
\mathcal C(x)\subseteq\mathcal C(y)
\Longrightarrow
H(x)\le H(y).
\]
\end{definition}

Thus entropy measures admissible continuation.

It is not fundamentally probabilistic.

Probability enters only after realizations of continuation have been chosen.

%%%%%%%%%%%%%%%%%%%%%%%%%%%%%%%%%%%%%%%%%%%%%%%%%%%%%%%%%%%%%%%%%%%%%%%%%%%%%%%
\section{Recovery of Shannon Entropy}
%%%%%%%%%%%%%%%%%%%%%%%%%%%%%%%%%%%%%%%%%%%%%%%%%%%%%%%%%%%%%%%%%%%%%%%%%%%%%%%

Suppose a realization interprets admissible continuations as mutually exclusive
symbols with probabilities

\[
p_1,\ldots,p_n.
\]

The continuation entropy specializes to

\[
H
=
-\sum_i p_i\log p_i.
\]

Thus Shannon entropy appears as one particular realization of structural
continuation entropy.

The logarithm is therefore not foundational.

It is a coordinate expression arising after probabilistic realization.

%%%%%%%%%%%%%%%%%%%%%%%%%%%%%%%%%%%%%%%%%%%%%%%%%%%%%%%%%%%%%%%%%%%%%%%%%%%%%%%
\section{Mutual Information}
%%%%%%%%%%%%%%%%%%%%%%%%%%%%%%%%%%%%%%%%%%%%%%%%%%%%%%%%%%%%%%%%%%%%%%%%%%%%%%%

Two continuation systems may constrain one another.

\begin{definition}[Mutual continuation]
The mutual continuation between systems \(X\) and \(Y\) is the common reduction
of admissible continuation produced by considering them jointly.
\end{definition}

Information shared between systems therefore measures shared continuation.

Classical mutual information becomes a numerical realization of this structural
intersection.

%%%%%%%%%%%%%%%%%%%%%%%%%%%%%%%%%%%%%%%%%%%%%%%%%%%%%%%%%%%%%%%%%%%%%%%%%%%%%%%
\section{Communication}
%%%%%%%%%%%%%%%%%%%%%%%%%%%%%%%%%%%%%%%%%%%%%%%%%%%%%%%%%%%%%%%%%%%%%%%%%%%%%%%

Communication may now be reconstructed.

A communication channel is simply a continuation morphism

\[
F:X\rightarrow Y
\]

that preserves admissible continuation.

Noise corresponds to morphisms that enlarge continuation frontiers.

Compression corresponds to equivalent continuation representations possessing
smaller realizations.

Error correction corresponds to recovery of lost continuation through structural
redundancy.

Thus the principal concepts of classical information theory become structural
properties of continuation morphisms.

%%%%%%%%%%%%%%%%%%%%%%%%%%%%%%%%%%%%%%%%%%%%%%%%%%%%%%%%%%%%%%%%%%%%%%%%%%%%%%%
\section{The Recovery Theorem}
%%%%%%%%%%%%%%%%%%%%%%%%%%%%%%%%%%%%%%%%%%%%%%%%%%%%%%%%%%%%%%%%%%%%%%%%%%%%%%%

\begin{theorem}[Recovery of Information Theory]
Classical information theory is a realization of continuation theory.

Specifically,

\begin{enumerate}
\item information is recovered from continuation classes;

\item distinguishability from continuation inequivalence;

\item entropy from measures of continuation complexity;

\item communication from continuation-preserving morphisms;

\item mutual information from shared continuation;

\item perfect information from completion.
\end{enumerate}

Consequently information theory is not primitive mathematics.

It is a realization of the deeper mathematics of continuation.
\end{theorem}

\begin{proof}
Each classical notion is obtained by choosing a realization that interprets
continuation classes as messages, admissible continuations as possible symbol
extensions, completion as perfect decoding, and structural entropy as Shannon's
entropy functional.

The familiar theory is therefore recovered functorially from continuation
mathematics rather than postulated independently.
\end{proof}

%%%%%%%%%%%%%%%%%%%%%%%%%%%%%%%%%%%%%%%%%%%%%%%%%%%%%%%%%%%%%%%%%%%%%%%%%%%%%%%
\section{Foundational Significance}
%%%%%%%%%%%%%%%%%%%%%%%%%%%%%%%%%%%%%%%%%%%%%%%%%%%%%%%%%%%%%%%%%%%%%%%%%%%%%%%

The recovery of information theory reveals that information is not fundamentally
about symbols, messages, or probability.

Those notions belong to particular realizations.

At the universal level, information is simply the mathematical structure carried
by admissible continuation.

Distinction, uncertainty, entropy, communication, and decoding all arise because
continuation possesses internal organization before any probabilistic or
computational interpretation is introduced.

Continuation Mathematics therefore identifies continuation—not information—as the
more primitive mathematical object.

\chapter{Recovery of Category Theory}

Category theory is often regarded as a foundational language of modern
mathematics. Objects are introduced together with morphisms between them,
composition is postulated, identities are assumed, and all subsequent structure
is developed from these primitives.

Continuation Mathematics reverses this order.

Objects, morphisms, composition, identity, limits, adjunctions, and functors
have already been constructed intrinsically from continuation. The purpose of
this chapter is therefore not to develop category theory, but to show that
ordinary category theory is recovered as a realization of continuation systems.

%%%%%%%%%%%%%%%%%%%%%%%%%%%%%%%%%%%%%%%%%%%%%%%%%%%%%%%%%%%%%%%%%%%%%%%%%%%%%%%
\section{Objects as Completion Classes}
%%%%%%%%%%%%%%%%%%%%%%%%%%%%%%%%%%%%%%%%%%%%%%%%%%%%%%%%%%%%%%%%%%%%%%%%%%%%%%%

The first primitive of category theory is the notion of object.

Within Continuation Mathematics no object exists in complete isolation.
Every mathematical entity exists only through its continuation structure.

Accordingly, the fundamental entities are not isolated objects but completion
classes.

\begin{definition}[Categorical realization object]
A \emph{categorical object} is the realization of a completion class
\[
[X]
\]
under a realization functor.
\end{definition}

Thus the objects of category theory are not primitive mathematical entities.

They arise from completed continuation structures.

%%%%%%%%%%%%%%%%%%%%%%%%%%%%%%%%%%%%%%%%%%%%%%%%%%%%%%%%%%%%%%%%%%%%%%%%%%%%%%%
\section{Morphisms from Continuation}
%%%%%%%%%%%%%%%%%%%%%%%%%%%%%%%%%%%%%%%%%%%%%%%%%%%%%%%%%%%%%%%%%%%%%%%%%%%%%%%

Continuation Morphisms were developed intrinsically in Part IV.

A realization identifies these morphisms with ordinary categorical arrows.

\begin{definition}[Categorical arrow]
Given completion classes
\[
[X],\qquad[Y],
\]
a categorical arrow
\[
[X]\longrightarrow[Y]
\]
is the realization of a continuation morphism preserving admissible
continuation.
\end{definition}

Thus arrows are inherited from continuation.

They are never postulated independently.

%%%%%%%%%%%%%%%%%%%%%%%%%%%%%%%%%%%%%%%%%%%%%%%%%%%%%%%%%%%%%%%%%%%%%%%%%%%%%%%
\section{Composition}
%%%%%%%%%%%%%%%%%%%%%%%%%%%%%%%%%%%%%%%%%%%%%%%%%%%%%%%%%%%%%%%%%%%%%%%%%%%%%%%

Continuation morphisms compose naturally.

If

\[
F:X\rightarrow Y,
\qquad
G:Y\rightarrow Z,
\]

preserve continuation, then so does

\[
G\circ F.
\]

Associativity has already been established in the intrinsic continuation
algebra.

Therefore categorical composition is recovered automatically.

%%%%%%%%%%%%%%%%%%%%%%%%%%%%%%%%%%%%%%%%%%%%%%%%%%%%%%%%%%%%%%%%%%%%%%%%%%%%%%%
\section{Identity Morphisms}
%%%%%%%%%%%%%%%%%%%%%%%%%%%%%%%%%%%%%%%%%%%%%%%%%%%%%%%%%%%%%%%%%%%%%%%%%%%%%%%

Every continuation system possesses the identity continuation.

Its realization becomes the identity morphism.

\[
\operatorname{id}_X:X\rightarrow X.
\]

The categorical identity law is therefore a realization of the intrinsic
identity continuation.

%%%%%%%%%%%%%%%%%%%%%%%%%%%%%%%%%%%%%%%%%%%%%%%%%%%%%%%%%%%%%%%%%%%%%%%%%%%%%%%
\section{Categories}
%%%%%%%%%%%%%%%%%%%%%%%%%%%%%%%%%%%%%%%%%%%%%%%%%%%%%%%%%%%%%%%%%%%%%%%%%%%%%%%

A continuation universe naturally determines a category.

\begin{definition}[Recovered category]
The category associated with a continuation universe consists of

\begin{enumerate}
\item realization objects;

\item continuation morphisms;

\item intrinsic composition;

\item intrinsic identities.
\end{enumerate}
\end{definition}

The category axioms are already consequences of continuation algebra.

Nothing new is added.

%%%%%%%%%%%%%%%%%%%%%%%%%%%%%%%%%%%%%%%%%%%%%%%%%%%%%%%%%%%%%%%%%%%%%%%%%%%%%%%
\section{Functors}
%%%%%%%%%%%%%%%%%%%%%%%%%%%%%%%%%%%%%%%%%%%%%%%%%%%%%%%%%%%%%%%%%%%%%%%%%%%%%%%

Functoriality also precedes category theory.

Continuation Functors were introduced in Part IV as mappings preserving
continuation structure.

Their realization produces ordinary categorical functors.

If

\[
F:\mathcal U_1\rightarrow\mathcal U_2
\]

preserves continuation, then

\[
F
\]

is simultaneously a categorical functor.

Thus functoriality is inherited from continuation.

%%%%%%%%%%%%%%%%%%%%%%%%%%%%%%%%%%%%%%%%%%%%%%%%%%%%%%%%%%%%%%%%%%%%%%%%%%%%%%%
\section{Natural Transformations}
%%%%%%%%%%%%%%%%%%%%%%%%%%%%%%%%%%%%%%%%%%%%%%%%%%%%%%%%%%%%%%%%%%%%%%%%%%%%%%%

Natural transformations likewise possess an intrinsic continuation meaning.

Whenever two realization functors preserve continuation in different but
compatible ways, the comparison between them forms a continuation natural
transformation.

Its realization is precisely the classical natural transformation.

Consequently naturality is not imposed externally.

It expresses coherence between continuation-preserving realizations.

%%%%%%%%%%%%%%%%%%%%%%%%%%%%%%%%%%%%%%%%%%%%%%%%%%%%%%%%%%%%%%%%%%%%%%%%%%%%%%%
\section{Universal Properties}
%%%%%%%%%%%%%%%%%%%%%%%%%%%%%%%%%%%%%%%%%%%%%%%%%%%%%%%%%%%%%%%%%%%%%%%%%%%%%%%

Perhaps the deepest insight of category theory is the notion of universal
property.

Continuation Mathematics explains why such properties appear.

Universal objects arise whenever continuation closure determines a unique
maximal completion satisfying a specified continuation constraint.

Thus universality is nothing more than uniqueness forced by completion.

This explains why universal constructions recur throughout mathematics.

They are manifestations of maximal continuation.

%%%%%%%%%%%%%%%%%%%%%%%%%%%%%%%%%%%%%%%%%%%%%%%%%%%%%%%%%%%%%%%%%%%%%%%%%%%%%%%
\section{Limits and Colimits}
%%%%%%%%%%%%%%%%%%%%%%%%%%%%%%%%%%%%%%%%%%%%%%%%%%%%%%%%%%%%%%%%%%%%%%%%%%%%%%%

Limits were developed intrinsically in Part IV.

A limit is the completion forced simultaneously by an entire compatible
continuation family.

Colimits arise dually by universal continuation extension.

Their categorical realizations are precisely the ordinary categorical limits
and colimits.

Thus diagrams, cones, cocones, products, coproducts, pullbacks, pushouts,
equalizers, coequalizers, inverse limits, and direct limits are recovered from
continuation closure.

%%%%%%%%%%%%%%%%%%%%%%%%%%%%%%%%%%%%%%%%%%%%%%%%%%%%%%%%%%%%%%%%%%%%%%%%%%%%%%%
\section{Adjunctions}
%%%%%%%%%%%%%%%%%%%%%%%%%%%%%%%%%%%%%%%%%%%%%%%%%%%%%%%%%%%%%%%%%%%%%%%%%%%%%%%

Adjunctions also admit a continuation interpretation.

Whenever two continuation processes represent opposite directions of structural
completion while preserving the same continuation information, they determine a
Continuation Adjunction.

Realization produces the ordinary categorical adjunction

\[
F\dashv G.
\]

Thus adjointness expresses structural reciprocity between two continuation
processes.

%%%%%%%%%%%%%%%%%%%%%%%%%%%%%%%%%%%%%%%%%%%%%%%%%%%%%%%%%%%%%%%%%%%%%%%%%%%%%%%
\section{Why Category Theory Works}
%%%%%%%%%%%%%%%%%%%%%%%%%%%%%%%%%%%%%%%%%%%%%%%%%%%%%%%%%%%%%%%%%%%%%%%%%%%%%%%

Category theory often appears surprisingly universal.

The continuation viewpoint explains this phenomenon.

Every sufficiently rich mathematical discipline possesses continuation,
completion, morphisms, and admissible extensions.

Once these exist, categorical structure appears automatically.

Category theory therefore succeeds not because categories are primitive, but
because continuation is universal.

Categories are the visible algebra of continuation.

%%%%%%%%%%%%%%%%%%%%%%%%%%%%%%%%%%%%%%%%%%%%%%%%%%%%%%%%%%%%%%%%%%%%%%%%%%%%%%%
\section{Recovery Theorem}
%%%%%%%%%%%%%%%%%%%%%%%%%%%%%%%%%%%%%%%%%%%%%%%%%%%%%%%%%%%%%%%%%%%%%%%%%%%%%%%

\begin{theorem}[Recovery of Category Theory]
Classical category theory is recovered from continuation mathematics.

Specifically,

\begin{enumerate}

\item objects are realizations of completion classes;

\item morphisms are realizations of continuation morphisms;

\item composition is inherited from continuation algebra;

\item identities are inherited from identity continuation;

\item functors are realization-preserving continuation functors;

\item natural transformations arise from continuation coherence;

\item universal constructions arise from maximal continuation;

\item limits and colimits arise from continuation closure;

\item adjunctions arise from reciprocal continuation processes.

\end{enumerate}

Accordingly category theory is a realization of continuation mathematics rather
than an independent mathematical foundation.
\end{theorem}

\begin{proof}
Every categorical notion listed above has already been constructed
intrinsically in Parts II, III, and IV.

Choosing a realization functor translates those intrinsic continuation
structures into their classical categorical counterparts.

Thus the entire categorical framework is recovered without introducing any new
primitive mathematical objects.
\end{proof}

%%%%%%%%%%%%%%%%%%%%%%%%%%%%%%%%%%%%%%%%%%%%%%%%%%%%%%%%%%%%%%%%%%%%%%%%%%%%%%%
\section{Foundational Significance}
%%%%%%%%%%%%%%%%%%%%%%%%%%%%%%%%%%%%%%%%%%%%%%%%%%%%%%%%%%%%%%%%%%%%%%%%%%%%%%%

Continuation Mathematics reveals that category theory occupies an intermediate
foundational level.

It is deeper than set theory because it captures relationships rather than
membership.

Yet it is not the deepest level.

Its true primitives—objects, arrows, composition, universality, and
functoriality—are themselves recoverable from continuation.

The hierarchy is therefore

\[
\text{Continuation}
\;\Longrightarrow\;
\text{Category Theory}
\;\Longrightarrow\;
\text{Classical Mathematics}.
\]

The remarkable unity of category theory is thereby explained.

It is the universal language of mathematics because continuation itself is the
universal structure from which mathematical realization proceeds.

\chapter{Quantum Continuation}

The preceding chapters have shown that many of the great branches of
mathematics are not primitive disciplines but realizations of continuation
structures. Set theory, topology, algebra, geometry, analysis, probability,
information theory, and category theory all arise by imposing different
realization principles upon continuation universes.

Quantum mathematics occupies a distinguished position among these
realizations. Unlike classical mathematics, which assumes that objects possess
determinate states, quantum mathematics begins with collections of admissible
possibilities constrained by composition.

Continuation Mathematics provides a more primitive viewpoint.

It does not begin with states.

It begins with admissible continuations.

%%%%%%%%%%%%%%%%%%%%%%%%%%%%%%%%%%%%%%%%%%%%%%%%%%%%%%%%%%%%%%%%%%%%%%%%%%%%%%%
\section{The Failure of Classical Determinism}
%%%%%%%%%%%%%%%%%%%%%%%%%%%%%%%%%%%%%%%%%%%%%%%%%%%%%%%%%%%%%%%%%%%%%%%%%%%%%%%

Classical mathematics is fundamentally deterministic.

A mathematical object is assumed to possess a unique state, and every
admissible operation produces another uniquely determined state.

Symbolically,

\[
X
\longrightarrow
x
\longrightarrow
f(x).
\]

This architecture is appropriate whenever continuation is already complete.

There exist, however, mathematical systems whose continuation has not yet
collapsed to a single realization.

For such systems the primitive object is not a state but a continuation
frontier.

%%%%%%%%%%%%%%%%%%%%%%%%%%%%%%%%%%%%%%%%%%%%%%%%%%%%%%%%%%%%%%%%%%%%%%%%%%%%%%%
\section{Quantum Continuation Systems}
%%%%%%%%%%%%%%%%%%%%%%%%%%%%%%%%%%%%%%%%%%%%%%%%%%%%%%%%%%%%%%%%%%%%%%%%%%%%%%%

\begin{definition}[Quantum continuation system]
A \emph{quantum continuation system} is a continuation system for which at
least one object possesses more than one admissible continuation.
\end{definition}

Thus there exist objects

\[
x
\]

for which

\[
\operatorname{Cont}(x)
=
\{x_1,x_2,\ldots,x_n\},
\qquad
n>1.
\]

No continuation is distinguished intrinsically.

The continuation structure itself is the primary mathematical object.

%%%%%%%%%%%%%%%%%%%%%%%%%%%%%%%%%%%%%%%%%%%%%%%%%%%%%%%%%%%%%%%%%%%%%%%%%%%%%%%
\section{Continuation Frontiers}
%%%%%%%%%%%%%%%%%%%%%%%%%%%%%%%%%%%%%%%%%%%%%%%%%%%%%%%%%%%%%%%%%%%%%%%%%%%%%%%

The admissible continuations determine a frontier.

\begin{definition}[Quantum frontier]
The continuation frontier of a quantum continuation system is the collection of
all simultaneously admissible continuations prior to realization.
\end{definition}

The frontier is therefore neither uncertainty nor ignorance.

It is an objective mathematical structure.

It records every continuation still compatible with the intrinsic continuation
laws.

%%%%%%%%%%%%%%%%%%%%%%%%%%%%%%%%%%%%%%%%%%%%%%%%%%%%%%%%%%%%%%%%%%%%%%%%%%%%%%%
\section{Observation}
%%%%%%%%%%%%%%%%%%%%%%%%%%%%%%%%%%%%%%%%%%%%%%%%%%%%%%%%%%%%%%%%%%%%%%%%%%%%%%%

Continuation Mathematics introduces a purely mathematical notion of
observation.

\begin{definition}[Observation]
An \emph{observation} is a continuation-completion operation that selects one
realization from a continuation frontier.
\end{definition}

Observation is therefore not a psychological process.

It is an operation within continuation mathematics.

No reference to consciousness is required.

%%%%%%%%%%%%%%%%%%%%%%%%%%%%%%%%%%%%%%%%%%%%%%%%%%%%%%%%%%%%%%%%%%%%%%%%%%%%%%%
\section{Measurement}
%%%%%%%%%%%%%%%%%%%%%%%%%%%%%%%%%%%%%%%%%%%%%%%%%%%%%%%%%%%%%%%%%%%%%%%%%%%%%%%

Classically, measurement is often treated as an additional physical axiom.

Within continuation mathematics it becomes a mathematical consequence.

\begin{definition}[Measurement]
A measurement is a realization map

\[
\mu :
\mathcal F
\longrightarrow
[X]
\]

from a continuation frontier to one of its completion classes.
\end{definition}

Thus measurement is simply completion.

%%%%%%%%%%%%%%%%%%%%%%%%%%%%%%%%%%%%%%%%%%%%%%%%%%%%%%%%%%%%%%%%%%%%%%%%%%%%%%%
\section{Classical Reality}
%%%%%%%%%%%%%%%%%%%%%%%%%%%%%%%%%%%%%%%%%%%%%%%%%%%%%%%%%%%%%%%%%%%%%%%%%%%%%%%

A classical mathematical object is therefore nothing more than a completed
continuation.

\begin{theorem}
Every deterministic mathematical object is the realization of a continuation
frontier whose completion has become unique.
\end{theorem}

\begin{proof}
If a continuation frontier possesses only one admissible completion, every
admissible realization selects the same object. The frontier therefore behaves
deterministically and coincides with the classical notion of state.
\end{proof}

Determinism is therefore a degenerate case of continuation.

%%%%%%%%%%%%%%%%%%%%%%%%%%%%%%%%%%%%%%%%%%%%%%%%%%%%%%%%%%%%%%%%%%%%%%%%%%%%%%%
\section{Emergence of Hilbert Structure}
%%%%%%%%%%%%%%%%%%%%%%%%%%%%%%%%%%%%%%%%%%%%%%%%%%%%%%%%%%%%%%%%%%%%%%%%%%%%%%%

Continuation Mathematics does not assume vector spaces, linear operators,
inner products, or complex numbers.

Instead one asks a realization question.

\emph{Which mathematical structures faithfully represent continuation
frontiers together with their admissible completion operations?}

Among all realizations, Hilbert spaces provide one particularly successful
answer.

They supply a linear geometry capable of encoding continuation frontiers while
preserving the composition laws governing admissible continuations.

Thus Hilbert space is not primitive.

It is a realization of continuation geometry.

%%%%%%%%%%%%%%%%%%%%%%%%%%%%%%%%%%%%%%%%%%%%%%%%%%%%%%%%%%%%%%%%%%%%%%%%%%%%%%%
\section{Superposition}
%%%%%%%%%%%%%%%%%%%%%%%%%%%%%%%%%%%%%%%%%%%%%%%%%%%%%%%%%%%%%%%%%%%%%%%%%%%%%%%

The classical notion of superposition receives a structural interpretation.

\begin{definition}[Continuation superposition]
A continuation superposition is the realization of a continuation frontier
whose admissible continuations have not yet undergone completion.
\end{definition}

Superposition therefore does not describe an object simultaneously occupying
multiple realized states.

It describes an object possessing multiple admissible continuations.

%%%%%%%%%%%%%%%%%%%%%%%%%%%%%%%%%%%%%%%%%%%%%%%%%%%%%%%%%%%%%%%%%%%%%%%%%%%%%%%
\section{Interference}
%%%%%%%%%%%%%%%%%%%%%%%%%%%%%%%%%%%%%%%%%%%%%%%%%%%%%%%%%%%%%%%%%%%%%%%%%%%%%%%

Continuation paths need not exist independently.

Distinct admissible continuations may constrain one another through the
continuation algebra.

The realization of these structural interactions appears as interference.

Interference is therefore not mysterious.

It is the visible manifestation of relationships already present within the
continuation frontier.

%%%%%%%%%%%%%%%%%%%%%%%%%%%%%%%%%%%%%%%%%%%%%%%%%%%%%%%%%%%%%%%%%%%%%%%%%%%%%%%
\section{Entanglement}
%%%%%%%%%%%%%%%%%%%%%%%%%%%%%%%%%%%%%%%%%%%%%%%%%%%%%%%%%%%%%%%%%%%%%%%%%%%%%%%

Consider two continuation systems

\[
X
\qquad\text{and}\qquad
Y.
\]

If admissible continuations of one system cannot be specified independently of
admissible continuations of the other, the systems share a common continuation
frontier.

\begin{definition}[Continuation entanglement]
Two continuation systems are \emph{continuation-entangled} whenever their
admissible continuation structures fail to decompose into independent
continuation products.
\end{definition}

Entanglement is therefore an intrinsic geometric property of continuation
spaces.

%%%%%%%%%%%%%%%%%%%%%%%%%%%%%%%%%%%%%%%%%%%%%%%%%%%%%%%%%%%%%%%%%%%%%%%%%%%%%%%
\section{Quantum Information}
%%%%%%%%%%%%%%%%%%%%%%%%%%%%%%%%%%%%%%%%%%%%%%%%%%%%%%%%%%%%%%%%%%%%%%%%%%%%%%%

Continuation entropy already measures uncertainty remaining within a frontier.

Quantum information is recovered by restricting continuation entropy to
quantum continuation systems.

Classical information theory appears when every frontier has collapsed.

Quantum information appears when frontiers remain open.

%%%%%%%%%%%%%%%%%%%%%%%%%%%%%%%%%%%%%%%%%%%%%%%%%%%%%%%%%%%%%%%%%%%%%%%%%%%%%%%
\section{Recovery Principle}
%%%%%%%%%%%%%%%%%%%%%%%%%%%%%%%%%%%%%%%%%%%%%%%%%%%%%%%%%%%%%%%%%%%%%%%%%%%%%%%

\begin{theorem}[Recovery of Quantum Mathematics]
Quantum mathematics is recovered as the realization theory of continuation
systems possessing non-trivial continuation frontiers.
\end{theorem}

\begin{proof}
Continuation Mathematics supplies the primitive notions of admissible
continuation, frontier, completion, morphism, entropy, geometry, and
realization.

Choosing realizations that preserve these structures yields the mathematical
architecture traditionally used in quantum theory, including state spaces,
measurement, superposition, interference, and entanglement.

These notions therefore arise from continuation rather than serving as
primitive axioms.
\end{proof}

%%%%%%%%%%%%%%%%%%%%%%%%%%%%%%%%%%%%%%%%%%%%%%%%%%%%%%%%%%%%%%%%%%%%%%%%%%%%%%%
\section{Foundational Consequences}
%%%%%%%%%%%%%%%%%%%%%%%%%%%%%%%%%%%%%%%%%%%%%%%%%%%%%%%%%%%%%%%%%%%%%%%%%%%%%%%

Continuation Mathematics proposes a shift in mathematical foundations.

Classical mathematics begins with realized objects.

Quantum mathematics begins with state spaces.

Continuation Mathematics begins earlier.

It begins with admissible continuation.

Every realized mathematical object is the completion of a continuation.

Every deterministic theory is a theory of completed continuation.

Every probabilistic theory is a theory of partially completed continuation.

Every quantum theory is a theory of non-trivial continuation frontiers.

Accordingly,

\[
\boxed{
\text{Continuation}
\Longrightarrow
\text{Quantum Mathematics}
\Longrightarrow
\text{Classical Mathematics}.
}
\]

Quantum mathematics is therefore not an exception to classical mathematics.

It is a deeper realization of the universal mathematics of continuation.

%%%%%%%%%%%%%%%%%%%%%%%%%%%%%%%%%%%%%%%%%%%%%%%%%%%%%%%%%%%%%%%%%%%%%%%%%%%%%%%
\chapter{Universal Continuation Theory}
%%%%%%%%%%%%%%%%%%%%%%%%%%%%%%%%%%%%%%%%%%%%%%%%%%%%%%%%%%%%%%%%%%%%%%%%%%%%%%%

The preceding chapters have established a sequence of recovery theorems for the
principal branches of mathematics. Each recovery began from continuation
structures rather than from the classical primitives of the discipline under
consideration. Nevertheless every recovery reconstructed an existing
mathematical theory.

This repeated phenomenon demands explanation.

It cannot be regarded as a coincidence that topology, algebra, geometry,
analysis, probability, information theory, category theory, and quantum
mathematics all emerge from one continuation framework.

The purpose of this chapter is to prove that this phenomenon is universal.

Continuation is not one mathematical theory among many.

It is the universal structural framework from which mathematical theories arise
as realizations.

%%%%%%%%%%%%%%%%%%%%%%%%%%%%%%%%%%%%%%%%%%%%%%%%%%%%%%%%%%%%%%%%%%%%%%%%%%%%%%%
\section{The Universal Recovery Phenomenon}
%%%%%%%%%%%%%%%%%%%%%%%%%%%%%%%%%%%%%%%%%%%%%%%%%%%%%%%%%%%%%%%%%%%%%%%%%%%%%%%

Every recovery theorem established in this book possesses the same logical
structure.

One first identifies the admissible continuations intrinsic to a mathematical
system.

One then derives the continuation structures forced by those admissible
continuations.

Finally one realizes those structures in a classical mathematical language.

Thus every recovery follows the universal pattern

\[
\boxed{
\text{Admissibility}
\Longrightarrow
\text{Continuation}
\Longrightarrow
\text{Structure}
\Longrightarrow
\text{Realization}.
}
\]

The classical theory appears only at the final stage.

%%%%%%%%%%%%%%%%%%%%%%%%%%%%%%%%%%%%%%%%%%%%%%%%%%%%%%%%%%%%%%%%%%%%%%%%%%%%%%%
\section{Continuation Before Realization}
%%%%%%%%%%%%%%%%%%%%%%%%%%%%%%%%%%%%%%%%%%%%%%%%%%%%%%%%%%%%%%%%%%%%%%%%%%%%%%%

The distinction between continuation and realization is fundamental.

Continuation concerns the intrinsic propagation permitted by a mathematical
system.

Realization concerns the language chosen to describe that propagation.

Accordingly two mathematical theories may differ completely in their classical
presentation while possessing identical continuation structures.

Such theories are not mathematically independent.

They are distinct realizations of one continuation system.

%%%%%%%%%%%%%%%%%%%%%%%%%%%%%%%%%%%%%%%%%%%%%%%%%%%%%%%%%%%%%%%%%%%%%%%%%%%%%%%
\section{Universal Continuation Systems}
%%%%%%%%%%%%%%%%%%%%%%%%%%%%%%%%%%%%%%%%%%%%%%%%%%%%%%%%%%%%%%%%%%%%%%%%%%%%%%%

\begin{definition}[Universal continuation system]
A continuation system is called \emph{universal} if every admissible realization
preserving continuation structure factors uniquely through it.
\end{definition}

Universality is therefore not attached to particular mathematical objects.

It is attached to continuation itself.

Every realization preserving continuation is determined by the same universal
continuation system.

%%%%%%%%%%%%%%%%%%%%%%%%%%%%%%%%%%%%%%%%%%%%%%%%%%%%%%%%%%%%%%%%%%%%%%%%%%%%%%%
\section{The Universal Factorization Principle}
%%%%%%%%%%%%%%%%%%%%%%%%%%%%%%%%%%%%%%%%%%%%%%%%%%%%%%%%%%%%%%%%%%%%%%%%%%%%%%%

The recovery theorems established throughout this work imply a common
factorization phenomenon.

Whenever a mathematical theory preserves admissible continuation, its structural
content factors through continuation mathematics before any classical
realization is chosen.

Consequently the apparent diversity of mathematical disciplines reflects only
the diversity of their realizations.

Their continuation structures remain universal.

%%%%%%%%%%%%%%%%%%%%%%%%%%%%%%%%%%%%%%%%%%%%%%%%%%%%%%%%%%%%%%%%%%%%%%%%%%%%%%%
\section{The Fundamental Universality Theorem}
%%%%%%%%%%%%%%%%%%%%%%%%%%%%%%%%%%%%%%%%%%%%%%%%%%%%%%%%%%%%%%%%%%%%%%%%%%%%%%%

\begin{theorem}[Universal Continuation Theorem]
Every mathematical theory preserving admissible continuation is a realization
of a universal continuation system.
\end{theorem}

\begin{proof}
Every mathematical theory determines a collection of admissible continuations.

These continuations generate continuation systems.

Continuation systems generate continuation spaces together with their intrinsic
relations, morphisms, observables, completion structures, and invariants.

Any faithful realization preserving these continuation structures therefore
factors through the corresponding continuation system before introducing its
particular mathematical language.

Consequently every continuation-preserving mathematical theory is a realization
of a universal continuation system.
\end{proof}

This theorem is independent of the particular realization chosen.

It applies equally to existing mathematical disciplines and to mathematical
theories not yet discovered.

%%%%%%%%%%%%%%%%%%%%%%%%%%%%%%%%%%%%%%%%%%%%%%%%%%%%%%%%%%%%%%%%%%%%%%%%%%%%%%%
\section{The Unity of Mathematics}
%%%%%%%%%%%%%%%%%%%%%%%%%%%%%%%%%%%%%%%%%%%%%%%%%%%%%%%%%%%%%%%%%%%%%%%%%%%%%%%

The preceding theorem establishes a precise sense in which mathematics is
structurally unified.

Topology, algebra, geometry, analysis, probability, information theory,
category theory, and quantum mathematics are not unrelated foundations.

They are distinct realizations of universal continuation structures.

The unity of mathematics therefore precedes its classical divisions.

%%%%%%%%%%%%%%%%%%%%%%%%%%%%%%%%%%%%%%%%%%%%%%%%%%%%%%%%%%%%%%%%%%%%%%%%%%%%%%%
\section{The Scope of Universal Continuation}
%%%%%%%%%%%%%%%%%%%%%%%%%%%%%%%%%%%%%%%%%%%%%%%%%%%%%%%%%%%%%%%%%%%%%%%%%%%%%%%

Universal continuation theory is not itself another branch of mathematics.

It is the structural framework within which mathematical branches are generated.

Accordingly,

\[
\begin{aligned}
\text{Topology}
&=\text{Spatial Realization},\\
\text{Algebra}
&=\text{Operational Realization},\\
\text{Geometry}
&=\text{Metric Realization},\\
\text{Analysis}
&=\text{Limiting Realization},\\
\text{Probability}
&=\text{Statistical Realization},\\
\text{Information Theory}
&=\text{Informational Realization},\\
\text{Category Theory}
&=\text{Functorial Realization},\\
\text{Quantum Mathematics}
&=\text{Frontier Realization}.
\end{aligned}
\]

None of these theories precedes continuation.

Each is generated by choosing a realization of continuation structures.

%%%%%%%%%%%%%%%%%%%%%%%%%%%%%%%%%%%%%%%%%%%%%%%%%%%%%%%%%%%%%%%%%%%%%%%%%%%%%%%
\section{A Universal Framework for Mathematical Investigation}
%%%%%%%%%%%%%%%%%%%%%%%%%%%%%%%%%%%%%%%%%%%%%%%%%%%%%%%%%%%%%%%%%%%%%%%%%%%%%%%

The universality established here extends beyond the recovery of existing
mathematics.

Whenever a new mathematical problem is encountered, the primary task is no
longer to invent auxiliary constructions or heuristic invariants.

Instead one first determines the admissible continuations intrinsic to the
problem.

The continuation structures generated by those admissible continuations then
determine the observables, invariants, completion phenomena, and realizations
appropriate to the investigation.

Thus universal continuation theory provides a single mathematical framework for
the investigation of arbitrary mathematical systems.

%%%%%%%%%%%%%%%%%%%%%%%%%%%%%%%%%%%%%%%%%%%%%%%%%%%%%%%%%%%%%%%%%%%%%%%%%%%%%%%
\section{Transition}
%%%%%%%%%%%%%%%%%%%%%%%%%%%%%%%%%%%%%%%%%%%%%%%%%%%%%%%%%%%%%%%%%%%%%%%%%%%%%%%

Universal continuation theory establishes that every continuation-preserving
mathematical theory is a realization of one universal structural framework.

The remaining question is therefore no longer one of continuation.

It is one level deeper.

What mathematical principle generates admissibility itself?

The final chapter answers this question by identifying admissibility as the
primitive mathematical notion from which continuation, structure, realization,
and ultimately mathematics itself are generated.

%%%%%%%%%%%%%%%%%%%%%%%%%%%%%%%%%%%%%%%%%%%%%%%%%%%%%%%%%%%%%%%%%%%%%%%%%%%%%%%
\chapter{The Theory of Admissible Mathematics}

%%%%%%%%%%%%%%%%%%%%%%%%%%%%%%%%%%%%%%%%%%%%%%%%%%%%%%%%%%%%%%%%%%%%%%%%%%%%%%%
% Introduction
%%%%%%%%%%%%%%%%%%%%%%%%%%%%%%%%%%%%%%%%%%%%%%%%%%%%%%%%%%%%%%%%%%%%%%%%%%%%%%%

The preceding chapters established that continuation is universal.

They demonstrated that the principal branches of classical mathematics are not
independent mathematical worlds, but faithful realizations of one underlying
continuation structure.

A deeper question nevertheless remains.

If continuation explains the emergence of mathematical structure, what explains
continuation itself?

The purpose of this chapter is to answer this question.

The answer completes the constitutional development of this work.

Continuation is not primitive.

Continuation is generated.

Its source is admissibility.

Accordingly, the primitive mathematical notion is neither object nor operation,
neither language nor realization.

It is admissibility.

Every mathematical structure developed throughout this work is therefore shown
to arise from one constitutional principle.

%%%%%%%%%%%%%%%%%%%%%%%%%%%%%%%%%%%%%%%%%%%%%%%%%%%%%%%%%%%%%%%%%%%%%%%%%%%%%%%
\section{The Primitive Mathematical Question}
%%%%%%%%%%%%%%%%%%%%%%%%%%%%%%%%%%%%%%%%%%%%%%%%%%%%%%%%%%%%%%%%%%%%%%%%%%%%%%%

The historical development of mathematics has typically begun by postulating
primitive mathematical objects.

Sets.

Numbers.

Functions.

Relations.

Spaces.

Categories.

Each foundation then investigates the structures generated by its chosen
primitive.

Continuation Mathematics reverses this order.

Objects are not primitive.

Operations are not primitive.

Even continuation is not primitive.

All are recovered.

The first mathematical question is therefore neither

\[
\text{What objects exist?}
\]

nor

\[
\text{What operations may be performed?}
\]

It is

\[
\boxed{
\text{What is intrinsically admissible?}
}
\]

Every subsequent mathematical construction is forced by the answer.

%%%%%%%%%%%%%%%%%%%%%%%%%%%%%%%%%%%%%%%%%%%%%%%%%%%%%%%%%%%%%%%%%%%%%%%%%%%%%%%
\section{Admissibility Before Mathematics}
%%%%%%%%%%%%%%%%%%%%%%%%%%%%%%%%%%%%%%%%%%%%%%%%%%%%%%%%%%%%%%%%%%%%%%%%%%%%%%%

Admissibility precedes every mathematical realization.

Before one can speak of objects, one must know which objects are admissible.

Before one can speak of operations, one must know which operations preserve
admissibility.

Before one can speak of spaces, one must know which extensions are admissible.

Before one can speak of proofs, one must know which deductions preserve
admissibility.

Thus admissibility is logically prior to every mathematical construction.

Mathematics therefore begins neither with objects nor with operations.

It begins with admissibility.

%%%%%%%%%%%%%%%%%%%%%%%%%%%%%%%%%%%%%%%%%%%%%%%%%%%%%%%%%%%%%%%%%%%%%%%%%%%%%%%
\section{The Constitutional Hierarchy}
%%%%%%%%%%%%%%%%%%%%%%%%%%%%%%%%%%%%%%%%%%%%%%%%%%%%%%%%%%%%%%%%%%%%%%%%%%%%%%%

The preceding parts of this work establish a single constitutional hierarchy.

\[
\boxed{
\begin{aligned}
\text{Admissibility}
&\Longrightarrow
\text{Constraint}\\
&\Longrightarrow
\text{Relation}\\
&\Longrightarrow
\text{Propagation}\\
&\Longrightarrow
\text{Continuation}\\
&\Longrightarrow
\text{Completion}\\
&\Longrightarrow
\text{Distinguishability}\\
&\Longrightarrow
\text{Information}\\
&\Longrightarrow
\text{Realization}\\
&\Longrightarrow
\text{Mathematics.}
\end{aligned}
}
\]

Nothing in this hierarchy is introduced independently.

Each level is generated by the preceding one.

Consequently every mathematical structure possesses a constitutional origin.

%%%%%%%%%%%%%%%%%%%%%%%%%%%%%%%%%%%%%%%%%%%%%%%%%%%%%%%%%%%%%%%%%%%%%%%%%%%%%%%
\section{The Fundamental Generation Principle}
%%%%%%%%%%%%%%%%%%%%%%%%%%%%%%%%%%%%%%%%%%%%%%%%%%%%%%%%%%%%%%%%%%%%%%%%%%%%%%%

\begin{principle}[Fundamental Generation Principle]
Every mathematical structure is generated by the propagation of admissibility.
\end{principle}

The content of mathematics therefore lies not in isolated mathematical objects,
but in the admissibility structures from which those objects necessarily arise.

Objects are consequences.

Admissibility is constitutional.

%%%%%%%%%%%%%%%%%%%%%%%%%%%%%%%%%%%%%%%%%%%%%%%%%%%%%%%%%%%%%%%%%%%%%%%%%%%%%%%
\section{The Fundamental Theorem of Admissible Mathematics}
%%%%%%%%%%%%%%%%%%%%%%%%%%%%%%%%%%%%%%%%%%%%%%%%%%%%%%%%%%%%%%%%%%%%%%%%%%%%%%%

\begin{theorem}[Fundamental Theorem of Admissible Mathematics]
Every mathematical theory is a realization of an admissibility structure.
\end{theorem}

\begin{proof}
Every mathematical theory determines which constructions are mathematically
admissible.

Those admissibility conditions determine allowable constraints.

Constraints generate relations.

Relations generate propagation.

Propagation generates continuation.

Continuation generates completion together with distinguishability,
information, and realization.

Consequently every mathematical theory is generated from an admissibility
structure.

Every mathematical theory is therefore a realization of admissibility.
\end{proof}

%%%%%%%%%%%%%%%%%%%%%%%%%%%%%%%%%%%%%%%%%%%%%%%%%%%%%%%%%%%%%%%%%%%%%%%%%%%%%%%
\section{The Recovery of Classical Mathematics}
%%%%%%%%%%%%%%%%%%%%%%%%%%%%%%%%%%%%%%%%%%%%%%%%%%%%%%%%%%%%%%%%%%%%%%%%%%%%%%%

The recovery theorems established throughout this work now acquire a unified
interpretation.

Classical mathematics is not replaced.

It is recovered.

Set theory is recovered as one realization of admissible completion.

Topology is recovered as one realization of admissible continuity.

Algebra is recovered as one realization of admissible operations.

Geometry is recovered as one realization of admissible spatial structure.

Analysis is recovered as one realization of admissible limiting behaviour.

Probability is recovered as one realization of admissible statistical
continuation.

Information theory is recovered as one realization of admissible
distinguishability.

Category theory is recovered as one realization of admissible structural
composition.

Quantum mathematics is recovered as one realization of admissible continuation
frontiers.

Their diversity reflects different realizations rather than different
foundations.

%%%%%%%%%%%%%%%%%%%%%%%%%%%%%%%%%%%%%%%%%%%%%%%%%%%%%%%%%%%%%%%%%%%%%%%%%%%%%%%
\section{The Universality of Admissibility}
%%%%%%%%%%%%%%%%%%%%%%%%%%%%%%%%%%%%%%%%%%%%%%%%%%%%%%%%%%%%%%%%%%%%%%%%%%%%%%%

The preceding recovery theorems establish more than a collection of independent
equivalences.

They establish universality.

Whenever a mathematical theory possesses internally coherent admissibility
conditions, continuation structures necessarily emerge.

Whenever continuation structures emerge, completion, distinguishability,
information, and realization follow.

Accordingly, admissibility is not one mathematical principle among many.

It is the constitutional source from which every coherent mathematical theory
is generated.

%%%%%%%%%%%%%%%%%%%%%%%%%%%%%%%%%%%%%%%%%%%%%%%%%%%%%%%%%%%%%%%%%%%%%%%%%%%%%%%
\section{The Constitution of Mathematics}
%%%%%%%%%%%%%%%%%%%%%%%%%%%%%%%%%%%%%%%%%%%%%%%%%%%%%%%%%%%%%%%%%%%%%%%%%%%%%%%

The purpose of mathematics may now be stated precisely.

Mathematics is not fundamentally the study of objects.

Neither is it fundamentally the study of operations.

Nor is it merely the manipulation of formal symbolic systems.

Rather, mathematics is the investigation of admissibility together with the
structures that admissibility necessarily generates.

Objects acquire meaning only through admissibility.

Relations acquire meaning only through admissibility.

Continuation exists only because admissibility permits continuation.

Completion exists only because admissibility permits completion.

Every mathematical realization therefore derives its meaning from the
admissibility structure that generates it.

%%%%%%%%%%%%%%%%%%%%%%%%%%%%%%%%%%%%%%%%%%%%%%%%%%%%%%%%%%%%%%%%%%%%%%%%%%%%%%%
\section{Beyond Recovery}
%%%%%%%%%%%%%%%%%%%%%%%%%%%%%%%%%%%%%%%%%%%%%%%%%%%%%%%%%%%%%%%%%%%%%%%%%%%%%%%

The objective of the present work has been to recover mathematics from
admissibility.

That objective has now been achieved.

A new question therefore presents itself.

If mathematical structures are generated rather than invented,

can mathematical discovery itself be generated rather than invented?

Can there exist a canonical theory governing the discovery of mathematics?

Can mathematical investigation itself become an object of mathematical study?

These questions do not concern particular mathematical objects.

They concern the process by which mathematical objects are systematically
recovered from admissibility.

The constitutional development established throughout this work now makes such
a theory possible.

%%%%%%%%%%%%%%%%%%%%%%%%%%%%%%%%%%%%%%%%%%%%%%%%%%%%%%%%%%%%%%%%%%%%%%%%%%%%%%%
\section{Closing Theorem}
%%%%%%%%%%%%%%%%%%%%%%%%%%%%%%%%%%%%%%%%%%%%%%%%%%%%%%%%%%%%%%%%%%%%%%%%%%%%%%%

\begin{theorem}[The Constitution of Mathematics]
Every mathematical theory is the realization of an admissibility structure, and
every mathematical object derives its meaning from the admissible structures
that generate it.
\end{theorem}

\begin{proof}
The preceding parts established that continuation is generated by admissibility,
that completion is generated by continuation, that realization is generated by
completion, and that the principal branches of mathematics are recoverable as
faithful realizations of these structures.

Consequently, mathematical objects are not foundational.

They are generated.

The common constitutional source underlying every coherent mathematical theory
is therefore admissibility.

Accordingly, mathematics is the study of admissibility together with the
structures that admissibility necessarily generates.
\end{proof}

\vspace{1em}

\begin{center}
\Large
\emph{The foundations have now been established.}

\vspace{0.75em}

\emph{The remaining question is no longer what mathematics is.}

\vspace{0.75em}

\emph{The remaining question is how mathematics discovers itself.}
\end{center}


\part{The Mathematics of Mathematical Discovery}

%%%%%%%%%%%%%%%%%%%%%%%%%%%%%%%%%%%%%%%%%%%%%%%%%%%%%%%%%%%%%%%%%%%%%%%%%%%%%%%
\chapter{The Theory of Canonical Mathematical Investigation}
%%%%%%%%%%%%%%%%%%%%%%%%%%%%%%%%%%%%%%%%%%%%%%%%%%%%%%%%%%%%%%%%%%%%%%%%%%%%%%%

The preceding parts established a new constitutional foundation for
mathematics.

Beginning from admissibility alone, continuation structures were generated.
Continuation generated completion.
Completion generated realization.
The principal branches of classical mathematics were then recovered as faithful
realizations of this common constitutional hierarchy.

One question nevertheless remained.

How should new mathematics be discovered?

Traditionally this question has been regarded as lying outside mathematics
itself.

Discovery has been attributed to intuition, ingenuity, experimentation,
experience, or heuristic insight.

The mathematical theory begins only after discovery has occurred.

The present part rejects this distinction.

If mathematical objects are generated by necessity rather than invention, then
the process by which those objects are discovered must itself possess
mathematical structure.

Mathematical investigation therefore becomes a legitimate mathematical object.

The purpose of this part is to develop that object.

%%%%%%%%%%%%%%%%%%%%%%%%%%%%%%%%%%%%%%%%%%%%%%%%%%%%%%%%%%%%%%%%%%%%%%%%%%%%%%%
\section{The Classical View of Mathematical Discovery}
%%%%%%%%%%%%%%%%%%%%%%%%%%%%%%%%%%%%%%%%%%%%%%%%%%%%%%%%%%%%%%%%%%%%%%%%%%%%%%%

Classical mathematical practice has historically separated two fundamentally
different activities.

The first is discovery.

The second is proof.

Discovery has traditionally been regarded as informal.

Proof has been regarded as mathematical.

Consequently the logical development of mathematics begins only after the
essential mathematical ideas have already been found.

This separation has proved remarkably successful.

Nevertheless it leaves unanswered one fundamental question.

Why were the successful ideas discovered rather than infinitely many possible
alternatives?

The traditional methodology supplies no general answer.

%%%%%%%%%%%%%%%%%%%%%%%%%%%%%%%%%%%%%%%%%%%%%%%%%%%%%%%%%%%%%%%%%%%%%%%%%%%%%%%
\section{The Structural View of Investigation}
%%%%%%%%%%%%%%%%%%%%%%%%%%%%%%%%%%%%%%%%%%%%%%%%%%%%%%%%%%%%%%%%%%%%%%%%%%%%%%%

Continuation Mathematics suggests a different perspective.

Throughout the preceding parts, mathematical objects were never introduced
because they appeared useful.

Instead they were generated because the preceding structures forced their
existence.

Investigation therefore proceeded by necessity rather than invention.

This observation admits a profound generalization.

If every mathematical object is generated by structural necessity, then every
successful mathematical investigation must likewise follow a structural path.

The task of mathematical investigation is therefore not to invent.

It is to determine.

%%%%%%%%%%%%%%%%%%%%%%%%%%%%%%%%%%%%%%%%%%%%%%%%%%%%%%%%%%%%%%%%%%%%%%%%%%%%%%%
\section{Canonical Investigation}
%%%%%%%%%%%%%%%%%%%%%%%%%%%%%%%%%%%%%%%%%%%%%%%%%%%%%%%%%%%%%%%%%%%%%%%%%%%%%%%

The purpose of mathematical investigation is to determine the structures
intrinsically forced by the mathematical system under consideration.

Nothing more.

Nothing less.

This principle immediately distinguishes two fundamentally different forms of
research.

The first introduces external mathematical constructions because they appear
useful.

The second derives every construction from the intrinsic structure of the
system itself.

Only the second will be called canonical.

\begin{definition}[Canonical mathematical investigation]
A mathematical investigation is called \emph{canonical} if every mathematical
object introduced during the investigation is generated by structures already
intrinsic to the system under investigation.
\end{definition}

Canonical investigation therefore forbids arbitrary mathematical invention.

Every object must possess a constitutional justification.

%%%%%%%%%%%%%%%%%%%%%%%%%%%%%%%%%%%%%%%%%%%%%%%%%%%%%%%%%%%%%%%%%%%%%%%%%%%%%%%
\section{The Canonical Investigation Principle}
%%%%%%%%%%%%%%%%%%%%%%%%%%%%%%%%%%%%%%%%%%%%%%%%%%%%%%%%%%%%%%%%%%%%%%%%%%%%%%%

The preceding definition immediately yields the governing principle.

\begin{principle}[Canonical Investigation Principle]
No mathematical object may be introduced unless its existence is forced by the
preceding structural development.
\end{principle}

This principle governs every subsequent chapter of the present part.

It replaces heuristic construction by structural generation.

%%%%%%%%%%%%%%%%%%%%%%%%%%%%%%%%%%%%%%%%%%%%%%%%%%%%%%%%%%%%%%%%%%%%%%%%%%%%%%%
\section{The Structure of Investigation}
%%%%%%%%%%%%%%%%%%%%%%%%%%%%%%%%%%%%%%%%%%%%%%%%%%%%%%%%%%%%%%%%%%%%%%%%%%%%%%%

Every canonical investigation possesses the same logical architecture.

One begins by identifying the mathematical system itself.

One then determines its admissibility structure.

Admissibility generates propagation.

Propagation generates continuation.

Continuation generates observables.

Observables generate structural quantities.

Structural quantities generate compatibility.

Compatibility generates structural obstructions.

Finally the desired theorem emerges as the unique admissible conclusion.

Accordingly every successful canonical investigation possesses the same
constitutional form.

%%%%%%%%%%%%%%%%%%%%%%%%%%%%%%%%%%%%%%%%%%%%%%%%%%%%%%%%%%%%%%%%%%%%%%%%%%%%%%%
\section{The Canonical Investigation Operator}
%%%%%%%%%%%%%%%%%%%%%%%%%%%%%%%%%%%%%%%%%%%%%%%%%%%%%%%%%%%%%%%%%%%%%%%%%%%%%%%

The preceding observations suggest that investigation itself admits an
operator-theoretic description.

Given a mathematical system

\[
\mathcal{S},
\]

there exists an investigation operator

\[
\mathfrak{I},
\]

whose purpose is to generate every mathematical structure forced by
\(\mathcal{S}\).

Accordingly

\[
\boxed{
\mathfrak{I}(\mathcal{S})
=
\text{the canonical structural development of }\mathcal{S}.
}
\]

The investigation operator is not algorithmic.

It is mathematical.

Its output is the complete hierarchy of structures constitutionally generated
by the system.

%%%%%%%%%%%%%%%%%%%%%%%%%%%%%%%%%%%%%%%%%%%%%%%%%%%%%%%%%%%%%%%%%%%%%%%%%%%%%%%
\section{Structural Necessity}
%%%%%%%%%%%%%%%%%%%%%%%%%%%%%%%%%%%%%%%%%%%%%%%%%%%%%%%%%%%%%%%%%%%%%%%%%%%%%%%

The introduction of the investigation operator fundamentally changes the role
of mathematical creativity.

Creativity no longer consists in inventing mathematical objects.

Creativity consists in recognizing the structures already forced by the
mathematical system.

Accordingly mathematical progress becomes the progressive elimination of
arbitrary choices.

The more canonical an investigation becomes, the fewer external decisions
remain.

The ideal investigation contains none.

%%%%%%%%%%%%%%%%%%%%%%%%%%%%%%%%%%%%%%%%%%%%%%%%%%%%%%%%%%%%%%%%%%%%%%%%%%%%%%%
\section{Canonical and Heuristic Mathematics}
%%%%%%%%%%%%%%%%%%%%%%%%%%%%%%%%%%%%%%%%%%%%%%%%%%%%%%%%%%%%%%%%%%%%%%%%%%%%%%%

The distinction between heuristic and canonical investigation is therefore
precise.

A heuristic investigation permits external constructions whose usefulness is
established retrospectively.

A canonical investigation admits only constructions whose necessity is
established prospectively.

Both approaches may discover correct mathematics.

Only one explains why the discovered mathematics had to exist.

%%%%%%%%%%%%%%%%%%%%%%%%%%%%%%%%%%%%%%%%%%%%%%%%%%%%%%%%%%%%%%%%%%%%%%%%%%%%%%%
\section{The Fundamental Theorem of Canonical Investigation}
%%%%%%%%%%%%%%%%%%%%%%%%%%%%%%%%%%%%%%%%%%%%%%%%%%%%%%%%%%%%%%%%%%%%%%%%%%%%%%%

\begin{theorem}[Fundamental Theorem of Canonical Investigation]
Every mathematical theorem possesses a canonical investigation whose
mathematical objects are generated solely from the intrinsic structure of the
underlying system.
\end{theorem}

\begin{proof}
Every mathematical theorem concerns a mathematical system.

Every mathematical system possesses admissibility conditions.

By the constitutional hierarchy established in the preceding parts,
admissibility generates propagation, continuation, completion,
distinguishability, information, and realization.

Every mathematical object required for the investigation is therefore generated
from structures already intrinsic to the system.

Consequently the investigation itself may be organized canonically.

The resulting mathematical development contains no constitutionally
unjustified constructions.

Therefore every mathematical theorem admits a canonical investigation.
\end{proof}

%%%%%%%%%%%%%%%%%%%%%%%%%%%%%%%%%%%%%%%%%%%%%%%%%%%%%%%%%%%%%%%%%%%%%%%%%%%%%%%
\section{Toward Theorem Spaces}
%%%%%%%%%%%%%%%%%%%%%%%%%%%%%%%%%%%%%%%%%%%%%%%%%%%%%%%%%%%%%%%%%%%%%%%%%%%%%%%

Canonical investigation determines mathematical structures.

The next question concerns the structures generated by the resulting
mathematical theorems themselves.

Do theorems possess organization?

Can one theorem generate another?

Can families of theorems exhibit structural dependence, completion,
compatibility, or fixed-point behaviour?

These questions concern not mathematical objects, but mathematics itself.

To answer them we must regard collections of theorems as mathematical objects
having their own intrinsic structure.

The development of this theory forms the subject of the next chapter.

%%%%%%%%%%%%%%%%%%%%%%%%%%%%%%%%%%%%%%%%%%%%%%%%%%%%%%%%%%%%%%%%%%%%%%%%%%%%%%%
\chapter{The Theory of Canonical Theorem Spaces}
%%%%%%%%%%%%%%%%%%%%%%%%%%%%%%%%%%%%%%%%%%%%%%%%%%%%%%%%%%%%%%%%%%%%%%%%%%%%%%%

The preceding chapter established that mathematical investigation is itself a
legitimate mathematical object.

Investigation was shown to consist not in the invention of mathematical
constructions, but in the canonical determination of structures already
intrinsic to the mathematical system under consideration.

One question now becomes unavoidable.

What is the mathematical object generated by an investigation?

The answer cannot simply be "a proof."

A proof is only one realization of mathematical understanding.

The true output of a canonical investigation is the collection of mathematical
theorems together with the structural relations that exist among them.

Accordingly, mathematical theorems must themselves become mathematical
objects.

This chapter develops their theory.

%%%%%%%%%%%%%%%%%%%%%%%%%%%%%%%%%%%%%%%%%%%%%%%%%%%%%%%%%%%%%%%%%%%%%%%%%%%%%%%
\section{The Mathematical Status of Theorems}
%%%%%%%%%%%%%%%%%%%%%%%%%%%%%%%%%%%%%%%%%%%%%%%%%%%%%%%%%%%%%%%%%%%%%%%%%%%%%%%

Traditionally a theorem is regarded as an isolated mathematical statement
together with its proof.

This viewpoint is sufficient for verification.

It is insufficient for discovery.

A successful investigation rarely produces a single theorem.

Instead it generates families of mutually dependent results.

Definitions produce propositions.

Propositions produce lemmas.

Lemmas produce structural principles.

Structural principles generate major theorems.

Thus theorems are not isolated.

They possess organization.

They possess dependency.

They possess propagation.

Consequently they admit mathematical structure.

%%%%%%%%%%%%%%%%%%%%%%%%%%%%%%%%%%%%%%%%%%%%%%%%%%%%%%%%%%%%%%%%%%%%%%%%%%%%%%%
\section{Theorem Spaces}
%%%%%%%%%%%%%%%%%%%%%%%%%%%%%%%%%%%%%%%%%%%%%%%%%%%%%%%%%%%%%%%%%%%%%%%%%%%%%%%

The preceding observation motivates the introduction of the fundamental object
of the present part.

\begin{definition}[Theorem space]
A \emph{theorem space} is the mathematical structure consisting of every
theorem generated by a canonical investigation together with every structural
dependency existing among those theorems.
\end{definition}

The theorem space generated by an investigation will be denoted

\[
\mathcal T.
\]

Theorems therefore cease to be isolated logical statements.

They become points of a mathematical structure.

%%%%%%%%%%%%%%%%%%%%%%%%%%%%%%%%%%%%%%%%%%%%%%%%%%%%%%%%%%%%%%%%%%%%%%%%%%%%%%%
\section{Structural Dependency}
%%%%%%%%%%%%%%%%%%%%%%%%%%%%%%%%%%%%%%%%%%%%%%%%%%%%%%%%%%%%%%%%%%%%%%%%%%%%%%%

Every theorem possesses hypotheses.

Every proof depends upon preceding mathematical results.

Accordingly every theorem determines structural dependencies.

These dependencies are intrinsic.

They are independent of exposition.

Different presentations of the same mathematics produce the same dependency
structure.

\begin{definition}[Dependency relation]
For theorems
\[
A,B\in\mathcal T,
\]
we write
\[
A\prec B
\]
whenever the proof of \(B\) constitutionally depends upon \(A\).
\end{definition}

The dependency relation partially orders the theorem space.

%%%%%%%%%%%%%%%%%%%%%%%%%%%%%%%%%%%%%%%%%%%%%%%%%%%%%%%%%%%%%%%%%%%%%%%%%%%%%%%
\section{The Dependency Graph}
%%%%%%%%%%%%%%%%%%%%%%%%%%%%%%%%%%%%%%%%%%%%%%%%%%%%%%%%%%%%%%%%%%%%%%%%%%%%%%%

The dependency relation naturally generates a directed graph.

\begin{definition}[Canonical dependency graph]
The canonical dependency graph of a theorem space is the directed graph whose

\begin{itemize}
\item vertices are the theorems of \(\mathcal T\),
\item directed edges represent structural dependency.
\end{itemize}
\end{definition}

The dependency graph is intrinsic.

It records mathematical necessity rather than historical order.

%%%%%%%%%%%%%%%%%%%%%%%%%%%%%%%%%%%%%%%%%%%%%%%%%%%%%%%%%%%%%%%%%%%%%%%%%%%%%%%
\section{Canonical Generation}
%%%%%%%%%%%%%%%%%%%%%%%%%%%%%%%%%%%%%%%%%%%%%%%%%%%%%%%%%%%%%%%%%%%%%%%%%%%%%%%

Not every theorem contributes equally to the development of mathematics.

Some theorems terminate local arguments.

Others generate entire mathematical theories.

This distinction is structural.

\begin{definition}[Generator theorem]
A theorem is called a \emph{generator theorem} if its introduction
constitutionally forces the existence of further theorems.
\end{definition}

Generator theorems occupy privileged positions within the dependency graph.

%%%%%%%%%%%%%%%%%%%%%%%%%%%%%%%%%%%%%%%%%%%%%%%%%%%%%%%%%%%%%%%%%%%%%%%%%%%%%%%
\section{The Rank of a Theorem}
%%%%%%%%%%%%%%%%%%%%%%%%%%%%%%%%%%%%%%%%%%%%%%%%%%%%%%%%%%%%%%%%%%%%%%%%%%%%%%%

Dependency immediately induces hierarchy.

\begin{definition}[Structural rank]
The structural rank of a theorem is the minimal dependency depth separating it
from the primitive assumptions of the investigation.
\end{definition}

Primitive constitutional principles therefore possess minimal rank.

Terminal realization theorems possess maximal rank.

Rank measures structural generation rather than logical complexity.

%%%%%%%%%%%%%%%%%%%%%%%%%%%%%%%%%%%%%%%%%%%%%%%%%%%%%%%%%%%%%%%%%%%%%%%%%%%%%%%
\section{The Geometry of Theorem Spaces}
%%%%%%%%%%%%%%%%%%%%%%%%%%%%%%%%%%%%%%%%%%%%%%%%%%%%%%%%%%%%%%%%%%%%%%%%%%%%%%%

The dependency graph endows theorem spaces with a natural geometry.

Distance is no longer numerical.

Instead it measures structural separation.

Two theorems lying close together possess nearly identical constitutional
origins.

Widely separated theorems require many successive generations.

Entire mathematical disciplines therefore appear as regions within one
underlying theorem space.

%%%%%%%%%%%%%%%%%%%%%%%%%%%%%%%%%%%%%%%%%%%%%%%%%%%%%%%%%%%%%%%%%%%%%%%%%%%%%%%
\section{The Investigation Map}
%%%%%%%%%%%%%%%%%%%%%%%%%%%%%%%%%%%%%%%%%%%%%%%%%%%%%%%%%%%%%%%%%%%%%%%%%%%%%%%

The investigation operator introduced in the preceding chapter may now be
described precisely.

For every mathematical system

\[
\mathcal S,
\]

canonical investigation produces its theorem space

\[
\mathcal T.
\]

Thus

\[
\boxed{
\mathfrak I:
\mathcal S
\longrightarrow
\mathcal T.
}
\]

The investigation operator maps mathematical systems to the theorem spaces they
constitutionally generate.

%%%%%%%%%%%%%%%%%%%%%%%%%%%%%%%%%%%%%%%%%%%%%%%%%%%%%%%%%%%%%%%%%%%%%%%%%%%%%%%
\section{Structural Equivalence}
%%%%%%%%%%%%%%%%%%%%%%%%%%%%%%%%%%%%%%%%%%%%%%%%%%%%%%%%%%%%%%%%%%%%%%%%%%%%%%%

Different mathematical presentations may produce identical theorem spaces.

Accordingly the true mathematical object is not a particular exposition.

It is the dependency structure generated by the exposition.

\begin{definition}[Structural equivalence]
Two investigations are structurally equivalent if they generate isomorphic
theorem spaces.
\end{definition}

Canonical investigation therefore studies theorem spaces rather than
presentations.

%%%%%%%%%%%%%%%%%%%%%%%%%%%%%%%%%%%%%%%%%%%%%%%%%%%%%%%%%%%%%%%%%%%%%%%%%%%%%%%
\section{The Fundamental Theorem of Theorem Spaces}
%%%%%%%%%%%%%%%%%%%%%%%%%%%%%%%%%%%%%%%%%%%%%%%%%%%%%%%%%%%%%%%%%%%%%%%%%%%%%%%

\begin{theorem}[Fundamental Theorem of Theorem Spaces]
Every canonical mathematical investigation generates a unique theorem space up
to structural equivalence.
\end{theorem}

\begin{proof}
Canonical investigation introduces no arbitrary mathematical constructions.

Every theorem is generated from structures already intrinsic to the
mathematical system.

Consequently the dependency relations among the resulting theorems are likewise
intrinsic.

Different presentations may alter notation, exposition, or order of
development, but they cannot alter the constitutional dependency structure.

Therefore the theorem space generated by a canonical investigation is unique up
to structural equivalence.
\end{proof}

%%%%%%%%%%%%%%%%%%%%%%%%%%%%%%%%%%%%%%%%%%%%%%%%%%%%%%%%%%%%%%%%%%%%%%%%%%%%%%%
\section{Toward Closure}
%%%%%%%%%%%%%%%%%%%%%%%%%%%%%%%%%%%%%%%%%%%%%%%%%%%%%%%%%%%%%%%%%%%%%%%%%%%%%%%

The theorem space is not static.

A newly established theorem frequently forces the existence of further
theorems.

Theorem spaces therefore possess their own dynamics.

A theorem generates new dependencies.

Those dependencies generate further theorems.

The theorem space expands.

This process continues until no further theorems are constitutionally forced.

The resulting process is not heuristic.

It is structural.

Its mathematical study requires the introduction of a closure operator acting
upon theorem spaces.

The development of this operator forms the subject of the next chapter.

%%%%%%%%%%%%%%%%%%%%%%%%%%%%%%%%%%%%%%%%%%%%%%%%%%%%%%%%%%%%%%%%%%%%%%%%%%%%%%%
\chapter{Canonical Theorem Closure}
%%%%%%%%%%%%%%%%%%%%%%%%%%%%%%%%%%%%%%%%%%%%%%%%%%%%%%%%%%%%%%%%%%%%%%%%%%%%%%%

The preceding chapter established that every canonical mathematical
investigation generates a theorem space together with its intrinsic dependency
structure.

A theorem space, however, is not a static mathematical object.

The establishment of one theorem frequently forces the existence of additional
theorems that were not initially apparent. These newly generated theorems may
themselves generate further consequences, producing an expanding structural
development.

Accordingly, theorem spaces possess an intrinsic dynamics.

The purpose of the present chapter is to determine this dynamics and to show
that it is governed by a canonical closure operator.

%%%%%%%%%%%%%%%%%%%%%%%%%%%%%%%%%%%%%%%%%%%%%%%%%%%%%%%%%%%%%%%%%%%%%%%%%%%%%%%
\section{The Dynamics of Theorem Generation}
%%%%%%%%%%%%%%%%%%%%%%%%%%%%%%%%%%%%%%%%%%%%%%%%%%%%%%%%%%%%%%%%%%%%%%%%%%%%%%%

Mathematical discovery rarely terminates after the proof of a single theorem.

Instead, every theorem alters the structural landscape of the investigation.

Previously unrelated results become connected.

New definitions become meaningful.

Additional propositions become unavoidable.

Entire families of consequences emerge.

The theorem space therefore evolves.

This evolution is not heuristic.

It is generated by the dependency structure already present within the theorem
space itself.

%%%%%%%%%%%%%%%%%%%%%%%%%%%%%%%%%%%%%%%%%%%%%%%%%%%%%%%%%%%%%%%%%%%%%%%%%%%%%%%
\section{Forced Theorem Generation}
%%%%%%%%%%%%%%%%%%%%%%%%%%%%%%%%%%%%%%%%%%%%%%%%%%%%%%%%%%%%%%%%%%%%%%%%%%%%%%%

The essential observation is that some theorems are not independently
discovered.

They are structurally forced.

Once the prerequisite mathematical structures exist, the theorem follows by
necessity.

Accordingly, theorem generation possesses two distinct forms.

The first is primitive.

The second is forced.

Only primitive results require genuine investigation.

Forced results are generated automatically by the structural organization of
the theorem space.

%%%%%%%%%%%%%%%%%%%%%%%%%%%%%%%%%%%%%%%%%%%%%%%%%%%%%%%%%%%%%%%%%%%%%%%%%%%%%%%
\section{Canonical Closure}
%%%%%%%%%%%%%%%%%%%%%%%%%%%%%%%%%%%%%%%%%%%%%%%%%%%%%%%%%%%%%%%%%%%%%%%%%%%%%%%

These observations motivate the introduction of the fundamental operator of
the present part.

\begin{definition}[Canonical Theorem Closure]
Let
\[
\mathcal T
\]
be a theorem space.

The \emph{canonical theorem closure} of
\(\mathcal T\)
is the smallest theorem space containing
\(\mathcal T\)
that is closed under every theorem constitutionally forced by the dependency
structure of
\(\mathcal T\).

The canonical closure will be denoted

\[
\operatorname{CTCA}(\mathcal T).
\]
\end{definition}

The operator
\(\operatorname{CTCA}\)
does not invent mathematics.

It merely completes the mathematics already latent within the theorem space.

%%%%%%%%%%%%%%%%%%%%%%%%%%%%%%%%%%%%%%%%%%%%%%%%%%%%%%%%%%%%%%%%%%%%%%%%%%%%%%%
\section{Closure Stability}
%%%%%%%%%%%%%%%%%%%%%%%%%%%%%%%%%%%%%%%%%%%%%%%%%%%%%%%%%%%%%%%%%%%%%%%%%%%%%%%

Repeated application of canonical closure cannot continue indefinitely by
creating genuinely new mathematics.

Instead, each application enlarges the theorem space until no further forced
theorems remain.

At this point structural equilibrium is reached.

\begin{definition}[Closure-complete theorem space]
A theorem space
\(\mathcal T\)
is called \emph{closure-complete} if

\[
\operatorname{CTCA}(\mathcal T)
=
\mathcal T.
\]
\end{definition}

Closure-complete theorem spaces therefore represent mathematically stable
investigations.

%%%%%%%%%%%%%%%%%%%%%%%%%%%%%%%%%%%%%%%%%%%%%%%%%%%%%%%%%%%%%%%%%%%%%%%%%%%%%%%
\section{The Closure Operator}
%%%%%%%%%%%%%%%%%%%%%%%%%%%%%%%%%%%%%%%%%%%%%%%%%%%%%%%%%%%%%%%%%%%%%%%%%%%%%%%

The canonical theorem closure possesses the characteristic properties of a
classical closure operator.

\begin{theorem}[Extensiveness]
For every theorem space,

\[
\mathcal T
\subseteq
\operatorname{CTCA}(\mathcal T).
\]
\end{theorem}

\begin{proof}

Canonical closure enlarges a theorem space only by adjoining theorems already
forced by its dependency structure.

Consequently no theorem is removed.

Therefore

\[
\mathcal T
\subseteq
\operatorname{CTCA}(\mathcal T).
\]

\end{proof}

\begin{theorem}[Monotonicity]

If

\[
\mathcal T_1
\subseteq
\mathcal T_2,
\]

then

\[
\operatorname{CTCA}(\mathcal T_1)
\subseteq
\operatorname{CTCA}(\mathcal T_2).
\]

\end{theorem}

\begin{proof}

Every dependency present in
\(\mathcal T_1\)
also exists in
\(\mathcal T_2\).

Consequently every theorem forced by
\(\mathcal T_1\)
is likewise forced by
\(\mathcal T_2\).

Therefore closure preserves inclusion.

\end{proof}

\begin{theorem}[Idempotence]

For every theorem space,

\[
\operatorname{CTCA}
(
\operatorname{CTCA}
(\mathcal T)
)
=
\operatorname{CTCA}
(\mathcal T).
\]

\end{theorem}

\begin{proof}

Once every structurally forced theorem has been generated, no additional
applications of canonical closure can enlarge the theorem space.

Consequently the closure operator stabilizes after completion.

\end{proof}

%%%%%%%%%%%%%%%%%%%%%%%%%%%%%%%%%%%%%%%%%%%%%%%%%%%%%%%%%%%%%%%%%%%%%%%%%%%%%%%
\section{Canonical Completion}
%%%%%%%%%%%%%%%%%%%%%%%%%%%%%%%%%%%%%%%%%%%%%%%%%%%%%%%%%%%%%%%%%%%%%%%%%%%%%%%

The preceding theorems establish that canonical theorem closure is a genuine
closure operator in the classical mathematical sense.

Completion therefore appears once again.

Just as continuation spaces admitted completion,

theorem spaces admit completion.

The completion of a theorem space is precisely its canonical closure.

This parallel is not accidental.

Completion is a universal consequence of admissibility.

%%%%%%%%%%%%%%%%%%%%%%%%%%%%%%%%%%%%%%%%%%%%%%%%%%%%%%%%%%%%%%%%%%%%%%%%%%%%%%%
\section{Minimal Generators}
%%%%%%%%%%%%%%%%%%%%%%%%%%%%%%%%%%%%%%%%%%%%%%%%%%%%%%%%%%%%%%%%%%%%%%%%%%%%%%%

Closure naturally distinguishes two kinds of theorems.

Some theorems merely belong to the closure.

Others generate the closure.

\begin{definition}[Minimal generator]
A subset

\[
G
\subseteq
\mathcal T
\]

is called a \emph{minimal generating family} if

\[
\operatorname{CTCA}(G)
=
\mathcal T,
\]

and no proper subset of
\(G\)
possesses the same property.

\end{definition}

Minimal generators constitute the constitutional core of a mathematical
theory.

%%%%%%%%%%%%%%%%%%%%%%%%%%%%%%%%%%%%%%%%%%%%%%%%%%%%%%%%%%%%%%%%%%%%%%%%%%%%%%%
\section{Closure Complexity}
%%%%%%%%%%%%%%%%%%%%%%%%%%%%%%%%%%%%%%%%%%%%%%%%%%%%%%%%%%%%%%%%%%%%%%%%%%%%%%%

Different theorem spaces require different amounts of structural generation
before stabilization.

This motivates the following invariant.

\begin{definition}[Closure depth]

The \emph{closure depth} of a theorem space is the minimum number of successive
applications of canonical closure required to produce a closure-complete
theorem space.

\end{definition}

Closure depth measures structural rather than computational complexity.

%%%%%%%%%%%%%%%%%%%%%%%%%%%%%%%%%%%%%%%%%%%%%%%%%%%%%%%%%%%%%%%%%%%%%%%%%%%%%%%
\section{The Fundamental Closure Theorem}
%%%%%%%%%%%%%%%%%%%%%%%%%%%%%%%%%%%%%%%%%%%%%%%%%%%%%%%%%%%%%%%%%%%%%%%%%%%%%%%

\begin{theorem}[Fundamental Closure Theorem]

Every canonical theorem space possesses a unique canonical closure.

\end{theorem}

\begin{proof}

Canonical investigation determines a unique theorem space up to structural
equivalence.

The dependency relations intrinsic to that theorem space uniquely determine
every structurally forced theorem.

Consequently canonical closure generates one and only one closure-complete
extension.

Therefore every theorem space possesses a unique canonical closure.

\end{proof}

%%%%%%%%%%%%%%%%%%%%%%%%%%%%%%%%%%%%%%%%%%%%%%%%%%%%%%%%%%%%%%%%%%%%%%%%%%%%%%%
\section{Toward Structural Fixed Points}
%%%%%%%%%%%%%%%%%%%%%%%%%%%%%%%%%%%%%%%%%%%%%%%%%%%%%%%%%%%%%%%%%%%%%%%%%%%%%%%

Canonical closure completes mathematical investigations.

Completion, however, is not the final phenomenon.

Every closure operator possesses fixed points.

These fixed points represent theorem spaces that are structurally stable under
further generation.

The study of these fixed points reveals the deepest organizational principles
of mathematical discovery.

They determine which concepts are fundamental, which structures are
constitutionally inseparable, and where genuine mathematical obstructions
reside.

The mathematical investigation of these fixed points forms the subject of the
next chapter.

%%%%%%%%%%%%%%%%%%%%%%%%%%%%%%%%%%%%%%%%%%%%%%%%%%%%%%%%%%%%%%%%%%%%%%%%%%%%%%%
\chapter{Structural Fixed Point Theory}
%%%%%%%%%%%%%%%%%%%%%%%%%%%%%%%%%%%%%%%%%%%%%%%%%%%%%%%%%%%%%%%%%%%%%%%%%%%%%%%

The preceding chapter established that canonical theorem closure is a genuine
closure operator acting upon theorem spaces.

Every theorem space therefore evolves through successive generations of
structurally forced theorems until no further mathematical consequences remain.

Completion, however, is not the final phenomenon.

Closure naturally gives rise to fixed points.

Within the present theory these fixed points possess a remarkable
interpretation.

They represent investigations that have become structurally complete.

Nothing further can be generated without introducing genuinely new
mathematical assumptions.

Accordingly, structural fixed points are not merely algebraic objects.

They are complete mathematical theories.

%%%%%%%%%%%%%%%%%%%%%%%%%%%%%%%%%%%%%%%%%%%%%%%%%%%%%%%%%%%%%%%%%%%%%%%%%%%%%%%
\section{Structural Stability}
%%%%%%%%%%%%%%%%%%%%%%%%%%%%%%%%%%%%%%%%%%%%%%%%%%%%%%%%%%%%%%%%%%%%%%%%%%%%%%%

Canonical theorem closure enlarges theorem spaces only by adjoining theorems
already forced by the existing dependency structure.

Eventually this process must stabilize.

At stabilization the theorem space reproduces itself.

No additional structurally forced theorem exists.

The investigation has therefore become constitutionally complete.

%%%%%%%%%%%%%%%%%%%%%%%%%%%%%%%%%%%%%%%%%%%%%%%%%%%%%%%%%%%%%%%%%%%%%%%%%%%%%%%
\section{Structural Fixed Points}
%%%%%%%%%%%%%%%%%%%%%%%%%%%%%%%%%%%%%%%%%%%%%%%%%%%%%%%%%%%%%%%%%%%%%%%%%%%%%%%

\begin{definition}[Structural fixed point]

A theorem space

\[
\mathcal T
\]

is called a \emph{structural fixed point} whenever

\[
\operatorname{CTCA}(\mathcal T)
=
\mathcal T.
\]

\end{definition}

Structural fixed points therefore coincide precisely with closure-complete
theorem spaces.

%%%%%%%%%%%%%%%%%%%%%%%%%%%%%%%%%%%%%%%%%%%%%%%%%%%%%%%%%%%%%%%%%%%%%%%%%%%%%%%
\section{The Meaning of Structural Stability}
%%%%%%%%%%%%%%%%%%%%%%%%%%%%%%%%%%%%%%%%%%%%%%%%%%%%%%%%%%%%%%%%%%%%%%%%%%%%%%%

The preceding definition admits an important interpretation.

A structural fixed point does not signify the end of mathematical activity.

It signifies the exhaustion of one mathematical constitution.

Every theorem forced by the existing admissibility structure has already been
generated.

Further mathematical development is possible only by enlarging the underlying
admissibility structure itself.

Consequently every structural fixed point represents a mathematically complete
investigation.

%%%%%%%%%%%%%%%%%%%%%%%%%%%%%%%%%%%%%%%%%%%%%%%%%%%%%%%%%%%%%%%%%%%%%%%%%%%%%%%
\section{The Fixed Point Principle}
%%%%%%%%%%%%%%%%%%%%%%%%%%%%%%%%%%%%%%%%%%%%%%%%%%%%%%%%%%%%%%%%%%%%%%%%%%%%%%%

\begin{principle}[Structural Fixed Point Principle]

Every canonical mathematical investigation converges toward a structural fixed
point.

\end{principle}

Canonical investigation is therefore not an open-ended process of invention.

It is a convergent process of structural completion.

%%%%%%%%%%%%%%%%%%%%%%%%%%%%%%%%%%%%%%%%%%%%%%%%%%%%%%%%%%%%%%%%%%%%%%%%%%%%%%%
\section{Fixed Point Basins}
%%%%%%%%%%%%%%%%%%%%%%%%%%%%%%%%%%%%%%%%%%%%%%%%%%%%%%%%%%%%%%%%%%%%%%%%%%%%%%%

Different investigations may begin from different primitive assumptions while
nevertheless converging toward the same structural fixed point.

This motivates the following definition.

\begin{definition}[Fixed point basin]

The \emph{basin} of a structural fixed point consists of every theorem space
whose canonical closure converges to that fixed point.

\end{definition}

Basins classify mathematical investigations according to their ultimate
structural completion rather than their initial presentation.

%%%%%%%%%%%%%%%%%%%%%%%%%%%%%%%%%%%%%%%%%%%%%%%%%%%%%%%%%%%%%%%%%%%%%%%%%%%%%%%
\section{Structural Equivalence of Investigations}
%%%%%%%%%%%%%%%%%%%%%%%%%%%%%%%%%%%%%%%%%%%%%%%%%%%%%%%%%%%%%%%%%%%%%%%%%%%%%%%

Two mathematical investigations may differ substantially during their early
development.

Definitions may differ.

Notation may differ.

Intermediate propositions may differ.

Yet if both investigations converge toward the same structural fixed point,
their constitutional mathematics is identical.

\begin{definition}[Canonical equivalence]

Two canonical investigations are canonically equivalent whenever they converge
to the same structural fixed point.

\end{definition}

Canonical equivalence therefore classifies mathematics according to completion
rather than exposition.

%%%%%%%%%%%%%%%%%%%%%%%%%%%%%%%%%%%%%%%%%%%%%%%%%%%%%%%%%%%%%%%%%%%%%%%%%%%%%%%
\section{Generator Cores}
%%%%%%%%%%%%%%%%%%%%%%%%%%%%%%%%%%%%%%%%%%%%%%%%%%%%%%%%%%%%%%%%%%%%%%%%%%%%%%%

Every structural fixed point possesses a smallest generating family.

\begin{definition}[Generator core]

The \emph{generator core} of a structural fixed point is the smallest family
of theorems whose canonical closure produces the entire fixed point.

\end{definition}

Generator cores represent the constitutional nucleus of mathematical theories.

Everything else follows necessarily.

%%%%%%%%%%%%%%%%%%%%%%%%%%%%%%%%%%%%%%%%%%%%%%%%%%%%%%%%%%%%%%%%%%%%%%%%%%%%%%%
\section{Structural Rigidity}
%%%%%%%%%%%%%%%%%%%%%%%%%%%%%%%%%%%%%%%%%%%%%%%%%%%%%%%%%%%%%%%%%%%%%%%%%%%%%%%

Once a structural fixed point has been reached, every theorem belongs to one of
two classes.

Either it belongs to the generator core,

or it is generated by the generator core.

Nothing else exists.

Accordingly every structural fixed point possesses intrinsic rigidity.

No theorem may be removed without destroying completion.

No theorem may be added unless the admissibility structure itself is enlarged.

%%%%%%%%%%%%%%%%%%%%%%%%%%%%%%%%%%%%%%%%%%%%%%%%%%%%%%%%%%%%%%%%%%%%%%%%%%%%%%%
\section{Structural Obstructions}
%%%%%%%%%%%%%%%%%%%%%%%%%%%%%%%%%%%%%%%%%%%%%%%%%%%%%%%%%%%%%%%%%%%%%%%%%%%%%%%

Not every attempted mathematical investigation converges.

Some investigations repeatedly generate mutually incompatible structural
requirements.

Others fail because no closure-complete theorem space exists under the assumed
admissibility conditions.

These failures are not accidental.

They are themselves mathematical objects.

\begin{definition}[Structural obstruction]

A structural obstruction is an intrinsic incompatibility preventing a theorem
space from reaching a structural fixed point.

\end{definition}

Structural obstructions therefore become canonical objects of investigation.

%%%%%%%%%%%%%%%%%%%%%%%%%%%%%%%%%%%%%%%%%%%%%%%%%%%%%%%%%%%%%%%%%%%%%%%%%%%%%%%
\section{The Fundamental Fixed Point Theorem}
%%%%%%%%%%%%%%%%%%%%%%%%%%%%%%%%%%%%%%%%%%%%%%%%%%%%%%%%%%%%%%%%%%%%%%%%%%%%%%%

\begin{theorem}[Fundamental Fixed Point Theorem]

Every closure-complete theorem space is a structural fixed point, and every
structural fixed point represents a constitutionally complete mathematical
theory.

\end{theorem}

\begin{proof}

By definition,

\[
\operatorname{CTCA}(\mathcal T)
=
\mathcal T
\]

holds precisely when canonical closure produces no additional structurally
forced theorems.

Accordingly every closure-complete theorem space is a structural fixed point.

Conversely, every structural fixed point contains every theorem forced by its
dependency structure.

Its admissibility structure is therefore constitutionally complete.

Hence the corresponding mathematical investigation is complete.

\end{proof}

%%%%%%%%%%%%%%%%%%%%%%%%%%%%%%%%%%%%%%%%%%%%%%%%%%%%%%%%%%%%%%%%%%%%%%%%%%%%%%%
\section{Toward Mathematical Universes}
%%%%%%%%%%%%%%%%%%%%%%%%%%%%%%%%%%%%%%%%%%%%%%%%%%%%%%%%%%%%%%%%%%%%%%%%%%%%%%%

The structural fixed points established in this chapter complete individual
mathematical investigations.

A deeper question nevertheless remains.

Different structural fixed points need not be unrelated.

They themselves may possess structural dependencies.

Entire mathematical theories may therefore generate larger mathematical
landscapes whose elements are complete investigations rather than individual
theorems.

The study of these landscapes requires one final abstraction.

Instead of studying theorem spaces,

we study spaces whose points are complete mathematical theories.

The development of these spaces forms the subject of the next chapter.

%%%%%%%%%%%%%%%%%%%%%%%%%%%%%%%%%%%%%%%%%%%%%%%%%%%%%%%%%%%%%%%%%%%%%%%%%%%%%%%
\chapter{Canonical Discovery Theory}
%%%%%%%%%%%%%%%%%%%%%%%%%%%%%%%%%%%%%%%%%%%%%%%%%%%%%%%%%%%%%%%%%%%%%%%%%%%%%%%

The preceding chapters established that mathematical investigation generates
theorem spaces, that theorem spaces admit canonical closure, and that closure
naturally converges toward structural fixed points representing complete
mathematical theories.

These results determine the architecture of mathematical knowledge.

One question nevertheless remains.

How does mathematical knowledge evolve?

This question concerns neither individual theorems nor complete theories.

It concerns the process by which one transforms into the other.

The purpose of this chapter is to develop a mathematical theory of discovery
itself.

Discovery will be shown to be neither heuristic nor mysterious.

It is a canonical structural evolution generated by admissibility.

%%%%%%%%%%%%%%%%%%%%%%%%%%%%%%%%%%%%%%%%%%%%%%%%%%%%%%%%%%%%%%%%%%%%%%%%%%%%%%%
\section{Discovery as Structural Evolution}
%%%%%%%%%%%%%%%%%%%%%%%%%%%%%%%%%%%%%%%%%%%%%%%%%%%%%%%%%%%%%%%%%%%%%%%%%%%%%%%

Every canonical investigation begins with a mathematical system.

Successive stages of investigation reveal structures that were previously
implicit.

Each newly established theorem enlarges the structural organization of the
theorem space.

This enlargement generates further mathematical consequences.

Discovery therefore proceeds by successive structural evolution.

Mathematical knowledge grows because mathematical structure propagates.

%%%%%%%%%%%%%%%%%%%%%%%%%%%%%%%%%%%%%%%%%%%%%%%%%%%%%%%%%%%%%%%%%%%%%%%%%%%%%%%
\section{The Discovery Operator}
%%%%%%%%%%%%%%%%%%%%%%%%%%%%%%%%%%%%%%%%%%%%%%%%%%%%%%%%%%%%%%%%%%%%%%%%%%%%%%%

The entire process may be represented by a single operator.

\begin{definition}[Discovery operator]

The \emph{canonical discovery operator}

\[
\mathfrak D
\]

maps a theorem space to the theorem space obtained after one complete stage of
canonical structural generation.

Accordingly,

\[
\boxed{
\mathfrak D :
\mathcal T
\longrightarrow
\mathcal T.
}
\]

\end{definition}

The discovery operator does not invent mathematics.

It reveals mathematics already implicit within the existing theorem space.

%%%%%%%%%%%%%%%%%%%%%%%%%%%%%%%%%%%%%%%%%%%%%%%%%%%%%%%%%%%%%%%%%%%%%%%%%%%%%%%
\section{Discovery Trajectories}
%%%%%%%%%%%%%%%%%%%%%%%%%%%%%%%%%%%%%%%%%%%%%%%%%%%%%%%%%%%%%%%%%%%%%%%%%%%%%%%

Repeated application of the discovery operator generates a sequence

\[
\mathcal T_0,
\mathcal T_1,
\mathcal T_2,
\ldots,
\]

where

\[
\mathcal T_{n+1}
=
\mathfrak D(\mathcal T_n).
\]

This sequence will be called the discovery trajectory of the investigation.

Every stage represents a mathematically richer understanding of the same
underlying system.

%%%%%%%%%%%%%%%%%%%%%%%%%%%%%%%%%%%%%%%%%%%%%%%%%%%%%%%%%%%%%%%%%%%%%%%%%%%%%%%
\section{Discovery Before Proof}
%%%%%%%%%%%%%%%%%%%%%%%%%%%%%%%%%%%%%%%%%%%%%%%%%%%%%%%%%%%%%%%%%%%%%%%%%%%%%%%

A proof certifies an already discovered theorem.

Discovery precedes proof.

The present theory therefore distinguishes two complementary mathematical
processes.

The first determines which theorem must exist.

The second establishes that theorem rigorously.

Discovery generates necessity.

Proof certifies necessity.

Neither replaces the other.

%%%%%%%%%%%%%%%%%%%%%%%%%%%%%%%%%%%%%%%%%%%%%%%%%%%%%%%%%%%%%%%%%%%%%%%%%%%%%%%
\section{Canonical Discovery}
%%%%%%%%%%%%%%%%%%%%%%%%%%%%%%%%%%%%%%%%%%%%%%%%%%%%%%%%%%%%%%%%%%%%%%%%%%%%%%%

Not every mathematical discovery is canonical.

Some investigations proceed by experimentation.

Others rely upon analogy.

Others employ inspired conjecture.

These approaches may be extraordinarily successful.

They nevertheless contain mathematical choices not justified by the underlying
structure.

Canonical discovery excludes such choices.

\begin{definition}[Canonical discovery]

A discovery is called canonical whenever every newly introduced mathematical
object is structurally forced by the preceding theorem space.

\end{definition}

Canonical discovery therefore minimizes arbitrariness.

%%%%%%%%%%%%%%%%%%%%%%%%%%%%%%%%%%%%%%%%%%%%%%%%%%%%%%%%%%%%%%%%%%%%%%%%%%%%%%%
\section{Discovery Acceleration}
%%%%%%%%%%%%%%%%%%%%%%%%%%%%%%%%%%%%%%%%%%%%%%%%%%%%%%%%%%%%%%%%%%%%%%%%%%%%%%%

Different investigations need not evolve at the same structural rate.

Some newly established theorems generate very few additional consequences.

Others generate entire mathematical disciplines.

This motivates the introduction of a structural invariant.

\begin{definition}[Discovery rate]

The discovery rate of a theorem space is the amount of new mathematical
structure generated under one application of the discovery operator.

\end{definition}

Discovery rate measures structural productivity rather than computational
difficulty.

%%%%%%%%%%%%%%%%%%%%%%%%%%%%%%%%%%%%%%%%%%%%%%%%%%%%%%%%%%%%%%%%%%%%%%%%%%%%%%%
\section{Discovery Horizons}
%%%%%%%%%%%%%%%%%%%%%%%%%%%%%%%%%%%%%%%%%%%%%%%%%%%%%%%%%%%%%%%%%%%%%%%%%%%%%%%

At any stage of an investigation, certain mathematical structures have already
been generated.

Others remain constitutionally hidden.

The boundary separating these regions will be called the discovery horizon.

\begin{definition}[Discovery horizon]

The discovery horizon consists of those mathematical structures that are not
yet present within the current theorem space but become structurally accessible
after one further stage of canonical discovery.

\end{definition}

Discovery therefore proceeds by successive expansion of the discovery horizon.

%%%%%%%%%%%%%%%%%%%%%%%%%%%%%%%%%%%%%%%%%%%%%%%%%%%%%%%%%%%%%%%%%%%%%%%%%%%%%%%
\section{Discovery Frontiers}
%%%%%%%%%%%%%%%%%%%%%%%%%%%%%%%%%%%%%%%%%%%%%%%%%%%%%%%%%%%%%%%%%%%%%%%%%%%%%%%

The discovery horizon is local.

A broader notion concerns the global boundary of mathematical knowledge.

The collection of all presently inaccessible but constitutionally admissible
structures forms the discovery frontier.

Unlike the horizon, which changes after every stage of investigation, the
frontier characterizes the entire investigation.

%%%%%%%%%%%%%%%%%%%%%%%%%%%%%%%%%%%%%%%%%%%%%%%%%%%%%%%%%%%%%%%%%%%%%%%%%%%%%%%
\section{The Principle of Structural Discovery}
%%%%%%%%%%%%%%%%%%%%%%%%%%%%%%%%%%%%%%%%%%%%%%%%%%%%%%%%%%%%%%%%%%%%%%%%%%%%%%%

\begin{principle}[Principle of Structural Discovery]

Every canonical mathematical discovery is generated by the progressive
elimination of structural incompleteness.

\end{principle}

Discovery therefore does not add mathematics.

It removes ignorance.

%%%%%%%%%%%%%%%%%%%%%%%%%%%%%%%%%%%%%%%%%%%%%%%%%%%%%%%%%%%%%%%%%%%%%%%%%%%%%%%
\section{The Fundamental Discovery Theorem}
%%%%%%%%%%%%%%%%%%%%%%%%%%%%%%%%%%%%%%%%%%%%%%%%%%%%%%%%%%%%%%%%%%%%%%%%%%%%%%%

\begin{theorem}[Fundamental Discovery Theorem]

Every canonical mathematical investigation determines a unique discovery
trajectory converging toward a structural fixed point.

\end{theorem}

\begin{proof}

Canonical investigation generates a unique theorem space.

Canonical theorem closure determines the structurally forced enlargement of
that theorem space.

Repeated closure therefore generates a unique sequence of theorem spaces.

By the Structural Fixed Point Principle established in the preceding chapter,
this sequence converges toward a closure-complete theorem space.

Consequently every canonical investigation determines a unique discovery
trajectory terminating at a structural fixed point.

\end{proof}

%%%%%%%%%%%%%%%%%%%%%%%%%%%%%%%%%%%%%%%%%%%%%%%%%%%%%%%%%%%%%%%%%%%%%%%%%%%%%%%
\section{Discovery as a Mathematical Object}
%%%%%%%%%%%%%%%%%%%%%%%%%%%%%%%%%%%%%%%%%%%%%%%%%%%%%%%%%%%%%%%%%%%%%%%%%%%%%%%

The present chapter completes a profound transition.

Discovery is no longer regarded as lying outside mathematics.

It has become an intrinsic mathematical object.

It possesses operators.

It possesses trajectories.

It possesses horizons.

It possesses frontiers.

It possesses invariants.

It possesses completion.

It possesses fixed points.

Accordingly the mathematics of discovery now stands alongside algebra,
geometry, topology, and analysis as a legitimate branch of classical
mathematics.

%%%%%%%%%%%%%%%%%%%%%%%%%%%%%%%%%%%%%%%%%%%%%%%%%%%%%%%%%%%%%%%%%%%%%%%%%%%%%%%
\section{Toward Canonical Structural Cryptography}
%%%%%%%%%%%%%%%%%%%%%%%%%%%%%%%%%%%%%%%%%%%%%%%%%%%%%%%%%%%%%%%%%%%%%%%%%%%%%%%

Canonical discovery determines how mathematical knowledge grows.

A deeper question nevertheless remains.

Why do some mathematical truths remain hidden despite already being
constitutionally present?

Equivalently,

what mathematical structures conceal theorem generation?

These questions concern not discovery itself but the encoding and decoding of
mathematical structure.

Investigation therefore proceeds from discovery to mathematical cryptography.

The next chapter develops the theory of Canonical Structural Cryptography,
where concealment, revelation, and structural decoding become intrinsic
mathematical phenomena.

%%%%%%%%%%%%%%%%%%%%%%%%%%%%%%%%%%%%%%%%%%%%%%%%%%%%%%%%%%%%%%%%%%%%%%%%%%%%%%%
\chapter{The Constitution of Mathematical Discovery}
%%%%%%%%%%%%%%%%%%%%%%%%%%%%%%%%%%%%%%%%%%%%%%%%%%%%%%%%%%%%%%%%%%%%%%%%%%%%%%%

The preceding chapters have established a new branch of classical mathematics.

They have shown that mathematical discovery is itself a mathematical object,
that investigations generate theorem spaces, that theorem spaces admit
canonical closure, and that mathematical theories converge toward structural
fixed points determined entirely by admissibility.

One question nevertheless remains.

What principles govern every future mathematical investigation?

The purpose of this final chapter is not to introduce further mathematical
structures.

Its purpose is to state the constitutional laws that every canonical
mathematical investigation must satisfy.

These laws do not describe one particular mathematical discipline.

They govern mathematics itself.

%%%%%%%%%%%%%%%%%%%%%%%%%%%%%%%%%%%%%%%%%%%%%%%%%%%%%%%%%%%%%%%%%%%%%%%%%%%%%%%
\section{The Constitutional View of Mathematics}
%%%%%%%%%%%%%%%%%%%%%%%%%%%%%%%%%%%%%%%%%%%%%%%%%%%%%%%%%%%%%%%%%%%%%%%%%%%%%%%

Mathematics is not fundamentally the study of numbers.

Nor is it fundamentally the study of sets.

Nor of spaces.

Nor of algebraic operations.

Nor of geometric configurations.

Nor of formal deductions.

The preceding parts of this work have established a different conclusion.

Mathematics is the study of admissible structure together with the canonical
processes by which admissible structures generate further admissible
structures.

Objects therefore do not constitute mathematics.

Generation does.

%%%%%%%%%%%%%%%%%%%%%%%%%%%%%%%%%%%%%%%%%%%%%%%%%%%%%%%%%%%%%%%%%%%%%%%%%%%%%%%
\section{The First Constitutional Law}
%%%%%%%%%%%%%%%%%%%%%%%%%%%%%%%%%%%%%%%%%%%%%%%%%%%%%%%%%%%%%%%%%%%%%%%%%%%%%%%

\begin{principle}[The Law of Admissibility]

No mathematical object may be introduced unless its existence is
constitutionally admissible.

\end{principle}

Admissibility therefore precedes construction.

%%%%%%%%%%%%%%%%%%%%%%%%%%%%%%%%%%%%%%%%%%%%%%%%%%%%%%%%%%%%%%%%%%%%%%%%%%%%%%%
\section{The Second Constitutional Law}
%%%%%%%%%%%%%%%%%%%%%%%%%%%%%%%%%%%%%%%%%%%%%%%%%%%%%%%%%%%%%%%%%%%%%%%%%%%%%%%

\begin{principle}[The Law of Canonical Generation]

Every admissible mathematical structure generates all subsequent mathematical
structures by necessity rather than invention.

\end{principle}

Mathematics therefore grows through structural propagation.

%%%%%%%%%%%%%%%%%%%%%%%%%%%%%%%%%%%%%%%%%%%%%%%%%%%%%%%%%%%%%%%%%%%%%%%%%%%%%%%
\section{The Third Constitutional Law}
%%%%%%%%%%%%%%%%%%%%%%%%%%%%%%%%%%%%%%%%%%%%%%%%%%%%%%%%%%%%%%%%%%%%%%%%%%%%%%%

\begin{principle}[The Law of Structural Investigation]

Every mathematical investigation must determine intrinsic mathematical
structure before introducing derived mathematical constructions.

\end{principle}

Canonical investigation therefore precedes theorem proving.

%%%%%%%%%%%%%%%%%%%%%%%%%%%%%%%%%%%%%%%%%%%%%%%%%%%%%%%%%%%%%%%%%%%%%%%%%%%%%%%
\section{The Fourth Constitutional Law}
%%%%%%%%%%%%%%%%%%%%%%%%%%%%%%%%%%%%%%%%%%%%%%%%%%%%%%%%%%%%%%%%%%%%%%%%%%%%%%%

\begin{principle}[The Law of Canonical Observability]

Every mathematical quantity employed within an investigation must arise from
the intrinsic observable structure of the mathematical system itself.

\end{principle}

Heuristic quantities may suggest mathematics.

Canonical observables determine mathematics.

%%%%%%%%%%%%%%%%%%%%%%%%%%%%%%%%%%%%%%%%%%%%%%%%%%%%%%%%%%%%%%%%%%%%%%%%%%%%%%%
\section{The Fifth Constitutional Law}
%%%%%%%%%%%%%%%%%%%%%%%%%%%%%%%%%%%%%%%%%%%%%%%%%%%%%%%%%%%%%%%%%%%%%%%%%%%%%%%

\begin{principle}[The Law of Structural Closure]

Every canonical investigation continues until every structurally forced theorem
has been generated.

\end{principle}

Completion is therefore not optional.

It is constitutionally required.

%%%%%%%%%%%%%%%%%%%%%%%%%%%%%%%%%%%%%%%%%%%%%%%%%%%%%%%%%%%%%%%%%%%%%%%%%%%%%%%
\section{The Sixth Constitutional Law}
%%%%%%%%%%%%%%%%%%%%%%%%%%%%%%%%%%%%%%%%%%%%%%%%%%%%%%%%%%%%%%%%%%%%%%%%%%%%%%%

\begin{principle}[The Law of Structural Stability]

Every completed mathematical investigation terminates at a structural fixed
point.

\end{principle}

Mathematical theories are therefore equilibrium states of admissible
generation.

%%%%%%%%%%%%%%%%%%%%%%%%%%%%%%%%%%%%%%%%%%%%%%%%%%%%%%%%%%%%%%%%%%%%%%%%%%%%%%%
\section{The Seventh Constitutional Law}
%%%%%%%%%%%%%%%%%%%%%%%%%%%%%%%%%%%%%%%%%%%%%%%%%%%%%%%%%%%%%%%%%%%%%%%%%%%%%%%

\begin{principle}[The Law of Structural Obstruction]

Every impossibility theorem arises from an intrinsic structural obstruction.

\end{principle}

Mathematical impossibility is never accidental.

It is generated.

%%%%%%%%%%%%%%%%%%%%%%%%%%%%%%%%%%%%%%%%%%%%%%%%%%%%%%%%%%%%%%%%%%%%%%%%%%%%%%%
\section{The Eighth Constitutional Law}
%%%%%%%%%%%%%%%%%%%%%%%%%%%%%%%%%%%%%%%%%%%%%%%%%%%%%%%%%%%%%%%%%%%%%%%%%%%%%%%

\begin{principle}[The Law of Canonical Discovery]

Every mathematical discovery consists in revealing mathematical structure that
was already constitutionally implicit.

\end{principle}

Discovery therefore creates no mathematics.

It uncovers mathematics.

%%%%%%%%%%%%%%%%%%%%%%%%%%%%%%%%%%%%%%%%%%%%%%%%%%%%%%%%%%%%%%%%%%%%%%%%%%%%%%%
\section{The Ninth Constitutional Law}
%%%%%%%%%%%%%%%%%%%%%%%%%%%%%%%%%%%%%%%%%%%%%%%%%%%%%%%%%%%%%%%%%%%%%%%%%%%%%%%

\begin{principle}[The Law of Mathematical Universality]

Every branch of mathematics preserving admissibility belongs to one common
constitutional mathematics.

\end{principle}

The apparent diversity of mathematics reflects diversity of realization rather
than diversity of foundation.

%%%%%%%%%%%%%%%%%%%%%%%%%%%%%%%%%%%%%%%%%%%%%%%%%%%%%%%%%%%%%%%%%%%%%%%%%%%%%%%
\section{The Tenth Constitutional Law}
%%%%%%%%%%%%%%%%%%%%%%%%%%%%%%%%%%%%%%%%%%%%%%%%%%%%%%%%%%%%%%%%%%%%%%%%%%%%%%%

\begin{principle}[The Law of Future Mathematics]

Every future mathematical theory must itself be investigable through canonical
mathematical investigation.

\end{principle}

The present work therefore establishes not merely a new mathematical theory but
a permanent constitutional framework within which future mathematics may be
developed.

%%%%%%%%%%%%%%%%%%%%%%%%%%%%%%%%%%%%%%%%%%%%%%%%%%%%%%%%%%%%%%%%%%%%%%%%%%%%%%%
\section{The Constitution of Mathematical Discovery}
%%%%%%%%%%%%%%%%%%%%%%%%%%%%%%%%%%%%%%%%%%%%%%%%%%%%%%%%%%%%%%%%%%%%%%%%%%%%%%%

The preceding constitutional laws possess a common origin.

Every one of them expresses a different manifestation of the same underlying
principle.

Admissibility generates continuation.

Continuation generates structure.

Structure generates observability.

Observability generates theorem spaces.

Theorem spaces generate closure.

Closure generates fixed points.

Fixed points determine complete mathematical theories.

Accordingly,

\[
\boxed{
\text{Admissibility}
\Longrightarrow
\text{Continuation}
\Longrightarrow
\text{Structure}
\Longrightarrow
\text{Discovery}
\Longrightarrow
\text{Mathematics}.
}
\]

This chain is not one mathematical construction among many.

It is the constitutional architecture from which every admissible mathematics
may be recovered.

%%%%%%%%%%%%%%%%%%%%%%%%%%%%%%%%%%%%%%%%%%%%%%%%%%%%%%%%%%%%%%%%%%%%%%%%%%%%%%%
\section*{The Constitutional Theorem}
%%%%%%%%%%%%%%%%%%%%%%%%%%%%%%%%%%%%%%%%%%%%%%%%%%%%%%%%%%%%%%%%%%%%%%%%%%%%%%%

\begin{theorem}[The Constitution of Mathematical Discovery]

Every admissible mathematical theory is generated through canonical structural
investigation and is completed by canonical structural closure.

\end{theorem}

\begin{proof}

The preceding parts established that admissibility generates continuation and
that continuation generates every subsequent mathematical structure.

Part VIII established that canonical investigation determines theorem spaces,
that theorem spaces admit canonical closure, and that closure converges toward
structural fixed points representing complete mathematical theories.

Consequently every admissible mathematical theory arises through canonical
structural investigation and reaches completion through canonical structural
closure.

Therefore mathematical discovery is itself governed by constitutional laws.

\end{proof}

\vspace{2em}

\begin{center}

\Large

\emph{Mathematics is not invented.}

\vspace{0.5em}

\emph{It is constitutionally discovered.}

\end{center}

\end{document}
